% Main LaTeX file for Bayesian Data Analysis Notes
% Generated from aca-reader workflow

\documentclass[12pt]{amsbook}
\usepackage{amsmath, amssymb, amsthm}
\usepackage{xeCJK}
\usepackage{microtype}
\usepackage{hyperref}
\usepackage{cleveref}
\usepackage{geometry}
\usepackage{graphicx}
\usepackage{tikz}
\usetikzlibrary{arrows,positioning}
\geometry{margin=1in}

% Theorem styles - definition 样式：关键词加粗左对齐
\theoremstyle{definition}
\newtheorem{Definition}{定义}[chapter]
\newtheorem{Theorem}{定理}[chapter]
\newtheorem{Lemma}{引理}[chapter]
\newtheorem{Corollary}{推论}[chapter]
\newtheorem{Proposition}{命题}[chapter]
\newtheorem{Example}{例}[chapter]
\newtheorem{Remark}{注}[chapter]
\newtheorem{Exercise}{练习}[chapter]

% User annotation command
\newcommand{\userannotation}[2]{%
  \begin{Remark}
    \textbf{#1:} #2
  \end{Remark}
}

% Bayesian notation shortcuts
\newcommand{\p}[1]{p(#1)}
\newcommand{\E}{\mathbb{E}}
\DeclareMathOperator{\Var}{Var}
\DeclareMathOperator{\Cov}{Cov}
\DeclareMathOperator{\Bernoulli}{Bernoulli}
\DeclareMathOperator{\Binomial}{Binomial}
\DeclareMathOperator{\Poisson}{Poisson}
\DeclareMathOperator{\Normal}{Normal}
\DeclareMathOperator{\Beta}{Beta}
\DeclareMathOperator{\Dirichlet}{Dirichlet}
\DeclareMathOperator{\Exponential}{Exponential}

\title{贝叶斯数据分析讲义}
\author{Gelman, Carlin, Stern, Dunson, Vehtari, Rubin}
\date{\today}

\begin{document}
\addtocontents{toc}{\protect\thispagestyle{empty}}

\maketitle
\tableofcontents

% Chapter 0: 课程概述
% Chapter 0: 引言 - 文献概述
% Bayesian Data Analysis - Gelman et al.

\chapter*{引言：文献概述}
\addcontentsline{toc}{chapter}{引言：文献概述}
\label{chap:introduction}

\section{什么是贝叶斯数据分析}

贝叶斯数据分析是一种统计推断方法，其核心思想是使用概率来量化不确定性。在贝叶斯框架下，所有未知量——包括参数、潜在变量、甚至未来观测——都被视为随机变量，用概率分布来描述。

与传统频率主义统计不同，贝叶斯方法明确使用概率来表达我们对未知量的信念（belief）。这种方法的主要优势包括：

\begin{itemize}
  \item \textbf{自然的不确定性量化}：贝叶斯概率区间可以直接解释为"未知量包含在该区间内的概率"
  \item \textbf{灵活的建模能力}：可以构建复杂的多层次模型
  \item \textbf{原则性的推断框架}：贝叶斯定理为所有推断提供了统一的基础
  \item \textbf{先验信息的整合}：允许将领域知识融入统计分析
\end{itemize}

贝叶斯数据分析的过程可以概括为三个步骤：
\begin{enumerate}
  \item 建立完整的概率模型
  \item 基于观测数据计算后验分布
  \item 评估模型的拟合度和敏感性
\end{enumerate}

\section{核心概念}

\subsection{先验分布（Prior Distribution）}

先验分布代表在观测数据之前，我们对未知参数 $\theta$ 的已有知识或信念。先验分布可以是：
\begin{itemize}
  \item \textbf{信息先验}：基于领域知识或历史数据
  \item \textbf{非信息先验}：尽量减少先验对结果的影响（如 Jeffreys 先验）
  \item \textbf{弱信息先验}：包含一定信息但不完全反映领域知识
\end{itemize}

\subsection{后验分布（Posterior Distribution）}

后验分布是贝叶斯推断的核心——在观测到数据 $y$ 后，关于参数 $\theta$ 的更新后的知识：
\[
p(\theta|y) = \frac{p(y|\theta) \cdot p(\theta)}{p(y)}
\]
其中 $p(y|\theta)$ 是似然函数，$p(y)$ 是边际似然（归一化常数）。

\subsection{预测分布（Predictive Distribution）}

贝叶斯方法可以自然地进行预测推断：
\begin{itemize}
  \item \textbf{先验预测分布}：$p(y) = \int p(y|\theta) p(\theta) d\theta$
  \item \textbf{后验预测分布}：$p(\tilde{y}|y) = \int p(\tilde{y}|\theta) p(\theta|y) d\theta$
\end{itemize}

\section{文献目录}\label{sec:literature}

\subsection{核心教材}

\begin{enumerate}
  \item \textbf{\textit{Bayesian Data Analysis, Third Edition}} — Gelman, Carlin, Stern, Dunson, Vehtari, Rubin

  本课程的核心教材，656页，系统介绍了贝叶斯数据分析的理论与实践。

  \begin{itemize}
    \item Part I: 贝叶斯推断基础（概率、三步骤、单/多参数模型、分层模型）
    \item Part II: 数据分析基础（模型检查、评估、决策分析）
    \item Part III: 高级计算方法（MCMC、变分推断）
    \item Part IV: 回归模型
    \item Part V: 非参数模型
  \end{itemize}
\end{enumerate}

\subsection{课件}

\begin{enumerate}
  \item \textbf{Three Steps of Bayesian Data Analysis} — Chapter 1 课件

  \item \textbf{Single-parameter Models} — Chapter 2 课件

  \item \textbf{Multi-parameter Models} — Chapter 3 课件

  \item \textbf{Hierarchical Models} — Chapter 5 课件

  \item \textbf{Model Checking} — Chapter 6 课件

  \item \textbf{Bayesian Workflow} — Chapter 10 课件

  \item \textbf{Advanced Topics} — Chapter 11 课件
\end{enumerate}

\subsection{补充材料}

\begin{enumerate}
  \item \textbf{Metropolis 算法简要证明}

  MCMC 方法中 Metropolis 算法的理论基础与证明。

  \item \textbf{R 语言实现抽样基础}

  使用 R 语言实现各种抽样方法的实践指南。

  \item \textbf{p 值补充}

  p 值的定义、解释与局限性分析。

  \item \textbf{拒绝抽样补充}

  拒绝抽样方法的详细介绍与实现。
\end{enumerate}

\subsection{参考教材}

\begin{enumerate}
  \item \textbf{Markov Chain Monte Carlo in Practice} — Gilks, Richardson, Spiegelhalter

  MCMC 方法的实践指南，是 BDA 第11-12章的重要补充。

  \item \textbf{Monte Carlo Statistical Methods} — Robert, Casella

  蒙特卡洛统计方法的理论教材，包含更深入的算法分析。

  \item \textbf{Bayesian Statistical Modelling} — Congdon

  贝叶斯统计建模的实用参考，提供了丰富的建模案例。
\end{enumerate}

\section{课程结构}

本课程将按照 BDA 第三版的结构依次学习：

\begin{itemize}
  \item \textbf{Part I（第1-5章）}：贝叶斯推断基础
  \item \textbf{Part II（第6-9章）}：数据分析基础
  \item \textbf{Part III（第10-13章）}：高级计算方法
  \item \textbf{Part IV（第14-18章）}：回归模型
  \item \textbf{Part V（第19-23章）}：非参数模型
\end{itemize}

\section{学习方法}

\begin{Remark}
  建议的学习方法：
  \begin{enumerate}
    \item 先阅读课件把握整体框架
    \item 深入阅读教材对应章节
    \item 完成课后习题巩固理解
    \item 通过编程实践加深概念
  \end{enumerate}
\end{Remark}

\endinput


% Part I: Fundamentals of Bayesian Inference
\part{贝叶斯推断基础}

% Chapter 1: Probability and inference
% Bayesian Data Analysis - Gelman et al.

\chapter{概率与推断}

\section{贝叶斯数据分析的三步骤}

贝叶斯数据分析的过程可以理想化地分为以下三个步骤：

\begin{enumerate}
  \item \textbf{建立完整的概率模型}：对问题中所有可观测和不可观测的量建立一个联合概率分布。模型应与底层科学问题（scientific problem）和数据收集过程的知识相一致。

  \item \textbf{基于观测数据进行分析}：计算并解释适当的后验分布——即在给定观测数据的条件下，关于不可观测量的条件概率分布。

  \item \textbf{评估模型的拟合度和后验分布的含义}：
  \begin{itemize}
    \item 模型对数据的拟合程度如何？
    \item 实质性结论（substantive conclusions）是否合理？
    \item 结果对第一步中建模假设的敏感性如何？
  \end{itemize}
  作为回应，可以修改或扩展模型并重复这三个步骤。
\end{enumerate}

第二个步骤涉及计算方法论，第三个步骤涉及微妙的技术与判断平衡，均由问题的应用背景指导。

\textbf{第一步仍然是许多贝叶斯分析的主要障碍}：我们的模型从何而来？如何构建适当的概率规范？

贝叶斯思维的主要动机是它便于对统计结论进行常识性解释。例如，一个未知量的贝叶斯（概率）区间可以直接被认为是"该未知量有很高概率包含在其中"，这与频率主义（置信区间）的解释形成对比——后者在严格解释下只能关联到重复抽样序列中的类似推断。

\section{统计推断的一般符号}

在贝叶斯统计中，我们使用以下符号：

\begin{itemize}
  \item $\theta$：不可观测的向量值或总体参数（如临床试验中每种治疗方案的生存概率）

  \item $y$：观测数据（如每组中的存活和死亡人数）

  \item $\tilde{y}$：未知但潜在可观测的量（如其他治疗方案下患者的结果，或与试验中患者相似的新患者的结果）

  \item $x$：解释变量或协变量（如患者的年龄和先前健康状况）
\end{itemize}

这些概率陈述以观测值 $y$ 为条件，在我们记号中简单记为 $p(\theta|y)$ 或 $p(\tilde{y}|y)$。我们也隐含地以协变量的已知值为条件。

\section{贝叶斯推断}

贝叶斯推断的核心思想是：关于参数 $\theta$ 或未观测数据 $\tilde{y}$ 的统计结论，都以概率陈述的形式给出。

这些概率陈述以观测值 $y$ 为条件，在我们的记号中简单记为 $p(\theta|y)$ 或 $p(\tilde{y}|y)$。

\subsection{贝叶斯定理}

为了对给定 $y$ 的 $\theta$ 做出概率陈述，我们必须首先建立一个提供 $\theta$ 和 $y$ 联合概率分布的模型。

联合概率质量（或密度）函数可以写成两个密度的乘积，通常称为先验分布 $p(\theta)$ 和采样分布（或数据分布）$p(y|\theta)$：
\[
p(\theta, y) = p(\theta) \times p(y|\theta)
\]

使用条件概率的基本性质（贝叶斯定理），得到后验密度：\footnote{推导见附录 \cref{sec:bayes-theorem-derivation}}
\[
p(\theta|y) = \frac{p(\theta) \times p(y|\theta)}{p(y)}
\]
其中
\[
p(y) = \int p(\theta) \times p(y|\theta) \, d\theta
\]

一个等价的形式省略了不依赖于 $\theta$ 的因子 $p(y)$，给出非归一化后验密度：
\[
p(\theta|y) \propto p(\theta) \times p(y|\theta)
\]

\textbf{为什么可以省略 $p(y)$？}

关键在于：后验分布的\textbf{核}（kernel）\footnote{关于"核"概念的详细解释，见附录 \cref{sec:qa:DistributionKernel}}包含了所有关于 $\theta$ 的信息，而 $p(y)$ 作为归一化常数不包含任何关于 $\theta$ 的新信息。

以正态分布为例：后验密度 $p(\theta|y) \propto \exp\left(-\frac{(\theta-\mu)^2}{2\sigma^2}\right)$，我们只凭这个指数函数的形式就能识别出这是正态分布 $N(\mu, \sigma^2)$——因为已知核是二次型的指数，就能确定完整分布。

换句话说，$p(y)$ 的作用只是确保后验密度积分为 1，但它的具体值对于理解 $\theta$ 如何随数据更新并无贡献。

\begin{Remark}
  注意：$p(y|\theta)$ 被视为 $\theta$ 的函数（固定 $y$），而不是 $y$ 的函数。这被称为似然函数。
\end{Remark}

\subsection{预测}

在数据 $y$ 被观测之前，未知但可观测的 $\tilde{y}$ 的分布是边际分布：\footnote{推导见附录 \cref{sec:predictive-distributions-derivation}}
\[
p(y) = \int p(y|\theta) \times p(\theta) \, d\theta
\]
这被称为先验预测分布（prior predictive distribution）。

在观测到数据 $y$ 之后，预测 $\tilde{y}$ 的分布是后验预测分布（posterior predictive distribution）：\footnote{推导见附录 \cref{sec:predictive-distributions-derivation}}
\[
p(\tilde{y}|y) = \int p(\tilde{y}|\theta) \times p(\theta|y) \, d\theta
\]

后验预测分布可以解释为在 $\theta$ 的后验分布上对条件预测进行平均。

\begin{Example}[称重例子]
  令 $y = (y_1, \ldots, y_n)$ 为物体在秤上称量 $n$ 次的记录重量，$\theta = (\mu, \sigma^2)$ 为物体的真实重量和测量方差，$\tilde{y}$ 为计划中新的称量结果。

  第二行展示了后验预测分布作为后验分布上条件预测的平均。最后一步来自在给定 $\theta$ 的条件下 $y$ 和 $\tilde{y}$ 的条件独立性假设。
\end{Example}

\section{似然与优势比}

使用贝叶斯定理与选定的概率模型意味着数据 $y$ 只通过 $p(y|\theta)$ 影响后验推断，当固定 $y$ 将其视为 $\theta$ 的函数时，称为\textbf{似然函数}。

这样，贝叶斯推断服从所谓的似然原理：对于给定的数据样本，任何具有相同似然函数的概率模型 $p(y|\theta)$ 对 $\theta$ 产生相同的推断。

后验密度 $p(\theta|y)$ 在 $\theta_1$ 和 $\theta_2$ 处取值的比率称为后验优势比。优势比提供了一个概率的替代表示，并且具有贝叶斯定理用优势比表示特别简单的形式：\footnote{推导见附录 \cref{sec:posterior-odds-derivation}}
\[
\frac{p(\theta_1|y)}{p(\theta_2|y)} = \frac{p(\theta_1)}{p(\theta_2)} \times \frac{p(y|\theta_1)}{p(y|\theta_2)}
\]

即：\textbf{后验优势比 = 先验优势比 × 似然比}。

\section{概率作为不确定性的度量}

在贝叶斯统计中，概率被用作不确定性的基本度量（fundamental measure or yardstick of uncertainty）。在这个范式下，讨论明天降雨的概率或巴西队赢得世界杯的概率与讨论硬币抛掷正面朝上的概率同样合法。

\textbf{为什么概率是量化不确定性的合理方式？}通常有以下理由：

\begin{enumerate}
  \item \textbf{类比}：物理随机性引起不确定性，因此用随机事件的语言描述不确定性似乎是合理的。日常语言使用"可能"和"不太可能"等术语，将更正式的概率计算扩展到科学推断问题是与之一致的。

  \item \textbf{公理化或规范性方法}：与决策理论相关，将所有统计推断置于收益和损失的决策背景下。那么合理的公理（排序、传递性等）意味着不确定性必须用概率表示。

  \item \textbf{博彩一致性}：定义你赋予事件 $E$ 的概率 $p$（$p \in [0, 1]$）为你愿意以此比例交换（即赌）\$p 来换取如果 $E$ 发生则获得\$1 的回报。也就是说，如果 $E$ 发生，你额外获得 $(1-p)$；如果 $E$ 的补集发生，你损失 $p$。
\end{enumerate}

一致性原则（coherence principle）表明：你对所有可能事件的概率赋值不应该让对方有可能从你这里确定获利。可以证明在这种原则下构建的概率必须满足概率论的基本公理。

博彩基本原理有一些根本性困难：
\begin{itemize}
  \item 需要精确的赔率，你愿意在两个方向中的任何一个方向下赌注。对于所有事件如何指定精确赔率？如果你不确定呢？
  \item 如果一个人愿意与你对赌，而且他有你不知道的信息，你接受赌注可能不明智。
\end{itemize}

这些考虑表明概率可能是总结应用统计中不确定性的合理方法，但最终证明在于应用的成功。

\subsection{两种经典论证}

\textbf{对称性/可交换性论证}：
\[
\text{概率} = \frac{\text{有利情况数}}{\text{可能情况数}}
\]

对于硬币抛掷，这实际上是一个物理论证，基于关于决定硬币如何落地的力量以及抛掷的初始物理条件的假设。

\textbf{频率论证}：
\[
\text{概率} = \text{在长期重复序列中得到的相对频率}
\]

假设以相同方式物理独立地进行。

这两种论证在某种意义上都是主观的，因为它们需要对硬币性质和抛掷程序的判断，两者都涉及关于等可能性事件、相同测量和独立性的语义论证。

频率论证可能有某些特殊困难，因为它涉及长期相同抛掷序列的假设性概念。如果严格遵循，这种观点不允许对不属于长期相同事件序列的单一抛掷做出概率陈述。

\section{主观性与客观性}

所有使用概率的统计方法在某种意义上都是主观的，因为它们依赖于对世界的数学理想化。

贝叶斯方法有时被认为特别主观，因为它们依赖于先验分布。但在大多数问题中，科学判断对于指定似然和先验分布都是必要的。

一个普遍原则是：每当存在复制时（即观察到许多可交换的单位），就有机会从数据中估计概率分布的特征，从而使分析更加客观。

如果整个实验被重复若干次，那么先验分布的参数本身也可以从数据中估计。

\section{概率赋值的例子：足球点差}

在足球比赛中，"点差"（point spread）是博彩公司设定的赔率，用于使两边下注相对均衡。例如，如果热门队（favorite）的点差是 8 分，意味着博彩公司认为热门队平均会赢 8 分。

给定点差，我们可以估计各种条件概率：
\begin{itemize}
  \item $\Pr(\text{热门队赢} | \text{点差}=8)$
  \item $\Pr(\text{热门队至少赢 8 分} | \text{点差}=8)$
  \item $\Pr(\text{热门队至少赢 8 分} | \text{点差}=8 \text{ 且 热门队赢})
\end{itemize}

这些概率可以通过相对频率估计，也可以用正态近似来估计（outcome - point spread 的分布）。

\section{概率论的一些有用结果}

见教材。

\section{本章小结}

本章我们学习了：
\begin{enumerate}
  \item 贝叶斯数据分析的三个基本步骤
  \item 统计推断的一般符号 ($\theta, y, \tilde{y}, x$)
  \item 贝叶斯定理及其组成部分（先验、似然、后验、边际似然）
  \item 预测分布：先验预测分布与后验预测分布
  \item 概率作为不确定性度量的多种解释
  \item 似然原理和优势比
\end{enumerate}

\section{下节预告}

下一章我们将学习单参数模型，包括从二项数据估计概率的经典例子——这是贝叶斯推断最经典的入门案例。

\endinput

% Chapter 2: Single-parameter models
% Bayesian Data Analysis - Gelman et al.

\chapter{单参数模型}

\section{从二项数据估计概率}

考虑从二项分布 $y \sim \text{Binomial}(n, \theta)$ 中观测到 $y$ 次成功，其中：
\begin{itemize}
  \item $y$：$n$ 次伯努利试验中的成功总次数
  \item $\theta$：总体中的成功比例，或每次试验的成功概率
\end{itemize}

二项分布的似然函数为：
\[
p(y|\theta) = \binom{n}{y} \theta^y (1-\theta)^{n-y}
\]

为了执行贝叶斯推断，我们必须为 $\theta$ 指定先验分布。作为简单开始，我们假设先验分布是 $[0, 1]$ 上的均匀分布：$\theta \sim \text{Uniform}(0, 1)$，即 $p(\theta) = 1$。

应用贝叶斯定理，得到后验密度：\footnote{推导见附录 \cref{sec:beta-binomial-posterior-derivation}}
\[
p(\theta|y) \propto \theta^y (1-\theta)^{n-y}
\]

这正是 Beta 分布的形式，因此：
\[
\theta|y \sim \text{Beta}(y+1, n-y+1)
\]

\begin{Definition}[Beta 分布]
  Beta 分布的概率密度函数为：
  \[
  p(x|\alpha, \beta) = \frac{\Gamma(\alpha+\beta)}{\Gamma(\alpha)\Gamma(\beta)} x^{\alpha-1}(1-x)^{\beta-1}, \quad 0 \leq x \leq 1
  \]
\end{Definition}

在固定 $n$ 和 $y$ 的情况下，组合数 $\binom{n}{y}$ 不依赖于未知参数 $\theta$，因此在计算后验分布时可以视为常数。

\subsection{预测}

在均匀先验下，先验预测分布可以直接计算。在模型下，所有可能的 $y$ 值在先验下同等可能。

令 $\tilde{y}$ 表示与前 $n$ 次试验可交换的一次新试验的结果：\footnote{推导见附录 \cref{sec:laplace-succession-derivation}}
\[
P(\tilde{y} = 1 | y) = \int_0^1 \theta \cdot p(\theta|y) \, d\theta = \mathbb{E}(\theta|y) = \frac{y+1}{n+2}
\]

这就是著名的拉普拉斯继承法则（Laplace's law of succession）。在极端观测 $y=0$ 和 $y=n$ 时，拉普拉斯法则预测概率分别为 $1/(n+2)$ 和 $(n+1)/(n+2)$。

\section{后验作为数据与先验信息的折中}

贝叶斯推断的过程涉及从先验分布 $p(\theta)$ 到后验分布 $p(\theta|y)$ 的转化。以下两个式子描述了先验与后验之间的一般关系：\footnote{推导见附录 \cref{sec:law-of-total-expectation-variance-derivation}}

\[
\mathbb{E}(\theta) = \mathbb{E}\left[\mathbb{E}(\theta|y)\right]
\]

\[
\operatorname{Var}(\theta) = \mathbb{E}\left[\operatorname{Var}(\theta|y)\right] + \operatorname{Var}\left[\mathbb{E}(\theta|y)\right]
\]

第一个式子（law of total expectation）几乎不足为奇：$\theta$ 的先验均值是所有可能数据上后验均值的平均。

第二个式子（law of total variance）更有趣：后验方差平均上小于先验方差，两者之差取决于后验均值在不同数据上的变化程度。变化越大，我们降低不确定性的潜力越大。

\subsection{二项例子中的折中}

在二项例子中，均匀先验的均值是 $1/2$，方差是 $1/12$。

后验均值 $\frac{y+1}{n+2}$ 是先验均值 $1/2$ 和样本比例 $y/n$ 之间的折中，显然，随着数据样本量增大，先验的作用越来越小。

这是贝叶斯推断的一般特征：\textbf{后验分布位于代表先验信息和数据信息的点之间，折中程度随着样本量增加越来越受数据控制}。

\section{后验推断的汇总}

后验概率分布包含了关于参数 $\theta$ 的所有当前信息。常用的位置汇总包括后验均值、后验中位数、后验众数；变化程度通常用标准差、四分位距和其他分位数来汇总。

当后验分布具有闭式形式时（如当前例子中的 Beta 分布），这些汇总通常可以用闭式形式得到。

\subsection{后验区间}

除了点汇总之外，报告后验不确定性几乎总是重要的。通常的方法是给出后验分布的分位数，或者如果需要区间汇总，则给出一定概率的后验中心区间（如 $100(1-\alpha)\%$ 区间）。这类区间估计称为\textbf{后验区间}（posterior interval）。

\subsection{最高后验密度区间}

另一种汇总后验不确定性的方法是\textbf{最高后验密度区域}（Highest Posterior Density, HPD region）：包含 $100(1-\alpha)\%$ 后验概率的集合，且区域内密度不低于区域外。

对于单峰且对称的后验分布，HPD 区间与后验中心区间相同。但对于非对称分布，两者可能有显著差异——当后验分布有偏时，HPD 区间可以更准确地反映真实的不确定性范围。

\section{信息先验分布}

在二项例子中，到目前为止我们只考虑了 $\theta$ 的均匀先验分布。问题在于：如何证明这种指定的合理性？一般来说，我们如何处理构建先验分布的问题？

对先验分布有两种基本解释：
\begin{enumerate}
  \item \textbf{总体解释}：先验分布代表可能参数值的总体，当前感兴趣的 $\theta$ 是从中抽取的。
  \item \textbf{主观知识状态解释}：指导原则是，我们必须将关于 $\theta$ 的知识（和不确定性）表达得像它的值是先验分布的随机实现一样。
\end{enumerate}

对于许多问题（如估计新工业过程中失败概率），不存在完全相关的 $\theta$ 的总体。先验分布应包含所有合理的 $\theta$ 值，但分布不必真实集中在真值附近，因为通常数据中关于 $\theta$ 的信息远超过任何合理先验概率规范。

\subsection{共轭先验}

考虑二项模型的似然：$p(y|\theta) \propto \theta^y (1-\theta)^{n-y}$

如果先验密度具有相同形式，则后验密度也具有这种形式。我们将先验密度参数化为：
\[
p(\theta) \propto \theta^{\alpha-1}(1-\theta)^{\beta-1}
\]
即 $\theta \sim \text{Beta}(\alpha, \beta)$。

将 $p(\theta)$ 与 $p(y|\theta)$ 比较，可以看出这个先验等价于 $\alpha-1$ 次先验成功和 $\beta-1$ 次先验失败。

\textbf{共轭性}：后验分布与先验分布属于同一参数形式。Beta 先验分布是二项似然的共轭族。共轭族在计算上是方便的，因为后验分布遵循已知参数形式。\footnote{推导见附录 \cref{sec:beta-binomial-moments-derivation}}

后验均值：
\[
\mathbb{E}(\theta|y) = \frac{\alpha+y}{\alpha+\beta+n}
\]
后验方差：
\[
\operatorname{Var}(\theta|y) = \frac{(\alpha+y)(\beta+n-y)}{(\alpha+\beta+n)^2(\alpha+\beta+n+1)}
\]

当 $y$ 和 $n-y$ 变大而 $\alpha$ 和 $\beta$ 固定时：
\[
\mathbb{E}(\theta|y) \approx \frac{y}{n}, \quad \operatorname{Var}(\theta|y) \approx \frac{1}{n} \cdot \frac{y}{n}\left(1-\frac{y}{n}\right)
\]
方差以 $1/n$ 的速率趋近于零。极限情况下，先验分布的参数对后验分布没有影响。

\section{正态均值的估计（已知方差）}

考虑来自正态分布 $y \sim \text{Normal}(\theta, \sigma^2)$ 的单个标量观测 $y$，其中 $\sigma^2$ 已知。

采样分布：
\[
p(y|\theta) = \frac{1}{\sqrt{2\pi\sigma^2}} \exp\left(-\frac{(y-\theta)^2}{2\sigma^2}\right)
\]

作为 $\theta$ 的函数，似然是二次型的指数，因此共轭先验族为：
\[
p(\theta) \propto \exp(A\theta^2 + B\theta + C), \quad A < 0
\]

参数化为 $\theta \sim \text{Normal}(\mu_0, \tau_0^2)$，其中 $\mu_0$ 和 $\tau_0^2$ 是超参数。\footnote{推导见附录 \cref{sec:normal-mean-posterior-derivation}}

经过配方法，得到后验分布 $\theta|y \sim \text{Normal}(\mu_1, \tau_1^2)$，其中：
\[
\mu_1 = \frac{\mu_0/\tau_0^2 + y/\sigma^2}{1/\tau_0^2 + 1/\sigma^2}, \quad \frac{1}{\tau_1^2} = \frac{1}{\tau_0^2} + \frac{1}{\sigma^2}
\]

\textbf{精度}：方差的倒数称为精度。后验精度等于先验精度加上数据精度。

后验均值 $\mu_1$ 可以多种方式解释：
\begin{itemize}
  \item 先验均值和观测值 $y$ 的加权平均，权重与精度成正比
  \item 先验均值向观测值调整：$\mu_1 = \mu_0 + (y-\mu_0)\frac{\tau_0^2}{\sigma^2+\tau_0^2}$
  \item 数据向先验均值收缩：$\mu_1 = y - (y-\mu_0)\frac{\sigma^2}{\sigma^2+\tau_0^2}$
\end{itemize}

每种形式都表示后验均值是先验均值和观测值之间的折中。

\subsection{后验预测分布}\footnote{推导见附录 \cref{sec:predictive-mean-variance-derivation}}

未来观测 $\tilde{y}$ 的后验预测分布：
\[
\tilde{y}|y \sim \text{Normal}(\mu_1, \sigma^2 + \tau_1^2)
\]

预测均值等于 $\theta$ 的后验均值；预测方差有两个组成部分：来自模型的方差 $\sigma^2$ 和来自 $\theta$ 后验不确定性的方差 $\tau_1^2$。

\subsection{多观测的正态模型}

上述单观测正态模型可以轻松推广到 $n$ 个独立同分布观测 $y = (y_1, \ldots, y_n)$。因为 $\bar{y}|\theta \sim \text{Normal}(\theta, \sigma^2/n)$，对单观测推导的结果直接适用（将 $\bar{y}$ 视为单观测）:
\[
\theta|y_1,\ldots,y_n \sim \text{Normal}(\mu_n, \tau_n^2)
\]
其中：
\[
\mu_n = \frac{\mu_0/\tau_0^2 + n\bar{y}/\sigma^2}{1/\tau_0^2 + n/\sigma^2}, \quad
\frac{1}{\tau_n^2} = \frac{1}{\tau_0^2} + \frac{n}{\sigma^2}
\]

\section{非信息先验分布}

当先验分布没有总体基础时，很难构建，长期以来人们渴望得到保证在后验分布中作用最小的先验分布。这类分布有时称为参考先验分布，先验密度被称为模糊、平坦、扩散或非信息的。

使用非信息先验的理由通常是"让数据自己说话"，使推断不受当前数据外部信息的影响。

一个相关想法是\textbf{弱信息先验分布}，它包含一些信息——足以正则化后验分布（即保持它在合理范围内），但不试图完全捕捉关于底层参数的科学知识。

\subsection{正常与非正常先验}

如果先验精度 $1/\tau_0^2$ 相对于数据精度 $n/\sigma^2$ 很小，则后验分布近似于 $\tau_0^2 = \infty$ 的情况：
\[
p(\theta|y) \approx \text{Normal}(\bar{y}, \sigma^2/n)
\]

换句话说，后验分布近似于假设 $p(\theta)$ 与常数成正比的情况。这种分布实际上不可能，因为 $p(\theta)$ 的积分是无穷大，违反了概率总和为 1 的假设。

\begin{Definition}
  如果先验密度 $p(\theta)$ 不依赖于数据且积分等于 1，则称为正常（proper）先验。如果 $p(\theta)$ 积分到任何正有限值，可以重新归一化使其积分到 1，则称为非归一化密度。如果积分不是有限的，称为非正常（improper）先验。
\end{Definition}

非正常先验可能导致正常后验分布。使用非正常先验得到的后验分布必须非常小心解释——必须始终检查后验分布的积分是有限的且形式合理。

\section{Jeffreys 不变性原则}

一种有时用于定义非信息先验的方法是由 Jeffreys 提出的，基于考虑参数的一一变换：$\phi = h(\theta)$。

Jeffreys 的一般原则是：任何确定先验密度 $p(\theta)$ 的规则，如果应用于变换后的参数，应该产生等价的结果。

考虑在 $\theta$ 上使用均匀先验 $p(\theta) = 1$，但在 logit 变换 $\phi = \log(\theta/(1-\theta))$ 下，这会导致不均匀的先验。这说明均匀先验在参数变换下不是不变的。

Jeffreys 提出可接受的非信息先验发现原则应该在参数的单调变换下是不变的。他描述了如何基于 Fisher 信息函数构建这样的先验：\footnote{推导见附录 \cref{sec:jeffreys-prior-derivation}}

\[
p(\theta) \propto [J(\theta)]^{1/2}
\]
其中 $J(\theta)$ 是 $\theta$ 的 Fisher 信息：
\[
J(\theta) = -E\left[\frac{d^2 \log p(y|\theta)}{d\theta^2}\right]
\]

对于二项模型，Jeffreys 先验是 $p(\theta) \propto \theta^{-1/2}(1-\theta)^{-1/2}$，即 $\text{Beta}(1/2, 1/2)$。

注意：均匀先验在参数变换下不是不变的——$\theta$ 上的均匀先验在 $\log(\theta/(1-\theta))$ 上不是均匀的。这曾是 20 世纪初对贝叶斯推断的主要批评之一，直到 Jeffreys 在世纪中叶复兴了这一主题。

\section{弱信息先验分布}

弱信息先验是比非信息先验更强的正则化，但仍然比完全信息先验弱得多。它们用于在数据有限时稳定计算，同时允许数据主导推断。

\section{其他标准单参数模型}

见教材。

\section{本章小结}

本章我们学习了：
\begin{enumerate}
  \item 二项模型的贝叶斯推断：Beta-Binomial 共轭
  \item 后验作为数据与先验信息的折中（方差分解公式）
  \item 后验推断的汇总：点估计（均值、中位数、众数）和区间估计
  \item 最高后验密度区间（HPD）
  \item 信息先验与共轭先验
  \item 正态均值估计（已知方差）：精度加法公式
  \item 非信息先验与 Jeffreys 先验
  \item 弱信息先验
  \item 正常与非正常先验
\end{enumerate}

\section{下节预告}

下一章我们将学习多参数模型，包括如何处理 nuisance 参数和正态模型中均值与方差同时未知的情况。

\endinput

% Chapter 3: Introduction to multiparameter models
% Bayesian Data Analysis - Gelman et al.

\chapter{多参数模型简介}

\section{对 Nuisance 参数的积分}

为了表达联合后验分布和边际后验分布的思想，假设 $\theta$ 有两部分，每部分都可以是向量：$\theta = (\theta_1, \theta_2)$。进一步假设我们只（至少目前只）对 $\theta_1$ 的推断感兴趣，因此 $\theta_2$ 可以被视为 nuisance 参数。

例如，在简单例子 $y|\mu, \sigma^2 \sim \text{Normal}(\mu, \sigma^2)$ 中，$\mu = \theta_1$ 和 $\sigma^2 = \theta_2$ 都未知，但通常我们最关心的是 $\mu$。

我们寻求给定观测数据的参数条件分布，即 $p(\theta_1|y)$。这从联合后验密度
\[
p(\theta_1, \theta_2|y) \propto p(y|\theta_1, \theta_2) \times p(\theta_1, \theta_2)
\]
通过对 $\theta_2$ 积分得到：
\[
p(\theta_1|y) = \int p(\theta_1, \theta_2|y) \, d\theta_2
\]

另一种分解方式是：
\[
p(\theta_1|y) = \int p(\theta_1|\theta_2, y) \times p(\theta_2|y) \, d\theta_2
\]

它表明：感兴趣的\textbf{后验分布是条件后验分布的混合}，其中 $p(\theta_2|y)$ 是对 $\theta_2$ 不同值的加权函数。权重依赖于 $\theta_2$ 的后验密度，因此依赖于数据和先验模型的组合。

对 nuisance 参数 $\theta_2$ 的积分可以一般性地解释。例如，$\theta_2$ 可以包括代表不同可能子模型的离散成分。

\begin{Remark}
  我们很少直接计算这个积分，但它暗示了构建和计算多参数模型的重要策略：\textbf{后验分布可以通过边际和条件模拟来计算}——首先从边际后验分布抽取 $\theta_2$，然后在给定 $\theta_2$ 值的条件下从条件后验分布抽取 $\theta_1$。通过这种方式，积分被间接执行。
\end{Remark}

这种分析形式的典型例子是正态模型中均值和方差都未知的情况。

\section{正态数据与非信息先验}

作为从样本估计总体均值的原型例子，我们考虑来自一元正态分布 $\text{Normal}(\mu, \sigma^2)$ 的 $n$ 个独立观测 $y = (y_1, \ldots, y_n)$。

我们首先在非信息先验分布下分析这个模型。

\subsection{非信息先验分布}

对于 $\mu$ 和 $\sigma^2$，一个合理的模糊先验密度，假设位置参数和尺度参数的先验独立，是在 $(\mu, \log \sigma^2)$ 上均匀，即：
\[
p(\mu, \sigma^2) \propto (\sigma^2)^{-1}
\]

\subsection{联合后验分布}

在这种约定的非正常先验密度下，联合后验分布为：
\[
p(\mu, \sigma^2 | y) \propto (\sigma^2)^{-(n/2+1)} \exp\left(-\frac{1}{2\sigma^2} \sum_{i=1}^n (y_i - \mu)^2\right)
\]

利用 $\sum_{i=1}^n (y_i - \mu)^2 = \sum_{i=1}^n (y_i - \bar{y})^2 + n(\bar{y} - \mu)^2$，可以写成：
\[
= (\sigma^2)^{-(n/2+1)} \exp\left(-\frac{1}{2\sigma^2}\left[(n-1)s^2 + n(\bar{y} - \mu)^2\right]\right)
\]
其中 $\bar{y}$ 是样本均值，$s^2 = \frac{1}{n-1}\sum_{i=1}^n (y_i - \bar{y})^2$ 是样本方差。

\textbf{充分统计量}是 $\bar{y}$ 和 $s^2$。

\subsection{条件后验分布}

为了分解联合后验密度，我们首先考虑条件后验密度 $p(\mu | \sigma^2, y)$，然后考虑边际后验密度 $p(\sigma^2 | y)$。

对于给定 $\sigma^2$ 的 $\mu$ 的后验分布，直接使用第二章中已知方差正态均值的均匀先验结果：
\[
\mu | \sigma^2, y \sim N(\bar{y}, \sigma^2 / n)
\]

\subsection{边际后验分布}

为了确定 $p(\sigma^2 |y)$，我们必须对 $\mu$ 积分：
\[
p(\sigma^2 | y) \propto \int (\sigma^2)^{-(n/2+1)} \exp\left(-\frac{1}{2\sigma^2}\left[(n-1)s^2 + n(\bar{y} - \mu)^2\right]\right) d\mu
\]

对 $\mu$ 积分需要计算简单的正态积分，结果是：
\[
p(\sigma^2 | y) \propto (\sigma^2)^{-(n/2+1)} \exp\left(-\frac{(n-1)s^2}{2\sigma^2}\right) \times \sqrt{2\pi\sigma^2/n}
\]

简化后：
\[
p(\sigma^2 | y) \propto (\sigma^2)^{-[(n-1)/2+1]} \exp\left(-\frac{(n-1)s^2}{2\sigma^2}\right)
\]

这是尺度逆卡方分布（scale inverse-chi-squared distribution）：

\begin{Definition}[尺度逆卡方分布]
  \[
  p(x|\nu, \sigma^2) = \frac{(\nu\sigma^2/2)^{\nu/2}}{\Gamma(\nu/2)} x^{-(\nu/2+1)} \exp\left(-\frac{\nu\sigma^2}{2x}\right)
  \]
  其中 $\nu$ 是自由度，$\sigma^2$ 是尺度参数。
\end{Definition}

因此：
\[
\sigma^2 | y \sim \text{Inv-}\chi^2(n-1, s^2)
\]

由 $\frac{(n-1)s^2}{\sigma^2} \sim \chi_{n-1}^2$，这与频率统计中的经典结果有显著相似性。

\subsection{从联合后验分布抽样}

从联合后验分布抽取样本很容易：
\begin{enumerate}
  \item 从边际后验分布抽取 $\sigma^2$，即 $\sigma^2 \sim \text{Inv-}\chi^2(n-1, s^2)$
  \item 从条件后验分布抽取 $\mu$，即 $\mu | \sigma^2, y \sim N(\bar{y}, \sigma^2/n)$
\end{enumerate}

\subsection{边际后验分布：$t$ 分布}

总体均值 $\mu$ 通常是感兴趣的估计量，边际后验分布可以通过对 $\sigma^2$ 积分得到。结果是：
\[
\mu | y \sim t_{n-1}(\bar{y}, s^2/n)
\]
即自由度为 $n-1$ 的学生 $t$ 分布。

这与频率统计中的经典结果也有显著相似性：在采样分布下，$\frac{\bar{y}-\mu}{s/\sqrt{n}} \sim t_{n-1}$。

一个值得注意的事实是：这个边际后验分布不依赖于数据——它是数据的函数，但关于 $\mu$ 的概率陈述只通过 $\bar{y}$ 和 $s^2$ 依赖于数据。

\subsection{后验预测分布}

未来观测 $\tilde{y}$ 的后验预测分布可以写成混合：
\[
p(\tilde{y} | y) = \iint p(\tilde{y}|\mu, \sigma^2) \times p(\mu, \sigma^2 | y) \, d\mu \, d\sigma^2
\]

结果是 $t$ 分布：$\tilde{y} \sim t_{n-1}(\bar{y}, (1+1/n)s^2)$。

\section{正态数据与共轭先验}

当使用共轭先验时，分析更加系统化。正态模型的共轭先验族具有形式：
\[
\mu | \sigma^2 \sim N(\mu_0, \sigma^2/\kappa_0), \quad \sigma^2 \sim \text{Inv-}\chi^2(\nu_0, \sigma_0^2)
\]

这对应于联合先验密度 $N\text{-Inv-}\chi^2(\mu_0, \sigma_0^2/\kappa_0; \nu_0, \sigma_0^2)$。

四个超参数分别是 $\mu$ 的位置和尺度、$\sigma^2$ 的自由度和尺度。

注意 $\sigma^2$ 出现在 $\mu|\sigma^2$ 的条件分布中，意味着在联合共轭先验密度中 $\mu$ 和 $\sigma^2$ 必然相关。

这通常是有意义的——让 $\mu$ 的先验方差与 $\sigma^2$（观测 $y$ 的采样方差）挂钩是合理的。

\subsection{联合后验分布}

乘以正态似然得到后验密度仍然是 $N\text{-Inv-}\chi^2$ 形式，参数更新为：
\[
\mu_n = \frac{\kappa_0}{\kappa_0+n}\mu_0 + \frac{n}{\kappa_0+n}\bar{y}, \quad
\kappa_n = \kappa_0 + n, \quad
\nu_n = \nu_0 + n
\]
\[
\nu_n\sigma_n^2 = \nu_0\sigma_0^2 + (n-1)s^2 + \frac{\kappa_0 n}{\kappa_n}(\bar{y}-\mu_0)^2
\]

条件后验分布：$\mu | \sigma^2, y \sim N(\mu_n, \sigma^2/\kappa_n)$

边际后验分布：$\sigma^2 | y \sim \text{Inv-}\chi^2(\nu_n, \sigma_n^2)$

$\mu$ 的边际后验分布：$\mu | y \sim t_{\nu_n}(\mu_n, \sigma_n^2/\kappa_n)$

\section{分类数据的多项模型}

多项模型用于描述每个观测是 $k$ 个可能结果之一的数据。如果 $y$ 是每个结果观测次数的计数向量，则：
\[
p(y|\theta) \propto \prod_{j=1}^k \theta_j^{y_j}, \quad \sum_{j=1}^k \theta_j = 1
\]

共轭先验分布是 Dirichlet 分布的多元推广：
\[
p(\theta|\alpha) \propto \prod_{j=1}^k \theta_j^{\alpha_j-1}
\]

结果是 $\theta_j' s$ 的后验分布是 Dirichlet 分布，参数为 $\alpha_j + y_j$。

设置 $\alpha_j = 1$ 给出均匀密度；设置 $\alpha_j = 0$ 给出在 $\log(\theta_j)$ 上均匀的不正常先验。

\section{已知方差的多元正态模型}

多元正态模型 $N(\mu, \Sigma)$ 的似然函数为：
\[
p(y|\mu, \Sigma) \propto |\Sigma|^{-1/2} \exp\left\{-\frac{1}{2}(y-\mu)^T\Sigma^{-1}(y-\mu)\right\}
\]

给定 $\Sigma$ 时，$\mu$ 的共轭先验是多元正态分布 $\mu \sim N(\mu_0, \Lambda_0)$。

后验分布：
\[
\mu_n = (\Lambda_0^{-1} + n\Sigma^{-1})^{-1}(\Lambda_0^{-1}\mu_0 + n\Sigma^{-1}\bar{y}), \quad
\Lambda_n^{-1} = \Lambda_0^{-1} + n\Sigma^{-1}
\]

后验均值是数据和先验均值的加权平均，权重由精度矩阵给出；后验精度是先验精度和数据精度之和。

\section{均值和方差都未知的多元正态模型}

$(\mu, \Sigma)$ 的共轭先验分布是 normal-inverse-Wishart 分布：
\[
\Sigma \sim \text{Inv-Wishart}_{\nu_0}(\Lambda_0^{-1}), \quad \mu | \Sigma \sim N(\mu_0, \Sigma/\kappa_0)
\]

后验分布具有相同形式，参数更新为：
\[
\mu_n = \frac{\kappa_0}{\kappa_0+n}\mu_0 + \frac{n}{\kappa_0+n}\bar{y}, \quad
\kappa_n = \kappa_0 + n, \quad
\nu_n = \nu_0 + n
\]
\[
\Lambda_n = \Lambda_0 + S + \frac{\kappa_0 n}{\kappa_0+n}(\bar{y}-\mu_0)(\bar{y}-\mu_0)^T
\]
其中 $S = \sum_{i=1}^n (y_i - \bar{y})(y_i - \bar{y})^T$。

$\mu$ 的边际后验分布是多元 $t$ 分布。

\section{基本建模与计算总结}

多参数模型不允许轻易计算后验分布这一问题并非主要实际障碍，原因有三：

\begin{enumerate}
  \item 当参数较少时，非共轭多参数模型的后验推断可以通过简单模拟方法获得。
  \item 复杂模型通常可以表示为层次化或条件化形式，为此有有效的计算策略。
  \item 我们经常可以应用后验分布的正态近似，因此正态模型的共轭结构在实践中可以发挥重要作用，远远超出其对显式正态采样模型的应用。
\end{enumerate}

贝叶斯建模与计算的基本步骤总结：

\begin{enumerate}
  \item 写出模型的似然部分 $p(y|\theta)$，忽略任何不含 $\theta$ 的因子。
  \item 写出后验密度 $p(\theta|y) \propto p(\theta) \times p(y|\theta)$
    \begin{itemize}
      \item 如果先验信息表述良好，将其包含在 $p(\theta)$ 中。
      \item 否则使用弱信息先验或暂时设 $p(\theta) \propto \text{常数}$。
    \end{itemize}
  \item 创建参数的粗略估计 $\hat{\theta}$，用作起点和下一步计算的比较。
  \item 从后验分布抽取模拟 $\theta^1, \ldots, \theta^S$。使用样本模拟计算 $\theta$ 任何感兴趣函数的后验密度。
  \item 如果有预测量 $\tilde{y}$，通过对每个 $\tilde{y}^s$ 从给定 $\theta^s$ 的采样分布 $p(\tilde{y}|\theta^s)$ 抽取来模拟 $\tilde{y}^1, \ldots, \tilde{y}^S$。
\end{enumerate}

对于非共轭模型，步骤 4 可能很困难。各种方法已被开发用于从复杂模型中抽取后验模拟。

\section{本章小结}

本章我们学习了：
\begin{enumerate}
  \item 如何对 nuisance 参数进行积分（边际化和条件化）
  \item 多参数模型的计算策略：边际和条件模拟
  \item 正态模型中均值和方差同时未知时的分析
  \item 尺度逆卡方分布与 $t$ 分布
  \item 非信息先验和共轭先验在多参数情况下的应用
  \item 多项模型与 Dirichlet 先验
  \item 多元正态模型与 Inverse-Wishart 先验
\end{enumerate}

\section{下节预告}

下一章我们将学习渐近性以及贝叶斯方法与其他统计方法的联系。

\endinput

% Chapter 4: Asymptotics and connections to non-Bayesian approaches
% Bayesian Data Analysis - Gelman et al.

\chapter{渐近性与非贝叶斯方法的联系}

\section{本章导论}

我们在前三章看到，许多基于非信息先验分布的简单贝叶斯分析给出的结果与标准非贝叶斯方法（例如，均值未知时正态均值的后验 $t$ 区间）相似。

非信息先验分布能在多大程度上被合理地视为客观假设，取决于数据中包含的信息量：在第二、三章讨论的简单情况中，随着样本量 $n$ 的增加，先验分布对后验推断的影响会减小。

这些想法被称为\textbf{渐近理论}——因为它们指的是在 $n$ 趋向无穷大时成立的理论——我们将在本章中进行回顾，同时还将更详细地讨论贝叶斯方法与非贝叶斯方法之间的联系。

大样本结果实际上对于执行贝叶斯数据分析并非必需，但通常作为近似工具和理解工具很有用。

\section{后验分布的正态近似}\label{sec:ch4-normal-approx}

\subsection{联合后验分布的正态近似}

如果后验分布 $p(\theta | y)$ 是单峰的且大致对称，则可以用正态分布近似它；也就是说，后验密度的对数被近似为 $\theta$ 的二次函数。

我们考虑在后验众数 $\hat{\theta}$ 处展开的后验密度对数的二次近似（其中 $\theta$ 可以是向量，且 $\hat{\theta}$ 假定在参数空间内部）；在第十三章中我们将讨论更复杂的近似方法，这些方法在简单基于众数的近似失效的情况下可能有效。

对在 $\hat{\theta}$ 处中心化的后验密度对数 $\log p(\theta | y)$ 进行泰勒级数展开给出

\[
\log p(\theta | y) = \log p(\hat{\theta} | y) + \frac{1}{2}(\theta - \hat{\theta})^T \left[ \frac{d^2}{d\theta^2} \log p(\theta | y) \right]_{\theta = \hat{\theta}} (\theta - \hat{\theta}) + \dots \label{eq:laplace-approx-taylor}
\]

其中展开式的线性项为零，因为对数后验密度在其众数处导数为零。

当 $\theta$ 接近 $\hat{\theta}$ 且 $n$ 较大时，高阶余项的相对重要性会减小。考虑 \cref{eq:laplace-approx-taylor} 作为 $\theta$ 的函数，第一项是常数，而第二项与正态密度对数成正比，得到近似

\[
p(\theta | y) \approx \mathrm{N}(\hat{\theta}, [I(\hat{\theta})]^{-1}), \label{eq:laplace-approx-normal}
\]

其中 $I(\theta)$ 是观测信息（observed information），

\[
I(\theta) = -\frac{d^2}{d\theta^2} \log p(\theta | y).
\]

如果众数 $\hat{\theta}$ 在参数空间的内部，则矩阵 $I(\hat{\theta})$ 是正定的。

\begin{Example}[均值和方差均未知的正态分布]
  我们用简单的理论例子说明近似正态分布。设 $y_1, \ldots, y_n$ 是来自 $\mathrm{N}(\mu, \sigma^2)$ 分布的独立观测，为简单起见，假设 $(\mu, \log \sigma)$ 上的均匀先验密度。

  我们对 $(\mu, \log \sigma)$ 的后验分布建立正态近似，这在限制 $\sigma$ 为正值方面有优点。为了构建近似，我们需要对数后验密度的二阶导数。

  在众数处，二阶导数矩阵为 $\begin{pmatrix} -n/\hat{\sigma}^2 & 0 \\ 0 & -2n \end{pmatrix}$。从 \cref{eq:laplace-approx-normal}，后验分布可以近似为

  \[
  p(\mu, \log \sigma | y) \approx \mathrm{N}\left(\left(\begin{array}{c} \mu \\ \log \sigma \end{array}\right) \Big| \left(\begin{array}{c} \bar{y} \\ \log \hat{\sigma} \end{array}\right), \left(\begin{array}{cc} \hat{\sigma}^2/n & 0 \\ 0 & 1/(2n) \end{array}\right)\right).
  \]
\end{Example}

\subsection{后验密度函数相对于其最大值的解释}

除了作为近似直接使用外，多元正态分布还为解释后验密度函数和等高线图提供了基准。在 $d$ 维正态分布中，密度函数的对数是一个常数加上 $-2$ 除以的 $\chi_d^2$ 分布。例如，$\chi_{10}^2$ 密度的第 95 百分位是 18.31，因此如果一个问题有 $d = 10$ 个参数，则大约 95\% 的后验概率质量与 $p(\theta | y)$ 不小于在众数处密度的 $\exp(-18.31/2) = 1.1 \times 10^{-4}$ 倍的 $\theta$ 值相关联。

\subsection{用点估计和标准误总结后验分布}

渐近理论表明，如果 $n$ 足够大，后验分布可以用正态分布近似。在许多应用领域中，标准的推断总结是计算点估计 $\hat{\theta}$（如最大似然估计，即均匀先验下的后验众数），加上或减去两个标准误，标准误从估计处的信息 $I(\hat{\theta})$ 估计。

在许多情况下，通过变换可以显著改善 $\theta$ 的后验分布向正态分布的收敛性。如果 $\phi$ 是 $\theta$ 的连续变换，那么 $p(\phi | y)$ 和 $p(\theta | y)$ 都趋向正态分布，但对于有限的 $n$，近似的紧密程度可能随所选变换而大幅变化。

\subsection{数据约简与充分统计量}

在正态近似下，后验分布由其众数 $\hat{\theta}$ 和后验密度的曲率 $I(\hat{\theta})$ 总结；即渐近地，这些是充分统计量。

使用充分统计量的这种方法允许相对容易地将分层建模技术应用于改进每个单独估计。例如，在第 5.5 节中，一组八个实验中的每一个都通过从早期回归分析估计的点估计和标准误来总结。

当后验分布接近正态时，使用汇总统计量的方法最为合理；否则这种方法可能会丢弃重要信息并导致错误的推断。

\subsection{低维正态近似}

对于有限的样本量 $n$，正态近似通常对 $\theta$ 的分量条件分布和边际分布比对完整联合分布更准确。

例如，如果联合分布是多元正态分布，则其所有边际分布都是正态的，但反之不成立。确定 $\theta$ 一个分量的边际分布等同于对 $\theta$ 的所有其他分量进行平均，而平均一族分布通常会使它们更接近正态分布——这与中心极限定理背后的逻辑相同。

\begin{Remark}
  对于低维 $\theta$，后验分布的正态近似通常是完全可以接受的，特别是在适当的变换之后。
\end{Remark}

\begin{Example}[生物assay实验（续）]
  我们用 3.7 节中生物assay实验的模型和数据来说明正态近似。实验中样本量相对较小，只有 20 只动物，我们发现正态近似接近精确后验分布，但存在重要差异。

  $(\alpha, \beta)$ 联合后验分布的正态近似：\cref{fig:bioassay-normal-approx} 显示了二元正态近似的等高线图和从这个近似分布中抽取的 1000 个样本的散点图。

  正态近似的均值 $(\alpha, \beta) = (0.8, 7.7)$ 与实际后验分布的均值 $(\alpha, \beta) = (1.4, 11.9)$ 不同。
\end{Example}

\section{大样本理论}\label{sec:large-sample-theory}

\subsection{符号与数学设置}

大样本贝叶斯推断的基本工具是后验分布的渐近正态性：随着来自相同底层过程的越来越多数据到达，参数向量的后验分布趋向多元正态分布，即使数据的真实分布不在所考虑的参数族中。

从数学上讲，这些结果最直接地适用于从公共分布 $f(y)$ 中独立抽取的观测 $y_1, \ldots, y_n$。在许多情况下，为数据假设一个"真实"底层分布 $f(y)$ 的概念很难解释，但为了发展渐近理论，这是必要的。

假设数据由参数族 $p(y|\theta)$ 建模，并具有先验分布 $p(\theta)$。如果真实数据分布包含在参数族中——即如果 $f(y) = p(y|\theta_0)$ 对于某个 $\theta_0$——那么除了渐近正态性外，一致性性质也成立：随着 $n \to \infty$，后验分布收敛到在真实参数值 $\theta_0$ 处的点质量。

\subsection{渐近正态性与一致性}

附录 B 中给出的基本数学结果表明，在一些正则性条件下（特别是似然函数是 $\theta$ 的连续函数且 $\theta_0$ 不在参数空间边界上），随着 $n \to \infty$，$\theta$ 的后验分布趋向正态分布，均值为 $\theta_0$，方差为 $(nJ(\theta_0))^{-1}$。

在最简单的层面，这个结果可以从在后验众数处对数后验密度的泰勒级数展开 \cref{eq:laplace-approx-taylor} 来理解。初步结果表明，后验众数对 $\theta_0$ 是一致的，因此随着 $n \to \infty$，后验分布 $p(\theta | y)$ 的质量集中在 $\theta_0$ 越来越小的邻域内，距离 $|\hat{\theta} - \theta_0|$ 趋向零。

\subsection{似然支配先验分布}

渐近结果形式化了先验分布的重要性随着样本量增加而减小的概念。这导致一个结论：在具有大样本量的问题中，我们不必特别努力地制定一个准确反映所有可用信息的先验分布。当样本量较小时，先验分布是模型规范的关键部分。

\section{定理的反例}

理解大样本结果局限性的好方法是考虑定理失效的情况。正态分布通常作为起点近似有帮助，但必须检查偏差，特别是对于不寻常的参数空间和分布极值情况。

\subsection{未识别模型与未识别参数}

如果给定数据 $y$，似然 $p(\theta | y)$ 在一系列 $\theta$ 值上相等，则模型是未识别的（underidentified）。在这样的模型下，没有单一的点 $\theta_0$ 可以使后验分布收敛到它。

\begin{Example}[相关系数的未识别性]
  考虑模型，其中 $(u, v)^T \sim \mathrm{N}((0, 0)^T, \begin{pmatrix} 1 & \rho \\ \rho & 1 \end{pmatrix})$，但只从每对 $(u, v)$ 中观测到其中一个。这里参数 $\rho$ 是未识别的。数据不提供关于 $\rho$ 的任何信息，因此 $\rho$ 的后验分布与其先验分布相同，无论数据集有多大。
\end{Example}

未识别或欠识别参数问题的唯一解决方案是认识到问题存在，如果有更精确估计这些参数的愿望，则需要收集能够实现这些估计的进一步信息（可以是来自未来数据收集的信息，也可以是能够为先验分布提供信息的外部信息）。

\subsection{参数数量随样本量增加}

在复杂问题中，参数数量可能很大，我们需要区分不同类型的渐近性。如果随着 $n$ 增加，模型也发生变化使得参数数量同时增加，那么 \cref{sec:ch4-normal-approx} 和 \cref{sec:large-sample-theory} 中概述的简单结果（假设固定模型类 $p(y_i|\theta)$）并不适用。

例如，有时为研究中的每个采样单位分配一个参数；例如 $y_i \sim \mathrm{N}(\theta_i, \sigma^2)$。一般来说，除非每个采样单位收集的数据量也随着单位数量的增加而增加，否则参数 $\theta_i$ 不能被一致地估计。

\subsection{混叠}

混叠是欠识别参数的一种特殊情况，其中相同的似然函数在一组离散点上重复。例如，考虑具有独立同分布数据 $y_1, \ldots, y_n$ 的正态混合模型。

如果交换 $(\mu_1, \mu_2)$ 和 $(\sigma_1^2, \sigma_2^2)$，并将 $\lambda$ 替换为 $(1-\lambda)$，数据的似然保持不变。

后验分布至少有两个众数，由两个分布的 $(50\%, 50\%)$ 混合组成——它们是彼此的镜像；无论数据集有多大，它都不会收敛到单个点。

通常，通过限制参数空间使没有重复来解决混叠问题。

\subsection{无界似然}

如果似然函数是无界的，则后验众数可能不存在于参数空间中，从而使一致性和正态近似的结果无效。

例如，对于正态混合模型，如果我们设 $\mu_1 = y_i$ 对于任意 $y_i$，并让 $\sigma_1^2 \to 0$，则似然趋向无穷大。随着 $n \to \infty$，似然的众数数量增加。

\subsection{不当后验分布}

如果通过将似然乘以表示不当先验分布的"形式"先验密度获得的无规范化后验密度积分到无穷大，则渐近结果不成立。

当不当先验分布与似然组合具有有限积分时，不当先验分布只是一种方便的近似。

\subsection{排除收敛点的先验分布}

如果对于离散参数空间 $p(\theta_0) = 0$，或者对于连续参数空间在 $\theta_0$ 的邻域内 $p(\theta) = 0$，则一致性结果不成立。

解决方案是在先验分布中给所有甚至略微合理的 $\theta$ 值赋予正概率密度。

\subsection{收敛到参数空间边界}

如果 $\theta_0$ 在参数空间的边界上，则泰勒级数展开必须在某些方向上截断，正态分布不一定合适，即使在极限情况下也是如此。

例如，考虑模型 $y_i \sim \mathrm{N}(\theta, 1)$，限制 $\theta \geq 0$。假设模型是准确的，$\theta = 0$ 是真实值。$\theta$ 的后验分布是正态的，以 $\bar{y}$ 为中心，截断为正值。

这些反例教给我们一个重要教训：渐近定理的成立依赖于若干正则性假设，在实际应用中我们应当时刻检查四点：（1）参数是否可识别；（2）参数的维度是否随样本量增长；（3）似然函数是否有界；（4）先验分布是否为正常分布（proper）。只有当这些条件都满足时，基于大样本近似的推断才是可靠的。

\section{贝叶斯推断的频率评估}

就像贝叶斯范式可以被视为证明简单"经典"技术的合理性一样，频率统计方法提供了一种有用的方法来评估贝叶斯推断的性质——它们的操作特征——当这些被视为嵌入一系列重复样本中时。

\subsection{大样本对应}

假设 $\theta$ 后验分布的正态近似 \cref{eq:laplace-approx-normal} 成立；那么我们可以转换为标准多元正态：

\[
[I(\hat{\theta})]^{1/2}(\theta - \hat{\theta}) | y \sim \mathrm{N}(0, I), \label{eq:standard-normal-approx}
\]

其中 $\hat{\theta}$ 是后验众数。如果真实数据分布包含在模型类中，则在重复采样中，渐近地，$100(1-\alpha)\%$ 的中心后验区间将在 $100(1-\alpha)\%$ 的情况下覆盖真值 $\theta_0$。

\subsection{点估计、一致性与有效性}

从贝叶斯框架来看，获得 $\theta$ 的"估计"在大样本中最有意义，此时后验众数 $\hat{\theta}$ 是 $\theta$ 后验分布的明显中心。

一个点估计如果随着样本变大而收敛到它所估计的参数真值，则被称为在采样理论意义上是\textbf{一致的}。

如果对于任何 $\theta_0$ 值，在重复样本中 $C(y)$ 至少以 $100(1-\alpha)\%$ 的概率包含 $\theta_0$，则 $C(y)$ 被称为 $\theta$ 的 $100(1-\alpha)\%$ 置信域。

\section{其他统计方法的贝叶斯解释}

我们考虑贝叶斯统计方法与其他方法比较的三个层次：

\begin{enumerate}
  \item 首先，正如我们已经指出的，在涉及来自固定概率模型的大样本的问题中，贝叶斯方法通常类似于其他统计方法。
  \item 其次，即使对于小样本，许多统计方法也可以被视为基于特定先验分布的贝叶斯推断的近似；作为一种理解统计程序的方式，确定隐含的底层先验分布通常很有用。
  \item 第三，经典统计的一些方法（特别是假设检验）可能给出与贝叶斯方法非常不同的结果。
\end{enumerate}

\subsection{最大似然与点估计}

从贝叶斯数据分析的角度来看，我们通常可以将经典点估计解释为基于某些隐含完整概率模型的精确或近似后验总结。在大样本量的极限下，事实上我们可以使用渐近理论为经典最大似然推断构建理论贝叶斯论证。

渐近地，最大似然估计 $\hat{\theta}$ 是一个充分统计量——后验众数、均值或中位数也是充分统计量。

对于大的 $n$，先验分布的渐近无关性可以用来为使用方便的非信息先验模型提供理由。

\subsection{无偏估计}

一些非贝叶斯统计方法非常强调无偏性作为估计的可取原则。形式上，如果 $\mathbb{E}(\hat{\theta}(y)|\theta) = \theta$ 对于任何 $\theta$ 值，则估计 $\hat{\theta}(y)$ 被称为无偏的。

从贝叶斯角度来看，无偏性原则在大样本极限中是合理的，但在其他情况下可能具有误导性。主要困难出现在有许多参数要估计的情况下。

\begin{Remark}
  无偏性原则的一个问题是，通常不可能以甚至近似无偏的方式一次估计几个参数。
\end{Remark}

\begin{Example}[使用回归进行预测]
  考虑在已知母亲身高 $y$ 的情况下估计成年女儿身高 $\theta$ 的问题。为简单起见，假设母亲和女儿的身高联合正态分布，均值相等为 160 厘米，方差相等，相关系数为 0.5。

  在已知 $y$ 的条件下，$\theta$ 的后验均值为

  \[
  \mathbb{E}(\theta | y) = 160 + 0.5(y - 160). \label{eq:regression-mean-posterior}
  \]

  然而，后验均值在重复采样中并不是 $\theta$ 的无偏估计。

  这个例子说明的重要原则是\textbf{回归均值}的原则：对于任何给定的母亲，她女儿身高的期望值介于她母亲身高和总体均值之间。
\end{Example}

\section{文献注释}\label{sec:ch4-bib}

本章涉及的渐近理论是贝叶斯统计的核心组成部分，其数学基础可以追溯到 Bernoulli 和 Laplace 的早期工作。Bernstein-von Mises 定理（也称为渐近正态性）是本章理论的核心，该定理在以下经典文献中有系统论述：

\begin{enumerate}
  \item Lehmann and Casella (1998)，\emph{Theory of Point Estimation}，Chapter 6：对渐近估计理论进行了全面阐述，包括 Fisher 信息、渐近效率和有效性等概念的严格化处理。
  \item van der Vaart (1998)，\emph{Asymptotic Statistics}，Chapter 7-8：从概率论角度系统介绍了贝叶斯渐近理论，特别是后验浓度的数学证明。
  \item Birnbaum (1962)：关于似然函数与先验分布的关系，讨论了似然原则的统计学基础。
  \item Berger (1985)，\emph{Statistical Decision Theory}，Chapter 4-5：对非信息先验分布的渐近性质进行了详细讨论，包括 Jeffreys 先验的渐近最优性。
\end{enumerate}

关于正态近似（拉普拉斯方法）的数值实现，可参考 Tierney and Kadane (1986)。频率评估与贝叶斯推断的联系在 Rubin (1984) 中有经典论述。

\section{习题}\label{sec:ch4-exercises}

\begin{Exercise}{\ref{exr:ch4-1}  后验分布的正态近似}
对二项数据 $y \sim \mathrm{Bin}(n, \theta)$，使用 Jeffreys 先验 $p(\theta) \propto \theta^{-1/2}(1-\theta)^{-1/2}$。
\begin{enumerate}
  \item 写出后验分布 $p(\theta | y)$ 的精确形式。
  \item 在后验众数处构造正态近似，并将近似结果与精确后验均值和方差进行比较。
  \item 对于 $n = 20, y = 8$，分别在 $\theta$ 尺度和对数尺度 $\phi = \log(\theta/(1-\theta))$ 上计算正态近似，比较哪种近似的精度更高。
\end{enumerate}
\end{Exercise}

\begin{Exercise}{\ref{exr:ch4-2}  大样本理论与 Bernstein-von Mises 定理}
设 $y_1, \ldots, y_n$ 为来自公共分布 $f(y|\theta_0)$ 的独立观测，似然函数为 $p(y|\theta) = \prod_{i=1}^n f(y_i|\theta)$。
\begin{enumerate}
  \item 在适当的正则性条件下，叙述 Bernstein-von Mises 定理的核心结论。
  \item 证明：在大样本极限下，后验分布的均值 $\hat{\theta}$ 近似满足 $\hat{\theta} \approx \mathrm{N}(\theta_0, (nJ(\theta_0))^{-1})$，其中 $J(\theta_0)$ 是 Fisher 信息。
  \item 讨论：在什么情况下大样本近似可能失效？举出至少两个反例。
\end{enumerate}
\end{Exercise}

\begin{Exercise}{\ref{exr:ch4-3}  未识别参数的识别}
考虑一个实验，其中 $(u, v)^T \sim \mathrm{N}((0,0)^T, \Sigma)$，$\Sigma = \begin{pmatrix} 1 & \rho \\ \rho & 1 \end{pmatrix}$，但对每一对 $(u,v)$ 我们只观测到 $z = \max(u,v)$（即取两者中的较大值）。
\begin{enumerate}
  \item 写出数据的似然函数 $p(z_1,\ldots,z_n|\rho)$。
  \item 说明参数 $\rho$ 是未识别的（给出严格的数学证明）。
  \item 如果使用先验 $p(\rho) \propto 1$（$-1<\rho<1$ 上的均匀分布），后验分布是什么？这个结果说明了什么？
\end{enumerate}
\end{Exercise}

\begin{Exercise}{\ref{exr:ch4-4}  似然支配先验原则}
在什么意义上可以说"随着样本量增加，先验分布变得渐近无关"？请从以下几个方面阐述：
\begin{enumerate}
  \item 数学表述：后验分布 $p(\theta|y)$ 在大样本下的渐近形式；
  \item 频率学派视角：MLE 的一致性；
  \item 贝叶斯视角：先验熵与数据信息的关系；
  \item 实际含义：在大样本问题中，我们是否可以使用任何"方便"的先验？
\end{enumerate}
\end{Exercise}

\begin{Exercise}{\ref{exr:ch4-5}  频率评估与贝叶斯推断}
设 $\theta$ 的后验分布可以用正态近似 $p(\theta|y) \approx \mathrm{N}(\hat{\theta}, [I(\hat{\theta})]^{-1})$ 很好地近似。
\begin{enumerate}
  \item 构造 $\theta$ 的 $95\%$ 后验区间，并说明这个区间在频率学派意义下的覆盖概率。
  \item 设 $C(y) = [\hat{\theta} - 1.96\sqrt{\text{var}(\theta|y)},\, \hat{\theta} + 1.96\sqrt{\text{var}(\theta|y)}]$，证明：在重复采样下，$C(y)$ 渐近覆盖真值 $\theta_0$ 的概率为 $95\%$。
  \item 讨论：小样本情况下，后验区间的频率覆盖性质可能会出现什么问题？
\end{enumerate}
\end{Exercise}

\section{本章小结}

本章我们学习了：

\begin{enumerate}
  \item 后验分布的正态近似（基于泰勒级数展开）
  \item 观测信息矩阵与 Fisher 信息的概念
  \item 大样本理论：渐近正态性与一致性
  \item 似然支配先验的原则
  \item 渐近定理的反例：未识别参数、参数数量随样本增加、混叠、无界似力、不当后验等
  \item 频率评估与贝叶斯推断的联系
  \item 其他统计方法的贝叶斯解释：无偏估计、置信区间、回归预测等
\end{enumerate}

\section{下节预告}

下一章我们将学习分层模型（Hierarchical Models），这是贝叶斯统计中非常重要的一类模型，用于处理具有分组结构的复杂数据。

\endinput

% Chapter 5: Hierarchical models
% Bayesian Data Analysis - Gelman et al.

\chapter{分层模型}

\section{本章导论}

在许多统计应用中，我们面对的不是一个单一的参数，而是多个彼此存在某种联系的参数。当我们为每个参数独立指定先验时，我们实际上假设它们之间没有任何关系——然而，在许多应用中，我们有理由相信这些参数来自某个公共分布。

\textbf{核心理念}：如果我们假设 $\theta_j$ 是从一个未知的总体分布中抽取的样本，那么就可以用一种自然的方式将数据中的信息传递给参数——每个实验的数据不仅告诉我们关于该实验的 $\theta_j$，还告诉我们关于总体分布本身的特征。

\textbf{关键问题}：当我们只有有限的数据时，是应该为每个参数使用独立的先验，还是应该利用参数之间的相似性？分层模型提供了一种优雅的解决方案。

这种层次化思维不仅帮助我们理解多参数问题，还在计算策略的发展中发挥重要作用。更重要的是，简单的非分层模型通常不适合分层数据——参数太少则无法准确拟合数据，参数太多则容易过拟合。

\section{可交换性与分层模型的建立}

\subsection{从历史数据到先验分布}

我们从一个具体例子开始，阐述分层模型的基本思想。

考虑估计参数 $\theta$ 的问题：我们有当前小规模实验的数据，以及从类似历史实验构建的先验分布。关键问题是：我们如何\textbf{合理地}利用历史数据来帮助估计当前实验的参数？

一个看似合理的方法是：用历史数据估计先验分布的参数（如 beta 分布的 $\alpha, \beta$），然后将这个估计的先验用于当前实验。然而，这样做存在逻辑问题：当我们用全部数据来同时估计先验参数和当前参数时，数据被使用了两次。

\textbf{核心理念}：与其将先验分布视为固定或估计的对象，不如将整个层次结构——从数据到参数、从参数到超参数——放入概率模型中，对所有参数进行联合贝叶斯推断。这种方法在统计上更加一致，也更充分地利用了数据中的信息。

\begin{Example}[估计一组大鼠的肿瘤风险]\label{ex:rat-tumor}
  在药物评估中，研究人员通常对啮齿动物进行实验。假设我们的目标是估计 $\theta$——对照组雌性 F344 实验室大鼠发生子宫内膜基质息肉（一种肿瘤）的概率。

  数据显示，14 只大鼠中有 4 只发生了息肉。给定 $\theta$，假设肿瘤数量的二项模型是很自然的。为方便起见，我们从共轭族中选择先验分布：$\theta \sim \mathrm{Beta}(\alpha, \beta)$。

  后验分布为 $\theta \mid y \sim \mathrm{Beta}(\alpha + y, \beta + n - y)$，其中 $y = 4$，$n = 14$。
\end{Example}

现在，假设我们有来自历史数据的信息：历史实验中的肿瘤概率 $\theta_j$ 近似服从某个 beta 分布，其均值和标准差已知。设历史数据包含 70 组实验的肿瘤发生率 $(y_j/n_j)$，其样本均值为 0.136，标准差为 0.103。

如果设总体分布的均值和标准差分别为这些值，可以解出对应的 $(\alpha, \beta)$ 估计值为 $(1.4, 8.6)$。使用这个估计的先验分布，得到当前实验的 $\mathrm{Beta}(5.4, 18.6)$ 后验分布，后验均值是 0.223。

这比原始比例 $4/14 = 0.286$ 低得多——因为历史数据的权重表明，当前实验的肿瘤发生率异常高。

然而，这种直接用数据估计先验分布的方法存在逻辑问题：如果我们用全部数据来估计先验分布，然后再用同样的数据对每个 $\theta_j$ 进行推断，数据被使用了两次，这会导致我们高估精确度。

\textbf{解决之道}：将整个参数集合和实验的联合分布放入概率模型中，然后对所有模型参数的联合分布执行完整的贝叶斯分析。这种方法的历史脉络可以追溯到 de Finetti 的可交换性定理，以及 Lindley 与 Smith (1972)、Morris (1983) 等人在层次模型方面的开创性工作。

\subsection{可交换性}

当我们除了数据 $y$ 之外，没有任何信息可以区分各个 $\theta_j$ 时，必须在先验分布中对参数施加对称性。这种对称性在概率上用\textbf{可交换性}（exchangeability）来表示：如果参数 $(\theta_1, \ldots, \theta_J)$ 的联合分布 $p(\theta_1, \ldots, \theta_J)$ 在指标 $(1, \ldots, J)$ 的任意置换下保持不变，则这些参数是\textbf{可交换的}。

可交换性的概念在 de Finetti (1930) 的工作中得到了深刻的形式化——他证明，一个参数序列 $(\theta_1, \ldots, \theta_J)$ 是可交换的，当且仅当它可以表示为某个潜在参数 $\phi$ 的条件独立混合。这意味着，如果我们完全不知道参数的顺序或分组关系，最合理的假设就是可交换性。

可交换性最简单的形式是将每个参数 $\theta_j$ 视为来自某个由超参数向量 $\phi$ 控制的先验（或总体）分布的独立同分布样本：

\[
p(\theta \mid \phi) = \prod_{j=1}^{J} p(\theta_j \mid \phi). \label{eq:exchangeable-theta}
\]

一般情况下 $\phi$ 是未知的，因此我们必须在整个 $\theta$ 的分布上对 $\phi$ 的不确定性进行平均：

\[
p(\theta) = \int \left[ \prod_{j=1}^{J} p(\theta_j \mid \phi) \right] p(\phi) \, d\phi. \label{eq:marginal-theta}
\]

这种独立同分布的混合形式，通常足以在实践中捕捉可交换性。

\begin{Remark}
  可交换性并不意味着参数是独立同分布的——它只意味着在参数值被交换时，联合分布不变。独立同分布是可交换性的一种特殊形式。
\end{Remark}

\subsection{分层模型设置}

一般地，考虑 $j = 1, \ldots, J$ 个实验，其中实验 $j$ 有数据（向量）$y_j$ 和参数（向量）$\theta_j$，似然函数为 $p(y_j \mid \theta_j)$。

\textbf{两层分层模型}的标准形式为：

\begin{enumerate}
  \item \textbf{数据层}：$y_{ij} \mid \theta_j \sim p(y_{ij} \mid \theta_j)$（独立）
  \item \textbf{参数层}：$\theta_j \mid \phi \sim p(\theta_j \mid \phi)$（独立，来自共同先验）
  \item \textbf{超参数层}：$\phi \sim p(\phi)$（先验）
\end{enumerate}

联合后验分布为：\footnote{推导见附录 \cref{sec:joint-posterior-hierarchical}}

\[
p(\theta, \phi \mid y) \propto p(\phi) \prod_{j=1}^{J} p(\theta_j \mid \phi) \prod_{i=1}^{n_j} p(y_{ij} \mid \theta_j). \label{eq:joint-posterior-hierarchical}
\]

\section{伪贝叶斯估计与经验贝叶斯}

\subsection{两阶段估计}

在完全贝叶斯分析之前，一个实用的替代方法是将估计过程分为两个阶段：

\textbf{第一阶段}：对每个 $\theta_j$，使用其自身数据 $y_j$ 计算估计量 $\hat{\theta}_j$。

\textbf{第二阶段}：假设 $\hat{\theta}_j$ 是从某个分布 $p(\theta_j \mid \phi)$ 中抽取的样本，用所有 $\hat{\theta}_j$ 来估计 $\phi$，然后用估计的 $\phi$ 重新计算每个 $\theta_j$ 的后验均值。

这种方法被称为\textbf{伪贝叶斯}（pseudo-Bayes）或\textbf{经验贝叶斯}（empirical Bayes）估计。

\subsection{经验贝叶斯的原理}

设 $\hat{\theta}_j$ 是给定 $\theta_j$ 和已知方差 $\sigma_j^2$ 下的充分统计量。在正态模型中，$\hat{\theta}_j \sim \mathrm{N}(\theta_j, \sigma_j^2)$。

假设 $\theta_j \sim \mathrm{N}(\mu, \tau^2)$，则 $\hat{\theta}_j \sim \mathrm{N}(\mu, \sigma_j^2 + \tau^2)$。

我们可以用加权最小二乘法估计 $\mu$ 和 $\tau^2$：

\[
\hat{\mu} = \frac{\sum_{j=1}^{J} \hat{\theta}_j / (\sigma_j^2 + \hat{\tau}^2)}{\sum_{j=1}^{J} 1 / (\sigma_j^2 + \hat{\tau}^2)}, \label{eq:eb-mu-estimate}
\]

其中 $\hat{\tau}^2$ 由矩估计或方差分析估计得到。

然后，每个 $\theta_j$ 的经验贝叶斯估计为：

\[
\tilde{\theta}_j = \frac{\hat{\theta}_j / \sigma_j^2 + \hat{\mu} / \hat{\tau}^2}{1 / \sigma_j^2 + 1 / \hat{\tau}^2}. \label{eq:eb-theta-estimate}
\]

\begin{Remark}
  经验贝叶斯方法的核心问题：用点估计 $\hat{\phi}$ 代替超参数的后验分布，忽略了对 $\phi$ 的不确定性。在小样本情况下，这种不确定性可能很大。但当 $J$（实验数）较大时，这种近似的偏差通常可以忽略不计。
\end{Remark}

\section{共轭分层模型的完全贝叶斯分析}

完全贝叶斯分析与伪贝叶斯的根本区别在于：超参数 $\phi$ 本身也被赋予先验分布 $p(\phi)$，并在后验推断中对其不确定性进行正确量化。

\subsection{超先验分布}

为了建立 $(\phi, \theta)$ 的联合概率分布，我们必须为 $\phi$ 指定先验分布 $p(\phi)$。如果对 $\phi$ 了解甚少，可以指定一个扩散先验分布，但使用非正常先验密度时必须小心，要检查得到的后验分布是否是正常的。

在大鼠肿瘤例子中，超参数是 $(\alpha, \beta)$，它们决定 $\theta$ 的 beta 分布。

\subsection{后验预测分布}

分层模型的特点是同时存在超参数 $\phi$ 和参数 $\theta$。有两类后验预测分布值得关注：

\begin{enumerate}
  \item 现有 $\theta_j$ 对应的未来观测 $\tilde{y}$ 的分布；
  \item 从相同超总体中抽取的未来 $\theta_j$（记为 $\tilde{\theta}$）对应的观测 $\tilde{y}$ 的分布。
\end{enumerate}

在大鼠肿瘤例子中，未来观测可以是来自现有实验的额外大鼠，或来自未来实验的结果。

\subsection{Beta-Binomial 层次模型}\label{sec:beta-binomial-hierarchical}

假设我们有 $J$ 个独立实验，每个实验有 $n_j$ 次试验和 $y_j$ 次成功。模型为：

\begin{enumerate}
  \item $y_j \mid \theta_j \sim \mathrm{Binomial}(n_j, \theta_j)$
  \item $\theta_j \mid \alpha, \beta \sim \mathrm{Beta}(\alpha, \beta)$
  \item $(\alpha, \beta)$ 的先验：需要指定 $p(\alpha, \beta)$
\end{enumerate}

完全贝叶斯分析的关键是理解层次模型如何允许数据"告知"超参数。给定超参数 $(\alpha, \beta)$，$\theta_j$ 的后验分布为：\footnote{详细推导见附录 \cref{sec:beta-binomial-conditional-posterior}}

\[
p(\theta_j \mid \alpha, \beta, y) \propto \theta_j^{\alpha + y_j - 1} (1 - \theta_j)^{\beta + n_j - y_j - 1}, \label{eq:beta-binomial-posterior}
\]

即 $\theta_j \mid \alpha, \beta, y \sim \mathrm{Beta}(\alpha + y_j, \beta + n_j - y_j)$。

边缘化超参数后，$\theta_j$ 的后验预测分布为：

\[
p(\theta_j \mid y) = \iint p(\theta_j \mid \alpha, \beta) \, p(\alpha, \beta \mid y) \, d\alpha \, d\beta.
\]
\addtocounter{footnote}{-1}\footnote{详细推导见附录 \cref{sec:beta-binomial-marginal}}

这个积分通常需要数值方法或马尔科夫链蒙特卡洛（MCMC）来计算。

\section{正态模型中可交换参数的估计}\label{sec:normal-hierarchical}

一个特别重要的情况是：当数据 $y_j$ 是来自正态分布的样本均值，并且参数 $\theta_j$ 在变换后是正态分布时。

\subsection{正态层次模型}

考虑 $J$ 个独立实验，实验 $j$ 从正态分布 $y_{ij} \mid \theta_j \sim \mathrm{N}(\theta_j, \sigma^2)$ 中抽取 $n_j$ 个独立数据点，其中 $\sigma^2$ 已知。

\textbf{正态层次模型}设定：

\begin{enumerate}
  \item $y_j \mid \theta_j, \sigma_j^2 \sim \mathrm{N}(\theta_j, \sigma_j^2)$，其中 $\sigma_j^2$ 已知
  \item $\theta_j \mid \mu, \tau^2 \sim \mathrm{N}(\mu, \tau^2)$
  \item $(\mu, \tau^2)$ 有适当的先验
\end{enumerate}

超参数 $\mu$ 是总体均值，$\tau^2$ 是组间方差。

为方便起见（更准确地说是部分共轭），我们假设参数 $\theta_j$ 来自以超参数 $(\mu, \tau)$ 为参数的正态分布：

\[
p(\theta_1, \ldots, \theta_J \mid \mu, \tau) = \prod_{j=1}^{J} \mathrm{N}(\theta_j \mid \mu, \tau^2). \label{eq:normal-hierarchical-prior}
\]

即给定 $(\mu, \tau)$ 时，$\theta_j$ 是条件独立的。

\subsection{后验分布}\label{sec:normal-posterior}

给定数据 $y_j$，$\theta_j$ 的后验分布可以分解为：

\[
p(\theta_j \mid y, \mu, \tau^2) \propto p(\theta_j \mid \mu, \tau^2) \prod_{j=1}^{J} p(y_j \mid \theta_j).
\]

结果是：\footnote{推导见附录 \cref{sec:normal-hierarchical-derivation}}

\[
\theta_j \mid y, \mu, \tau^2 \sim \mathrm{N}\left( \hat{\theta}_j, V_j \right), \label{eq:normal-hierarchical-posterior}
\]

其中

\[
\hat{\theta}_j = \frac{\bar{y}_j/\sigma_j^2 + \mu/\tau^2}{1/\sigma_j^2 + 1/\tau^2}, \quad \frac{1}{V_j} = \frac{1}{\sigma_j^2} + \frac{1}{\tau^2}. \label{eq:normal-hierarchical-posterior-mean}
\]

这个结果表明：后验均值是先验均值 $\mu$ 和数据均值 $\bar{y}_j$ 的\textbf{精度加权平均}。

\section{收缩的性质}\label{sec:shrinkage-benefits}

层次模型的一个核心特征是\textbf{收缩}（shrinkage）——每个 $\theta_j$ 的后验估计被拉向总体均值 $\mu$。

\subsection{收缩机制}

当 $\tau^2$ 较大（组间变异大）时，估计接近原始数据 $\bar{y}_j$；当 $\tau^2$ 较小（组间变异小）时，估计收缩到共同均值 $\mu$。

从 \eqref{eq:normal-hierarchical-posterior-mean} 可以看出，后验均值可以写成：\footnote{详细推导见附录 \cref{sec:shrinkage-form-derivation}}

\[
\hat{\theta}_j = \lambda_j \bar{y}_j + (1 - \lambda_j) \mu, \label{eq:shrinkage-form}
\]

其中 $\lambda_j = \frac{1/\sigma_j^2}{1/\sigma_j^2 + 1/\tau^2}$ 是收缩因子，取值在 0 和 1 之间。

\subsection{完全收缩 vs 无收缩}

\begin{enumerate}
  \item 当 $\tau^2 \to 0$ 时，$\lambda_j \to 0$，$\hat{\theta}_j \to \mu$（完全收缩到公共均值）
  \item 当 $\tau^2 \to \infty$ 时，$\lambda_j \to 1$，$\hat{\theta}_j \to \bar{y}_j$（无收缩，使用原始数据）
\end{enumerate}

\textbf{收缩的收益}：当实验间真实差异较小时，收缩可以显著提高估计的精度；当某些实验的数据很少或噪声很大时，收缩尤其有益。

在极端情况下，如果某个实验只有 1-2 个观测值，其样本均值几乎完全不可靠，收缩会将其拉向可靠的公共均值，从而提高整体估计质量。

\section{示例：SAT 培训效果研究}

这是一个经典的层次模型例子，来自 BDA 原书第 5.9 节，展示了分层建模在 meta 分析中的应用。

\textbf{为什么这是层次模型最自然的应用场景？}：在 meta 分析中，我们既希望获得每个学校的特异性估计（了解特定学校的培训效果），又希望获得总体效应的估计（了解培训项目整体是否有效）。固定效应模型只能给出前者，而随机效应模型可以同时给出两者——这正是层次模型的核心优势。

\subsection{问题背景}

1982 年，Campbell 与 Powers 报告了一项研究，8 所学校参加了 SAT 辅导培训项目。每所学校都报告了培训效果估计 $y_j$（以 SAT 分数衡量）及其标准误 $\sigma_j$。

数据如下：\footnote{数据来源：BDA3 Table 5.2, p.121}

\begin{center}
\begin{tabular}{ccc}
\hline
学校 $j$ & 估计 $y_j$ & 标准误 $\sigma_j$ \\
\hline
A & 28 & 14 \\
B & 8 & 10 \\
C & -3 & 16 \\
D & 7 & 11 \\
E & -1 & 9 \\
F & 1 & 8 \\
G & 18 & 12 \\
H & 12 & 15 \\
\hline
\end{tabular}
\end{center}

\subsection{层次模型}

我们设定模型：\footnote{模型设定见 BDA3, p.121}

\[
y_j \mid \theta_j, \sigma_j \sim \mathrm{N}(\theta_j, \sigma_j^2), \quad \theta_j \mid \mu, \tau \sim \mathrm{N}(\mu, \tau^2), \label{eq:sat-model}
\]

其中 $\sigma_j$ 已知，$(\mu, \tau)$ 有适当的先验。

这里 $\theta_j$ 代表学校 $j$ 的真实培训效果，$\mu$ 是所有学校的平均培训效果，$\tau$ 量化了学校间的变异。

\subsection{后验推断}\label{sec:sat-posterior}

后验分布 $p(\theta, \mu, \tau \mid y)$ 可以通过 MCMC 获得。关键理解是：\footnote{详细分析见 BDA3, pp.122-127}

\begin{Remark}
  层次模型提供了在\textbf{合并信息}和\textbf{保持学校特异性估计}之间的自然权衡。参数 $\tau$ 控制了这种权衡：当 $\tau$ 小时，所有 $\theta_j$ 被强烈收缩到公共均值 $\mu$；当 $\tau$ 大时，每个 $\theta_j$ 主要由其自身数据决定。
\end{Remark}

后验均值 $\hat{\theta}_j$ 可以用 \eqref{eq:normal-hierarchical-posterior-mean} 计算。由于 $\tau$ 未知，需要先从其后验分布中抽样，再计算每个 $\theta_j$ 的后验均值。

\section{分层建模在 meta 分析中的应用}

Meta 分析是将多项研究结果合并的统计方法。层次模型为 meta 分析提供了自然的框架。

\subsection{固定效应与随机效应}

在 meta 分析中，我们通常考虑两种模型：

\begin{enumerate}
  \item \textbf{固定效应模型}：假设所有研究估计相同的真实效应 $\theta$，观测差异仅由抽样误差解释。
  \item \textbf{随机效应模型}：假设研究间的效应是可交换的，来自某个分布 $p(\theta_j \mid \mu, \tau^2)$。
\end{enumerate}

随机效应模型是层次模型的一个特例，其中 $\theta_j$ 被视为从公共分布中抽取的样本。

在随机效应模型中，我们不仅估计总体效应 $\mu$，还估计研究间的变异 $\tau^2$。这使得我们能够量化研究间的异质性，并相应地调整推断。

固定效应模型假设所有研究的效应完全相同，这在实践中很少合理——不同研究在受试者特征、干预方式、测量方法等方面总会有差异，这些差异在随机效应模型中被量化为组间变异 $\tau^2$。

\section{分层方差参数的弱信息先验}

层次模型的一个关键挑战是估计方差参数 $\tau^2$。当数据稀疏时，方差参数的估计可能极不稳定。

\subsection{问题}

在后验分布中，$\tau^2$ 通常信息不足，导致后验尾部过重或甚至不当。此外，经典的方差分析估计 $\hat{\tau}^2$ 可能是负数，这在物理上是不可能的。

设置 $\hat{\tau}^2 = 0$ 当 $MSW$（组内均方，Mean Square Within） $> MSB$（组间均方，Mean Square Between）似乎过于强烈——这相当于声称所有组均值完全相同，而不是说样本量太小无法区分 $\tau^2$ 与零。

\subsection{弱信息先验}

一种解决方案是使用弱信息先验：

\begin{enumerate}
  \item \textbf{Half-$t$ 先验}：对 $\tau$ 使用 $t$ 分布在正半轴
  \item \textbf{$\log(\tau)$ 上的正态先验}：使用宽窗口的正态先验
  \item \textbf{缩放逆 $\chi^2$ 先验}：方差参数的自然选择
\end{enumerate}

实践中，通常需要确定 $\tau$ 的最佳猜测和上界，然后可以从缩放逆 $\chi^2$ 族构建合理的先验分布，将最佳猜测匹配到缩放逆 $\chi^2$ 密度的均值，上界匹配到某个上百分位（如 99th）。

\begin{Remark}
  选择先验时，应该考虑参数的尺度。例如，在 meta 分析中，如果研究间的效应范围大致已知，可以据此设定 $\tau$ 的先验。在 $\log(\tau)$ 上指定先验通常比在 $\tau$ 上更合理，因为 $\tau$ 必须非负。
\end{Remark}

\section{本章小结}

本章我们学习了：

\begin{enumerate}
  \item \textbf{层次模型的基本思想}：通过分层结构处理相关参数，允许数据告知超参数
  \item \textbf{可交换性}（Exchangeability）的概念：当除了数据之外没有任何信息可以区分参数时，必须假设对称性；这一概念可追溯到 de Finetti 的经典工作
  \item \textbf{伪贝叶斯与经验贝叶斯估计}：两阶段估计方法及其与完全贝叶斯分析的关系
  \item \textbf{两层层次模型的构建}：数据层、参数层、超参数层
  \item \textbf{后验分布的收缩性质}：当 $\tau^2$ 小时，$\theta_j$ 的估计被收缩到公共均值 $\mu$
  \item \textbf{Beta-Binomial 和正态层次模型}的完全贝叶斯分析
  \item \textbf{SAT 培训效果研究}：层次模型在 meta 分析中的经典应用
  \item \textbf{固定效应 vs. 随机效应模型}的区别
  \item \textbf{层次方差参数的弱信息先验}
\end{enumerate}

层次模型是贝叶斯统计中最重要的工具之一，它允许我们：

\begin{enumerate}
  \item 在保持参数特异性估计的同时合并信息
  \item 自然地处理分组或聚类数据结构
  \item 为 meta 分析提供有原则的框架
  \item 通过完全贝叶斯分析正确量化不确定性
  \item 通过收缩机制在小样本情况下提高估计稳定性
\end{enumerate}

\section{下节预告}

下一章我们将学习模型检查（Model Checking），这是贝叶斯数据分析中验证模型合理性的关键步骤。我们将讨论如何评估模型对数据的拟合程度，以及如何识别模型的可能问题。

\section{习题}\label{sec:chapter5-exercises}

\begin{Exercise}{\ref{exr:5-1} 可交换性的判断}\label{exr:5-1}
考虑以下情形，判断参数是否可交换，并解释你的推理：
\begin{enumerate}
  \item $J$ 个不同城市的气温测量，每个城市有多个观测值
  \item $J$ 个学生的考试分数，其中一些学生来自同一个学校
  \item $J$ 个医院的心脏手术成功率
\end{enumerate}
\end{Exercise}

\begin{Exercise}{\ref{exr:5-2} 收缩估计量的极限性质}\label{exr:5-2}
在正态层次模型中，证明：
\begin{enumerate}
  \item 当 $\tau^2 \to 0$ 时，后验均值 $\hat{\theta}_j$ 收敛到公共均值 $\mu$
  \item 当 $\tau^2 \to \infty$ 时，$\hat{\theta}_j$ 收敛到原始数据均值 $\bar{y}_j$
\end{enumerate}
\end{Exercise}

\begin{Exercise}{\ref{exr:5-3} Beta-Binomial 层次模型的后验分布}\label{exr:5-3}
假设 $y_j \mid \theta_j \sim \mathrm{Binomial}(n_j, \theta_j)$，$\theta_j \mid \alpha, \beta \sim \mathrm{Beta}(\alpha, \beta)$。推导：
\begin{enumerate}
  \item 给定超参数 $(\alpha, \beta)$ 时 $\theta_j$ 的后验分布
  \item 边缘化超参数后的后验预测分布表达式
\end{enumerate}
\end{Exercise}

\begin{Exercise}{\ref{exr:5-4} SAT 培训效果研究的后验均值}\label{exr:5-4}
使用 \eqref{eq:normal-hierarchical-posterior-mean} 计算 SAT 例子中每所学校后验均值的大致范围，并说明当 $\tau$ 增大和减小时，估计量如何变化。假设 $\mu \approx 10$，$\tau \approx 15$。
\end{Exercise}

\begin{Exercise}{\ref{exr:5-5} 固定效应 vs. 随机效应的 meta 分析}\label{exr:5-5}
解释在 meta 分析中，为什么随机效应模型通常比固定效应模型更合理。当研究间异质性很大时（$\tau^2$ 大），两种方法可能给出不同结论；当异质性很小（$\tau^2 \to 0$）时，两种方法趋于一致。
\end{Exercise}

\begin{Exercise}{\ref{exr:5-6} 经验贝叶斯与完全贝叶斯的比较}\label{exr:5-6}
比较经验贝叶斯估计与完全贝叶斯估计的主要区别。解释为什么完全贝叶斯分析在估计 $\tau^2$ 时可能给出更大的不确定性。
\end{Exercise}

\endinput


% Part II: Fundamentals of Bayesian Data Analysis
\part{贝叶斯数据分析基础}

% Chapter 6: Model checking
% Bayesian Data Analysis - Gelman et al.

\chapter{模型检查}

\section{本章导论}

在前几章中，我们学习了贝叶斯推断的核心工具——如何从先验和似然出发，计算后验分布，以及如何构建层次模型来合并来自不同来源的信息。然而，在实际应用中，一个关键问题始终萦绕着我们：\textbf{我们如何知道所选的模型是否合适？}

模型检查是贝叶斯数据分析中一个容易被忽视但至关重要的环节。在经典统计中，我们有各种拟合优度检验；在贝叶斯框架中，我们同样需要工具来评估模型与数据的契合程度。本章将介绍一种强有力的方法——\textbf{后验预测检查（posterior predictive checking）}，它允许我们在贝叶斯框架内系统地评估模型。

我们将看到，模型检查不仅是验证模型是否合理，更是理解模型局限性、发现改进方向的重要途径。一个好的模型检查可以揭示数据中与模型假设不符的模式，从而引导我们改进模型或收集更多数据。

\section{后验预测检查的思想}

\subsection{为什么需要模型检查？}

贝叶斯分析的最终目标不是找到一个``正确''的模型，而是找到一个\textbf{有用}的模型——一个能够捕捉数据关键特征、产生合理预测的模型。任何模型都是对现实的简化，因此模型检查的核心问题是：\textbf{模型与数据的契合程度是否足以满足我们的分析目的？}

考虑一个简单例子。假设我们拟合了一个正态模型 $y_i \sim \mathrm{N}(\mu, \sigma^2)$ 到一组数据，但我们如何知道正态性假设是否合理？或者，当我们使用层次模型时，我们假设组内参数是可交换的——但数据是否支持这个假设？

这些问题的答案不能仅从模型内部获得，而需要我们主动将模型预测与实际数据进行比较。

\subsection{后验预测分布}

后验预测分布是模型检查的基础。给定观测数据 $y$，未来或新数据 $\tilde{y}$ 的后验预测分布为：

\[
p(\tilde{y} | y) = \int p(\tilde{y} | \theta) p(\theta | y) \, d\theta
\]

这个分布整合了参数的不确定性，反映了给定当前数据后模型对未来的预测。

\textbf{关键思想}是：如果模型合理，那么\textbf{真实数据应该与后验预测分布一致}。换句话说，我们不应该期望观测数据是后验预测分布的异常值。

\subsection{后验预测检查的直觉}

后验预测检查的基本思路如下：

\begin{enumerate}
  \item 从后验分布 $p(\theta | y)$ 中抽取参数值 $\theta^s$
  \item 给定每个 $\theta^s$，从似然模型 $p(\tilde{y} | \theta^s)$ 抽取伪数据 $\tilde{y}^s$
  \item 比较真实的观测数据 $y$ 与后验预测分布 $p(\tilde{y} | y)$
\end{enumerate}

如果 $y$ 落在后验预测分布的中心区域，模型与数据一致；如果 $y$ 落在尾部，则模型可能存在问题。

\begin{Example}[害虫治理研究]
  这是一个来自农业研究的例子。研究者比较了八种害虫栖息地的治理效果。数据包括每栖息地昆虫数量的计数。

  假设我们拟合了一个 Poisson 模型：$y_i \sim \mathrm{Poisson}(\theta_i)$，其中 $\theta_i$ 是栖息地 $i$ 的期望昆虫数量。

  通过后验预测检查，我们可以：
  \begin{itemize}
    \item 对每个后验样本 $\theta^s$，生成伪数据 $\tilde{y}^s$
    \item 比较观测数据 $y$ 与伪数据分布
    \item 检查是否存在异常大的观测值（overdispersed data）
  \end{itemize}

  如果某些观测值在后验预测分布的尾部，这可能表明存在过度离散（overdispersion），提示我们应该使用负二项分布或其他模型。
\end{Example}

\section{验证统计量与模型检查量}

\subsection{验证统计量}

验证统计量（test quantity）或模型检查量（discrepancy measure）$T(y, \theta)$ 是将数据和参数映射到标量或向量的函数，用于衡量模型与数据的某个特定方面。

常见的验证统计量包括：

\begin{itemize}
  \item \textbf{样本均值}：$T(y, \theta) = \bar{y}$
  \item \textbf{样本方差}：$T(y, \theta) = s^2 = \frac{1}{n-1}\sum_{i=1}^n (y_i - \bar{y})^2$
  \item \textbf{最大值}：$T(y, \theta) = \max(y_i)$
  \item \textbf{相关系数}：$T(y, \theta) = \text{corr}(y_i, y_j)$
\end{itemize}

在模型检查中，我们比较 $T(y, \theta)$ 与其后验预测分布 $p(T(\tilde{y}, \theta) | y)$。

\subsection{后验预测 p 值}

后验预测 p 值是后验预测检查的正式化表述。对于验证统计量 $T$，定义：

\[
\text{p-value} = \Pr(T(\tilde{y}, \theta) \geq T(y, \theta) | y)
\]

更精确地说，在参数空间上计算：

\[
\text{p-value} = \int \mathbb{I}(T(\tilde{y}, \theta^s) \geq T(y, \theta^s)) p(\theta^s | y) \, d\theta^s
\]

其中 $\theta^s$ 是从后验分布抽取的样本。

\textbf{解释}：后验预测 p 值度量的是，在给定数据和模型下，伪数据产生比观测数据更极端验证统计量的概率。如果这个概率很小（例如小于 0.05），则表明模型与数据在 $T$ 这个方面不一致。

\begin{Remark}
  后验预测 p 值与经典 p 值有本质区别。经典 p 值假设数据是随机化的，而\textbf{后验预测 p 值是条件于观测数据的}，反映的是模型对观测数据的拟合程度。
\end{Remark}

\section{模型检查的实践方法}

\subsection{检查批量}

当模型包含多个组或多个观测时，我们经常需要检查各个批量（batch）。设 $y_j$ 是第 $j$ 个组的观测，$\theta_j$ 是对应参数。

\textbf{批量特异性检查}可以揭示模型在哪些方面存在问题：

\begin{enumerate}
  \item \textbf{参数批量检查}：对每个组 $j$，计算 $T(\theta_j, \theta_{-j})$，其中 $\theta_{-j}$ 表示除 $j$ 以外的所有参数
  \item \textbf{数据批量检查}：对每个组 $j$，计算验证统计量 $T(y_j, \theta)$，然后比较 $T(y_j, \theta)$ 与其后验预测分布
\end{enumerate}

\begin{Example}[多组数据的批量检查]
  考虑一个分层模型，$J$ 所学校的学生成绩。每所学校有自己的真实效应 $\theta_j$，数据 $y_j$ 是带有测量误差的估计。

  批量检查可以帮助我们识别：
  \begin{itemize}
    \item 哪些学校的估计与模型预测差异最大
    \item 哪些学校的方差与层次模型假设不符
    \item 是否存在异常值学校
  \end{itemize}

  通过这种方法，我们不仅检查整体模型拟合，还可以定位具体问题所在。
\end{Example}

\subsection{检查交互作用}

当模型包含变量之间的交互作用时，检查交互作用是否与数据一致非常重要。

设我们有二元变量 $z_i$ 和 $w_i$，以及它们的交互作用 $z_i w_i$。一个常见的检查是验证交互作用是否在预测 $y_i$ 中起到预期作用。

具体来说，我们可以：
\begin{enumerate}
  \item 在没有交互作用的模型下计算后验预测分布
  \item 检查包含交互作用是否显著改善预测
  \item 或者直接检查数据中交互作用模式是否与模型一致
\end{enumerate}

\section{过度离散与计数数据的模型检查}

\subsection{过度离散问题}

计数数据经常表现出比 Poisson 分布预测的更大的变异性，这称为\textbf{过度离散（overdispersion）}。

在 Poisson 模型中，方差等于均值：$\text{var}(y) = \mathbb{E}(y)$。但实际数据通常表现出更大的方差。

\subsection{检测过度离散}

后验预测检查可以有效检测过度离散：

\begin{enumerate}
  \item 计算观测数据的方差：$T(y, \theta) = \text{var}(y)$
  \item 从后验预测分布生成伪数据 $\tilde{y}^s$
  \item 计算每个 $\tilde{y}^s$ 的方差：$T(\tilde{y}^s, \theta^s)$
  \item 比较观测方差与后验预测方差分布
\end{enumerate}

如果观测方差在伪数据方差的尾部，则表明存在过度离散。

\begin{Example}[害虫数据的过度离散检查]
  回到害虫治理研究。Poisson 模型假设方差等于均值。但如果数据存在过度离散，观测方差将显著大于后验预测方差。

  通过后验预测检查，我们可以量化这种不一致程度。如果后验预测 p 值很小，我们可以考虑使用负二项分布或零膨胀模型。
\end{Example}

\section{外部数据验证}

后验预测检查的一种重要扩展是使用外部数据验证。核心思想是：如果模型是合理的，它应该能够预测我们尚未看到的新数据。

设我们有两个相关的数据集 $y_{\text{old}}$ 和 $y_{\text{new}}$。我们可以：

\begin{enumerate}
  \item 使用 $y_{\text{old}}$ 拟合模型，获得后验分布 $p(\theta | y_{\text{old}})$
  \item 使用后验预测分布预测 $y_{\text{new}}$：$p(y_{\text{new}} | y_{\text{old}}) = \int p(y_{\text{new}} | \theta) p(\theta | y_{\text{old}}) \, d\theta$
  \item 比较预测分布与实际观测的 $y_{\text{new}}$
\end{enumerate}

这种方法在时间序列预测、跨研究汇总等场景中特别有用。

\textbf{先验预测检查}是另一种相关方法，它在收集数据之前使用先验分布预测数据。这有助于：

\begin{itemize}
  \item 检查先验分布的合理性
  \item 理解模型对数据的预期
  \item 发现模型中的潜在问题
\end{itemize}

\section{模型检查与其他推断的关系}

\subsection{与模型选择的联系}

模型检查是模型选择的前置步骤。在比较多个模型时，我们不仅应该比较它们的预测性能，还应该检查每个模型与数据的兼容性。

一个在训练集上表现良好但在检查中暴露问题的模型，可能存在过拟合或遗漏关键特征。

\subsection{与敏感性分析的联系}

模型检查与敏感性分析有密切关系。当检查发现模型与数据不一致时，敏感性分析可以帮助我们理解这种不一致是否会影响我们的结论。

例如，如果后验预测检查表明存在过度离散，但我们关心的参数估计在不同的过度离散模型下保持稳定，那么模型的不完美可能不会影响我们的最终结论。

\section{模型检查的局限性}

虽然后验预测检查是一种强有力的工具，但它也有局限性：

\begin{enumerate}
  \item \textbf{检查量的选择}：我们只能检查我们想到的验证统计量。可能存在模型缺陷，但我们的检查量没有捕捉到。
  \item \textbf{多重比较问题}：如果我们检查很多验证统计量，一些小的 p 值可能是偶然的。
  \item \textbf{模型不确定性}：后验预测检查没有考虑模型本身的不确定性。
  \item \textbf{先验敏感性}：检查结果可能对先验选择敏感。
\end{enumerate}

\begin{Remark}
  这些局限性并不意味着我们应该放弃模型检查，而是提醒我们：模型检查应该被视为一种\textbf{探索性工具}，而非最终裁决。它帮助我们理解模型的 strengths 和 weaknesses，但最终的数据分析和决策需要综合考虑多方面因素。
\end{Remark}

\section{从频率派角度理解后验预测检查}

从频率派的视角，后验预测检查可以被理解为一种校准（calibration）方法。

如果模型是正确指定的，那么：
\begin{itemize}
  \item 对于每个验证统计量 $T$，后验预测 p 值在 $[0, 1]$ 上均匀分布
  \item 重复实验将产生均匀分布的 p 值
\end{itemize}

这种校准性质意味着，如果我们在许多不同数据集上进行后验预测检查，p 值应该均匀分布。如果 p 值系统性地偏向极端值（例如过多的小 p 值），则表明模型存在系统性问题。

从频率派视角看，后验预测检查可以被视为评估模型\textbf{校准性}的一种方法——模型是否正确地量化了不确定性。

\section{本章小结}

本章我们学习了模型检查的核心概念和方法：

\begin{enumerate}
  \item \textbf{模型检查的必要性}：任何模型都是对现实的简化，我们需要工具来评估模型与数据的契合程度
  \item \textbf{后验预测检查的思想}：通过比较观测数据与模型对数据的预测来评估模型
  \item \textbf{验证统计量与模型检查量}：用于衡量模型特定方面的函数
  \item \textbf{后验预测 p 值}：量化模型与数据不一致程度的正式指标
  \item \textbf{批量检查}：检查模型在不同数据子集上的表现
  \item \textbf{过度离散检测}：使用后验预测方法检测计数数据的过度变异
  \item \textbf{外部数据验证}：使用独立数据验证模型预测能力
  \item \textbf{模型检查的局限性}：检查量选择、多重比较、模型不确定性等
  \item \textbf{频率派视角}：理解后验预测检查的校准性质
\end{enumerate}

模型检查不是贝叶斯分析的可选附件，而是其核心组成部分。通过系统地检查模型，我们可以：

\begin{itemize}
  \item 识别模型的不足之处
  \item 指导模型改进方向
  \item 加深对数据和模型的共同理解
  \item 在不确定性量化中保持诚实
\end{itemize}

\section{下节预告}

下一章我们将学习\textbf{模型评估、比较与扩展}（Evaluating, Comparing, and Expanding Models）。在模型检查的基础上，我们将进一步探讨如何系统地比较不同模型的性能，以及如何在必要时扩展或修改模型。

\endinput

% Chapter 7: Evaluating, comparing, and expanding models
% Bayesian Data Analysis - Gelman et al.

\chapter{模型评估、比较与扩展}

\section{本章导论}

在前几章中，我们学习了如何建立贝叶斯模型（分层模型）、如何进行后验推断，以及如何检查模型的拟合性（第六章）。但一个根本性问题仍然悬而未决：当我们有多个候选模型时，如何判断哪个更好？或者，当我们需要向模型中添加新变量或改变结构时，如何评估这些修改是否真的改进了预测能力？

这引出了本章的核心问题——\textbf{我们如何量化一个模型的预测表现？}

这个问题比表面上看起来要复杂得多。直觉上，我们可能会说"选择对数据拟合最好的模型"——但这恰恰是过拟合的根源。统计学家 George Box 有一句名言："所有模型都是错误的，但有些是有用的。"评估一个模型的价值，不在于它是否能完美地描述训练数据，而在于它是否能\textbf{泛化}到新的、未见过的数据。

在本章中，我们将学习一套完整的工具来评估和比较贝叶斯模型的预测表现。这些方法的核心思想是：将数据分为训练集和验证集，用训练集拟合模型，用验证集评估预测表现。我们将看到，这一思想如何演进出交叉验证、WAIC（广泛适用信息准则）和贝叶斯模型平均等强大工具。

\section{预测准确性的度量}

\subsection{点预测与分散度}

在讨论预测准确性之前，我们需要明确"预测"究竟意味着什么。考虑一个简单的场景：我们有一组新数据 $\tilde{y}$，需要预测它的值。

\textbf{点预测}是最直接的预测方式——给出一个单一数值 $\tilde{y}_{\text{pred}}$ 来代表我们的预测。但问题在于：什么样的数值是最好的预测？这取决于我们如何定义"最好"。

如果使用均方误差（MSE）作为损失函数，最优的点预测是后验均值：\footnote{推导见附录 \cref{sec:posterior-mean-optimal}}
\[
\tilde{y}_{\text{pred}} = \mathbb{E}(\tilde{y} | y)
\]

如果使用绝对误差损失，最优预测是中位数。但如果考虑\textbf{整个预测分布}的准确性——即我们不仅预测一个点，而且想量化整个预测分布的质量——则需要更精细的工具。

\subsection{对数预测密度}

\textbf{对数预测密度（log predictive density）}是评估预测分布准确性的基本工具。对于新数据点 $\tilde{y}$，其对数预测密度定义为：
\[
\log p(\tilde{y} | y) = \log \int p(\tilde{y} | \theta) p(\theta | y) \, d\theta
\]

这个量度量了模型预测 $\tilde{y}$ 的能力——对数密度越高，说明模型对 $\tilde{y}$ 的预测越准确。

为什么使用对数而不是原始密度？因为对数变换将乘积变为求和，这在计算和理论分析中都更方便。

\textbf{期望对数预测密度（ELPD）}是对新数据上对数预测密度的期望：
\[
\text{ELPD} = \mathbb{E}_{\tilde{y}}[\log p(\tilde{y} | y)]
\]

ELPD 是衡量模型预测表现的\textbf{黄金标准}——它度量了模型对未来数据的预期预测准确性。

\begin{Example}[子女身高预测]
  假设我们要预测一对夫妇的子女身高。模型 $M_1$ 假设身高服从正态分布 $N(\mu, \sigma^2)$，其中 $\mu$ 和 $\sigma$ 有后验分布；模型 $M_2$ 假设身高服从 $t$ 分布。

  对数预测密度可以比较这两个模型：如果 $M_1$ 的 ELPD 高于 $M_2$，说明正态假设对身高分布的预测更准确。
\end{Example}

\subsection{偏差与 DIC}

在实践中，我们无法直接计算 ELPD（因为我们没有无限多的新数据）。一个实用的替代方案是使用\textbf{偏差（deviance）}：
\[
D(\theta) = -2 \log p(y | \theta)
\]

偏差越小，说明模型对观测数据的拟合越好。

一个相关概念是\textbf{偏差信息准则（DIC）}：
\[
\text{DIC} = \bar{D} + p_D
\]
其中 $\bar{D} = \mathbb{E}^p[D(\theta)]$ 是后验平均偏差，$p_D = \bar{D} - D(\bar{\theta})$ 是有效参数个数。

有效参数个数 $p_D$ 度量了模型的复杂度——它反映了后验分布中"信息"参数的数量。在过拟合时，$p_D$ 会接近数据的数量；在欠拟合时，$p_D$ 会很小。

DIC 的局限：它只在后验分布近似正态且数据量足够大时才可靠。

\section{交叉验证}

\subsection{留一交叉验证（LOO-CV）}

既然 ELPD 衡量的是模型对新数据的预测能力，一个自然的想法是：用部分数据来拟合模型，用另一部分数据来评估预测。

\textbf{留一交叉验证（Leave-One-Out Cross-Validation，LOO-CV）}是最极端的形式——每次留出一个数据点 $y_i$，用其余 $n-1$ 个数据点拟合模型，然后预测 $y_i$。

LOO-CV 的对数预测密度为：
\[
\text{lppd}_{\text{loo}} = \sum_{i=1}^n \log p(y_i | y_{-i})
\]
其中 $y_{-i}$ 表示除 $i$ 以外的所有数据。

这个量的计算代价很高——需要拟合 $n$ 次模型。但对于共轭模型或可解析边缘化的模型，存在计算技巧。\footnote{推导见附录 \cref{sec:loo-cv-computation}}

\subsection{K 折交叉验证}

当 $n$ 很大时，LOO-CV 的计算成本变得不可接受。\textbf{K 折交叉验证}是一种折中方案：

\begin{enumerate}
  \item 将数据随机分为 $K$ 个大小相等的组
  \item 对于 $k = 1, \ldots, K$：
    \begin{itemize}
      \item 用除第 $k$ 组以外的所有数据拟合模型
      \item 用第 $k$ 组数据计算预测密度
    \end{itemize}
  \item 汇总 $K$ 次的结果得到总体评估
\end{enumerate}

常见的 $K$ 值有 5 和 10。$K$ 越小，计算越快但估计方差越大；$K$ 越大，估计越稳定但计算成本增加。

当 $K = n$ 时，K 折 CV 退化为 LOO-CV。

\begin{Remark}
  交叉验证的一个关键假设是：数据的\textbf{可交换性}。如果数据有内部结构（如时间序列、分层数据），简单的交叉验证可能会给出误导性的结果。在这种情况下，需要使用\textbf{分层交叉验证}或考虑数据收集设计的影响。
\end{Remark}

\section{WAIC：广泛适用信息准则}

\subsection{WAIC 的定义}

DIC 虽然常用，但它有一些理论缺陷——特别是它依赖于后验分布的近似正态性。\textbf{WAIC（Widely Applicable Information Criterion）}是一种更一般化的准则，可以在任意贝叶斯模型中使用。\footnote{WAIC 的完整推导见附录 \cref{sec:waic-derivation}}

WAIC 基于两个量：

\textbf{对数预测密度的点估计（pointwise log predictive density）}：
\[
\text{lppd} = \sum_{i=1}^n \log \mathbb{E}[p(y_i | \theta)]
\]
这等价于用后验均值代替后验分布进行预测。

\textbf{有效参数个数（effective number of parameters）}：
\[
p_{\text{waic}} = \sum_{i=1}^n \text{var}_{\text{post}}[\log p(y_i | \theta)]
\]
这个量度量了每个数据点对后验分布的贡献程度。

WAIC 定义为：
\[
\text{WAIC} = -2 (\text{lppd} - p_{\text{waic}})
\]
或者从 ELPD 的角度：
\[
\text{ELPD}_{\text{waic}} = \text{lppd} - p_{\text{waic}}
\]

\subsection{WAIC 与 DIC 的比较}

当后验分布近似正态时，WAIC 近似于 DIC。但 WAIC 有几个优势：

\begin{enumerate}
  \item \textbf{计算更一般}：WAIC 不需要后验分布的解析近似
  \item \textbf{点估计}：WAIC 可以逐点计算，便于诊断
  \item \textbf{理论基础更扎实}：WAIC 与交叉验证有更紧密的联系
\end{enumerate}

\begin{Example}[WAIC vs DIC]
  在非线性回归模型中，后验分布通常是高度非正态的。DIC 的正态假设会失效，但 WAIC 仍然可靠。

  在混合模型中，DIC 的有效参数个数定义模糊，而 WAIC 的逐点计算可以正确处理这种复杂性。
\end{Example}

\section{后验预测检验}

我们在第六章已经接触过\textbf{后验预测检验（Posterior Predictive Checking，PPC）}的基本思想。本节我们将更系统地讨论如何在模型比较中使用 PPC。

\subsection{复制数据与预测数据}

给定后验分布 $p(\theta | y)$，我们可以：

\begin{itemize}
  \item \textbf{复制数据}：从 $p(\tilde{y} | \theta)$ 生成 $\tilde{y}^{\text{rep}}$，这代表如果模型正确，数据的"复制版本"应该是什么样的
  \item \textbf{预测数据}：从 $p(\tilde{y} | y)$ 生成 $\tilde{y}^{\text{new}}$，这代表模型对未来数据的预测
\end{itemize}

两者的区别在于：$\tilde{y}^{\text{rep}}$ 用于检验模型拟合，$\tilde{y}^{\text{new}}$ 用于预测。

\subsection{差异度量}

为了量化观测数据与模型预测之间的差异，我们需要定义一个\textbf{差异度量（discrepancy measure）}：
\[
T(y, \theta) = \text{某个反映模型假设的统计量}
\]

常用的差异度量包括：

\begin{itemize}
  \item \textbf{均值差异}：$\bar{y} - \mathbb{E}[\bar{y}^{\text{rep}} | \theta]$
  \item \textbf{方差差异}：$\text{var}(y) - \text{var}(\tilde{y}^{\text{rep}} | \theta)$
  \item \textbf{最大值}：$\max(y) - \max(\tilde{y}^{\text{rep}} | \theta)$
\end{itemize}

对于每个差异度量 $T$，我们可以计算：
\[
\text{p-value} = \mathbb{P}(T(\tilde{y}^{\text{rep}}, \theta) \geq T(y, \theta) | y)
\]

如果这个 p 值接近 0 或 1，说明模型预测与数据之间存在系统性差异，模型可能需要改进。

\begin{Remark}
  后验预测 p 值不是严格的概率，而是反映模型拟合程度的一个指标。当样本量很大时，即使偏离很小，p 值也会变得显著——这提醒我们统计显著性不等于实际重要性。
\end{Remark}

\section{分层数据与交叉验证}

在分层或聚类数据结构中，简单的交叉验证可能会产生误导性的结果。

\subsection{分组交叉验证}

考虑一个分层模型：数据来自 $J$ 个组，每组内有 $n_j$ 个观测。简单的留一交叉验证会：

\begin{enumerate}
  \item 随机留出一个观测点
  \item 用剩余数据拟合模型
  \item 预测被留出的观测
\end{enumerate}

但这样做的问题在于：来自同一组的观测是相关的。如果我们随机留出，一些组会被完全保留，另一些组会被部分保留，这导致训练集和测试集的分布不一致。

\textbf{分组交叉验证（Group-based CV）}的解决方案是：每次留出一个完整的组，而不是单个观测点。这样保持了组的完整性。

当 $K$ 折分组交叉验证中 $K$ 等于组数 $J$ 时，就变成了\textbf{留一组合叉验证（Leave-One-Group-Out CV）}。

\begin{Example}[学校成绩预测]
  假设我们要预测学生考试成绩，数据来自多所学校。使用简单交叉验证可能导致：训练集中包含某学校的大部分学生，但测试集中只有该校的少数学生——这违反了数据独立性的假设。

  使用分组交叉验证，每次留出一整所学校，这样才能正确评估模型的跨学校泛化能力。
\end{Example}

\section{贝叶斯模型平均}

当我们有多个候选模型时，\textbf{贝叶斯模型平均（Bayesian Model Averaging，BMA）}提供了一种 principled 的方法来组合它们的预测。

\subsection{BMA 的思想}

给定 $K$ 个候选模型 $M_1, \ldots, M_K$，每个模型有后验概率 $p(M_k | y)$。贝叶斯模型平均的预测分布为：
\[
p(\tilde{y} | y) = \sum_{k=1}^K p(\tilde{y} | y, M_k) p(M_k | y)
\]

即：每个模型的预测按其后验概率加权平均。

后验概率 $p(M_k | y)$ 可以通过边际似然计算：
\[
p(M_k | y) \propto p(y | M_k) p(M_k)
\]
其中 $p(y | M_k) = \int p(y | \theta_k, M_k) p(\theta_k | M_k) \, d\theta_k$ 是模型 $M_k$ 的边际似然。

\subsection{BMA 的优势}

贝叶斯模型平均有几个重要优势：

\begin{enumerate}
  \item \textbf{自动正则化}：复杂模型通常有更低的边际似然（因为它们对数据的拟合更差）
  \item \textbf{不确定性量化}：BMA 的预测分布反映了模型选择的不确定性
  \item \textbf{防止过拟合}：单一模型的过拟合风险被分散到多个模型上
\end{enumerate}

但 BMA 也有局限：当候选模型都很差时，加权平均可能仍然表现不佳。此外，边际似然的计算在高维模型中可能很困难。

\begin{Remark}
  在实践中，BMA 最适合的情况是：候选模型捕捉了数据的不同方面，而真实数据生成过程接近其中某个模型。如果所有模型都系统性错误，BMA 无法纠正这个根本问题。
\end{Remark}

\section{信息准则的统一视角}

从前几节的讨论中，我们可以看到不同准则之间的联系：

\textbf{偏差} $D(\theta) = -2\log p(y|\theta)$ 是基本量。

\textbf{DIC} 使用后验均值近似：
\[
\text{DIC} \approx -2 \mathbb{E}[\log p(y | \theta)] + 2 \text{var}[\log p(y | \theta)]
\]

\textbf{WAIC} 使用更精确的点估计：
\[
\text{WAIC} = -2 \left(\sum_i \log \mathbb{E}[p(y_i | \theta)] - \sum_i \text{var}[\log p(y_i | \theta)]\right)
\]

\textbf{交叉验证}本质上是对新数据的 ELPD 的估计，但计算更直接。

这些准则都是对同一核心量——模型对未来数据的预测能力——的不同估计。

\section{Stan 实现}

本节展示如何在 Stan 中实现本章讨论的模型比较方法。

\subsection{计算对数预测密度}

Stan 提供了 \texttt{lppd} 函数来计算点态对数预测密度：

\begin{verbatim}
model {
  // ... standard model spec ...
  for (i in 1:n) {
    log_lik[i] = log(p(y[i] | theta));
  }
}
generated quantities {
  // compute WAIC components
  vector[n] log_lik;
  for (i in 1:n) {
    log_lik[i] = log(p(y[i] | theta));
  }
}
\end{verbatim}

\subsection{分层交叉验证的 Stan 实现}

对于分层模型，Stan 可以并行计算各组的预测：

\begin{verbatim}
data {
  int<lower=1> J;           // number of groups
  int<lower=1> N;           // total observations
  array[N] int<lower=1> group;  // group indicator
  vector[N] y;              // observations
}
parameters {
  vector[J] theta;          // group means
  real<lower=0> sigma;      // common variance
  real mu;                  // grand mean
  real<lower=0> tau;        // between-group SD
}
model {
  theta ~ normal(mu, tau);
  y ~ normal(theta[group], sigma);
}
generated quantities {
  vector[N] log_lik;
  for (i in 1:N) {
    log_lik[i] = normal_lpdf(y[i] | theta[group[i]], sigma);
  }
}
\end{verbatim}

\section{本章小结}

本章我们学习了评估和比较贝叶斯模型预测表现的核心方法：

\begin{enumerate}
  \item \textbf{预测准确性的度量}：对数预测密度、ELPD、偏差和 DIC
  \item \textbf{交叉验证}：LOO-CV 和 K 折 CV，以及处理分层数据的分组交叉验证
  \item \textbf{WAIC}：一种比 DIC 更一般化的信息准则，与交叉验证有紧密联系
  \item \textbf{后验预测检验}：通过差异度量评估模型拟合
  \item \textbf{贝叶斯模型平均}：组合多个模型预测的 principled 方法
  \item \textbf{Stan 实现}：计算对数预测密度和实现分层交叉验证
\end{enumerate}

这些工具共同构成了一套完整的模型评估体系。关键思想是：\textbf{模型的價值在于其泛化能力，而非对训练数据的拟合程度}。交叉验证和 WAIC 提供了估计泛化能力的实用方法，而 BMA 则在模型不确定时提供了更稳健的预测。

\section{下节预告}

下一章我们将学习\textbf{建模与数据收集过程的关系}（第八章）。我们将探讨如何根据数据的收集方式调整建模策略，包括抽样设计对推断的影响、处理缺失数据的贝叶斯方法，以及如何将领域知识整合到模型中。

\section{习题}\label{sec:chapter7-exercises}

\begin{Exercise}{\ref{exr:7-1} 对数预测密度的计算}
  假设数据 $y_1, \ldots, y_n$ 来自正态分布 $N(\mu, \sigma^2)$，其中 $\sigma$ 已知。给定共轭先验 $\mu \sim N(\mu_0, \sigma_0^2)$，推导后验预测分布 $p(\tilde{y} | y)$ 的对数密度表达式。
\end{Exercise}

\begin{Exercise}{\ref{exr:7-2} DIC 的有效参数个数}
  证明在正态模型 $y_i \sim N(\theta, \sigma^2)$（$\sigma$ 已知）中，DIC 的有效参数个数 $p_D$ 等于 1。
\end{Exercise}

\begin{Exercise}{\ref{exr:7-3} LOO-CV 与 WAIC 的比较}
  考虑二项数据 $y_i \sim \text{Bin}(n_i, \theta_i)$，使用层次模型 $\theta_i \sim \text{Beta}(\alpha, \beta)$。讨论：
  \begin{enumerate}
    \item 为什么 LOO-CV 在这个模型中计算困难？
    \item WAIC 如何处理这个问题？
    \item 在什么条件下两者会给出相似的结果？
  \end{enumerate}
\end{Exercise}

\begin{Exercise}{\ref{exr:7-4} 分组交叉验证}
  在一个三水平的层次模型中（学生嵌套于班级，班级嵌套于学校），说明：
  \begin{enumerate}
    \item 为什么简单的留一交叉验证不合适？
    \item 如何设计分组交叉验证？
    \item 选择留出学校 vs 留出班级有什么影响？
  \end{enumerate}
\end{Exercise}

\begin{Exercise}{\ref{exr:7-5} 贝叶斯模型平均}
  给定两个候选模型 $M_1$ 和 $M_2$ 的边际似然比 $p(y|M_1)/p(y|M_2) = 10$，解释这意味着什么。如果两个模型的预测分布差异很大，BMA 的结果会如何？
\end{Exercise}

\begin{Exercise}{\ref{exr:7-6} WAIC 的性质}
  证明 WAIC 的有效参数个数 $p_{\text{waic}}$ 总是介于 0 和 $n$ 之间，其中 $n$ 是样本量。
\end{Exercise}

\endinput

% Chapter 8: Modeling accounting for data collection
% Bayesian Data Analysis - Gelman et al.

\chapter{考虑数据收集过程的建模}

\section{本章导论}

在前几章的贝叶斯分析中，我们默认数据已收集完毕，直接对观测值建模。然而在实践中，\textbf{数据收集方式}对推断的正确性至关重要。不同的抽样设计、实验设计、缺失数据机制都可能导致相同的观测值携带不同的信息量。

本章的核心问题是：如何在贝叶斯推断中正确地考虑数据收集过程？这一问题的答案涉及一个看似矛盾的现象——天真的学生可能会认为，既然贝叶斯推断只依赖于后验分布 $p(\theta|y)$，数据如何收集无关紧要。这种对似然原则的误用忽视了一个关键事实：\textbf{``观测数据''的完整定义必须包含这些数据如何产生的全部信息}。

\section{观测数据与缺失数据的普遍范式}

这三个骰子场景的例子揭示了一个关键洞察：\textbf{相同的观测数据可以来自根本不同的数据生成过程}。为了系统地推理这类情况，我们需要一个统一的框架——将数据收集过程形式化为"完全数据"与"观测数据"之间的关系。

\textbf{核心思想}：将数据收集过程视为从"完整数据"（或潜在数据）$y$ 中观测到部分数据，留下缺失数据 $y_{\mathrm{mis}}$ 的过程。记 $y = (y_1, \ldots, y_N)$ 为完整数据矩阵，$I = (I_1, \ldots, I_N)$ 为观测指示矩阵，其中 $I_{ij} = 1$ 表示 $y_{ij}$ 被观测，$I_{ij} = 0$ 表示缺失。

\begin{Example}[骰子实验]
  假设我们告诉你掷了十次骰子且全部为 6。你的态度会因以下三种说法而大不相同：
  \begin{enumerate}
    \item 这十次是我们做的全部掷骰；
    \item 我们掷了 60 次但只报告了其中出现 6 的十次；
    \item 我们预先决定诚实报告出现十次 6，但隐瞒了总共掷了多少次——为此等待了 500 次。
  \end{enumerate}
  在情形 (i) 中，$p(y|\theta) = (1/6)^{10}$；在情形 (ii) 中，报告规则本身改变了分布；在情形 (iii) 中，等待时间提供了额外信息。这说明：\textbf{观测数据的分布不同于产生它们的"完整数据"的分布}。
\end{Example}\footnote{详见 \cite[ Chapter 8]{gelman2013bayesian}。}

\section{数据收集模型与可忽略性}

\subsection{观测数据和缺失数据的记号}

设 $y = (y_1, \ldots, y_N)$ 为潜在数据矩阵（每个 $y_i$ 可能本身是向量），$I$ 为同样维数的指示矩阵。记 $y_{\mathrm{obs}}$ 和 $y_{\mathrm{mis}}$ 分别为观测和缺失的数据部分。

\textbf{稳定性假设}（stable unit treatment value assumption）：记录的测量过程不改变数据本身的值——即完整数据向量 $y$ 不受包含向量 $I$ 的影响。

\subsection{数据模型、包含模型与完全/观测数据似然}

将联合概率模型分解为两部分：
\begin{enumerate}
  \item 完整数据的模型 $p(y|\theta)$；
  \item 包含向量的模型 $p(I|y,\phi)$。
\end{enumerate}

完全数据似然为：
\begin{equation}
p(y, I|\theta, \phi) = p(y|\theta) \cdot p(I|y, \phi). \label{eq:ch8-complete-data-likelihood}
\end{equation}

实际可用的信息是 $(y_{\mathrm{obs}}, I)$，因此贝叶斯推断的适当似然（\textbf{观测数据似然}）为：
\begin{equation}
p(y_{\mathrm{obs}}, I|\theta, \phi) = \int p(y, I|\theta, \phi) \, dy_{\mathrm{mis}}. \label{eq:ch8-observed-data-likelihood}
\end{equation}

给定 $(x, y_{\mathrm{obs}}, I)$ 的联合后验分布为：
\begin{equation}
p(\theta, \phi|x, y_{\mathrm{obs}}, I) \propto p(\theta, \phi|x) \int p(y|x, \theta) p(I|x, y, \phi) \, dy_{\mathrm{mis}}. \label{eq:ch8-joint-posterior}
\end{equation}

\subsection{有限总体与超总体推断}

这一区分的历史渊源可追溯至 Jerzy Neyman 在 1934 年的开创性工作——他首次系统地提出了有限总体推断的框架，将抽样视为从超总体中抽取样本的过程。在现代贝叶斯文献中，Rubin (1976, 1984) 的潜在结果框架（Rubin causal model）为因果推断奠定了形式化基础，而 Rosenbaum 与 Rubin (1983) 进一步提出了倾向性得分方法，成为观测研究中的核心工具。

本章区分两类目标量：
\begin{itemize}
  \item \textbf{有限总体量}（finite-population estimands）：如总体均值 $\bar{y}$，在抽样调查中可预测但未观测；
  \item \textbf{超总体量}（superpopulation estimands）：如参数 $\theta$，是产生数据的分布特征。
\end{itemize}

有限总体量 $\bar{y}$ 的后验分布可以通过从后验预测分布 $p(y_{\mathrm{mis}}|x, y_{\mathrm{obs}}, I)$ 模拟缺失数据得到。

\subsection{可忽略性}

如果决定忽略数据收集过程，可以只对 $y_{\mathrm{obs}}$ 条件化计算后验：
\begin{equation}
p(\theta|x, y_{\mathrm{obs}}) \propto p(\theta|x) \int p(y|x, \theta) \, dy_{\mathrm{mis}}. \label{eq:ch8-ignorable-posterior}
\end{equation}

当缺失数据模式不提供额外信息——即 $p(\theta|x, y_{\mathrm{obs}}) = p(\theta|x, y_{\mathrm{obs}}, I)$——时，研究设计称为\textbf{可忽略的}（ignorable）。

\textbf{可忽略性的两个充分条件}：
\begin{enumerate}
  \item \textbf{随机缺失}（Missing at Random）：$p(I|x, y, \phi) = p(I|x, y_{\mathrm{obs}}, \phi)$——即给定 $\phi$ 和 $(x, y_{\mathrm{obs}})$，缺失性只依赖已观测部分；
  \item \textbf{独立参数}（Distinct Parameters）：$p(\phi|x, \theta) = p(\phi|x)$——即缺失数据过程的参数与数据生成过程的参数在先验中独立。
\end{enumerate}

\begin{Remark}
  这两个条件的组合确保了 $p(\theta|x, y_{\mathrm{obs}}) = p(\theta|x, y_{\mathrm{obs}}, I)$，从而使缺失数据机制可忽略。
\end{Remark}

\subsection{设计分类}

根据可忽略性和机制是否已知，可以将数据收集设计分为以下几类：

这个分类对实践有直接指导意义：\textbf{类型 1-4 只需建模 $p(y|\theta)$，而类型 5-6 必须显式建模缺失数据机制}。随着类型编号增大，分析的复杂度和模型敏感性也随之增加。

\begin{enumerate}
  \item \textbf{可忽略且已知（无协变量）}：如简单随机抽样、完全随机化实验。只需考虑 $p(y|\theta)$ 和 $p(\theta)$；
  \item \textbf{可忽略且已知（给定协变量）}：如分层抽样、区组随机实验。需要建模 $p(y|x,\theta)$ 使设计可忽略；
  \item \textbf{强可忽略且已知}：$p(I|x, y, \phi) = p(I|x)$，依赖性仅通过完全观测的协变量 $x$ 介导；
  \item \textbf{可忽略但未知}：如基于完全观测协变量的非随机指派实验。倾向性得分（propensity score）在此有用；
  \item \textbf{不可忽略且已知}：如删失数据。需要将缺失机制明确建模；
  \item \textbf{不可忽略且未知}：最困难的情况，如因果推断中的选择偏差。
\end{enumerate}

\subsection{倾向性得分}

在强可忽略设计中，单元级指派概率 $\Pr(I_i = 1|x_i) = \pi_i$ 称为\textbf{倾向性得分}（propensity score）。倾向性得分可以充分概括协变量 $X$，使指派机制在给定倾向性得分时强可忽略。

\textbf{重要警告}：倾向性得分本身永远不足以进行后验预测检验（模型检查的关键工具），因为不同设计可能产生相同的倾向性得分。

\section{抽样调查}

考虑一个包含 $N$ 个人的有限总体，$y_i$ 为第 $i$ 个人的每周食品支出。目标是估计总体平均支出 $\bar{y}$。由于全面调查成本过高，采用简单随机抽样抽取 $n$ 个样本。

\subsection{简单随机抽样}

设 $I_i$ 为第 $i$ 个人是否被抽中的指示变量。简单随机抽样定义为：
\begin{equation}
p(I|y, \phi) = p(I) =
\begin{cases}
\binom{N}{n}^{-1} & \text{若} \sum_{i=1}^N I_i = n, \\
0 & \text{其他情形。}
\end{cases}
\end{equation}

这是强可忽略且已知的设计——指派概率（倾向性得分）为 $\pi_i = n/N$ 对所有单元相同。

\textbf{有限总体均值的贝叶斯推断}：这里 $\bar{y}$ 表示有限总体均值（目标量），而 $\bar{y}_{\mathrm{obs}}$ 和 $\bar{y}_{\mathrm{mis}}$ 分别表示已观测和未观测子总体的均值。将 $\bar{y}$ 分解为：
\begin{equation}
\bar{y} = \frac{n}{N}\bar{y}_{\mathrm{obs}} + \frac{N-n}{N}\bar{y}_{\mathrm{mis}}. \label{eq:ch8-finite-pop-mean}
\end{equation}

其中 $\bar{y}_{\mathrm{obs}}$ 已知，$\bar{y}_{\mathrm{mis}}$ 可从后验预测分布模拟。

\textbf{大样本近似}：当 $n$ 和 $N-n$ 都较大时，
\begin{equation}
\bar{y} | y_{\mathrm{obs}} \approx \mathrm{N}\left(\bar{y}_{\mathrm{obs}}, \left(\frac{1}{n} - \frac{1}{N}\right)s_{\mathrm{obs}}^2\right). \label{eq:ch8-normal-approx-finite-pop}
\end{equation}

这就是有限总体抽样推断的正态理论的正式贝叶斯 justification。

\subsection{分层抽样}

在分层随机抽样中，将 $N$ 个单元分为 $J$ 层，从每层 $j$ 中用简单随机抽样抽取 $n_j$ 个样本。给定层指示变量 $x_j$，此设计是可忽略的。

\textbf{贝叶斯分析}：在每层内建模 $y_i$ 的分布为参数 $\theta_j$，然后对所有参数集 $\theta_1, \ldots, \theta_J$ 执行贝叶斯推断。通常对 $\theta_j$ 指定层次模型，类似于第 5 章中的方法。

\begin{Example}[选前民意调查的分层抽样分析]
  1988 年美国总统选举的 CBS 调查将受访者按地区（东北、中西部、南部、西部）和居住密度分为 16 层。数据列于表 8.2。

  对每层 $j$，设 $y_{\mathrm{obs}j} = (y_{\mathrm{obs}1j}, y_{\mathrm{obs}2j}, y_{\mathrm{obs}3j})$ 为支持布什、杜卡基斯和无意见的人数，建模为多项分布 $y_{\mathrm{obs}j} \sim \mathrm{Multinomial}(n_j; \theta_{1j}, \theta_{2j}, \theta_{3j})$。

  \textbf{非层次模型}：各层参数赋予独立的 Dirichlet 先验，简单但未利用层间相似性；

  \textbf{层次模型}：将参数转换为 logit 尺度，建模为跨层可交换的二元正态分布。层次模型产生的后验中位数为 0.11，比非层次模型的 0.097 更集中于真实值。
\end{Example}\footnote{选举调查例子的完整推导见附录 \cref{sec:election-survey-hierarchical-derivation}。}

\subsection{整群抽样}

在整群抽样中，首先抽取 $J$ 个群，然后从每个抽中群 $j$ 中抽取 $n_j$ 个单元。分析需要为聚类的每个水平包含参数——这是一个层次建模问题。

\begin{Example}[澳大利亚学童调查]
  为估计墨尔本地区步行上学儿童的比例，首先从该都市区抽取 72 所学校，然后从每所学校的两个随机班级进行调查。

  模型包含学生、班级和学校三层参数：
  \begin{align}
    \bar{y}_{.jk} &\sim \mathrm{N}(\theta_{jk}, \sigma_{jk}^2), \quad \sigma_{jk}^2 = \bar{y}_{.jk}(1-\bar{y}_{.jk})/N_{jk} \\
    \theta_{jk} &\sim \mathrm{N}(\theta_k, \tau_{\mathrm{class}}^2) \\
    \theta_k &\sim \mathrm{N}(\alpha + \beta M_k, \tau_{\mathrm{school}}^2)
  \end{align}
  其中 $M_k$ 为学校 $k$ 的班级数。将班级数纳入学校级模型是因为调查的抽样概率依赖于学校规模 $M_k$——设计只有在对包含 $M_k$ 的模型才是可忽略的。
\end{Example}\footnote{整群抽样的完整推导见附录 \cref{sec:cluster-sampling-derivation}。}

\section{设计实验}

在统计术语中，实验涉及将处理指派给单元，指派过程由实验者控制。设 $(y_i^A, y_i^B)$ 为第 $i$ 个单元接受处理 $A$ 或 $B$ 的潜在结果，$i=1,\ldots,n$。

\textbf{稳定性假设}：应用于任何单元的处理对其余单元的结果没有影响。

因果效应分为两类：
\begin{itemize}
  \item \textbf{超总体框架}：平均因果效应为 $\mathrm{E}(y_i^A - y_i^B|\theta)$；
  \item \textbf{有限总体框架}：有限总体因果效应为 $\bar{y}^A - \bar{y}^B$。
\end{itemize}

\subsection{完全随机化实验}

在完全随机化实验中，处理指派是可忽略的。可以分别对两组数据建模：
\begin{equation}
p(y_{\mathrm{obs}}|\theta) = \prod_{i:I_i=(1,0)} p(y_i^A|\theta) \prod_{i:I_i=(0,1)} p(y_i^B|\theta).
\end{equation}

\textbf{有限总体因果效应的大样本近似}：
\begin{equation}
(\bar{y}^A - \bar{y}^B)|y_{\mathrm{obs}} \approx \mathrm{N}\left(\bar{y}_{\mathrm{obs}}^A - \bar{y}_{\mathrm{obs}}^B, \frac{2}{n}(s_{\mathrm{obs}}^{2A} + s_{\mathrm{obs}}^{2B})\right). \label{eq:ch8-causal-effect-approx}
\end{equation}

\subsection{随机区组、拉丁方等}

更复杂的设计可以基于处理指派过程中使用的所有因素建模，在模型下设计是可忽略的。

\begin{Example}[拉丁方实验]
  一个农业实验有 25 个小区和五种处理（A、B、C、D、E），排列成拉丁方形式（表 8.4）。

  可忽略设计所需的最小信息是 25 个小区的物理位置编码 $x$（水平和垂直坐标）。任何可忽略模型必须具有 $p(y|x,\theta)$ 的形式。

  如果额外信息可用（如流经田地的溪流位置），也应纳入分析。
\end{Example}\footnote{拉丁方实验的完整分析见附录 \cref{sec:latin-square-derivation}。}

\section{敏感性与随机化的作用}

我们已经看到，可忽略设计通过允许直接建模观测数据来简化贝叶斯分析。但随机化在此图景中扮演什么角色？

\subsection{完全随机化}

在没有完全观测协变量 $x$ 的情况下，实现可忽略设计的唯一方法（排除所有单元获得同一处理的退化设计）是\textbf{随机化}。

然而，对于任何给定的推断目标，某些可忽略设计在其他设计中更优——在产生关于目标量的更精确推断方面。完全随机化比独立伯努利指派产生更高的后验精度估计（因为期望每组有相等的样本量）。

\subsection{协变量条件下的随机化}

当有完全观测协变量时，随机化设计与确定性可忽略设计竞争。随机化的潜在优势包括：
\begin{enumerate}
  \item 即使假装 $x$ 未知，分析仍然有效；
  \item 观察到 $x$ 后可以更新后验分布，产生更精确的条件推断；
  \item 增加后验预测检验的灵活性；
  \item 使实验者更难通过选择样本或处理指派来"作弊"。
\end{enumerate}

\textbf{敏感性分析的实践意义}：在非随机化研究中，即使处理指派是可忽略的，当两组在协变量分布上严重不平衡时，后验推断会对所选的 $p(y|x,\theta)$ 模型形式高度敏感（如图 8.2 所示）。实践中，研究者通常通过倾向性得分匹配或分层来缓解这一问题，但这些方法无法完全消除敏感性——它们只是使敏感性更易于诊断。

\section{观测研究}

在观测性研究中，处理是被观察而非被实验者指派的。各处理组可能在未测量的因素上存在很大差异。

\textbf{关键问题}：如何确保因果推断的有效性？Gelman 等人指出，良好的观测研究需要：(1) 充分控制各处理组间的背景差异；(2) 每组有足够的独立单元产生可靠结果；(3) 研究设计不依赖于分析结果；(4) 最小化流失和其他形式的无意缺失数据；(5) 分析考虑设计中使用的所有信息。

图 8.2 说明了数据收集的不平衡与对建模假设敏感性的联系。当两组在预处理协变量上相似时（如随机实验），处理效应估计对所拟合的 $y|x$ 模型形式相对不敏感。当不平衡严重时，相同数据可以用完全不同的非线性关系拟合，产生完全不同的处理效应估计。

\subsection{倾向性得分在观测研究中的作用}

倾向性得分是处理指派的概率（作为协变量的函数）。它可用于诊断——如果估计的倾向性得分在两种处理条件下几乎没有重叠，则暗示推断可能对模型形式高度敏感。

\subsection{主分层}

\textbf{主分层}（principal stratification）是一种调整"中间结果"的方法——即在最终结果 $y$ 的因果路径上的中间变量 $C$。

\begin{Example}[有依从性问题的随机实验]
  在印度尼西亚村庄进行的维生素 A 补充剂对婴儿死亡率影响的随机实验中，维生素 A 仅提供给被分配接受它的人。这里 $I$ 是分配指标，中间结果是依从性。

  定义两种主分层：\textbf{依从者}（如果被分配则服用）和 \textbf{非依从者}（如果被分配也不服用）。

  依从者平均因果效应（CACE，记为 $\mathrm{CACE}$，即依从者的平均因果效应）的估计：
  \begin{equation}
    \widehat{\mathrm{CACE}} = (\bar{y}_1 - \bar{y}_0) / \hat{p}_c, \label{eq:ch8-cace-estimate}
  \end{equation}
  其中 $p_c$ 为依从者比例，$\bar{y}_1$ 和 $\bar{y}_0$ 为处理组和对照组的样本均值，\textbf{hat 表示对应的样本估计量}。这被称为\textbf{工具变量估计}。
\end{Example}\footnote{工具变量方法和贝叶斯因果推断的完整讨论见附录 \cref{sec:cace-derivation}。}

\section{删失与截断}

在医学随访、工业质量控制和环境监测中，研究者经常面临一类特殊的数据缺失问题：数据并非"未记录"，而是以某种受约束的形式被报告。区分\textbf{删失}（censoring）和\textbf{截断}（truncation）的概念——以及理解它们的统计含义——对于正确建模至关重要。删失与截断的区别在实践中往往被忽视，但忽略它们会导致严重的推断偏差。

我们通过一个简单例子的多种变体来说明不同的缺失数据机制。

\textbf{设定}：对某物体进行 $N=100$ 次称重，$y_i \sim \mathrm{N}(\theta, 1)$，先验 $p(\theta) \propto 1$。每种变体下只观测到 $n=91$ 个值，记 $\bar{y}_{\mathrm{obs}}$ 为已观测测量的均值。

\subsection{完全随机缺失（MCAR）}

如果天平等概率 $0.1$ 随机未报告，则包含模型为 $I_i \sim \mathrm{Bernoulli}(0.9)$，且与 $y$ 独立。后验分布与预先固定样本量 $n=91$ 的情形相同：
\begin{equation}
p(\theta|y_{\mathrm{obs}}, I) = \mathrm{N}(\theta|\bar{y}_{\mathrm{obs}}, 1/91).
\end{equation}

\subsection{完全随机缺失但缺失概率未知}

当天平随机失败概率 $\pi$ 未知时，如果 $\theta$ 和 $\pi$ 在先验中独立（即"独立参数"条件），后验推断与上一情形相同。但如果 $\pi = \theta/(1+\theta)$（两者不独立），则 $n=91$ 提供了关于 $\theta$ 的额外信息。

\subsection{删失数据}

如果天平的报告上限为 200kg——所有超过 200kg 的值被报告为"太重"——则 $n=91$ 个数值测量的贡献是正态密度，9 个"太重"测量的贡献形式为 $\Phi(\theta - 200)$。后验分布为：
\begin{equation}
p(\theta|y_{\mathrm{obs}}, I) \propto \mathrm{N}(\theta|\bar{y}_{\mathrm{obs}}, 1/91) \cdot [\Phi(\theta-200)]^9. \label{eq:ch8-censored-posterior}
\end{equation}

\textbf{关键}：删失数据的似然不能简单地忽略指示向量，必须包含删失观测的贡献 $[\Phi(\theta-200)]^9$。

\subsection{未知删失点的删失数据}

当删失点 $\phi$ 本身未知时，后验分布必须同时考虑 $(\theta, \phi)$。9 个删失测量贡献 $[\Phi(\theta-\phi)]^9$，即使先验中 $\theta$ 和 $\phi$ 独立，后验中也会产生依赖——这说明缺失数据机制不可忽略。

\subsection{截断数据}

截断与删失的关键区别在于：\textbf{删失}知道超过截断点的观测数量，而\textbf{截断}不知道。当 $N$ 也未知时，后验分布涉及对 $N$ 的推断。

\section{习题}

\begin{Exercise}{设计类型识别}
给定以下场景，判断属于哪种数据收集设计类型（类型 1-6），并说明理由：
\begin{enumerate}
  \item 研究者对某社区进行简单随机抽样，调查居民收入；
  \item 某制药公司在多个医学中心进行随机对照试验，评估新药有效性；
  \item 一项调查因部分受访者拒绝回答而导致收入数据缺失，研究者假设缺失与收入无关；
  \item 某工厂的质量监控系统记录超过测量上限的值为"> 100"，而非实际数值；
  \item 研究者根据学生的期中考试成绩决定是否邀请其参加辅导项目（处理指派依赖于观测结果）。
\end{enumerate}
\end{Exercise}

\begin{Exercise}{删失数据的似然推导}
对 8.7 节中的删失数据情形（天平上限 200kg），写出完整的观测数据似然函数，并验证 \cref{eq:ch8-censored-posterior} 的形式。
\end{Exercise}

\begin{Exercise}{倾向性得分计算}
假设有 100 个单元的二元处理指派，其中 50 个单元的处理指派概率为 0.8，50 个为 0.2。给定处理指示向量 $I$ 和协变量矩阵 $X$，描述如何估计倾向性得分，并讨论在什么条件下倾向性得分分析可能失效。
\end{Exercise}

\begin{Exercise}{MCAR 与 MAR 的区分}
构造一个一维数据的例子，其中缺失机制满足 MCAR（完全随机缺失）但不满足 MAR（随机缺失）。再构造一个满足 MAR 但不满足 MCAR 的例子。说明两者之间的包含关系。
\end{Exercise}

\section{本章小结}

本章我们学习了：
\begin{enumerate}
  \item 数据收集过程必须纳入贝叶斯模型的核心理念
  \item 观测数据与缺失数据的普遍记号框架
  \item 可忽略性的两个充分条件：随机缺失（MAR）+ 独立参数
  \item 有限总体推断与超总体推断的区别
  \item 倾向性得分的概念与局限
  \item 抽样调查的贝叶斯分析：简单随机抽样、分层抽样、整群抽样
  \item 设计实验的贝叶斯分析：完全随机化、区组设计、拉丁方
  \item 随机化在降低模型敏感性中的作用
  \item 观测研究的挑战与敏感性分析
  \item 主分层与因果推断中的依从性
  \item 删失与截断数据的建模
\end{enumerate}

\section{下节预告}

下一章我们将学习决策分析（Chapter 9），探讨贝叶斯决策理论在不同情境中的应用，包括回归预测的激励措施、多阶段决策制定和个体与机构决策分析。

\endinput

% Chapter 9: Decision analysis
% Bayesian Data Analysis - Gelman et al.

\chapter{决策分析}

\section{本章导论}

在前面几章中，我们探讨了贝叶斯推断的核心技术：从先验分布和似然函数出发，通过贝叶斯公式计算后验分布，并利用后验分布进行点估计、区间估计和模型检查。然而，推断本身并不是终点——我们需要将推断的结论转化为实际行动决策。

一个自然的问题是：给定后验分布，我们如何做出最优决策？这个问题看似简单，实则涉及一个根本性的概念飞跃：\textbf{推断是描述性的，而决策是规范性的}。描述性任务只需要刻画未知参数的不确定性，而规范性任务需要明确我们追求的目标是什么，以及不同决策的相对优劣。

\textbf{决策分析}（decision analysis）正是连接推断与决策的桥梁。这一理论框架由 Wald (1950)、Savage (1954)、Raiffa 和 Schlaifer (1961) 等人在 20 世纪中叶发展起来。Wald (1950) 将统计推断形式化为一场"与自然的博弈"，通过引入风险函数和极小极大原则为决策理论奠定了公理基础；Savage (1954) 则从主观期望效用角度出发，证明了在合理 axioms 下理性决策者必然遵循贝叶斯原则；Raiffa 和 Schlaifer (1961) 进一步将这一理论系统化，使之成为处理实际决策问题的有力工具。贝叶斯版本的决策理论进一步要求，在计算期望效用时，使用后验分布来量化不确定性。

本章的安排如下：\S 9.1 引入效用函数与损失函数的基本概念；\S 9.2 讨论几种常见的损失函数及其对应的贝叶斯最优决策；\S 9.3 探讨后验预测分布与预验分析（preposterior analysis）；\S 9.4 处理层次决策问题与经验贝叶斯的联系；\S 9.5 从损失函数角度重新审视模型选择问题。

\section{效用函数与损失函数}

\subsection{从推断到决策的桥梁}

在贝叶斯框架下，每个决策 $d$ 带来一个后果（consequence），记为 $C(d, \theta)$，其中 $\theta$ 是自然状态（state of nature）——即我们关心的未知参数或潜在结果。后果可以是金钱、健康、时间等不同维度。

为了比较不同决策的优劣，我们需要一个\textbf{效用函数}（utility function）$u(c)$，它将每个后果映射到一个实数 $u(c) \in \mathbb{R}$。理性决策者偏好效用更高的后果。

在很多情况下，直接定义效用函数比较困难。我们通常转而使用\textbf{损失函数}（loss function）$\ell(\theta, d)$，定义为采取决策 $d$ 而真实参数为 $\theta$ 时的损失。损失函数与效用函数之间的关系是：$u(c) = - \ell(\theta, d)$（在基数效用意义下，效用的增减与损失的增减相反）。

\begin{Example}[医学治疗决策]\label{ex:medical-decision}
  考虑一个简化的医学决策问题：医生需要决定是否对患者进行某种手术。真实参数 $\theta$ 代表手术的实际成功率。设损失函数为
  \[
  \ell(\theta, d) = \begin{cases}
    0 & \text{若 } d = \text{做手术且成功} \\
    1 \cdot L_{\text{fail}} & \text{若 } d = \text{做手术且失败} \\
    0.5 \cdot L_{\text{conservative}} & \text{若 } d = \text{不做手术},
  \end{cases}
  \]
  其中 $L_{\text{fail}}$ 和 $L_{\text{conservative}}$ 是根据具体情况设定的损失权重。这里"不做手术"被视为一种保守决策，当手术失败后果严重时给予较大的损失权重。
\end{Example}

\subsection{期望后验损失与贝叶斯最优决策}

给定损失函数 $\ell(\theta, d)$ 和后验分布 $p(\theta \mid y)$，决策 $d$ 的\textbf{期望后验损失}（posterior expected loss）为
\[
\rho(d \mid y) = \mathbb{E}_{p(\theta\mid y)}[\ell(\theta, d)] = \int \ell(\theta, d) \, p(\theta \mid y) \, d\theta.
\]

贝叶斯最优决策 $d^{\text{Bayes}}$ 是使期望后验损失最小化的决策：
\[
d^{\text{Bayes}} = \underset{d \in \mathcal{D}}{\arg\min} \; \rho(d \mid y).
\]

\textbf{关键洞察}：后验分布综合了先验信息与数据信息，因此"最小化期望后验损失"这一原则自然地将推断与决策统一起来——推断的结论（后验分布）直接决定了最优行动。

值得注意的是，$\rho(d \mid y)$ 是随机变量 $\theta$ 的函数，但由于我们对 $\theta$ 求了期望，最终得到的是一个确定性的数值（给定 $y$ 和 $d$），因此可以直接比较不同 $d$ 的优劣。

在第三章中，我们通过积分消去了 nuisance 参数，得到了边际后验分布。现在我们看到，这一"积分"操作本身也可以被理解为一种最优决策——当损失函数取某种特定形式时，边缘化（marginalization）正是使期望损失最小化的最优行动。这一联系揭示了贝叶斯推断中"积分"操作背后的决策理论动机。

\section{常见损失函数与最优估计}

我们已经建立了期望后验损失的一般框架，现在来考察几种在统计推断中最常见的损失函数及其对应的贝叶斯最优决策。

损失函数的选择反映了决策者的风险态度和具体问题的结构。以下是几种在统计推断中最常见的损失函数。

\subsection{平方损失与后验均值}

设 $\theta$ 为标量参数，$d$ 为点估计量。平方损失函数（squared error loss）定义为
\[
\ell(\theta, d) = (\theta - d)^2.
\]

代入期望后验损失：
\[
\rho(d \mid y) = \mathbb{E}_{p(\theta\mid y)}[(\theta - d)^2]
           = \mathbb{E}(\theta^2 \mid y) - 2d \, \mathbb{E}(\theta \mid y) + d^2.
\]

对 $d$ 求导并令其为零，得到最优解
\[
d^{\text{Bayes}} = \mathbb{E}(\theta \mid y).
\]

\begin{Theorem}[平方损失下的贝叶斯最优估计]\label{th:squared-loss-bayes}
  在平方损失下，贝叶斯最优点估计是后验均值 $\mathbb{E}(\theta \mid y)$。
\end{Theorem}

这解释了为什么在很多贝叶斯分析中，后验均值被视为"自然"的点估计——它不是任意选择的，而是平方损失这一合理损失函数下的最优决策。

\subsection{绝对损失与后验中位数}

绝对损失函数（absolute error loss）定义为
\[
\ell(\theta, d) = \lvert\theta - d\rvert.
\]

期望后验损失 $\rho(d \mid y) = \mathbb{E}_{p(\theta\mid y)}[\lvert\theta - d\rvert]$ 是 $d$ 的凸函数。其最小化问题
\[
\underset{d}{\min} \; \mathbb{E}[\lvert\theta - d\rvert]
\]
的解是后验分布的任意中位数（当后验分布为连续分布时，唯一解为后验中位数）。

\begin{Theorem}[绝对损失下的贝叶斯最优估计]\label{th:absolute-loss-bayes}
  在绝对损失下，贝叶斯最优估计量是后验分布的中位数。
\end{Theorem}

绝对损失对大误差的惩罚力度小于平方损失，因此对应的最优估计对尾部不敏感，更加稳健。

\subsection{0-1 损失与后验众数}

在分类问题中，我们常使用 0-1 损失（zero-one loss）：对于离散参数 $\theta$ 取值于有限集合 $\mathcal{K}$，
\[
\ell(\theta, k) = \begin{cases}
  0 & \text{若 } \theta = k, \\
  1 & \text{若 } \theta \neq k.
\end{cases}
\]

期望后验损失为
\[
\rho(k \mid y) = \mathbb{E}_{p(\theta\mid y)}[\ell(\theta, k)] = 1 - p(\theta = k \mid y).
\]

最小化期望后验损失等价于最大化后验概率 $p(\theta = k \mid y)$，因此
\[
k^{\text{Bayes}} = \underset{k \in \mathcal{K}}{\arg\max} \; p(\theta = k \mid y),
\]
即后验众数（posterior mode）或 MAP 估计（maximum a posteriori estimate）。

\begin{Theorem}[0-1 损失下的贝叶斯最优分类]\label{th:zero-one-loss-bayes}
  在 0-1 损失下，贝叶斯最优分类决策是选择具有最大后验概率的类别。
\end{Theorem}

\subsection{损失函数的比较}

\begin{Remark}
  平方损失、绝对损失和 0-1 损失分别强调了后验分布的不同特征：均值、中位数和众数。这三个量在对称分布（如正态分布）中重合，但在偏态分布中可能差异显著。选择哪种损失函数，取决于具体问题的结构和决策者的偏好。
\end{Remark}

\section{预验分析：决定收集什么数据}

在许多实际问题中，我们不仅需要在给定数据后做出决策，还需要\textbf{在收集数据之前}决定收集哪些数据。这一问题被称为预验分析（preposterior analysis）。

\subsection{预验预期损失}

设我们考虑一个未来的观测 $y'$，其分布依赖于当前未知但通过先验分布 $p(\theta)$ 建模的参数 $\theta$：
\[
p(y') = \int p(y' \mid \theta) \, p(\theta) \, d\theta.
\]

值得注意的是，这里 $p(y')$ 正是我们在第三章中学习的\textbf{后验预测分布}（posterior predictive distribution）：在先验 $p(\theta)$ 下对未来数据 $y'$ 的预测分布。预验分析与后验预测分布之间的联系是：前者评估的是"收集新数据"这一行动本身的预期价值，而后者描述的是"在现有知识下"对未观测数据的预测。

在收集数据 $y'$ 之前，我们对参数 $\theta$ 仍处于先验分布的认知。预验预期损失（preposterior expected loss）定义为：对所有可能的未来数据和参数值取期望，
\[
\rho(d) = \mathbb{E}_{p(y', \theta)}[\ell(\theta, d(y'))] = \int \left[ \int \ell(\theta, d(y')) \, p(y' \mid \theta) \, dy' \right] p(\theta) \, d\theta.
\]

注意这里有两个期望层：内层关于 $y'$（给定 $\theta$），外层关于 $\theta$（先验分布）。通过交换积分次序，可以将预验预期损失写成
\[
\rho(d) = \mathbb{E}_{p(\theta)}[\rho(d \mid \theta)],
\]
其中 $\rho(d \mid \theta) = \mathbb{E}_{p(y'\mid\theta)}[\ell(\theta, d(y'))]$ 是在参数固定为 $\theta$ 时，决策 $d$ 的期望损失。

预验分析的核心问题是：如何选择实验设计 $y'$ 的采样方案（样本量、测量变量等），使得预验预期损失最小？

\begin{Example}[样本量确定]\label{ex:sample-size}
  设我们在一项临床试验中需要确定样本量 $n$。试验结束后，我们将使用后验均值 $\mathbb{E}(\theta \mid y')$ 作为治疗效果的点估计，采用平方损失。预验分析需要计算不同 $n$ 值下的预验预期损失
  \[
  \rho(n) = \mathbb{E}_{p(\theta)}[\text{var}(\theta \mid y')],
  \]
  并选择使 $\rho(n)$ 最小（或在成本约束下使 $\rho(n)$ 可接受）的 $n$。
  \footnote{详细推导见附录 \cref{sec:derivation-sample-size}.}
\end{Example}

\section{层次决策问题与经验贝叶斯}

在层次模型中，决策问题呈现出特殊的结构。考虑一个两层次层次模型：
\[
\theta_i \sim p(\theta \mid \phi), \quad y_i \mid \theta_i \sim p(y_i \mid \theta_i),
\]
其中 $\phi$ 是超参数。

\subsection{层次决策的两阶段特性}

给定超参数 $\phi$ 的后验分布 $p(\phi \mid y)$，层次决策问题可以分解为两个层次：

\textbf{第一层（条件最优决策）}：给定 $\phi$ 和数据 $y$，对每个 $i$ 选择最优 $\theta_i$ 的估计 $\hat{\theta}_i(\phi, y)$。

\textbf{第二层（全贝叶斯决策）}：在第一层的基础上，对超参数 $\phi$ 的选择或估计进行决策。

经验贝叶斯（empirical Bayes）方法采取一种务实的策略：用数据估计 $\phi$（例如使用边缘似然函数 $p(y \mid \phi)$ 的最大似然估计 $\hat{\phi}$），然后在估计得到的 $\phi = \hat{\phi}$ 条件下进行第一层决策。

从决策理论的角度看，经验贝叶斯是一种\textbf{可迭代替代}（iterative alternative）全贝叶斯方法，在计算上更为可行，但理论上不如全贝叶斯方法严格。

\begin{Example}[正态均值层次模型的 shrinkage 估计]\label{ex:normal-hierarchical}
  考虑一个经典的正态均值层次模型：设有 $J$ 个学校（或 $J$ 个实验），第 $j$ 个学校的真实效果为 $\theta_j$，观测数据为
  \[
  y_j \mid \theta_j \sim \text{Normal}(\theta_j, \sigma_j^2), \quad \theta_j \sim \text{Normal}(\mu, \tau^2),
  \]
  其中 $\sigma_j^2$ 已知。在平方损失下，层次贝叶斯的最优估计具有著名的\textbf{shrinkage}形式：
  \[
  \hat{\theta}_j^{\text{Bayes}} = \left(1 - \frac{\sigma_j^2}{\sigma_j^2 + \tau^2}\right) y_j + \frac{\sigma_j^2}{\sigma_j^2 + \tau^2} \mu.
  \]
  这里估计被拉向群体均值 $\mu$，拉动的程度取决于组间方差 $\tau^2$ 与组内方差 $\sigma_j^2$ 的比值。当 $\tau^2$ 很大时（组间差异大），估计接近原始观测 $y_j$；当 $\tau^2$ 很小时（组间差异小），估计强烈收缩向 $\mu$。如果 $\tau^2$ 未知而需要用数据估计（经验贝叶斯），则用估计值 $\hat{\tau}^2$ 替代。
\end{Example}

\begin{Remark}
  层次决策问题的核心张力在于：第一层决策（对 $\theta_i$ 的估计）依赖于超参数 $\phi$ 的认知，而 $\phi$ 本身的认知又依赖于所有 $\theta_i$ 的整体结构。这种相互依赖关系在层次模型中通过\textbf{ shrinkage 估计}得到解决——个体估计被拉向群体均值，拉动的程度由数据的置信度决定。
\end{Remark}

\section{损失函数视角下的模型选择}

决策理论还为模型选择提供了统一框架。设我们有两个候选模型 $M_1$ 和 $M_2$，对应的后验概率分别为 $p(M_1 \mid y)$ 和 $p(M_2 \mid y)$。

在 0-1 损失（选择正确模型损失为 0，选择错误模型损失为 1）下，最优模型选择规则是选择后验概率最大的模型：
\[
\hat{M} = \underset{M \in \{M_1, M_2\}}{\arg\max} \; p(M \mid y).
\]

这一规则与贝叶斯因子（Bayes factor）密切相关，因为
\[
\frac{p(M_1 \mid y)}{p(M_2 \mid y)} = \frac{p(y \mid M_1)}{p(y \mid M_2)} \cdot \frac{p(M_1)}{p(M_2)} = B_{12} \cdot \text{先验比}.
\]

当先验等可能时，模型选择等价于比较贝叶斯因子。

损失函数框架的优势在于：它允许我们根据具体问题的代价结构来调整决策规则。例如，如果假阳性（错误选择 $M_1$）的代价远大于假阴性（错误选择 $M_2$），我们可能需要在后验概率上设定一个更严格的阈值。

\section{本章小结}

本章建立了贝叶斯决策理论的基本框架，其核心思想可以概括为以下几条：

\begin{enumerate}
  \item \textbf{效用与损失的二元性}：决策问题可以通过效用函数或损失函数来描述，两者通过取负号相互转换。
  \item \textbf{贝叶斯最优决策原则}：在给定后验分布 $p(\theta \mid y)$ 的条件下，最优决策是使期望后验损失最小化的决策。
  \item \textbf{损失函数决定估计量的性质}：平方损失对应后验均值，绝对损失对应后验中位数，0-1 损失对应后验众数。
  \item \textbf{预验分析将决策前移到实验设计阶段}：通过预验预期损失，我们可以在收集数据之前评估不同数据收集方案的优劣。
  \item \textbf{层次决策体现了多层次不确定性的处理}：层次模型中的 shrinkage 估计可以从决策理论角度得到自然的解释。
\end{enumerate}

决策分析将贝叶斯推断从"描述不确定性"推进到"基于不确定性做最优决策"，是贝叶斯思想在实际问题中落地的关键步骤。

\section{习题}\label{sec:ch9-exercises}

\begin{Exercise}{平方损失的性质}\label{exr:ch9-1}
证明在平方损失下，后验均值 $\mathbb{E}(\theta \mid y)$ 是唯一最优估计量。
\end{Exercise}

\begin{Exercise}{绝对损失与后验中位数}\label{exr:ch9-2}
设后验分布 $p(\theta \mid y)$ 均匀分布于区间 $[a, b]$。证明在绝对损失下，任何 $d \in [a, b]$ 都是最优决策，并解释这一结论的直觉。
\end{Exercise}

\begin{Exercise}{0-1 损失与后验概率}\label{exr:ch9-3}
考虑一个三分类问题，参数 $\theta \in \{\text{A}, \text{B}, \text{C}\}$，后验概率分别为 $0.7, 0.2, 0.1$。在 0-1 损失下，贝叶斯最优分类是什么？如果 A 类代表"患病"，且假阴性（将患者误判为健康）的代价是假阳性的 5 倍，应如何修改损失函数？
\end{Exercise}

\begin{Exercise}{预验分析}\label{exr:ch9-4}
设 $\theta \sim \text{Normal}(0, 1)$，数据 $y \mid \theta \sim \text{Normal}(\theta, \sigma^2)$，其中 $\sigma^2$ 已知。证明在平方损失下，预验预期损失 $\rho(n) = 1 / (1 + n)$，并说明随着样本量 $n$ 增加，预验预期损失单调递减。
\end{Exercise}

\begin{Exercise}{经验贝叶斯与全贝叶斯}\label{exr:ch9-5}
考虑一个分层模型：$\theta_i \sim \text{Normal}(\mu, \tau^2)$，$y_i \mid \theta_i \sim \text{Normal}(\theta_i, \sigma^2)$，其中 $\sigma^2$ 已知，$\mu$ 和 $\tau^2$ 为未知超参数。比较以下两种方法的决策：
\begin{enumerate}
  \item 全贝叶斯方法：在 $(\mu, \tau^2)$ 的先验分布下，对 $\theta_i$ 进行全贝叶斯推断；
  \item 经验贝叶斯方法：用边缘似然函数 $p(y \mid \mu, \tau^2)$ 估计 $(\mu, \tau^2)$，然后在估计值下对 $\theta_i$ 进行条件最优决策。
\end{enumerate}
哪种方法在理论上更严格？哪种在计算上更可行？
\end{Exercise}

\section{附录：公式推导}\label{sec:appendix-ch9}

\subsection{平方损失最优解的推导}\label{sec:derivation-squared-loss}

\textbf{背景（Background）}：在 \S 9.2 中，我们声称平方损失 $(\theta - d)^2$ 下的贝叶斯最优估计是后验均值 $\mathbb{E}(\theta \mid y)$。本节给出这一结论的完整推导。

\textbf{参数定义（Parameter Definitions）}：
\begin{itemize}
  \item $\theta$：未知标量参数；
  \item $d$：点估计量（决策变量）；
  \item $p(\theta \mid y)$：给定数据 $y$ 后的后验分布密度；
  \item $\rho(d \mid y) = \mathbb{E}_{p(\theta\mid y)}[(\theta - d)^2]$：决策 $d$ 的期望后验损失。
\end{itemize}

\textbf{目标（Goal）}：证明 $\underset{d}{\arg\min} \; \rho(d \mid y) = \mathbb{E}(\theta \mid y)$。

\textbf{推导步骤（Derivation Steps）}：

对 $\rho(d \mid y)$ 关于 $d$ 求导：
\[
\frac{\partial}{\partial d} \rho(d \mid y)
= \frac{\partial}{\partial d} \int (\theta - d)^2 \, p(\theta \mid y) \, d\theta
= \int \frac{\partial}{\partial d} (\theta - d)^2 \, p(\theta \mid y) \, d\theta
= \int -2(\theta - d) \, p(\theta \mid y) \, d\theta.
\]

交换求导与积分次序的合法性由被积函数的光滑性保证（控制收敛定理）。

令导数为零：
\[
\int -2(\theta - d) \, p(\theta \mid y) \, d\theta = 0
\;\Longleftrightarrow\;
\int \theta \, p(\theta \mid y) \, d\theta = d \int p(\theta \mid y) \, d\theta.
\]

由于 $p(\theta \mid y)$ 是密度函数，$\int p(\theta \mid y) \, d\theta = 1$，因此
\[
d = \int \theta \, p(\theta \mid y) \, d\theta = \mathbb{E}(\theta \mid y).
\]

二阶条件：$\frac{\partial^2}{\partial d^2} \rho(d \mid y) = 2 > 0$，确认是最小值点。

\textbf{结论}：平方损失下的贝叶斯最优估计量是后验均值 $\mathbb{E}(\theta \mid y)$。\qed

\subsection{样本量确定的预验分析推导}\label{sec:derivation-sample-size}

\textbf{背景（Background）}：在 \cref{ex:sample-size} 中，我们声称在平方损失下，预验预期损失 $\rho(n) = \mathbb{E}_{p(\theta)}[\text{var}(\theta \mid y')]$。本节证明这一结论，并推导具体形式。

\textbf{参数定义（Parameter Definitions）}：
\begin{itemize}
  \item $\theta$：手术成功率（未知参数）；
  \item $y'$：临床试验的 $n$ 个独立观测，$y'_i \sim \text{Binomial}(1, \theta)$；
  \item 决策 $d(y') = \mathbb{E}(\theta \mid y')$（平方损失下的贝叶斯最优决策）；
  \item 损失函数：$\ell(\theta, d) = (\theta - d)^2$。
\end{itemize}

\textbf{目标（Goal）}：计算 $\rho(n) = \mathbb{E}_{p(y', \theta)}[(\theta - \mathbb{E}(\theta \mid y'))^2]$。

\textbf{推导步骤（Derivation Steps）}：

第一步：对给定 $\theta$ 的条件期望损失 $\rho(n \mid \theta) = \mathbb{E}_{p(y'\mid\theta)}[(\theta - \mathbb{E}(\theta \mid y'))^2]$。

由贝叶斯公式，后验均值 $\mathbb{E}(\theta \mid y') = \frac{a + \sum y'_i}{a + b + n}$，其中 $a, b$ 是先验 Beta 参数。

对 $y'_i \sim \text{Binomial}(1, \theta)$ 求期望时，$\sum y'_i \sim \text{Binomial}(n, \theta)$，故 $\mathbb{E}(\sum y'_i) = n\theta$，$\text{var}(\sum y'_i) = n\theta(1-\theta)$。

因此：
\[
\mathbb{E}(\theta \mid y') = \frac{a + \sum y'_i}{a + b + n}
= \frac{a + b}{a + b + n} \cdot \frac{a}{a+b} + \frac{n}{a + b + n} \cdot \frac{\sum y'_i}{n}.
\]

第二步：对 $\theta$ 求期望。注意到 $\mathbb{E}(\theta) = \frac{a}{a+b}$，以及条件方差分解公式：
\[
\text{var}(\theta) = \mathbb{E}[\text{var}(\theta \mid y')] + \text{var}[\mathbb{E}(\theta \mid y')].
\]

经过计算可得：
\[
\rho(n) = \mathbb{E}_{p(\theta)}[\text{var}(\theta \mid y')]
= \frac{1}{1 + n} \cdot \text{var}(\theta)
\propto \frac{1}{1 + n}.
\]

当先验方差 $\text{var}(\theta)$ 固定时，$\rho(n)$ 随 $n$ 单调递减。\qed

\endinput


% Part III: Advanced Computation
\part{高级计算方法}

% Chapter 10: Introduction to Bayesian Computation
% Bayesian Data Analysis - Gelman et al.

\chapter{贝叶斯计算入门}

\section{本章导论}

在前面几章中，我们已经建立了贝叶斯推理的基本框架：先验分布、数据模型、似然函数，通过贝叶斯定理将它们结合为后验分布。理论上，这个框架是完整且优雅的。然而，当我们试图将这一理论应用于实际问题时，一个根本性的障碍立刻显现出来：

\textbf{后验分布的计算往往涉及高维积分，而这些积分在解析上通常是不可解的。}

例如，在层次模型中，我们需要计算后验均值、预测分布、模型证据等，每一种都需要在数十甚至数百个参数构成的高维空间中进行积分。传统的解析方法对此束手无策。

这正是本章要解决的问题：我们介绍三种主要的贝叶斯计算方法——\textbf{数值积分}、\textbf{蒙特卡洛（Monte Carlo）积分}和\textbf{渐近近似}。这些方法各有优劣，在不同场景下发挥各自的作用。我们将在第十一章深入讨论第四种也是最强大的方法——马尔科夫链蒙特卡洛（MCMC）。

\section{计算挑战的来源}

\subsection{从点估计到后验分布}

在经典统计学中，我们通常只关心某些感兴趣的参数的点估计——如均值或众数。贝叶斯方法则更进一步：我们想要完整的后验分布，以便进行区间估计、预测推断和模型比较。

形式上，对于参数 $\theta$ 和数据 $y$，后验分布为

\[
p(\theta | y) = \frac{p(y | \theta) p(\theta)}{p(y)},
\]

其中 $p(y) = \int p(y | \theta) p(\theta) \, d\theta$ 是边缘似然（或称证据）。

\textbf{问题出在哪里？}边缘似然 $p(y)$ 是一个维度等于 $\theta$ 维度的积分。当 $\theta$ 是单个参数时，这个积分尚可数值求解；但当 $\theta$ 包含数十或数百个参数时，\textbf{维度灾难（curse of dimensionality）}使得传统数值方法彻底失效。

\subsection{后验特征量的计算需求}

贝叶斯分析中常用的特征量包括：

\begin{enumerate}
  \item \textbf{后验均值}：$\mathbb{E}(\theta | y) = \int \theta \, p(\theta | y) \, d\theta$
  \item \textbf{后验方差}：$\text{var}(\theta | y) = \int (\theta - \mathbb{E}\theta)^2 \, p(\theta | y) \, d\theta$
  \item \textbf{预测分布}：$p(y_{\text{new}} | y) = \int p(y_{\text{new}} | \theta) \, p(\theta | y) \, d\theta$
  \item \textbf{边缘后验}：对部分参数积分得到其余参数的后验
\end{enumerate}

每一种都需要在高维空间中进行积分运算。

\begin{Remark}
  值得注意的是，在许多实际问题中，我们并不需要完整的联合后验分布 $p(\theta | y)$，而只需要某些边缘分布或特征量。这就引出了两类主要策略：\textbf{直接近似后验分布}（如变分推断）和\textbf{直接近似感兴趣的特征量}（如蒙特卡洛方法）。
\end{Remark}

\section{数值积分方法}

\subsection{简单求积法}

最基本的思路是用离散求和逼近连续积分。例如，对于一维积分，我们可以使用矩形法则或 Simpson 法则：

\[
\int_a^b f(x) \, dx \approx \sum_{i=1}^n w_i f(x_i),
\]

其中 $x_i$ 是求积节点，$w_i$ 是相应权重。

\textbf{为什么这在贝叶斯计算中很快失效？}在 $d$ 维空间中，网格方法需要 $n^d$ 个节点。如果 $d = 10$ 且每维取 10 个节点，就需要 $10^{10}$ 个节点——这在任何实际计算中都是不可行的。这种指数增长就是维度灾难的本质。

\subsection{Latin 超立方采样}

一种改进策略是\textbf{分层采样（stratified sampling）}：将积分区域划分为若干子区域，在每个子区域中独立采样。

Latin 超立方采样（Latin Hypercube Sampling, LHS）是一种优雅的分层采样方法，它确保样本在每个维度上均匀分布，同时保持采样的随机性。

对于 $d$ 维空间和 $n$ 个样本，LHS 将每个维度划分为 $n$ 个等概率区间，然后通过随机置换确保样本点不重复。这种方法的优势在于：即使样本量较小，也能获得比纯随机采样更均匀的覆盖。

然而，LHS 本质上仍是一种\textbf{低维方法}。当维度增加到数十甚至数百时，其效果迅速下降。

\section{蒙特卡洛积分}

\subsection{基本思想}

蒙特卡洛方法的核心思想出奇地简单：用随机模拟代替确定性计算。\footnote{关于蒙特卡洛方法的历史，详见附录 \cref{sec:mc-history}。}

假设我们想要计算 $\theta = \int h(x) \, dx$。如果我们能够从 $x$ 的分布 $p(x)$ 中独立抽取大量样本 $x_1, \ldots, x_n$，则可以用

\[
\hat{\theta}_n = \frac{1}{n} \sum_{i=1}^n h(x_i)
\]

作为 $\theta$ 的估计。强大数定律保证 $\hat{\theta}_n \to \theta$（当 $n \to \infty$）。

更一般地，如果我们想计算 $\int h(x) p(x) \, dx$，其中 $p(x)$ 是某个分布的密度，我们只需从 $p(x)$ 中采样，然后用样本均值估计积分。

\subsection{重要性采样}

在贝叶斯计算中，我们经常需要从难以采样的分布中估计积分。\textbf{重要性采样（importance sampling）}是一种巧妙的解决方案：不是从目标分布 $p(\theta | y)$ 直接采样，而是从另一个更容易采样的分布 $g(\theta)$ 中采样。

设我们想估计

\[
I = \int h(\theta) \, p(\theta | y) \, d\theta = \int h(\theta) \frac{p(y | \theta) p(\theta)}{p(y)} \, d\theta.
\]

如果从 $g(\theta)$ 中采样，注意到

\[
I = \int h(\theta) \frac{p(y | \theta) p(\theta)}{p(y)} \cdot \frac{g(\theta)}{g(\theta)} \, d\theta
    = \int \frac{h(\theta) p(y | \theta) p(\theta)}{p(y) g(\theta)} \, g(\theta) \, d\theta.
\]

令 $w(\theta) = \frac{p(y | \theta) p(\theta)}{g(\theta)}$，则

\[
\hat{I} = \frac{1}{n} \sum_{i=1}^n \frac{h(\theta_i) w(\theta_i)}{\frac{1}{n} \sum_{j=1}^n w(\theta_j)}
        = \sum_{i=1}^n h(\theta_i) \tilde{w}(\theta_i),
\]

其中 $\tilde{w}(\theta_i) = \frac{w(\theta_i)}{\sum_j w(\theta_j)}$ 是归一化权重。

\begin{Example}[10 维积分的重要性采样估计]\label{ex:importance-sampling-10d}
  考虑估计 $d$ 维单位球体积：$I_d = \int_{[0,1]^d} 1 \, d\theta$。

  使用标准正态分布作为重要性分布 $g(\theta)$，并将积分区域扩展到全空间 $I = \int_{\mathbb{R}^d} \mathbb{I}_{[0,1]^d}(\theta) \, d\theta$。这实际上是用 $g(\theta)$ 的样本估计 $\mathbb{I}_{[0,1]^d}(\theta)$ 在 $\mathbb{R}^d$ 上的积分。

  由于正态分布的尾部概率很小，我们需要足够多的样本才能获得稳定的估计。\footnote{详细推导和代码实现见附录 \cref{sec:importance-sampling-derivation}。}
\end{Example}

重要性采样的核心挑战是选择合适的\textbf{重要性分布} $g(\theta)$。好的重要性分布应该：
\begin{enumerate}
  \item \textbf{支撑集覆盖}：$g(\theta) > 0$ whenever $p(\theta | y) > 0$
  \item \textbf{轻尾部}：避免在尾部区域浪费太多样本
  \item \textbf{形状相近}：$g(\theta)$ 的形状应尽量接近 $|h(\theta)| p(\theta | y)$
\end{enumerate}

\subsection{有效样本量}

在重要性采样中，权重 $w_i$ 的分布反映了估计的效率。如果所有权重几乎相等，说明 $g(\theta)$ 与目标分布匹配良好；如果权重分布极其不均，说明大部分样本贡献很小。

我们用\textbf{有效样本量（Effective Sample Size, ESS）}量化这一点：

\[
\text{ESS} = \frac{n}{1 + \text{var}(w)},
\]

其中 $\text{var}(w)$ 是（未归一化）权重的变异系数。ESS 的直观含义是：与直接采样相比，我们需要多少样本来达到当前的估计精度？

在实践中，我们用

\[
\text{ESS} \approx \frac{n}{1 + \frac{\text{var}(w)}{\mathbb{E}(w)^2}}
\]

来估计 ESS。如果 $\text{ESS} \ll n$，说明重要性采样效果不佳，我们需要调整 $g(\theta)$ 或尝试其他方法。

\begin{Remark}
  ESS 是一个有用的诊断工具，但不是绝对的。当权重高度非均匀时，即使 ESS 看起来还可以，估计量 $\hat{I}$ 的方差也可能被低估。这被称为\textbf{权重复制问题（weight degeneracy）}，是在实际使用重要性采样时需要警惕的。
\end{Remark}

\subsection{拒绝采样}

\textbf{拒绝采样（rejection sampling）}是另一种从困难分布中采样的通用方法。\footnote{关于拒绝采样的详细讨论，包括与 MCMC 的关系，见附录 \cref{sec:rejection-sampling-detail}。}

其基本思想是：找到一个容易采样的包络分布 $g(\theta)$，使得存在常数 $c$ 满足

\[
p(\theta | y) \leq c \cdot g(\theta), \quad \forall \theta.
\]

然后算法如下：

\begin{enumerate}
  \item 从 $g(\theta)$ 中抽取候选样本 $\theta^*$
  \item 计算接受概率 $\alpha = \frac{p(\theta^* | y)}{c \cdot g(\theta^*)}$
  \item 以概率 $\alpha$ 接受 $\theta^*$，否则拒绝并返回步骤 1
\end{enumerate}

被接受的样本恰好服从目标分布 $p(\theta | y)$。

\begin{Example}[截断正态分布的拒绝采样]\label{ex:truncated-normal-rejection}
  假设我们想从截断正态分布 $N(\mu, \sigma^2)$ 在区间 $[a, b]$ 上的限制采样。一个自然的选择是使用同样的正态分布作为包络分布，但需要确保 $c \cdot g(\theta) \geq p(\theta | y)$ 在整个区间上成立。

  具体而言，我们需要 $c \geq \frac{\max_{\theta \in [a,b]} p(\theta | y)}{g(\theta)}$，这可以通过分析 $g$ 和 $p$ 的密度比来确定。\footnote{详细推导和 R 实现代码见附录 \cref{sec:truncated-normal-rejection}。}
\end{Example}

拒绝采样的关键挑战是找到合适的包络分布 $g(\theta)$ 和常数 $c$。如果 $c$ 太大，接受率会很低，算法效率严重下降；如果 $g(\theta)$ 的形状与 $p(\theta | y)$ 差异太大，$c$ 也不得不设得很大。

\subsection{采样重要性重采样}

\textbf{采样重要性重采样（Sequential Importance Resampling, SIR）}，也称为\textbf{粒子滤波}的基础，将重要性采样与重采样结合起来。

当直接估计后验分布 $p(\theta | y)$ 困难时，SIR 提供了一种近似采样的方法：

\begin{enumerate}
  \item 从简单建议分布 $g(\theta)$ 中抽取 $n$ 个样本 $\theta_1, \ldots, \theta_n$
  \item 计算非归一化权重 $w_i \propto \frac{p(y | \theta_i) p(\theta_i)}{g(\theta_i)}$
  \item 归一化权重 $\tilde{w}_i = w_i / \sum_j w_j$
  \item 通过重采样（依据权重 $\tilde{w}_i$）得到新样本集
\end{enumerate}

重采样步骤是 SIR 与普通重要性采样的关键区别：通过消除低权重样本、复制高权重样本，我们得到了一个更接近目标分布的样本集。

\begin{Remark}
  SIR 的一个重要应用是状态空间模型中的粒子滤波。在这类模型中，潜在状态 $x_t$ 是隐含的，我们只能通过观测 $y_t$ 间接推断。SIR 通过序贯地处理每个时间步的观测，逐步更新状态的后验分布。
\end{Remark}

\section{渐近近似方法}

当样本量较大时，中心极限定理为后验分布提供了良好的近似。

\subsection{Laplace 方法}

Laplace 方法的核心思想是用正态分布逼近后验分布。设

\[
p(\theta | y) \propto p(y | \theta) p(\theta) = \exp\{-n \cdot h(\theta)\},
\]

其中 $h(\theta) = -\frac{1}{n} \log p(y | \theta) p(\theta)$。

将 $h(\theta)$ 在其后验众数 $\hat{\theta}$ 处 Taylor 展开：

\[
h(\theta) \approx h(\hat{\theta}) + \frac{1}{2} (\theta - \hat{\theta})^T H(\hat{\theta}) (\theta - \hat{\theta}),
\]

其中 $H(\hat{\theta}) = \nabla^2 h(\hat{\theta})$ 是 Hessian 矩阵。

代入得

\[
p(\theta | y) \approx \exp\left\{-n h(\hat{\theta})\right\} \cdot \exp\left\{-\frac{n}{2} (\theta - \hat{\theta})^T H(\hat{\theta}) (\theta - \hat{\theta})\right\}.
\]

这正是均值为 $\hat{\theta}$、协方差为 $\frac{1}{n} H(\hat{\theta})^{-1}$ 的正态分布。\footnote{关于 Laplace 方法的更严格推导和误差分析，见附录 \cref{sec:laplace-method-derivation}。}

\begin{Example}[逻辑回归的 Laplace 近似]\label{ex:logistic-laplace}
  在逻辑回归中，似然函数为
  \[
  p(y | \theta) = \prod_{i=1}^n \text{logit}^{-1}(\theta^T x_i)^{y_i} (1 - \text{logit}^{-1}(\theta^T x_i))^{1-y_i}.
  \]
  对数似然的 Hessian 矩阵为
  \[
  H(\theta) = \sum_{i=1}^n \text{logit}^{-1}(\theta^T x_i) (1 - \text{logit}^{-1}(\theta^T x_i)) x_i x_i^T,
  \]
  其中 $\text{logit}^{-1}(z) = \frac{1}{1 + e^{-z}}$。

  Laplace 近似给出 $\theta | y \approx N(\hat{\theta}, H(\hat{\theta})^{-1})$，其中 $\hat{\theta}$ 是后验众数。\footnote{关于如何在实践中使用这个近似的详细讨论，见附录 \cref{sec:logistic-laplace-extra}。}
\end{Example}

Laplace 方法的优势在于计算效率高：当 $n$ 足够大时，它提供了对后验分布的快速近似。其局限性在于它假设后验分布是单峰的且接近正态——这在许多实际问题是合理的，但绝非普遍成立。

\subsection{BIC 近似}

贝叶斯信息准则（BIC）是大样本下边缘似然的重要近似。设 $d$ 是参数维度，$n$ 是样本量，则

\[
\text{BIC} = -2 \log p(y | \hat{\theta}_{\text{MLE}}) + d \log n.
\]

更精确地，边缘似然的 BIC 近似为

\[
\log p(y) \approx \log p(y | \hat{\theta}_{\text{MLE}}) - \frac{d}{2} \log n + O(1).
\]

BIC 在模型选择中广泛应用，因为它提供了一个在似然性和模型复杂度之间权衡的准则。\footnote{关于 BIC 与其他模型选择准则（如 AIC、WAIC）的比较，见附录 \cref{sec:model-selection-criteria}。}

在贝叶斯计算中，BIC 的意义在于它提供了一种廉价的边缘似然近似，而无需进行昂贵的高维积分。这在需要快速筛选大量候选模型时尤其有用。

\section{变分推断初步}

变分推断（Variational Inference）是近年来发展迅速的一种贝叶斯计算方法。\footnote{变分推断的完整理论，包括平均场近似和信念传播，见附录 \cref{sec:variational-inference-full}。}

其核心思想是：\textbf{用一类简单的分布 $Q$ 去近似难以处理的后验分布 $p(\theta | y)$}。具体而言，我们寻找 $q(\theta) \in Q$ 使得

\[
q^*(\theta) = \arg\min_{q \in Q} \text{KL}(q(\theta) \| p(\theta | y)).
\]

由于 KL 散度涉及 $p(\theta | y)$，直接最小化 KL 不可行。但可以证明（通过简单的代数变形）：

\[
\log p(y) = \text{KL}(q \| p) + \mathcal{L}(q),
\]

其中 $\mathcal{L}(q) = \int q(\theta) \log \frac{p(y | \theta) p(\theta)}{q(\theta)} \, d\theta$ 是\textbf{变分下界（ELBO）}。

由于 $\log p(y)$ 是常数，最小化 $\text{KL}(q \| p)$ 等价于最大化 $\mathcal{L}(q)$。而 $\mathcal{L}(q)$ 只涉及 $q(\theta)$ 和对数似然/先验，因而是可以计算的。

变分推断的优势在于其\textbf{计算效率}：通过优化有限的变分参数，变分推断可以在百万级参数规模的模型中进行后验近似。这使其成为现代大规模贝叶斯模型（如主题模型、深度贝叶斯网络）的首选方法。

其局限性在于近似误差难以控制：我们无法知道 $q^*(\theta)$ 与真实后验 $p(\theta | y)$ 相差多少，除非进行昂贵的后验校验。

\begin{Remark}
  变分推断与 MCMC 是互补的关系。当我们需要\textbf{可扩展性}时（大规模数据、大参数空间），变分推断是首选；当我们需要\textbf{无偏估计}和\textbf{可证明收敛}时，MCMC 更适合。在实际应用中，两者经常结合使用：先用变分推断快速获得初始近似，再用 MCMC 进行精细采样。
\end{Remark}

\section{方法选择指南}

面对多样的贝叶斯计算方法，如何选择？以下是一些经验法则：

\begin{enumerate}
  \item \textbf{低维问题（$d \leq 5$）}：优先尝试数值积分或拒绝采样。网格方法或简单的自适应求积可能已经足够。
  \item \textbf{中等维度（$5 < d \leq 20$）}：重要性采样和 SIR 是合理的选择。关键是选择好的重要性分布——通常可以使用 Laplace 近似的正态分布作为建议分布。
  \item \textbf{高维问题（$d > 20$）}：MCMC（下一章主题）或变分推断是主要工具。简单的重要性采样在极高维度下效果很差。
  \item \textbf{模型结构}：如果模型具有共轭结构，利用条件分布的解析形式可以大大简化计算（如 EM 算法）。
  \item \textbf{计算资源}：当有 GPU 可用时，基于蒙特卡洛的方法可以并行化；当只有 CPU 时，变分推断的确定性优化可能更快。
\end{enumerate}

\section{本章小结}

本章我们介绍了贝叶斯计算的基本工具：

\begin{enumerate}
  \item \textbf{计算挑战的本质}：高维积分是贝叶斯计算的核心障碍
  \item \textbf{数值积分方法}：低维有效，但受维度灾难限制
  \item \textbf{蒙特卡洛积分}：用随机采样代替确定性积分
  \item \textbf{重要性采样}：从建议分布采样，加权修正
  \item \textbf{拒绝采样}：用包络分布控制接受率
  \item \textbf{采样重要性重采样（SIR）}：结合重要性采样与重采样
  \item \textbf{Laplace 方法}：用正态分布近似后验分布
  \item \textbf{变分推断}：用优化方法近似后验
\end{enumerate}

这些方法为下一章马尔科夫链蒙特卡洛（MCMC）奠定了基础。MCMC 是解决高维贝叶斯计算问题的最强大武器，它将在第十一章和第十二章中得到系统阐述。

\section{下节预告}

下一章我们将学习\textbf{马尔科夫链蒙特卡洛（MCMC）}的基本原理，包括 Metropolis-Hastings 算法和 Gibbs 采样。这两种方法是现代贝叶斯统计计算的核心工具。

\endinput

% Chapter 11: Basics of Markov chain simulation
% Bayesian Data Analysis - Gelman et al.

\chapter{马尔可夫链模拟基础}

\section{本章导论}

在前几章中，我们学习了如何通过共轭先验和解析计算来处理贝叶斯推断。然而，现实世界中的大多数贝叶斯模型都无法得到解析解——后验分布涉及高维积分，这些积分通常没有闭式形式。

这就引出了一个根本性的问题：\textbf{如何从无法解析表达的分布中抽取样本？}

答案在统计物理和概率论中早有先例。物理学家 Metropolis 和 Ulam 在 20 世纪 40 年代就发现，通过模拟随机过程可以从复杂分布中抽取样本。这一思想后来发展成为马尔可夫链蒙特卡洛（MCMC）方法，成为现代贝叶斯统计的计算基石。

本章我们将学习两类核心的模拟技术：
\begin{enumerate}
  \item \textbf{重要性抽样（Importance Sampling）}和\textbf{拒绝抽样（Rejection Sampling）}：直接从目标分布抽样的技术
  \item \textbf{马尔可夫链模拟（Markov Chain Simulation）}：通过构造马尔可夫链来探索目标分布的技术
\end{enumerate}

这些方法共同构成了贝叶斯计算的工具箱，使我们能够处理从前无法解决的复杂模型。

\section{重要性抽样}

\subsection{基本思想}

考虑一个简单的问题：我们希望计算某个函数 $h(\theta)$ 在后验分布下的期望
\begin{equation}\label{eq:posterior-expectation}
\mathbb{E}(h(\theta) | y) = \int h(\theta) \, p(\theta | y) \, d\theta
\end{equation}
但直接计算这个积分通常是不可行的。

重要性抽样的核心思想是：\textbf{不直接从 $p(\theta | y)$ 抽样，而是从一个更容易抽样的分布 $g(\theta)$ 中抽取样本，然后用权重调整来估计期望。}

具体地，如果我们从 $g(\theta)$ 中抽取样本 $\theta^{(1)}, \ldots, \theta^{(m)}$，则
\begin{equation}\label{eq:is-estimate}
\mathbb{E}(h(\theta) | y) \approx \frac{1}{m} \sum_{i=1}^{m} h(\theta^{(i)}) \cdot \frac{p(\theta^{(i)} | y)}{g(\theta^{(i)})}
\end{equation}
其中比值 $\frac{p(\theta | y)}{g(\theta)}$ 称为\textbf{重要性权重（importance weight）}。

\subsection{重要性权重的计算}

注意到 $p(\theta | y) \propto p(y | \theta) p(\theta)$，因此重要性权重可以写成
\begin{equation}\label{eq:importance-weight}
w(\theta) = \frac{p(y | \theta) p(\theta)}{g(\theta)}
\end{equation}
这个权重不需要归一化常数——只要 $g(\theta)$ 的支撑包含 $p(\theta | y)$ 的支撑即可。

\begin{Remark}
  重要性抽样的效果高度依赖于 $g(\theta)$ 和 $p(\theta | y)$ 的匹配程度。如果 $g(\theta)$ 在 $p(\theta | y)$ 重要的区域取值很小，重要性权重会变得极不稳定（方差可能趋于无穷大）。
\end{Remark}

\section{拒绝抽样}

\subsection{动机：为什么需要拒绝？}

重要性抽样虽然简单，但有一个根本缺陷：当 $g(\theta)$ 和 $p(\theta | y)$ 匹配不好时，重要性权重的方差会爆炸。

拒绝抽样（Rejection Sampling）提供了一种保证收敛精度的替代方案。其思想是：\textbf{找到一个可抽样的覆盖分布，并在接受率上引入随机拒绝机制。}

\subsection{算法}

设 $p(\theta)$ 是目标分布（可以是未归一化的），$g(\theta)$ 是我们能抽样的提议分布，$M$ 是一个常数使得对所有 $\theta$ 都有
\begin{equation}\label{eq:coverage-condition}
p(\theta) \leq M \cdot g(\theta)
\end{equation}

拒绝抽样的步骤如下：
\begin{enumerate}
  \item 从 $g(\theta)$ 中抽取候选样本 $\theta^*$
  \item 计算接受概率 $u = \frac{p(\theta^*)}{M \cdot g(\theta^*)}$
  \item 从 $\mathrm{Uniform}(0,1)$ 抽取 $u^*$
  \item 如果 $u^* \leq u$，则接受 $\theta^*$；否则拒绝
\end{enumerate}

被接受的样本 $\theta^*$ 确实来自目标分布 $p(\theta)$。\footnote{严格证明见附录 \cref{sec:derivation-rejection-sampling}}

\begin{Example}[拒绝抽样的几何直观]
  想象在 $p(\theta)$ 曲线下方投掷飞镖。如果我们的提议分布 $g(\theta)$ 是均匀的，那么 $M$ 就是 $p(\theta)$ 的最大值的倒数（归一化后）。只有落在 $p(\theta)$ 曲线下方的点才会被接受。
\end{Example}

\subsection{接受率的计算}

拒绝抽样的期望接受率为下式。\footnote{推导见附录 \cref{sec:derivation-rejection-acceptance}}
\begin{equation}\label{eq:rejection-acceptance-rate}
\text{accept rate} = \frac{1}{M} \int \min\bigl(p(\theta), M g(\theta)\bigr) \, d\theta
\end{equation}
当覆盖条件 \cref{eq:coverage-condition} 成立时，上式简化为 $\frac{1}{M} \int p(\theta) \, d\theta = \frac{1}{M}$。

因此，我们需要选择尽可能小的 $M$ 来提高接受率，但 $M$ 必须满足对所有 $\theta$ 的覆盖条件。

这个条件是拒绝抽样的主要限制：对于高维分布或支撑集不规则的分布，找到一个有效的覆盖分布 $g(\theta)$ 和常数 $M$ 往往非常困难。

\section{重要性重抽样（SIR）}

\subsection{动机：结合两种方法的优势}

拒绝抽样保证了精度，但接受率可能极低；重要性抽样效率高，但方差可能失控。

\textbf{重要性重抽样（SIR, Sampling Importance Resampling）}结合了两者的优点：首先用重要性抽样生成大量加权样本，然后通过重抽样来消除权重的不均匀性。

\subsection{算法}

SIR 的步骤：
\begin{enumerate}
  \item 从提议分布 $g(\theta)$ 抽取 $m$ 个候选样本 $\theta^{(1)}, \ldots, \theta^{(m)}$
  \item 计算每个样本的重要性权重 $w^{(i)} = \frac{p(\theta^{(i)} | y)}{g(\theta^{(i)})}$
  \item 归一化权重：$\tilde{w}^{(i)} = \frac{w^{(i)}}{\sum_{j=1}^{m} w^{(j)}}$
  \item 从 $\{ \theta^{(1)}, \ldots, \theta^{(m)} \}$ 中按概率 $\tilde{w}^{(i)}$ 有放回抽取 $n$ 个样本
\end{enumerate}

最终得到的 $n$ 个样本近似服从目标分布 $p(\theta | y)$。

\begin{Remark}
  SIR 的有效性依赖于 $m$ 足够大——$m$ 越大，重抽样得到的分布越接近真实目标分布。通常取 $m = 5n$ 到 $m = 10n$。
\end{Remark}

SIR 方法在实际应用中非常广泛，特别是在处理层次模型和混合模型时。它的缺点是仍然依赖于选择合适的提议分布 $g(\theta)$。

\section{马尔可夫链基础}

\subsection{从独立抽样到依赖抽样}

前述方法有一个共同特点：每次抽样都是独立的。但对于复杂的高维分布，找到一个好的提议分布本身就是一个难题。

马尔可夫链模拟提供了一种革命性的替代思路：\textbf{不追求独立同分布抽样，而是构造一个马尔可夫链，使其平稳分布恰好是我们想要的目标分布。}

这样，只需要运行马尔可夫链足够长时间，就可以从目标分布中抽样——即使我们永远无法写出这个分布的解析形式。

\subsection{马尔可夫链的定义}

回忆一下，\textbf{马尔可夫链}是一个随机过程 $\theta^{(1)}, \theta^{(2)}, \ldots$，满足马尔可夫性质：
\begin{equation}\label{eq:markov-property}
\theta^{(t+1)} | \theta^{(1)}, \ldots, \theta^{(t)} \sim P(\theta^{(t+1)} | \theta^{(t)})
\end{equation}
也就是说，下一个状态只依赖于当前状态，与历史无关。

转移核 $P(\theta' | \theta)$ 完全决定了链的行为。

\subsection{平稳分布}

如果存在分布 $\pi(\theta)$ 使得
\begin{equation}\label{eq:stationary-distribution}
\pi(\theta') = \int P(\theta' | \theta) \, \pi(\theta) \, d\theta
\end{equation}
则称 $\pi(\theta)$ 为该马尔可夫链的\textbf{平稳分布（stationary distribution）}。

我们的目标是：构造一个马尔可夫链，使其平稳分布 $\pi(\theta) = p(\theta | y)$（后验分布）。

\subsection{细致平衡条件}

一个充分条件是\textbf{细致平衡（detailed balance）}：\footnote{推导见附录 \cref{sec:derivation-detailed-balance}}
\begin{equation}\label{eq:detailed-balance}
P(\theta' | \theta) \, \pi(\theta) = P(\theta | \theta') \, \pi(\theta')
\end{equation}
如果转移核满足细致平衡，则平稳分布存在且就是 $\pi(\theta)$。

\subsection{收敛性}

对于大多数我们构造的马尔可夫链，在适当条件下，链会从任何初始值出发，在运行足够多步后，收敛到平稳分布。这意味着：

\begin{enumerate}
  \item 我们可以用链的样本近似目标分布
  \item 收敛所需的步数称为"\textbf{混合时间（mixing time）}"
  \item 我们通常需要"\textbf{预烧期（burn-in）}"——丢弃前若干步的样本
\end{enumerate}

理解马尔可夫链的收敛行为是有效使用 MCMC 方法的关键。

\section{Metropolis 算法}

\subsection{历史背景}

1953 年，Metropolis、Rosenbluth、Rosenbluth、Teller 和 Teller 在一项关于分子物理的工作中提出了著名的 Metropolis 算法。\cite{metropolis1953equation} 这一算法后来被 Hastings 在 1970 年推广为更一般的形式。\cite{hastings1970monte}

\textbf{毫不夸张地说，Metropolis 算法彻底改变了贝叶斯统计的计算格局。}

\subsection{算法思想}

Metropolis 算法的巧妙之处在于：它只需要能计算目标分布的\textbf{非归一化}形式 $p^*(\theta) = c \cdot p(\theta | y)$，而不需要知道归一化常数 $c$。

算法的步骤：
\begin{enumerate}
  \item 选择一个对称的提议分布 $q(\theta' | \theta) = q(\theta | \theta')$（例如，以当前值为均值的高斯分布）
  \item 从 $q(\cdot | \theta^{(t)})$ 抽取候选样本 $\theta^*$
  \item 计算接受比 $r = \frac{p^*(\theta^*)}{p^*(\theta^{(t)})}$
  \item 以概率 $\min(1, r)$ 接受候选值 $\theta^{(t+1)} = \theta^*$；否则 $\theta^{(t+1)} = \theta^{(t)}$
\end{enumerate}

\subsection{为什么有效？}

可以验证，Metropolis 算法构造的马尔可夫链满足细致平衡 \cref{eq:detailed-balance}，因此其平稳分布正是目标分布 $p(\theta | y)$。\footnote{详细证明见附录 \cref{sec:derivation-metropolis}}

关键的洞察是：即使我们不知道归一化常数，比值 $r = \frac{p^*(\theta^*)}{p^*(\theta^{(t)})}$ 也不依赖于归一化常数，因为它们被抵消了。

\begin{Example}[一维高斯目标分布]
  假设目标分布是标准正态分布 $p(\theta) = \phi(\theta)$。使用对称的随机游走提议分布 $q(\theta' | \theta) = \mathrm{N}(\theta', \theta, \sigma^2)$。

  接受比为
  \begin{equation}\label{eq:gaussian-acceptance-ratio}
  r = \frac{p(\theta^*)}{p(\theta^{(t)})} = \exp\!\left(-\frac{1}{2}(\theta^*)^2 + \frac{1}{2}(\theta^{(t)})^2\right)
  \end{equation}
  候选值以概率 $\min(1, r)$ 被接受。

  如果候选值在概率密度更高的区域（更接近 0），则 $r > 1$，必然被接受；如果在更低区域（$r < 1$），则以概率 $r$ 接受，以概率 $1 - r$ 拒绝。
\end{Example}

\subsection{提议分布的选择}

提议分布的选择直接影响 Metropolis 算法的效率：

\begin{enumerate}
  \item \textbf{步长（scale）}：如果步长太大，接受率会很低，链几乎不动；如果步长太小，混合会很慢
  \item \textbf{形状}：提议分布应该与目标分布的形状大致匹配
  \item 经验法则：目标接受率约为 0.25 到 0.5
\end{enumerate}

调参是使用 Metropolis 算法的重要技巧，需要根据具体问题摸索。

\section{Metropolis-Hastings 算法}

\subsection{推广到非对称提议}

Metropolis 算法要求提议分布对称，这在很多情况下限制了灵活性。Hastings 在 1970 年提出了一个更一般的版本，允许使用非对称的提议分布。

\subsection{算法}

Metropolis-Hastings 算法的步骤：
\begin{enumerate}
  \item 从提议分布 $q(\theta' | \theta^{(t)})$ 抽取候选样本 $\theta^*$
  \item 计算接受比
    \begin{equation}\label{eq:mh-acceptance-ratio}
    r = \frac{p^*(\theta^*)}{p^*(\theta^{(t)})} \cdot \frac{q(\theta^{(t)} | \theta^*)}{q(\theta^* | \theta^{(t)})}
    \end{equation}
  \item 以概率 $\min(1, r)$ 接受或拒绝
\end{enumerate}

接受比中的额外因子 $\frac{q(\theta^{(t)} | \theta^*)}{q(\theta^* | \theta^{(t)})}$ 纠正了非对称提议带来的偏差。

当 $q$ 是对称的时，这个因子等于 1，Metropolis-Hastings 算法就退化为普通的 Metropolis 算法。

\section{Gibbs 抽样}

\subsection{条件分布的神奇作用}

在多参数问题中，直接从联合后验分布抽样通常是不可能的。但如果我们可以从各个参数的\textbf{条件分布}中抽样，就会大大简化问题。

\textbf{Gibbs 抽样}正是利用了这一思想。它是 Metropolis-Hastings 算法的一种特殊形式，提议分布就是条件分布，因此接受率恒为 1。

这一方法由 Geman 和 Geman 在 1984 年提出，\cite{geman1984stochastic} 最初用于图像恢复，后来成为贝叶斯统计的核心计算工具。

\subsection{算法}

对于 $d$ 维参数 $\theta = (\theta_1, \ldots, \theta_d)$，Gibbs 抽样的步骤：

\begin{enumerate}
  \item 初始化 $\theta^{(0)} = (\theta_1^{(0)}, \ldots, \theta_d^{(0)})$
  \item 对于 $t = 1, 2, \ldots$：
    \begin{enumerate}
      \item 从 $p(\theta_1 | \theta_2^{(t)}, \ldots, \theta_d^{(t)}, y)$ 抽取 $\theta_1^{(t+1)}$
      \item 从 $p(\theta_2 | \theta_1^{(t+1)}, \theta_3^{(t)}, \ldots, \theta_d^{(t)}, y)$ 抽取 $\theta_2^{(t+1)}$
      \item $\ldots$
      \item 从 $p(\theta_d | \theta_1^{(t+1)}, \ldots, \theta_{d-1}^{(t+1)}, y)$ 抽取 $\theta_d^{(t+1)}$
    \end{enumerate}
\end{enumerate}

一轮更新完成后，$\theta^{(t+1)} = (\theta_1^{(t+1)}, \ldots, \theta_d^{(t+1)})$。

\subsection{为什么接受率为 1？}

考虑更新 $\theta_j$ 时，提议分布就是条件分布 $q(\theta_j' | \theta_{-j}) = p(\theta_j' | \theta_{-j}, y)$。\footnote{详细推导见附录 \cref{sec:derivation-gibbs-acceptance}}

接受比为
\begin{equation}\label{eq:gibbs-acceptance-ratio}
r = \frac{p^*(\theta_j', \theta_{-j})}{p^*(\theta_j, \theta_{-j})} \cdot \frac{p(\theta_j | \theta_{-j}, y)}{p(\theta_j' | \theta_{-j}, y)}
\end{equation}
由于 $p^*(\theta_j', \theta_{-j}) = p^*(\theta_{-j}) \cdot p(\theta_j' | \theta_{-j}, y)$，类似地 $p^*(\theta_j, \theta_{-j}) = p^*(\theta_{-j}) \cdot p(\theta_j | \theta_{-j}, y)$，于是两个因子互为倒数，$r = 1$。因此 Gibbs 抽样总是接受提议值。

\subsection{Gibbs 抽样的优势与局限}

\begin{Remark}
  Gibbs 抽样的优势在于：
  \begin{enumerate}
    \item 无需调参（没有接受率问题）
    \item 每步只需要从一元条件分布抽样，通常比较简单
    \item 自然地利用了问题的条件结构
  \end{enumerate}

  局限性在于：
  \begin{enumerate}
    \item 必须能够从所有条件分布有效抽样
    \item 当条件分布不是标准形式时，可能需要其他 MCMC 步骤（\textbf{混合 MCMC}）
    \item 参数高度相关时，收敛可能很慢
  \end{enumerate}
\end{Remark}

\section{从后验分布抽样的实践}

\subsection{构建 MCMC 策略}

在实际应用中，构建一个有效的 MCMC 抽样器通常需要结合多种技术：

\begin{enumerate}
  \item \textbf{识别条件结构}：首先分析模型，找出哪些条件分布是共轭的或容易抽样的
  \item \textbf{Gibbs 为主}：优先使用 Gibbs 抽样处理条件分布可解析处理的部分
  \item \textbf{Metropolis-Hastings 补充}：对难以抽样的条件分布，使用 Metropolis-Hastings 步骤
  \item \textbf{块更新}：有时需要对多个参数进行联合更新以加速收敛
\end{enumerate}

\subsection{收敛诊断}

MCMC 的一个核心问题是判断链是否已经收敛到目标分布。常用的诊断方法包括：

\begin{enumerate}
  \item \textbf{图形诊断}：绘制链的轨迹图（trace plot），检查是否"平稳"
  \item \textbf{Gelman-Rubin 统计量}：比较多条链的方差
  \item \textbf{自相关函数（ACF）}：检查样本之间的相关性
  \item \textbf{有效样本量（ESS）}：考虑自相关后的等效独立样本数
\end{enumerate}

需要强调的是，这些诊断方法只能检测\textbf{未收敛}，而不能保证收敛。良好的建模实践要求使用多种诊断并谨慎解释结果。

\subsection{预烧期和 thinning}

\begin{enumerate}
  \item \textbf{预烧期（burn-in）}：丢弃前 $b$ 次迭代，因为链在达到平稳分布之前可能处于"预热"状态
  \item \textbf{Thinning（稀疏化）}：每隔 $k$ 步取一个样本，以减少样本自相关
\end{enumerate}

这些技术虽然会减少有效样本量，但可以提高估计的独立性，是实践中常用的策略。

\section{示例：从分层模型抽样}

为了说明 MCMC 的实际应用，考虑一个简单的分层模型：

\begin{equation}\label{eq:hierarchical-model}
y_j | \theta_j \sim \mathrm{N}(\theta_j, \sigma^2), \quad \theta_j | \mu, \tau \sim \mathrm{N}(\mu, \tau^2)
\end{equation}

其中 $\sigma^2$ 已知，我们希望从后验分布 $p(\mu, \tau, \theta | y)$ 中抽样。

使用 Gibbs 抽样，我们可以依次从以下条件分布抽样：\footnote{具体推导见附录 \cref{sec:derivation-hierarchical-conditionals}}

\begin{enumerate}
  \item $p(\theta_j | \mu, \tau, y)$：这是正态分布
    \[
    \theta_j | \mu, \tau, y \sim \mathrm{N}\!\left(\frac{\frac{\mu}{\tau^2} + \frac{y_j}{\sigma^2}}{\frac{1}{\tau^2} + \frac{1}{\sigma^2}}, \left(\frac{1}{\tau^2} + \frac{1}{\sigma^2}\right)^{-1}\right)
    \]
  \item $p(\mu | \theta, \tau, y)$：$\mu$ 的条件后验也是正态的
    \[
    \mu | \theta, \tau, y \sim \mathrm{N}\!\left(\frac{1}{J}\sum_{j=1}^J \theta_j, \frac{\tau^2}{J}\right)
    \]
  \item $p(\tau | \theta, \mu, y)$：这个条件分布通常不是标准形式，需要 Metropolis 步骤
\end{enumerate}

这个例子展示了\textbf{混合 MCMC}的价值：利用 Gibbs 抽样的简便性，同时用 Metropolis 处理非标准条件分布。

\section{简单代码示例}

为完整起见，以下是 Metropolis 算法的 R 语言实现示例：

\begin{verbatim}
# Metropolis algorithm for sampling from N(0,1)
metropolis <- function(n, sigma_proposal = 1) {
  samples <- numeric(n)
  theta <- 0  # initial value

  for (i in 1:n) {
    # Propose new value
    theta_star <- rnorm(1, mean = theta, sd = sigma_proposal)

    # Compute acceptance ratio (unnormalized)
    log_r <- -0.5 * (theta_star^2 - theta^2)

    # Accept or reject
    if (log(runif(1)) < log_r) {
      theta <- theta_star
    }
    samples[i] <- theta
  }
  return(samples)
}

# Run the algorithm
set.seed(123)
samples <- metropolis(10000, sigma_proposal = 0.5)

# Check results
hist(samples, prob = TRUE, breaks = 50)
curve(dnorm(x), add = TRUE, col = "red")
\end{verbatim}

这个简单的例子展示了 Metropolis 算法的核心逻辑：提议 $\rightarrow$ 计算接受比 $\rightarrow$ 接受或拒绝。

\section{本章小结}

本章我们学习了贝叶斯计算的核心技术——马尔可夫链模拟：

\begin{enumerate}
  \item \textbf{重要性抽样}：通过提议分布和权重估计期望，但方差可能不稳定
  \item \textbf{拒绝抽样}：保证精度的覆盖抽样方法，但接受率可能很低
  \item \textbf{SIR}：结合重要性抽样和重抽样，实用性强
  \item \textbf{马尔可夫链基础}：平稳分布、细致平衡、收敛性概念
  \item \textbf{Metropolis 算法}：利用对称提议分布从目标分布抽样
  \item \textbf{Metropolis-Hastings 算法}：推广到非对称提议分布
  \item \textbf{Gibbs 抽样}：利用条件分布的特殊结构，接受率恒为 1
  \item \textbf{实践要点}：收敛诊断、预烧期、混合 MCMC 策略
\end{enumerate}

这些技术共同构成了现代贝叶斯统计的计算基础，使我们能够处理从前无法解决的高维复杂模型。

下一章我们将学习如何进一步提高 MCMC 的计算效率，包括切片抽样（slice sampling）、Hamiltonian Monte Carlo 等高级技术。

\section{参考文献}

\begin{thebibliography}{99}
  \bibitem{metropolis1953equation} Metropolis, N., Rosenbluth, A. W., Rosenbluth, M. N., Teller, A. H., \& Teller, E. (1953). Equation of state calculations by fast computing machines. \textit{The Journal of Chemical Physics}, 21(6), 1087--1092.

  \bibitem{hastings1970monte} Hastings, W. K. (1970). Monte Carlo sampling methods using Markov chains and their applications. \textit{Biometrika}, 57(1), 97--109.

  \bibitem{geman1984stochastic} Geman, S., \& Geman, D. (1984). Stochastic relaxation, Gibbs distributions, and the Bayesian restoration of images. \textit{IEEE Transactions on Pattern Analysis and Machine Intelligence}, (6), 721--741.
\end{thebibliography}

\endinput

% Chapter 12: Computationally Efficient Markov Chain Simulation
% Bayesian Data Analysis - Gelman et al.

\chapter{高效马尔可夫链模拟}

\section{本章导论}

在前两章中，我们介绍了马尔可夫链模拟的基本工具——Gibbs 采样器和 Metropolis 算法。这些方法具有广泛的适用性，但在处理复杂模型时可能面临效率问题。

本章将基本 Gibbs 采样器和 Metropolis 算法视为构建更高级马尔可夫链模拟算法的基石。\footnote{HMC 的开创性工作是 Duane et al. (1987) 在物理学期刊上提出的，Neal (1993, 2011) 是 HMC 和 NUTS 的权威参考文献。}我们依次讨论参数化变换和调优参数设置以提高效率的方法、切片采样（slice sampling）\footnote{切片采样源于 Neal (2003) 的工作。}、可逆跳转采样（reversible jump sampling）\footnote{可逆跳转采样由 Green (1995) 提出。}以及模拟回火（simulated tempering）\footnote{模拟回火由 Marinari \& Parisi (1992) 引入。}等扩展技术。核心内容是哈密顿蒙特卡洛（Hamiltonian Monte Carlo，HMC），这是一种通过引入"动量"变量使每次迭代能在参数空间中移动更远的 Metropolis 算法推广。最后，我们介绍 Stan——一个实现 HMC 的通用软件环境。

\section{高效 Gibbs 采样器}

\subsection{变换与参数化}

Gibbs 采样器在参数化为独立分量时效率最高。当参数高度相关时，收敛会变得缓慢。解决这一问题最简单的方法是对参数进行线性变换。

\begin{Remark}
  这一原则同样适用于 Metropolis 跳跃。在正态或近似正态的设定中，跳跃核的理想协方差结构应与目标分布相同。
\end{Remark}

\subsection{辅助变量与数据增广}

通过添加辅助变量往往可以简化或加速 Gibbs 采样器的计算。\footnote{数据增广的思想可追溯到 Tanner \& Wong (1987) 的工作。}例如，混合分布的指标变量就是一种常见的辅助变量。这一思想也称为数据增广（data augmentation），在概念上和计算上都是有用的工具。

\begin{Example}[$t$ 分布的正态混合表示]
  辅助变量的一个简单而重要的例子出现在 $t$ 分布的建模中。$t$ 分布可以表示为正态分布的混合，这一表示在金融和信号处理中广泛应用。考虑从 $t_\nu(\mu, \sigma^2)$ 分布中抽取 $n$ 个独立数据点，给定 $\mu, \sigma^2$ 的似然等价于以下模型：

  \[
  y_i \sim \mathrm{N}(\mu, V_i), \quad V_i \sim \chi^{-2}(\nu, \sigma^2),
  \]

  其中 $V_i$ 是不可直接观测的辅助变量。
\end{Example}

该模型的联合后验分布 $p(\mu, \sigma^2, V | y)$ 可以使用 Gibbs 采样器进行模拟，最终关于 $\mu, \sigma$ 的模拟样本将代表原始 $t$ 模型下的后验分布。

在 $t$ 模型的潜在参数形式中，如果 $\sigma$ 的模拟值接近零，收敛可能会很慢。通过添加一个额外参数——即参数扩展（parameter expansion）——可以打破 $\tau$ 与 $V_i$ 之间的相关性，从而改善收敛性。\footnote{参数扩展的详细讨论见 Liu \& Wu (1999) 和 van Dyk \& Meng (2001)。}

\section{高效 Metropolis 跳跃规则}

对于任何给定的后验分布，Metropolis-Hastings 算法都有无限多种实现方式。\footnote{Metropolis 跳跃规则的理论结果来自 Gelman, Roberts \& Gilks (1995)。}即使在重新参数化之后，跳跃规则 $J_t$ 仍有无穷多种选择。

\subsection{两类主要跳跃规则}

第一类是随机游走型跳跃规则，通常使用以当前参数值为中心、方差设定的正态跳跃核。第二类则构造与目标分布近似相同的提议分布，目的是尽可能多地接受样本，而 Metropolis-Hastings 接受步骤主要用于修正近似误差。

\subsection{最优跳跃尺度}

对于 $d$ 维后验分布（经适当变换后为多元正态），使用正态随机游走跳跃核时，最优尺度为

\[
c \approx 2.4 / \sqrt{d},
\]

其中效率是相对于从后验分布独立抽样定义的。\footnote{效率 $0.3/d$ 的推导见附录 \cref{sec:metropolis-efficiency}。}

对于多元正态目标分布，这一最优 Metropolis 跳跃规则的效率约为 $0.3 / d$。

最优接受率在一维约为 $0.44$，在高维（$d > 5$）下降至约 $0.23$。\footnote{关于随机游走 Metropolis 最优接受率的理论结果见 Gelman, Roberts \& Gilks (1995)。}这一结果提示了一种自适应模拟算法：在并行模拟开始时使用固定算法（如 Gibbs 采样器或正态随机游走 Metropolis 算法），经过一定次数的模拟后，根据接受率调整跳跃规则的尺度。

\section{切片采样}

切片采样的思想来自 Neal (2003)。从 $d$ 维目标分布 $p(\theta | y)$ 中随机抽取 $\theta$，等价于从该分布下的区域中均匀抽样。形式上，切片采样从 $(\theta, u)$ 的 $d+1$ 维分布中抽样，其中对任意 $\theta$，$p(\theta, u | y) \propto 1$ 当 $u \in [0, p(\theta | y)]$ 时，否则为 $0$。

切片采样在 Gibbs 采样结构中特别有用，可用于抽样一维条件分布。

\section{可逆跳转采样}

在许多问题中，需要进行跨维度的马尔可夫链模拟，即参数空间的维度在迭代之间可能发生变化。\footnote{可逆跳转 MCMC 的开创性工作是 Green (1995)。}例如模型平均和有限混合模型就属于这种情况。

可逆跳转采样仍可使用 Metropolis 算法实现。其核心思想是引入额外的随机变量来匹配不同模型间的参数空间维度。

\section{模拟回火与并行回火}

多峰分布对马尔可夫链模拟提出了特殊挑战。当两个（或多个）峰被极低后验密度的区域分隔时，链可能长时间停留在单个峰附近。\footnote{模拟回火由 Marinari \& Parisi (1992) 引入，并行回火见 Geyer \& Thompson (1993)。}

模拟回火是改善这类问题性能的一种策略。该算法使用一族分布 $p_k(\theta | y)$，其中 $p_0(\theta | y) = p(\theta | y)$，而 $p_1, \ldots, p_K$ 具有相同的基本形状但混合性能更好。并行回火是其变体，在该方法中 $K+1$ 条并行链同时模拟，每条链对应一个密度 $q_k$。

\section{哈密顿蒙特卡洛}

Gibbs 采样器和 Metropolis 算法的一个固有低效之处在于其随机游走行为——在参数空间中移动时可能需要很长时间才能穿过目标分布。

这一观察促使物理学家和统计学家将哈密顿动力学（描述行星运动、粒子碰撞等的确定性物理定律）引入蒙特卡洛模拟。\footnote{HMC 的开创性工作是 Duane et al. (1987)，Neal (2011) 是 HMC 和 NUTS 的权威参考文献。}哈密顿蒙特卡洛（Hamiltonian Monte Carlo，HMC）借鉴物理学中的思想来抑制 Metropolis 算法中的局部随机游走行为，从而使算法能够更快地穿过目标分布。HMC 也称为混合蒙特卡洛，因为它结合了马尔可夫链蒙特卡洛和确定性模拟方法。

\subsection{物理直觉}

想象一个粒子在势能 landscape 上运动。粒子的位置 $\theta$ 对应于我们要采样的参数，动量 $\phi$ 对应于粒子的运动速度。哈密顿动力学保证能量守恒：动能（与动量相关）和势能（与目标分布的对数密度相关）之和恒定。

当粒子处于后验密度平坦区域时，梯度为零，动量保持不变，粒子以恒定速度"滑行"穿过参数空间。当粒子进入低密度区域时，梯度为负，动量减小，最终可能改变方向。这种机制使 HMC 能够有效地探索后验分布的各个区域，而不像随机游走那样在局部徘徊。正是这种确定性的轨迹探索，使得 HMC 在高维空间中比随机游走 Metropolis 效率高出数个量级。

\subsection{动量分布}

HMC 为目标空间中的每个分量 $\theta_j$ 添加一个"动量"变量 $\phi_j$。通常将 $\phi$ 设为多元正态分布，$\phi \sim \mathrm{N}(0, M)$，其中 $M$ 是质量矩阵。

\subsection{HMC 迭代的三个步骤}

HMC 每次迭代包含以下三个部分：

\begin{enumerate}
  \item 从动量的后验分布 $\phi \sim \mathrm{N}(0, M)$ 中抽取新值。
  \item 通过 $L$ 个"蛙跳"步骤同时更新 $(\theta, \phi)$。
  \item 执行接受-拒绝步骤。
\end{enumerate}

蛙跳步骤的核心是：
\[
\phi \gets \phi + \frac{1}{2} \epsilon \frac{d \log p(\theta | y)}{d \theta}, \quad
\theta \gets \theta + \epsilon M^{-1} \phi, \quad
\phi \gets \phi + \frac{1}{2} \epsilon \frac{d \log p(\theta | y)}{d \theta}.
\]

接受-拒绝步骤计算
\[
r = \frac{p(\theta^* | y) p(\phi^*)}{p(\theta^{t-1} | y) p(\phi^{t-1})},
\]
并以概率 $\min(r, 1)$ 接受新值。

\begin{Remark}
  HMC 需要后验密度（只需计算到乘法常数）和对数后验密度的梯度。在实际编程中，我们建议同时用数值方法编程梯度作为检查，但实际部署时应使用自动微分以提高效率。\footnote{自动微分的介绍见 Griewank \& Walther (2008)。}
\end{Remark}

\subsection{调优参数设置}

HMC 可在三个地方调优：动量变量的概率分布 $p(\phi)$、蛙跳步骤的缩放因子 $\epsilon$ 以及每次迭代的蛙跳步数 $L$。

理论表明，当接受率约为 $65\%$ 时 HMC 效率最优。\footnote{最优接受率 $65\%$ 的推导见附录 \cref{sec:hmc-acceptance-rate}。}如果接受率过低，应减小 $\epsilon$ 并相应增加 $L$；如果接受率过高，则增大 $\epsilon$ 并减小 $L$。

\subsection{无 U 形转弯采样器}

无 U 形转弯采样器（no-U-turn sampler）\footnote{NUTS 由 Hoffman \& Gelman (2013) 提出。}自适应地确定每次迭代的步数 $L$，而不是使用固定值。其核心思想是沿轨迹运行直到轨迹开始转向（更准确地说，直到动量变量 $\phi$ 与从迭代开始时位置 $\theta$ 行驶的距离之间的点积达到负值）。

\subsection{Riemannian 适应}

另一种优化方法是 Riemannian 适应，\footnote{Riemannian 适应由 Girolami \& Calderhead (2011) 提出。}将质量矩阵 $M$ 设置为与每一步对数后验密度局部曲率相符的值。该方法可与无 U 形转弯采样器结合使用。

\subsection{教育测试示例}

我们用第 5 章中的教育测试实验来说明 HMC 的调优过程。该模型中，参数向量 $\theta$ 有 $d = 10$ 维，对应 $\alpha_1, \ldots, \alpha_8, \mu, \tau$。

HMC 需要计算对数后验密度关于每个参数的梯度。对于该正态层次模型，梯度可以解析地写出。\footnote{该例子的详细代码实现见附录 C.4。}在调优过程中，我们观察到不同链的接受率差异很大（从 $0.02$ 到 $0.87$），通过调整步长 $\epsilon$ 和步数 $L$（保持 $\epsilon_0 L_0 \approx 1$），最终达到约 $65\%$ 的目标接受率。

该示例说明，虽然对于这个特定问题 Gibbs 采样器已经足够有效，但在更大更复杂的问题中，HMC 能够更有效地穿越高维参数空间。

\section{Stan：计算环境的开发}

哈密顿蒙特卡洛需要一定的编程和调优努力。在更复杂的设置中，我们开发了一个计算机程序 Stan（通过自适应邻域采样）来自动应用 HMC。\footnote{Stan 的原始参考文献是 Gelman et al. (2015)。}

与 JAGS 和 OpenBUGS 等基于 Gibbs 采样的通用软件相比，Stan 的主要优势在于：
\begin{itemize}
  \item 使用 NUTS 采样器，自适应地确定轨迹长度
  \item 利用自动微分计算梯度，避免手动求导
  \item 在热身阶段自动调优步长和质量矩阵
\end{itemize}

Stan 的关键步骤包括：数据和模型的输入、对数后验密度（到任意常数）及其梯度的计算、热身阶段中调优参数的设置、使用无 U 形转弯采样器在参数空间中移动，以及最终的收敛监控和推理汇总。

Stan 使用自动解析微分来计算梯度。对于任意 $\mathrm{C}++$ 表达式，它解析表达式并应用微分的基本规则来构建梯度的 $\mathrm{C}++$ 程序。这种自动微分通常比通过有限差分计算数值梯度快得多。

Stan 还可以在热身阶段自适应地调整质量矩阵 $M$ 和步长 $\epsilon$。

\section{本章小结}

本章我们学习了提高马尔可夫链模拟效率的各种方法：

\begin{enumerate}
  \item \textbf{高效 Gibbs 采样器}：通过变换、参数化、添加辅助变量和参数扩展来加速收敛
  \item \textbf{高效 Metropolis 跳跃规则}：最优跳跃尺度 $c \approx 2.4 / \sqrt{d}$，最优接受率约为 $0.44$（一维）至 $0.23$（高维）
  \item \textbf{切片采样}：用于抽样一维条件分布的一般方法
  \item \textbf{可逆跳转采样}：处理跨维度问题的 Metropolis 扩展
  \item \textbf{模拟回火与并行回火}：处理多峰分布的技术
  \item \textbf{哈密顿蒙特卡洛（HMC）}：通过引入动量变量消除局部随机游走行为，特别适合高维问题
  \item \textbf{无 U 形转弯采样器}：HMC 的自适应扩展
  \item \textbf{Stan}：实现 HMC 的通用软件环境
\end{enumerate}

这些技术共同构成了现代贝叶斯计算的核心工具集，使得处理复杂层次模型和高维问题成为可能。

\section{下节预告}

下一章我们将学习模态和分布近似（Modal and Distributional Approximations），这是处理无法使用马尔可夫链模拟的复杂模型的替代方法。

\endinput

% Chapter 13: Modal and Distributional Approximations
% Bayesian Data Analysis - Gelman et al.

\chapter{模态与分布近似}

\section{本章导论}

在前几章中，我们探讨了蒙特卡洛模拟方法来逼近后验分布。对于低维问题，直接抽样或拒绝抽样往往已经足够；而对于更复杂的模型，马尔可夫链蒙特卡洛（MCMC）提供了通用的迭代模拟框架。然而，在超高维问题中，即便是 MCMC 也可能面临计算瓶颈——每步迭代需要遍历整个参数空间，而在高维曲面上的"随机行走"（random walk）行为会导致混合（mixing）极慢。

这就引出了本章的核心问题：\textbf{当我们不需要完整的后验分布，而只需要快速近似推断时，是否存在比迭代模拟更高效的方法？}

答案在于\textbf{模态与分布近似}。其基本思想是：首先找到后验分布的众数（mode）——即最可能出现的参数值——然后在众数附近用解析形式（如正态分布或其混合）来近似完整的后验。这种近似方法有三方面的价值：

\begin{enumerate}
\item \textbf{快速推断}：无需运行耗时的 MCMC 链，直接从近似分布抽样或计算矩；
\item \textbf{MCMC 的初始化}：找到的后验众数可以作为 MCMC 的起点，加速收敛；
\item \textbf{诊断工具}：比较近似分布与实际模拟结果，可以检验模型拟合质量。
\end{enumerate}

本章将依次介绍：寻找后验众数的算法（\cref{sec:finding-posterior-modes}）；避免边界估计的先验分布选择（\cref{sec:boundary-avoiding-priors}）；基于众数的正态及混合近似（\cref{sec:normal-approx}）；EM 算法及其变体用于寻找边缘众数（\cref{sec:em-algorithm}）；以及两种基于条件矩的近似方法——变分贝叶斯（\cref{sec:variational-bayes}）和期望传播（\cref{sec:expectation-propagation}）。

\section{寻找后验众数}\label{sec:finding-posterior-modes}

\subsection{动机：为什么关注众数？}

在后验分布 $p(\theta|y)$ 中，众数 $\hat{\theta}$ 是概率密度最大的点。当后验分布近似正态或仅有一个对称峰时，众数与后验均值、后验中位数三者重合，是良好的点估计。但在非对称或多峰分布中，众数可能远离分布中心——这一点在分层模型中方差参数的后验上尤为明显（见 \cref{sec:boundary-avoiding-priors}）。

更重要的是，在构建解析近似（如正态近似）时，我们需要知道：\textbf{在哪里近似？曲率多大？} 这两个信息都隐含在众数及其二阶导数中。因此，寻找后验众数是构建分布近似的第一步。

\subsection{条件最大化}\label{sec:conditional-maximization}

最直接的众数寻找方法是\textbf{条件最大化}（conditional maximization），又称逐步上升法（stepwise ascent）。其基本步骤如下：

\begin{enumerate}
\item 选择参数空间的某个初始点 $\theta^{(0)}$；
\item 固定部分参数，更新其余参数使对数后验密度上升；
\item 交替更新，直到对数后验密度变化小于预设阈值。
\end{enumerate}

在许多标准统计模型中，条件分布 $p(\theta_i | \theta_{-i}, y)$ 具有解析形式，最大化容易实现。例如在分层正态模型中，给定其他组参数后，单个组均值的后验条件分布是正态的，直接取其均值即可完成一次更新。

条件最大化的优点是简单且稳定，缺点是可能收敛较慢，且只能找到初始点附近的一个局部众数。要找到多个众数，需要从参数空间的不同位置多次启动算法。

\begin{Example}[EM 算法的特例——逐步上升]\label{ex:conditional-maximization}
考虑经典的两参数正态模型：$y_i \sim N(\mu, \sigma^2)$，先验取部分共轭形式。令 $\theta = (\mu, \sigma^2)$，对数后验为
\[
\log p(\mu, \sigma^2 | y) = -\frac{n}{2}\log\sigma^2 - \frac{1}{2\sigma^2}\sum_{i=1}^n (y_i - \mu)^2 + \text{常数}.
\]
固定 $\sigma^2$（取当前估计值 $\sigma_{\text{old}}^2$），对 $\mu$ 求导并令其为零，得
\[
\hat{\mu} = \frac{1}{n}\sum_{i=1}^n y_i,
\]
即样本均值——这正是逐步上升法的一次更新。反过来固定 $\mu$ 可得 $\sigma^2$ 的更新公式。交替迭代即条件最大化。
\end{Example}\footnote{推导见附录 \cref{sec:conditional-maximization-derivation}。}

\subsection{牛顿法}\label{sec:newtons-method}

\textbf{牛顿法}（Newton's method，或 Newton-Raphson 算法）是寻找后验众数的另一核心工具。其基本思想是对对数后验密度 $L(\theta) = \log p(\theta|y)$ 在当前点附近做二次泰勒展开：

\[
L(\theta) \approx L(\theta^{(t)}) + L'(\theta^{(t)})^\top (\theta - \theta^{(t)}) + \frac{1}{2}(\theta - \theta^{(t)})^\top L''(\theta^{(t)})(\theta - \theta^{(t)}).
\]

令梯度为零，求得迭代公式：
\begin{equation}\label{eq:newton-update}
\theta^{(t+1)} = \theta^{(t)} - [L''(\theta^{(t)})]^{-1} L'(\theta^{(t)}).
\end{equation}

其中 $L'(\theta)$ 是一阶导数向量，$L''(\theta)$ 是二阶导数矩阵（Hessian 矩阵）。当后验分布近似正态时——即在大样本极限下（见第四章）——二次近似非常准确，牛顿法展现出\textbf{二次收敛}速度：误差以平方速率下降。

\begin{Remark}[初始值的重要性]
牛顿法对初始值敏感。如果初始点选在 $L''$ 非正定的区域，迭代可能发散。实践中可先用条件最大化获得一个合理的初始点，再启动牛顿法加速收敛。
\end{Remark}

\begin{Example}[一维牛顿法]\label{ex:newton-1d}
设后验密度 $p(\theta|y) \propto \exp(-\frac{1}{2}\theta^2)$（标准正态），则 $L(\theta) = -\frac{1}{2}\theta^2 + \text{常数}$，$L'(\theta) = -\theta$，$L''(\theta) = -1$。代入 \cref{eq:newton-update} 得
\[
\theta^{(t+1)} = \theta^{(t)} - \frac{-\theta^{(t)}}{-1} = 0,
\]
一步即收敛到真值。这反映了正态后验下牛顿法的卓越表现。
\end{Example}\footnote{推导见附录 \cref{sec:newton-derivation}。}

\subsection{拟牛顿法与共轭梯度法}\label{sec:quasi-newton}

计算和存储完整的 Hessian 矩阵 $L''(\theta)$ 在高维问题中代价高昂。\textbf{拟牛顿法}（quasi-Newton methods）——如 BFGS（Broyden-Fletcher-Goldfarb-Shanno）算法——仅利用梯度信息迭代逼近 $-L''$，避免了二阶导数的直接计算。

\textbf{共轭梯度法}（conjugate gradient methods）同样只需一阶导数，但通过共轭方向公式构造优化方向，比最速下降法效率更高，在存储受限的高维问题中尤为有用。

\subsection{数值导数计算}\label{sec:numerical-derivatives}

当对数后验的解析导数难以求出时，可用有限差分近似：
\begin{equation}\label{eq:numerical-first-derivative}
L'_i(\theta) \approx \frac{L(\theta + \delta_i e_i) - L(\theta - \delta_i e_i)}{2\delta_i},
\end{equation}
其中 $e_i$ 是第 $i$ 个标准基向量，$\delta_i$ 是小量（通常取 $10^{-4}$）。二阶导数的数值估计为
\begin{equation}\label{eq:numerical-second-derivative}
L''_{ij}(\theta) \approx \frac{L(\theta + \delta_i e_i + \delta_j e_j) - L(\theta - \delta_i e_i + \delta_j e_j) - L(\theta + \delta_i e_i - \delta_j e_j) + L(\theta - \delta_i e_i - \delta_j e_j)}{4\delta_i \delta_j}.
\end{equation}
数值导数的缺点是计算量随维度平方增长，但实现简单，适合调试阶段与解析梯度交叉验证。

\section{避免边界估计的先验分布}\label{sec:boundary-avoiding-priors}

\subsection{问题：后验众数落在边界上}\label{sec:mode-on-boundary}

在后验分布的单峰接近参数空间边界时，以众数作为点估计会产生不合理的结论。最典型的例子是分层模型中方差参数 $\tau$（组间标准差）的后验分布。

考虑 \cref{fig:tau-posterior}（图 13.1）所示的 8 所学校例子。边缘后验 $p(\tau|y)$ 的众数恰好在 $\tau = 0$——这意味着点估计告诉我们：所有学校的辅导效果完全相同（即 $\tau = 0$ 对应零方差）。但从实质角度看，我们不相信效应真精确为零；后验分布集中在 $\tau$ 的一个小邻域内，只是似然函数在零附近较为平坦，导致后验密度在边界处达到峰值。

\begin{figure}[htbp]
\centering
\caption{8 所学校例子中 $\tau$ 的边缘后验密度 $p(\tau|y)$。若用众数概括此分布，将得到 $\hat{\tau} = 0$，即边界估计。}\label{fig:tau-posterior}
\end{figure}

问题在于：\textbf{从贝叶斯全后验角度看，$\tau = 0$ 本身无概率（它是连续参数空间中的一个零测集），边缘后验分布本身并无问题；但以众数作为点估计则选了一个边界值。}

\subsection{避免 $\tau$ 边界估计的先验}\label{sec:zero-avoiding-prior-tau}

如果我们的目标是用后验众数来概括分布（可能出于计算便利或快速近似的目的），则应选择使后验众数\textbf{远离边界}的先验分布。

对于组间标准差 $\tau > 0$，我们要求先验 $p(\tau)$ 在 $\tau = 0$ 处为零，从而保证后验众数不在边界上。几种常见选择：

\begin{itemize}
\item \textbf{对数正态}：$\log\tau \sim N(\mu_\tau, \sigma_\tau^2)$，但其在零附近衰减过快，可能排除太多合理的小值；
\item \textbf{逆 gamma}：$\tau^2 \sim \text{Inv-gamma}(\alpha_\tau, \beta_\tau)$，同样在零附近有有效下界；
\item \textbf{Gamma$(2, A^2)$}：\textbf{本书推荐}，其中 $A$ 是尺度参数。该先验在 $\tau = 0$ 处为零但线性增长，保证先验密度在零附近不会压倒似然函数。
\end{itemize}

Gamma$(2, A^2)$ 的优势在于：它在 $\tau = 0$ 处的导数为正（线性增长），而非对数正态或逆 gamma 那般迅速趋于零或无穷。这使得无论似然函数在零附近多平坦，后验密度都不会被先验"截断"。

\begin{Remark}[与全贝叶斯推断的对比]
若目标是后验模拟（而非众数点估计），则无必要使用边界回避先验。通常使用均匀先验或半 Cauchy 先验（见第五章），它们允许 $\tau$ 取任意小值而不产生边界效应。本节的边界回避先验专为\textbf{以众数为总结量}的场景设计。
\end{Remark}

\subsection{相关系数的边界回避先验}\label{sec:boundary-avoiding-correlation}

在多维分层模型中，相关矩阵的参数 $\rho$ 也有类似问题。若后验分布的众数落在 $\rho = \pm 1$（对应奇异协方差矩阵），则基于众数的点估计将给出退化的结论。

解决思路相同：用先验 $p(\rho) \propto (1-\rho^2)$（即 $\text{Beta}(2,2)$ 分布在 $\rho = \cos\pi U$ 下的像），该先验在 $\rho = \pm 1$ 处为零但线性增长，从而保证后验众数不会落在边界。

更一般地，对于 $d \times d$ 协方差矩阵，Wishart$(d+3, A\mathbb{I})$ 先验在协方差矩阵退化时趋于零，是协方差矩阵众数估计的默认边界回避选择。

\section{正态及混合近似}\label{sec:normal-approx}

\subsection{基于众数曲率的多元正态近似}\label{sec:normal-at-mode}

找到后验众数 $\hat{\theta}$ 后，可以利用二阶导数构建正态近似。对对数后验做二次展开：

\[
\log p(\theta|y) \approx \log p(\hat{\theta}|y) - \frac{1}{2}(\theta - \hat{\theta})^\top \underbrace{[-L''(\hat{\theta})]}_{V_\theta^{-1}} (\theta - \hat{\theta}),
\]
即
\begin{equation}\label{eq:normal-approx}
p(\theta|y) \approx N(\theta | \hat{\theta}, V_\theta), \quad V_\theta = [-L''(\hat{\theta})]^{-1}.
\end{equation}

这里的方差矩阵 $V_\theta$ 正是 Fisher 信息矩阵的近似——它度量了后验分布在众数附近的曲率。

\begin{Remark}[参数变换]
在应用 \cref{eq:normal-approx} 之前，通常需要对参数做变换（如取对数、logit 变换），使得变换后的后验分布接近对称且定义在整个实数轴上。变换后需乘以雅可比行列式的绝对值来修正密度。
\end{Remark}

\subsection{拉普拉斯方法}\label{sec:laplace-method}

\textbf{拉普拉斯方法}是逼近形式为 $\int h(\theta)p(\theta|y)\,d\theta$ 的积分的解析工具。在实际统计计算中，我们经常需要计算后验均值 $\mathbb{E}(h(\theta)|y)$ 或后验矩——这些本质上都是上述形式的积分。当 $h(\theta)$ 平滑且后验分布近似正态时，拉普拉斯方法提供了高效的近似。

其核心思想是：用 $p(\theta|y)$ 在其众数处的正态近似来代替 $p$ 本身。具体地，对 $u(\theta) = \log(h(\theta)p(\theta|y))$ 在其众数 $\theta_0$ 处做二次展开：

\[
\int h(\theta)p(\theta|y)\,d\theta \approx h(\theta_0)p(\theta_0|y)(2\pi)^{d/2} | -u''(\theta_0) |^{-1/2}.
\label{eq:laplace-integral-approx}
\]
其中 $u(\theta) = \log(h(\theta)p(\theta|y))$，$\theta_0$ 是其众数点。当 $h(\theta)$ 平滑且后验分布近似正态时（典型的大样本情形），该近似质量良好。拉普拉斯方法尤其适用于单峰后验分布；对于多峰情形，可对每个峰分别应用，然后组合。

从 \cref{eq:normal-approx} 的正态近似出发，可以进一步构造\textbf{多元 $t$ 近似}——用 $t$ 分布替代正态，引入额外的自由度参数 $\nu$ 以捕获更厚的尾部：
\[
p_t(\theta) \propto \left( \nu + (\theta - \hat{\theta})^\top V_\theta^{-1} (\theta - \hat{\theta}) \right)^{-(d+\nu)/2}.
\]

\subsection{多峰情形的混合近似}\label{sec:mixture-approx}

当后验分布有 $K$ 个分离良好的峰时，单一正态近似失效。此时可将后验近似为\textbf{多元正态混合分布}：

\begin{equation}\label{eq:normal-mixture}
p(\theta|y) \approx \sum_{k=1}^K \omega_k \, N(\theta | \hat{\theta}_k, V_k),
\end{equation}
其中 $\hat{\theta}_k$、$V_k$ 和 $\omega_k$ 分别是第 $k$ 个峰的位置、协方差矩阵和相对权重。权重可通过在各峰处匹配未归一化后验密度来估计：
\[
\omega_k \propto q(\hat{\theta}_k|y) \, |V_k|^{1/2},
\]
其中 $q(\theta|y)$ 是未归一化后验密度。

从混合近似中抽样也很直接：首先根据混合权重 $\{\omega_k\}$ 做多项式抽样选出某个成分，然后从相应的正态（或 $t$）分布中抽样。该样本可用于重要性抽样或重要性重抽样（见第十章），以获得更接近真实后验的样本。

\begin{Example}[双峰后验的混合近似]\label{ex:bimodal-mixture}
考虑一个分层模型，其边缘后验 $p(\tau|y)$ 具有双峰——一个在 $\tau$ 较小的区域（对应同质性假设），另一个在较大的区域（对应异质性假设）。设两峰分别为 $\hat{\tau}_1 = 0.1$ 和 $\hat{\tau}_2 = 4.5$，对应的曲率半径估计为 $V_1 = 0.05^2$、$V_2 = 1.2^2$。则
\[
p(\tau|y) \approx \omega_1 N(\tau | 0.1, 0.05^2) + \omega_2 N(\tau | 4.5, 1.2^2),
\]
其中 $\omega_1 + \omega_2 = 1$。若 $q(\hat{\tau}_1|y) = 0.3$、$q(\hat{\tau}_2|y) = 0.7$，且 $|V_1|^{1/2} = 0.05$、$|V_2|^{1/2} = 1.2$，则 $\omega_1 \propto 0.3 \times 0.05 = 0.015$，$\omega_2 \propto 0.7 \times 1.2 = 0.84$，归一化后 $\omega_1 \approx 0.018$、$\omega_2 \approx 0.982$。
\end{Example}

\section{EM 算法寻找边缘后验众数}\label{sec:em-algorithm}

\subsection{从联合众数到边缘众数}\label{sec:joint-to-marginal}

在多参数问题中，联合后验的众数往往不是关注参数的最好近似起点。例如在分层模型中，我们可能只关心组间方差 $\tau$ 的后验分布，而非所有学校效应 $\theta_j$ 的联合后验。此时应寻找 $p(\phi|y)$ 的众数（$\phi$ 是感兴趣参数），而非 $p(\theta|y)$ 的。

\textbf{EM 算法}（expectation-maximization algorithm）正是为这一场景设计的：它将隐变量（或缺失数据）$\gamma$ 引入模型，交替执行"求期望"（E步）和"最大化"（M步），逐步找到边缘后验 $p(\phi|y)$ 的众数。

\subsection{EM 算法的推导}\label{sec:em-derivation}

设我们观察数据为 $y$，完整数据为 $(y, \gamma)$，参数为 $\phi$。考虑恒等式
\[
\log p(\phi|y) = \log p(\gamma, \phi|y) - \log p(\gamma|\phi, y),
\]
对 $\gamma$ 取期望（使用当前估计 $\phi^{\text{old}}$ 下的条件分布 $p(\gamma|\phi^{\text{old}}, y)$）：

\[
\log p(\phi|y) = E_{\gamma|\phi^{\text{old}}, y}\!\left[ \log p(\gamma, \phi|y) \right] - E_{\gamma|\phi^{\text{old}}, y}\!\left[ \log p(\gamma|\phi, y) \right].
\label{eq:em-log-identity}
\]

第二项在 $\phi = \phi^{\text{old}}$ 处取最大值（因为条件分布的熵在 $\phi$ 变化时不减）。因此，若我们找到一个 $\phi^{\text{new}}$ 使得第一项增大，则 $\log p(\phi^{\text{new}}|y) > \log p(\phi^{\text{old}}|y)$。这正是 EM 算法的核心洞察。

定义 Q 函数：
\[
Q(\phi|\phi^{\text{old}}) = E_{\gamma|\phi^{\text{old}}, y}\!\left[ \log p(\gamma, \phi|y) \right],
\]
则 EM 算法的每步迭代为：
\begin{enumerate}
\item \textbf{E 步}：计算 $Q(\phi|\phi^{(t)})$；
\item \textbf{M 步}：求 $\phi^{(t+1)} = \arg\max_\phi Q(\phi|\phi^{(t)})$。
\end{enumerate}

由于对数后验在每步迭代中单调递增，EM 算法必收敛到局部众数。

\begin{Theorem}[EM 算法的单调性]\label{th:em-monotonicity}
若每步迭代满足 $Q(\phi^{(t+1)}|\phi^{(t)}) \geq Q(\phi^{(t)}|\phi^{(t)})$，则对数边缘后验密度单调不减：$\log p(\phi^{(t+1)}|y) \geq \log p(\phi^{(t)}|y)$。
\end{Theorem}\footnote{证明见附录 \cref{sec:em-monotonicity-proof}。}

GEM（广义 EM）算法放松了 M 步的要求——只需\textbf{增加} $Q$ 而非严格最大化。当 $Q$ 的最大化困难时（如参数有约束），GEM 通过在可行方向上做条件最大化来更新。

\begin{Example}[正态均值与方差的部分共轭模型]\label{ex:normal-mean-variance-em}
设 $y_1, \ldots, y_n \sim N(\mu, \sigma^2)$，先验 $\mu \sim N(\mu_0, \tau_0^2)$，$\log\sigma$ 上取无信息均匀先验。此模型非完全共轭，无闭合形式边缘后验。我们用 EM 找 $\mu$ 的边缘后验众数（对 $\sigma$ 积分）。

令 $\gamma = \sigma^2$（隐变量），$\phi = \mu$。对数联合后验为
\[
\log p(\mu, \sigma^2|y) = -\frac{n+1}{2}\log\sigma^2 - \frac{1}{2\sigma^2}\sum_{i=1}^n (y_i - \mu)^2 - \frac{(\mu - \mu_0)^2}{2\tau_0^2} + \text{常数}.
\label{eq:normal-joint-log}
\]

E 步中，需求 $E_{\sigma^2|\mu^{\text{old}}, y}[(\sigma^2)^{-1}]$。给定 $\mu$ 时，$\sigma^2|\mu, y \sim \text{Inv-}\chi^2(n, s_{\text{old}}^2)$，其中 $s_{\text{old}}^2 = \frac{1}{n}\sum_i (y_i - \mu^{\text{old}})^2$。于是 $E[(\sigma^2)^{-1}] = n / s_{\text{old}}^2$。

将 \cref{eq:normal-joint-log} 中的 $E_{\sigma^2|\mu^{\text{old}}, y}[(\sigma^2)^{-1}]$ 替换为 $n/s_{\text{old}}^2$，得 Q 函数的表达式。该 Q 函数关于 $\mu$ 是二次的（具有正态对数核的形式），直接对 $\mu$ 求导并令其为零，得 M 步更新公式：
\[
\mu^{\text{new}} = \frac{\frac{1}{\tau_0^2}\mu_0 + n\frac{\bar{y}}{s_{\text{old}}^2}}{\frac{1}{\tau_0^2} + \frac{n}{s_{\text{old}}^2}}.
\label{eq:mu-new}
\]
迭代此式即收敛到 $\mu$ 的边缘后验众数。
\end{Example}\footnote{详细推导见附录 \cref{sec:normal-em-derivation}。}

\subsection{EM 的变体}\label{sec:em-variants}

\begin{itemize}
\item \textbf{ECM 算法}：将 M 步替换为一系列条件最大化（CM 步）。每步只更新参数的一个子集，同时约束其他子集保持不变。这在参数有约束（如协方差矩阵必须正定）时特别有用。

\item \textbf{PX-EM（参数展开 EM）}：引入额外的扩展参数以加速收敛。核心思想是：如果原始参数的后验曲面曲率变化剧烈，EM 收敛缓慢；通过重新参数化（引入"伪数据"或"展开参数"）使 $Q$ 函数更接近二次，可显著加速。
\end{itemize}

\begin{Remark}[调试 EM]
计算每步迭代后的对数边缘后验密度 $\log p(\phi^{(t)}|y)$，检查其是否单调递增。若出现下降，说明 E 步或 M 步实现有误。这是验证 EM 实现正确性的标准检查。
\end{Remark}

\section{变分贝叶斯}\label{sec:variational-bayes}

\subsection{动机：为什么需要变分近似？}\label{sec:variational-motivation}

EM 算法找到的是边缘后验的众数，而非整个后验分布。在许多场景下，我们希望获得后验的完整近似——包括均值、方差以及边缘分布。但对于复杂模型，精确计算后验分布的积分是 intractable 的（维数灾难）。

\textbf{变分贝叶斯}（Variational Bayes，VB）提供了一种替代方案：它寻找一个形式简单的近似分布 $q(\theta)$ 来逼近真实后验 $p(\theta|y)$，同时保证计算可行。其核心是最小化 $q$ 与 $p$ 之间的 KL 散度：

\[
q^* = \arg\min_q \, \text{KL}\!\left( q(\theta) \, \| \, p(\theta|y) \right).
\]

与 MCMC 相比，变分贝叶斯的优点是\textbf{确定性}（无随机模拟误差）、\textbf{快速}（只需迭代优化少数变分参数）；缺点是近似质量依赖于近似分布的假设形式，无法保证捕获后验的所有特征。

\subsection{平均场变分近似}\label{sec:mean-field}

变分贝叶斯最常用的假设是\textbf{平均场}（mean-field）近似：设 $q(\theta)$ 可以分解为各参数子群因子的乘积，
\[
q(\theta) = \prod_{j=1}^J q_j(\theta_j),
\]
其中 $\theta = (\theta_1, \ldots, \theta_J)$ 是参数的一个划分，各子群内部可能包含多个原始参数。

在该假设下，最优因子 $q_j^*(\theta_j)$ 满足：
\[
\log q_j^*(\theta_j) = E_{-j}\!\left[ \log p(\theta, y) \right] + \text{常数}.
\label{eq:variational-update}
\]
其中 $E_{-j}[\cdot]$ 表示对除 $\theta_j$ 以外所有其他参数取期望（在当前变分分布 $q_{-j}$ 下）。

\cref{eq:variational-update} 表明：\textbf{每个因子的最优形式由对数联合密度的期望决定}。若该期望是某个标准分布的核（kernel），则 $q_j^*$ 即为该标准分布——这意味着变分更新具有解析形式。

\begin{Example}[8 学校模型的变分贝叶斯]\label{ex:vb-8schools}
在 8 学校分层模型中，参数分为三组：学校效应 $\alpha = (\alpha_1, \ldots, \alpha_8)$、总体均值 $\mu$ 和组间标准差 $\tau$。设平均场近似 $q(\alpha, \mu, \tau) = q_\alpha(\alpha) \, q_\mu(\mu) \, q_\tau(\tau)$，则各因子的最优更新为：

\begin{itemize}
\item $q_\alpha(\alpha_j) = N(M_{\alpha_j}, S_{\alpha_j}^2)$，其中均值和方差依赖于 $q_\mu$ 和 $q_\tau$ 的当前矩；
\item $q_\mu(\mu) = N(M_\mu, S_\mu^2)$；
\item $q_\tau(\tau)$ 取某非标准形式（难以解析积分），实践中用数值近似或重新参数化。
\end{itemize}

\cref{fig:vb-convergence} 展示了一个典型运行中的迭代收敛过程：变分参数在约 50 次迭代后趋于稳定，KL 散度逐步下降（由算法保证单调性）。
\end{Example}

\begin{figure}[htbp]
\centering
\caption{8 学校模型的变分贝叶斯迭代过程。各变分参数（均值和标准差）在约 50 次迭代后收敛。右下图为 KL 散度 $\text{KL}(g\|p)$ 随迭代的变化（任意常数偏移），可见其单调下降。}\label{fig:vb-convergence}
\end{figure}

\subsection{变分推断的局限性}\label{sec:variational-limitations}

变分近似假设后验可分解为独立因子的乘积，这在许多实际场景中并不成立——尤其是当参数间存在强相关时。平均场近似会\textbf{低估后验方差}（因为假设独立等价于忽略了负相关），在极端情况下可能给出严重误导的推断。

评估变分近似质量的一种方法是：比较从变分分布抽样的样本与从 MCMC 得到的真实后验样本在各阶矩上的差异。若差异明显，应考虑更灵活的近似形式（如正则化流、神经网络近似等）。

\section{期望传播}\label{sec:expectation-propagation}

\subsection{动机：变分贝叶斯的不足与 EP 的引入}\label{sec:ep-motivation}

变分贝叶斯通过最小化 $\text{KL}(q\|p)$ 逼近后验，其平均场假设虽然简化了计算，但也带来了根本性问题：\textbf{平均场分解假设各参数子群相互独立，这与许多实际模型中参数存在强相关的现实不符。}其后果是变分近似倾向于低估后验方差——因为独立假设强制忽略了负相关结构。

更重要的是，变分贝叶斯只保证近似分布的均值（通过 \cref{eq:variational-update} 更新）逼近真实后验均值，但并不直接约束高阶矩。在需要精确后验方差或协方差进行不确定性量化的场景中，这一缺陷可能是致命的。

\textbf{期望传播}（Expectation Propagation，EP）正是为解决这些问题而设计的。EP 的核心洞察是：\textbf{与其整体逼近后验，不如逐个因子地"提炼"近似。}每个因子（如某个似然项或先验）对应模型的一个局部结构；EP 在保持其余因子不动的同时，迭代地用真实因子的统计量（矩）来精化对应的近似因子，最终使得近似后验在各阶矩上与真实后验匹配。

EP 的另一优势是\textbf{模块化}：模型中的每个因子可以独立定义其近似形式（如用正态近似某个非正态因子），而不像变分贝叶斯那样要求全局的平均场分解。

\subsection{EP 的算法步骤}\label{sec:ep-algorithm}

设后验分布可分解为 $F$ 个因子的乘积：
\[
p(\theta|y) \propto \prod_{f=1}^F m_f(\theta),
\]
其中每个 $m_f(\theta)$ 对应模型的一个组件（如先验 $p(\theta)$ 或第 $i$ 个数据点的似然 $p(y_i|\theta)$）。

EP 用一组近似因子 $\{\tilde{m}_f(\theta)\}$ 替代真实因子，使近似后验
\[
q(\theta) \propto \prod_{f=1}^F \tilde{m}_f(\theta)
\]
具有正确的各阶矩。算法通过\textbf{迭代 refinement} 实现，具体步骤如下：

\begin{enumerate}
\item \textbf{初始化}：令所有近似因子 $\tilde{m}_f(\theta) = 1$（或某个简单的初始形式），得到初始近似 $q(\theta) \propto 1$。

\item \textbf{迭代精化}：对每个因子 $f = 1, \ldots, F$，重复以下步骤：
\begin{enumerate}
\item \textbf{删除}：从当前近似中移除因子 $f$ 的贡献，得到边缘近似
\[
q_{\setminus f}(\theta) \propto q(\theta) / \tilde{m}_f(\theta).
\]

\item \textbf{包含修正}：计算真实因子 $m_f$ 与当前近似 $q_{\setminus f}$ 的"商"在 $q_{\setminus f}$ 下的动差。具体地，定义非归一化分布
\[
r_f(\theta) \propto \frac{m_f(\theta)}{q_{\setminus f}(\theta)},
\]
在 $q_{\setminus f}$ 下计算 $r_f(\theta)$ 的各阶矩（如均值和协方差）。

\item \textbf{矩匹配}：用 $r_f$ 的矩来精化近似因子 $\tilde{m}_f$，使得新近似因子具有与 $r_f$ 相同的矩。通常假设 $\tilde{m}_f$ 属于某参数化族（如正态分布），通过匹配矩来确定其参数。

\item \textbf{更新}：将精化后的 $\tilde{m}_f$ 放回近似分布，更新 $q(\theta) \propto q_{\setminus f}(\theta) \cdot \tilde{m}_f(\theta)$。
\end{enumerate}
\item \textbf{收敛判断}：重复上述迭代，直到近似后验 $q(\theta)$ 的各阶矩变化小于预设阈值，或达到最大迭代次数。
\end{enumerate}

\begin{Example}[EP 对正态似然的精化]\label{ex:ep-normal-likelihood}
考虑一个正态模型 $y_i \sim N(\mu, \sigma^2)$，$\sigma^2$ 已知。先验 $p(\mu) \propto 1$。此时每个数据点贡献一个因子 $m_i(\mu) = N(y_i|\mu, \sigma^2)$，后验精确为 $N(\mu|\bar{y}, \sigma^2/n)$。

EP 初始化所有 $\tilde{m}_i(\mu) = 1$，则 $q(\mu) \propto 1$（不proper但可工作）。对第一个因子 $f=1$：
\begin{itemize}
\item $q_{\setminus 1}(\mu) \propto 1$；
\item $r_1(\mu) \propto m_1(\mu) / 1 = N(y_1|\mu, \sigma^2)$；
\item 在 $q_{\setminus 1}$（均匀分布）下，$r_1$ 的均值和方差即为 $y_1$ 和 $\sigma^2$；
\item 令 $\tilde{m}_1(\mu) = N(y_1|\mu, \sigma^2)$，则更新后 $q(\mu) \propto N(y_1|\mu, \sigma^2)$。
\end{itemize}
对所有因子迭代，最终 $q(\mu) = N(\bar{y}|\mu, \sigma^2/n)$，与精确后验完全一致——这体现了 EP 在正态模型中的优良性质。
\end{Example}\footnote{详细推导见附录 \cref{sec:ep-normal-derivation}。}

\begin{Remark}[EP 与变分贝叶斯的区别]
两者都通过迭代优化逼近后验，但优化目标不同：
\begin{itemize}
\item \textbf{变分贝叶斯}最小化 $\text{KL}(q\|p)$，等价于让 $q$ 覆盖 $p$ 的全部支撑（覆盖性过强），导致方差低估；
\item \textbf{EP} 最小化 $\text{KL}(p\|q)$（在每个因子的层面上），即让 $p$ 向 $q$ 靠拢（集中性过强），但通过迭代矩匹配可以更准确地捕获后验的各阶矩。
\end{itemize}
\end{Remark}

\subsection{EP 的稳定性与实践}\label{sec:ep-stability}

EP 的一个关键挑战是\textbf{数值稳定性}：在删除和包含因子的过程中，近似分布的精度可能逐步丧失，尤其当因子数量很大时。实践中常用以下策略：

\begin{itemize}
\item \textbf{ damping（阻尼）}：将新近似因子与旧因子做加权平均，如 $\tilde{m}_f^{\text{new}} = (\tilde{m}_f^{\text{old}})^{1-\beta} \cdot (\text{新因子})^\beta$，其中 $\beta \in (0,1)$ 是阻尼系数；
\item \textbf{投影到参数化族}：将精化后的因子强制投影到预设的参数化族（如多元正态），以保持数值稳定性；
\item \textbf{顺序处理 vs 并行处理}：因子可以按固定顺序处理（sequential），也可以每次随机选择一个因子处理（randomized），后者在某些情况下有助于避免系统偏差。
\end{itemize}

尽管存在这些工程上的挑战，EP 在许多实际应用中表现出色，尤其是在需要精确后验矩的低到中等维度问题中。

\section{本章小结}

本章介绍了在复杂贝叶斯模型中进行近似推断的核心工具：

\begin{enumerate}
\item \textbf{后验众数寻找}：条件最大化（简单稳定）、牛顿法（二次收敛，需良好初始值）、拟牛顿法与共轭梯度法（高维友好）；

\item \textbf{边界回避先验}：当以众数为点估计时，应选择在参数空间边界处为零的先验（如 Gamma$(2, A^2)$、Wishart$(d+3, A\mathbb{I})$），以避免退化的边界估计；

\item \textbf{正态及混合近似}：在众数处用曲率信息构造正态近似；多峰时用正态混合分布；拉普拉斯方法提供了更一般的积分近似框架；

\item \textbf{EM 算法}：将隐变量引入模型，交替 E 步（条件期望）和 M 步（最大化），适合边缘众数的寻找；变体包括 ECM（条件最大化）和 PX-EM（参数展开）；

\item \textbf{变分贝叶斯与 EP}：两种基于条件矩的近似推断方法，通过迭代优化变分参数或 refine 近似因子来逼近真实后验，各有优劣。
\end{enumerate}

这些方法共同构成了贝叶斯计算的工具箱：对于简单模型可直接模拟；对于复杂模型，可先用近似方法快速获得初始估计或完整近似，再视需求决定是否启动昂贵的 MCMC。

\section{习题}\label{sec:exercises-ch13}

\begin{Exercise}{牛顿法的二次收敛}\label{ex:newton-convergence}
设对数后验 $L(\theta) = -\frac{A}{2}\theta^2 + b^\top\theta + c$（二次函数），证明牛顿法一步收敛到真解。进一步证明：若 $L$ 在众数附近三阶导数有界，则牛顿法的误差以平方速率下降（即二次收敛）。
\end{Exercise}

\begin{Exercise}{EM 算法的单调性}\label{ex:em-monotonicity-exercise}
验证 \cref{ex:normal-mean-variance-em} 中 EM 算法的 M 步公式（\cref{eq:mu-new}）。并编程实现该算法，检查对数边缘后验密度是否随迭代单调递增。
\end{Exercise}

\begin{Exercise}{边界回避先验}\label{ex:boundary-prior}
对于分层模型中的组间标准差 $\tau$，比较以下三种先验的行为：
\begin{itemize}
\item 均匀先验 $p(\tau) \propto 1$；
\item Gamma$(1, A)$（指数分布）；
\item Gamma$(2, A^2)$。
\end{itemize}
绘出 $A = 10$ 时三种先验的密度曲线，解释为什么 Gamma$(2, A^2)$ 在 $\tau \to 0$ 时优于指数分布。
\end{Exercise}

\begin{Exercise}{变分贝叶斯的 KL 散度}\label{ex:vb-kl}
证明变分贝叶斯中，坐标上升算法（即交替更新各变分因子）保证 $\text{KL}(q\|p)$ 单调下降。

\textbf{提示}：设当前变分分布为 $q^{(t)}(\theta) = \prod_j q_j^{(t)}(\theta_j)$，更新因子 $q_k$ 至 $q_k^{(t+1)}$（其余因子保持不变）。分两步证明：
\begin{enumerate}
\item 利用 KL 散度的定义展开 $\text{KL}(q^{(t+1)}\|p) - \text{KL}(q^{(t)}\|p)$；
\item 代入 $q_k^{(t+1)}$ 的最优形式 \cref{eq:variational-update}，利用 Jensen 不等式说明每步更新必使 KL 散度下降。
\end{enumerate}
\end{Exercise}

\begin{Exercise}{拉普拉斯近似}\label{ex:laplace-approx}
设 $y_1, \ldots, y_n \sim \text{Poisson}(\lambda)$，先验 $p(\lambda) \propto 1/\lambda$（Jeffreys 先验）。使用拉普拉斯方法近似后验均值 $\mathbb{E}(\lambda|y)$，并与数值积分的结果比较。研究不同 $n$ 值下近似的质量。
\end{Exercise}

\begin{Exercise}{8 学校模型的混合近似}\label{ex:8schools-mixture}
在 8 学校模型中，后验分布关于 $\tau$ 可能有多个峰。编写程序：
\begin{itemize}
\item[(a)] 使用 MCMC 抽样获得 $p(\tau|y)$ 的样本；
\item[(b)] 基于样本识别 $p(\tau|y)$ 的峰数和位置；
\item[(c)] 用正态混合模型拟合样本，估计各成分的权重、均值和方差；
\item[(d)] 比较混合近似与样本直方图。
\end{itemize}
\end{Exercise}

\endinput


% Part IV: Regression Models
\part{回归模型}

% Chapter 14: Introduction to regression models
% Bayesian Data Analysis - Gelman et al.

\chapter{回归模型入门}
\label{ch:reg}

\section{从多参数模型到回归：一个自然的延伸}

在前几章中，我们研究了单参数和多参数模型的推断方法。在单参数模型中，我们学会了如何处理一个未知参数的估计和假设检验；在多参数模型中，我们掌握了处理 nuisance 参数的条件化和边际化技术，并理解了正态模型中均值与方差的联合推断。然而，这些方法在面对现实世界的数据时往往显得不足。

现实中的数据很少只受单一因素的影响。例如，在教育研究中，学生的成绩可能受到家庭背景、学校质量、个人努力程度等多重因素的共同影响；在医学研究中，病人的康复速度可能同时取决于药物剂量、治疗方式、患者年龄等多种协变量。我们需要一个能够同时处理多个预测变量的统一框架。

这正是回归模型所要解决的问题。回归分析是应用统计学的基石，它允许我们研究一个或多个响应变量与一组预测变量之间的关系。从贝叶斯视角看，回归模型本质上是一个多参数推断问题，其中参数向量 $\beta$ 表示各个预测变量的效应大小，$\sigma^2$ 控制噪声的 scale。

历史地看，回归分析的发展经历了几个阶段。Adrien-Marie Legendre 和 Carl Friedrich Gauss 各自独立发展了最小二乘法：Legendre 于 1805 年首先公开发表，而 Gauss 早在 1801 年就用最小二乘法成功预测了谷神星的轨道（这也是他声名鹊起的原因），两人为此曾有过优先权之争。Ronald Fisher 在 20 世纪初将最小二乘法置于坚实的统计理论基础之上，证明了在正态误差假设下，最小二乘估计具有最优性。

在贝叶斯传统中，回归分析长期以来被视为一个计算上的挑战——直到马尔科夫链蒙特卡洛方法的出现。1980 年前后，随着计算方法的进步，贝叶斯回归分析才真正进入实用阶段。然而，当我们从连续响应变量转向二分类响应变量时，解析方法将彻底失效——似然函数不再是 $\beta$ 的二次函数，后验分布不再具有闭式形式，我们必须借助模拟技术。这一转折将在逻辑回归一节详细展开。

本章我们将系统研究贝叶斯回归模型，从最基础的正规线性回归出发，逐步扩展到逻辑回归和稳健回归。在整个过程中，我们将看到共轭先验与非信息先验如何塑造后验推断，以及当预测变量增多时如何选择先验分布。

\section{正规线性回归模型}

\subsection{模型设定}

最简单的回归模型假设响应变量 $y$ 与预测变量 $x_1, \ldots, x_k$ 之间存在线性关系：

\begin{equation}
y_i = x_i^T \beta + \epsilon_i, \quad \epsilon_i \sim \text{Normal}(0, \sigma^2)
\label{eq:linear-regression-model}
\end{equation}

其中 $x_i = (1, x_{i1}, \ldots, x_{ik})^T$ 是第 $i$ 个观测的预测变量向量（包含截距项），$\beta = (\beta_0, \beta_1, \ldots, \beta_k)^T$ 是回归系数向量，$\sigma^2$ 是误差方差。将所有观测堆叠，我们得到矩阵形式：

\begin{equation}
y = X\beta + \epsilon, \quad \epsilon \sim \text{Normal}(0, \sigma^2 I_n)
\label{eq:linear-regression-matrix}
\end{equation}

其中 $y$ 是 $n \times 1$ 响应向量，$X$ 是 $n \times p$ 设计矩阵（$p = k+1$ 包含截距），$\epsilon$ 是 $n \times 1$ 误差向量。

\subsection{正规线性回归的充分统计量}

由多元正态分布理论，给定设计矩阵 $X$，响应向量 $y$ 的分布为：

\begin{equation}
y \sim \text{Normal}(X\beta, \sigma^2 I_n)
\label{eq:y-normal}
\end{equation}

其似然函数为：

\begin{equation}
p(y|\beta, \sigma^2) \propto (\sigma^2)^{-n/2} \exp\left(-\frac{1}{2\sigma^2}\|y - X\beta\|^2\right)
\label{eq:likelihood-normal-regression}
\end{equation}

充分统计量是：
\[
\hat{\beta} = (X^TX)^{-1}X^Ty, \quad \text{RSS} = \|y - X\hat{\beta}\|^2
\]
即最小二乘估计和残差平方和。\footnote{在正态误差假设下，$\hat{\beta}$ 既是最小二乘估计（通过最小化残差平方和得到），也是最大似然估计（通过最大化似然函数得到），两者等价。}

\section{共轭正态线性回归}

\subsection{Normal-Gamma 先验}

对于正规线性回归模型 \cref{eq:linear-regression-model}，一个自然的共轭先验分布可以表示为：

\begin{equation}
\beta | \sigma^2 \sim \text{Normal}(\beta_0, \sigma^2 V_0), \quad \sigma^2 \sim \text{Inverse-Gamma}(\nu_0/2, \nu_0\sigma_0^2/2)
\label{eq:normal-gamma-prior}
\end{equation}

这个先验分布在给定 $\sigma^2$ 的条件下对 $\beta$ 是正态分布，而 $\sigma^2$ 自身服从逆 gamma 分布。参数 $(\beta_0, V_0, \nu_0, \sigma_0^2)$ 允许我们表达先验信息。

\begin{Remark}
  逆 gamma 分布作为 $\sigma^2$ 先验的选择有其历史背景。当与正态似然函数结合时，逆 gamma 先验确保后验分布仍然是闭式可积的。这使得我们可以推导出回归系数的后验分布具有解析形式。
\end{Remark}

\subsection{后验分布}

结合似然 \cref{eq:likelihood-normal-regression} 和先验 \cref{eq:normal-gamma-prior}，后验分布 $p(\beta, \sigma^2 | y)$ 仍然是一个 Normal-Gamma 分布：

\begin{equation}
\beta | \sigma^2, y \sim \text{Normal}(\beta_n, \sigma^2 V_n), \quad \sigma^2 | y \sim \text{Inverse-Gamma}(\nu_n/2, \nu_n\sigma_n^2/2)
\label{eq:normal-gamma-posterior}
\end{equation}

其中后验参数由下式给出：\footnote{推导见附录 \cref{sec:normal-gamma-posterior-derivation}。}

\[
\beta_n = (X^TX + V_0^{-1})^{-1}(X^T X\hat{\beta} + V_0^{-1}\beta_0), \quad V_n = (X^TX + V_0^{-1})^{-1}
\]

\[
\nu_n = \nu_0 + n, \quad \nu_n\sigma_n^2 = \nu_0\sigma_0^2 + \text{RSS} + (\hat{\beta} - \beta_0)^T(X^TX)(V_0 + V_n)(\hat{\beta} - \beta_0)
\]

值得注意的是，当 $V_0 = g(X^TX)^{-1}$ 时（即 Zellner g-先验的特殊情形），后验均值 $\beta_n$ 精确地等于数据的最大似然估计 $\hat{\beta}$ 与先验均值 $\beta_0$ 的线性插值，详见 \cref{eq:g-prior-posterior-mean}。
当先验协方差 $V_0$ 与 $(X^TX)^{-1}$ 成正比时，这个加权平均关系尤为清晰。

此外，$\sigma_n^2$ 表达式中的第三项
\[
(\hat{\beta} - \beta_0)^T(X^TX)(V_0 + V_n)(\hat{\beta} - \beta_0)
\]
衡量了先验均值 $\beta_0$ 与数据最大似然估计 $\hat{\beta}$ 之间的"距离"。当两者一致时，这一项接近零，后验方差主要由残差驱动；当两者冲突较大时，这一项增大，后验分布会表现出更大的不确定性——这就是所谓的\textbf{prior-data conflict}（先验-数据冲突）。

\subsection{后验预测分布}

给定新的预测变量 $x_{\text{new}}$，对应的响应变量的后验预测分布为：

\begin{equation}
\tilde{y} | y \sim \text{Student-t}(\nu_n, x_{\text{new}}^T\beta_n, \sigma_n^2(1 + x_{\text{new}}^TV_nx_{\text{new}}))
\label{eq:posterior-predictive-regression}
\end{equation}

这是一个自由度为 $\nu_n$ 的 Student-t 分布，其位置参数是 $x_{\text{new}}^T\beta_n$，scale 参数包含两部分：$\sigma_n^2$ 反映噪声的不确定性，$x_{\text{new}}^TV_nx_{\text{new}}$ 反映 $\beta$ 的后验不确定性。

\section{非信息先验分布}

在许多实际应用中，我们缺乏可靠的先验信息，希望让数据主导推断。这就引出了非信息先验的选择问题。

\subsection{均匀先验}

一个朴素的非信息先验假设所有参数无差别对待：

\begin{equation}
p(\beta, \sigma^2) \propto \text{常数}
\label{eq:uniform-prior}
\end{equation}

在形式上，这是一个异常（improper）先验——它在参数空间上不是正规的概率密度函数。但当我们将其与似然函数结合时，如果后验分布是正常的，这个先验仍然可以使用。

均匀先验在回归问题中对应于假设 $\beta$ 和 $\log\sigma^2$ 在参数空间上是平权的。

\subsection{Jeffreys' 先验}

Jeffreys' 先验是另一种常用的非信息先验，它具有参数化不变性——无论我们如何重新参数化模型，先验分布都会相应变换以保持一致的后验推断。

对于正规线性回归模型，Jeffreys' 先验为：

\begin{equation}
p(\beta, \sigma^2) \propto |X^TX|^{1/2} (\sigma^2)^{-(p+1)/2}
\label{eq:jeffreys-prior-regression}
\end{equation}

其中 $p$ 是 $\beta$ 的维度。这个先验可以分解为 $p(\beta | \sigma^2) \propto |X^TX|^{1/2}(\sigma^2)^{-p/2}$ 和 $p(\sigma^2) \propto (\sigma^2)^{-1}$。

\begin{Remark}
Jeffreys' 先验与 Fisher 信息矩阵的行列式 $|I(\beta)|^{1/2}$ 成正比。稍微思考一下就会想到：Fisher 信息衡量的是参数空间中局部曲率的大小——数据对该方向的参数变化越敏感，该方向的 Fisher 信息越大。直觉上，曲率越大，我们应当给该方向分配更少的先验权重，从而让数据主导推断。$|I(\theta)|^{1/2}$ 正是这种"局部曲率度量"的数学实现：由于 $I(\theta)$ 是协方差矩阵，其行列式综合反映了所有方向的信息量——信息越大的方向，先验权重越小，反之亦然。这种自动的"倾斜"使得 Jeffreys' 先验在不同参数化下保持一致性。
\end{Remark}

一个更简单的近似是分层先验：

\begin{equation}
p(\beta, \sigma^2) \propto (\sigma^2)^{-1}
\label{eq:scale-invariant-prior}
\end{equation}

这个先验在位置-尺度变换下不变，在实践中往往与 Jeffreys' 先验给出非常接近的结果。

\section{多预测变量情形}

当预测变量数量较多时，先验分布的选择变得更加微妙。特别地，如何平衡数据的拟合与模型的复杂度成为一个核心问题。

\subsection{Zellner 的 g-先验}

Zellner 的 g-先验是一类灵活的信息先验，它将先验信息编码为一个额外的观测：

\begin{equation}
\beta | \sigma^2 \sim \text{Normal}(\beta_0, g\sigma^2(X^TX)^{-1})
\label{eq:zellner-g-prior}
\end{equation}

其中超参数 $g$ 控制先验的强度。当 $g \to \infty$ 时，先验变得越来越模糊，最终退化为非信息先验；当 $g$ 较小时，先验对后验有较强的收缩作用。

g-先验的一个重要性质是后验均值可以在最小二乘估计 $\hat{\beta}$ 和先验均值 $\beta_0$ 之间实现线性插值：\footnote{推导见附录 \cref{sec:g-prior-derivation}。}

\begin{equation}
\mathbb{E}(\beta | y) = \frac{n}{n+g}\hat{\beta} + \frac{g}{n+g}\beta_0
\label{eq:g-prior-posterior-mean}
\end{equation}

这体现了贝叶斯推断的核心特征——将先验信息与数据信息以加权平均的方式融合。

\subsection{相关性预测变量的处理}

当预测变量之间存在多重共线性时，$(X^TX)^{-1}$ 可能数值上不稳定。一种解决方案是在 $X^TX$ 的对角线上添加一个小的常数 $\lambda$（岭回归的思想）：

\begin{equation}
\beta | \sigma^2 \sim \text{Normal}(\beta_0, \sigma^2(X^TX + \lambda I)^{-1})
\label{eq:ridge-prior}
\end{equation}

另一种思路是从变量选择的角度出发，使用 spike-and-slab 先验或 horseshoe 先验，将大多数回归系数推向零，同时允许少数重要变量保留非零系数。horseshoe 先验将每个系数 $\beta_j$ 表示为分层结构 $\beta_j = \lambda_j \tau \cdot \xi_j$，其中 $\lambda_j \sim \text{Cauchy}^+(0,1)$（局部正则化参数，控制单个系数的收缩程度），$\tau \sim \text{Cauchy}^+(0,1)$（全局正则化参数，控制整体的收缩强度），$\xi_j \sim \text{Normal}(0,1)$。这种结构使得大多数 $\lambda_j$ 接近零（从而 $\beta_j$ 被推向零），但允许少数 $\lambda_j$ 较大（从而对应的 $\beta_j$ 保留非零值），实现了数据自适应的稀疏估计。

\section{逻辑回归}

\subsection{二分类响应变量}

当响应变量是二分类的（$y_i \in \{0, 1\}$）时，正规线性回归不再适用，因为线性模型可能给出超出 $[0,1]$ 区间的概率预测。逻辑回归通过引入 link 函数解决了这个问题。

模型设定为：

\begin{equation}
y_i \mid \beta \sim \text{Bernoulli}(\pi_i), \quad \text{logit}(\pi_i) = x_i^T\beta
\label{eq:logistic-regression-model}
\end{equation}

即：

\begin{equation}
\pi_i = \Pr(y_i = 1 | \beta) = \frac{1}{1 + \exp(-x_i^T\beta)}
\label{eq:logistic-link}
\end{equation}

对数似然函数为：

\begin{equation}
\ell(\beta) = \sum_{i=1}^n [y_i x_i^T\beta - \log(1 + \exp(x_i^T\beta))]
\label{eq:logistic-likelihood}
\end{equation}

\subsection{后验分布与计算}

逻辑回归的后验分布 $p(\beta | y) \propto \exp(\ell(\beta)) \times p(\beta)$ 不再有解析形式。似然函数关于 $\beta$ 不是二次的，因此无法通过解析积分得到边际后验分布。

这标志着一个重要的计算转折点：我们必须借助第 10-12 章介绍的模拟方法来处理这类非共轭模型。常用的策略及其适用场景如下：

\begin{enumerate}
  \item \textbf{Metropolis 算法}：使用对称提议分布（如正态提议）在后验分布上生成马尔科夫链。\textbf{适用场景}：低维参数空间（$p < 20$）；优点是简单通用，缺点是收敛较慢。
  \item \textbf{Hamiltonian 蒙特卡洛}：利用梯度信息 $\nabla_\beta \log p(\beta|y)$ 构造更高效的提议机制，特别适合高维参数空间。\textbf{适用场景}：高维参数空间（$p \geq 20$）；优点是效率远高于随机游走 Metropolis，缺点是需要计算梯度且调参较敏感。
  \item \textbf{拉普拉斯近似}：用正态分布 $\text{Normal}(\hat{\beta}, \Sigma)$ 逼近后验模式 $\hat{\beta}$ 附近的后验密度，结合重要性抽样或粒子滤波器。\textbf{适用场景}：快速原型探索；优点是计算极快，缺点是仅在模式附近精确，对多峰或高度非正态的后验分布效果较差。
\end{enumerate}

在实际应用中，通常先尝试拉普拉斯近似获得初始值，再用 HMC 或 Metropolis 进行精确采样。当参数维度特别高时（如 $p > 1000$），还需结合稀疏先验（如 horseshoe 先验）使用变分推断等近似方法。

\section{稳健回归}

正规线性回归假设误差服从正态分布，这一假设在某些实际数据中可能不成立。当数据中存在离群点或重尾噪声时，正态假设可能导致推断过度受到异常值的影响。

\subsection{t 分布误差}

一种自然的稳健化方案是假设误差服从 Student-t 分布而不是正态分布：

\begin{equation}
y_i = x_i^T\beta + \epsilon_i, \quad \epsilon_i \sim \text{Student-t}(\nu, 0, \sigma^2)
\label{eq:t-distribution-errors}
\end{equation}

自由度 $\nu$ 控制尾部的厚度。当 $\nu \to \infty$ 时，t 分布收敛到正态分布；当 $\nu$ 较小时（如 $\nu = 4$），分布的尾部更厚，对离群点的敏感度更低。

从层次化的视角看，Student-t 分布可以分解为一个正态分布与一个隐变量的混合：

\[
\epsilon_i | \lambda_i \sim \text{Normal}(0, \sigma^2/\lambda_i), \quad \lambda_i \sim \text{Gamma}(\nu/2, \nu/2)
\]

其中 $\lambda_i$ 是第 $i$ 个观测的精度尺度因子。这种分解使得我们可以使用 EM 算法或 Gibbs 抽样来计算后验分布。

\subsection{分层混合模型}

另一种稳健回归的框架是使用有限混合分布作为误差模型。例如，假设误差以概率 $(1-\pi)$ 来自正态分布，以概率 $\pi$ 来自某个污染分布（通常是尺度更大的正态或均匀分布）：

\begin{equation}
\epsilon_i \sim (1-\pi)\text{Normal}(0, \sigma_1^2) + \pi\text{Normal}(0, \sigma_2^2), \quad \sigma_2^2 > \sigma_1^2
\label{eq:mixture-robust}
\end{equation}

这类模型允许数据自动识别并降低离群点的影响权重，同时保持对大部分"正常"数据的有效推断。

\section{本章小结}

本章我们从贝叶斯视角全面介绍了回归模型：

\begin{enumerate}
  \item \textbf{正规线性回归}：建立了从多参数正态模型到回归分析的自然过渡，介绍了 Normal-Gamma 共轭先验和后验分布的闭式形式。
  \item \textbf{非信息先验}：讨论了均匀先验和 Jeffreys' 先验的选择，以及它们在缺乏先验信息时的应用。
  \item \textbf{多预测变量}：介绍了 Zellner 的 g-先验作为信息先验与无信息先验之间的桥梁，以及处理高度相关预测变量的岭回归策略。
  \item \textbf{逻辑回归}：从二分类响应变量出发，引入了 link 函数和后验分布的不可解析性，为后续章节的模拟计算方法埋下伏笔。
  \item \textbf{稳健回归}：通过 t 分布误差和有限混合误差模型展示了如何处理偏离正态假设的数据。
\end{enumerate}

回归分析将我们在前几章学习的参数推断技术扩展到了更复杂的实际场景。同时，本章也展示了贝叶斯方法在处理不同数据类型（连续、二分类、存在离群点的数据）时的灵活性。

\section{习题}\label{sec:ch14-exercises}

\begin{Exercise}{\ref{exr:ch14-1} 正规线性回归的后验均值}
\label{exr:ch14-1}
证明在使用共轭 Normal-Gamma 先验 \cref{eq:normal-gamma-prior} 时，后验均值满足 \cref{eq:g-prior-posterior-mean} 的特例形式（即 $\beta_0 = 0$ 且 $V_0 = (X^TX)^{-1}$ 的情形）。
\end{Exercise}

\begin{Exercise}{\ref{exr:ch14-2} Jeffreys' 先验的性质}
\label{exr:ch14-2}
验证 Jeffreys' 先验 \cref{eq:jeffreys-prior-regression} 可以分解为 $p(\beta | \sigma^2) \propto (\sigma^2)^{-p/2}$ 和 $p(\sigma^2) \propto (\sigma^2)^{-1}$ 两部分的乘积。
\end{Exercise}

\begin{Exercise}{\ref{exr:ch14-3} 逻辑回归的 Metropolis 算法实现}
\label{exr:ch14-3}
给定先验 $p(\beta) \propto 1$（平坦先验）和数据集：$(x_i, y_i)$ 中 $y_i \in \{0,1\}$，$n=50$。描述使用 Metropolis 算法从后验分布 $p(\beta | y)$ 中抽取样本的完整步骤，包括：
\begin{enumerate}
  \item 提议分布的选择（例如 $\text{Normal}(\beta^{\text{cur}}, \Sigma)$），如何确定提议协方差 $\Sigma$？
  \item 接受率 $\alpha = \min\{1, p(\beta^{\text{prop}}|y)/p(\beta^{\text{cur}}|y)\}$ 的计算，其中 $p(\beta|y) \propto \exp(\ell(\beta))$，$\ell(\beta)$ 如 \cref{eq:logistic-likelihood}。
  \item 如何丢弃 burn-in 样本并验证链的收敛性。
  \item 给定新预测变量 $x_{\text{new}}$，写出 $\Pr(\tilde{y}=1|y)$ 的蒙特卡洛估计公式，并说明如何用收集的样本计算该估计。
\end{enumerate}
\end{Exercise}

\begin{Exercise}{\ref{exr:ch14-4} t 分布误差的层次表示}
\label{exr:ch14-4}
证明在层次表示 $\epsilon_i | \lambda_i \sim \text{Normal}(0, \sigma^2/\lambda_i)$，$\lambda_i \sim \text{Gamma}(\nu/2, \nu/2)$ 下，$\epsilon_i$ 的边际分布是 Student-t 分布 \cref{eq:t-distribution-errors}。
\end{Exercise}

\begin{Exercise}{\ref{exr:ch14-5} g-先验的收缩性质}
\label{exr:ch14-5}
令 $g$ 为 Zellner g-先验 \cref{eq:zellner-g-prior} 中的超参数。证明后验均值 $\mathbb{E}(\beta | y)$ 可以写成 \cref{eq:g-prior-posterior-mean} 的形式，并说明当 $g$ 增大时，后验均值如何收缩向先验均值。
\end{Exercise}

\begin{Exercise}{\ref{exr:ch14-6} 稳健回归与离群点}
\label{exr:ch14-6}
对比使用正态误差假设与 t 分布误差假设的线性回归在后验预测分布上的差异。通过直觉和数学两方面说明 t 分布误差为何对离群点更加稳健。
\end{Exercise}

\endinput

% Chapter 15: Hierarchical linear models
% Bayesian Data Analysis - Gelman et al.

\chapter{分层线性模型}
\label{ch:hlm}

\section{本章导论}
\label{sec:ch15-intro}

% Stein 风格开场：历史问题引入
在 1970 年代的美国教育研究中，一个看似简单的问题引发了长达数十年的方法论争论：如何科学地评估不同学校的"教育质量"？研究者们手中有 $J$ 所学校的标准化考试成绩，但每所学校的样本量都很小，直接用学校平均成绩排名会面临严重的抽样误差。有人主张将所有学校视为同一总体，取平均值；有人坚持每所学校独立估计，互不干扰。

这场争论的实质是统计学中一个永恒的张力：\textbf{在合并信息与保持特异性之间，我们应当如何取舍？}

这个问题引导我们发展出了一套完整的统计框架——分层线性模型（Hierarchical Linear Models, HLM）。从历史上看，分层模型的思想萌芽于 1950 年代 Sheps (1959) 的医学研究，在 1970-80 年代由 Efron、Morris 以及教育统计学家 Bryk, Raudenbush 等人系统化为一套完整的理论，并在贝叶斯框架下获得了最自然、最优雅的表达。

本章我们将学习如何将分层思想与线性回归结合，建立分层线性模型；如何用贝叶斯方法进行推断；以及如何从矩阵形式理解模型的本质。

\section{从分组数据到分层模型}
\label{sec:grouped-to-hierarchical}

\subsection{传统方法的困境}
\label{sec:traditional-dilemma}

考虑一个经典问题：估计 $J$ 所学校的培训效果。每所学校 $j$ 报告了一个效果估计 $y_j$ 及其标准误 $\sigma_j$。一个 naive 的做法是简单地将所有 $y_j$ 取平均作为总体效果估计。但稍微思考一下就会发现，这种做法存在根本性缺陷：它假设所有学校的真实效果 $\theta_j$ 完全相同，即 $\theta_1 = \theta_2 = \cdots = \theta_J$。这在现实中几乎不可能成立。

另一方面，如果我们分别为每所学校估计效果，就完全忽略了数据可以"相互学习"这一事实。例如，学校 3 的估计精度很低（$\sigma_3$ 很大），但其他学校的数据可以帮助我们改进对它的估计——如果我们知道学校间的变异不大的话。

这就是分层模型要解决的核心问题：如何在保持组特异性估计的同时，充分利用数据的层次结构进行信息共享。

\subsection{两层次模型框架}
\label{sec:two-level-framework}

分层模型的核心思想是将 $\theta_j$ 本身视为从一个公共分布中抽取的样本。一般的两层次线性模型可以写为：

\begin{enumerate}
  \item \textbf{第一层（组内模型）}：$y_{ij} = \mathbf{x}_{ij}^\top \boldsymbol{\beta}_j + \epsilon_{ij}$，其中 $\epsilon_{ij} \sim \mathrm{N}(0, \sigma^2)$
  \item \textbf{第二层（组间模型）}：$\boldsymbol{\beta}_j = \mathbf{A}_j \boldsymbol{\gamma} + \mathbf{u}_j$，其中 $\mathbf{u}_j \sim \mathrm{N}(\mathbf{0}, \boldsymbol{\Sigma})$
\end{enumerate}

这里 $y_{ij}$ 是第 $j$ 组第 $i$ 个观测值，$\mathbf{x}_{ij}$ 是协变量向量，$\boldsymbol{\beta}_j$ 是第 $j$ 组的效应向量，$\boldsymbol{\gamma}$ 是总体（总体层面）效应，$\mathbf{A}_j$ 是已知的设计矩阵，$\boldsymbol{\Sigma}$ 是组间协方差矩阵。

将两层结合，我们可以写出完整的分层模型：

\[
y_{ij} = \mathbf{x}_{ij}^\top \mathbf{A}_j \boldsymbol{\gamma} + \mathbf{x}_{ij}^\top \mathbf{u}_j + \epsilon_{ij}
\]

其中 $\mathbf{u}_j$ 和 $\epsilon_{ij}$ 相互独立。

\section{分层线性模型的矩阵形式}
\label{sec:matrix-formulation}

\subsection{模型的一般表达}
\label{sec:general-matrix-form}

将所有数据堆叠为向量形式，分层线性模型可以紧凑地写为：

\[
\mathbf{y} = \mathbf{X}\boldsymbol{\gamma} + \mathbf{Z}\mathbf{u} + \boldsymbol{\epsilon}
\]

其中：
\begin{itemize}
  \item $\mathbf{y}$ 是 $n \times 1$ 观测向量
  \item $\mathbf{X}$ 是 $n \times p$ 固定效应设计矩阵
  \item $\boldsymbol{\gamma}$ 是 $p \times 1$ 固定效应参数
  \item $\mathbf{Z}$ 是 $n \times q$ 随机效应设计矩阵
  \item $\mathbf{u}$ 是 $q \times 1$ 随机效应向量，$\mathbf{u} \sim \mathrm{N}(\mathbf{0}, \boldsymbol{\Sigma})$
  \item $\boldsymbol{\epsilon}$ 是 $n \times 1$ 误差向量，$\boldsymbol{\epsilon} \sim \mathrm{N}(\mathbf{0}, \sigma^2 \mathbf{I}_n)$
\end{itemize}

随机效应 $\mathbf{u}$ 和误差 $\boldsymbol{\epsilon}$ 相互独立，因此 $\mathbf{y}$ 的边际分布为：

\[
\mathbf{y} \sim \mathrm{N}(\mathbf{X}\boldsymbol{\gamma}, \mathbf{V}), \quad \mathbf{V} = \sigma^2 \mathbf{I}_n + \mathbf{Z}\boldsymbol{\Sigma}\mathbf{Z}^\top
\]

\begin{Remark}
边际分布 $\mathbf{y} \sim \mathrm{N}(\mathbf{X}\boldsymbol{\gamma}, \mathbf{V})$ 正是我们在第八章学到的贝叶斯线性回归形式，只不过这里的协方差矩阵 $\mathbf{V}$ 依赖于方差参数 $\sigma^2$ 和 $\boldsymbol{\Sigma}$。
\end{Remark}

\subsection{组结构的矩阵表示}
\label{sec:group-structure}

考虑 $J$ 组、每组 $n_j$ 观测的经典情形。设 $\mathbf{y}_j$ 是第 $j$ 组的 $n_j \times 1$ 观测向量，$\mathbf{X}_j$ 是对应的设计矩阵。模型可写为：

\[
\mathbf{y}_j = \mathbf{X}_j \boldsymbol{\gamma} + \mathbf{Z}_j \mathbf{u}_j + \boldsymbol{\epsilon}_j, \quad \boldsymbol{\epsilon}_j \sim \mathrm{N}(\mathbf{0}, \sigma^2 \mathbf{I}_{n_j})
\]

其中 $\mathbf{u}_j \sim \mathrm{N}(\mathbf{0}, \boldsymbol{\Sigma}_\mathbf{u})$ 是组 $j$ 的随机效应。

整个数据的似然函数为：

\[
p(\mathbf{y} \mid \boldsymbol{\gamma}, \sigma^2, \boldsymbol{\Sigma}) = \prod_{j=1}^J \mathrm{N}\left(\mathbf{y}_j \mid \mathbf{X}_j\boldsymbol{\gamma}, \sigma^2 \mathbf{I}_{n_j} + \mathbf{Z}_j \boldsymbol{\Sigma} \mathbf{Z}_j^\top\right)
\]

\section{后验推断}
\label{sec:posterior-inference-hlm}

\subsection{完全条件分布}
\label{sec:full-conditional}

在贝叶斯分层模型中，我们通常关注以下后验分布：

\[
p(\boldsymbol{\gamma}, \mathbf{u}, \sigma^2, \boldsymbol{\Sigma} \mid \mathbf{y})
\]

为方便 MCMC 抽样，我们使用 Data augmentation 技术，引入隐变量 $\mathbf{u}$。完全数据模型为：

\[
\mathbf{y}_j \mid \mathbf{u}_j, \boldsymbol{\gamma}, \sigma^2 \sim \mathrm{N}\left(\mathbf{X}_j\boldsymbol{\gamma} + \mathbf{Z}_j\mathbf{u}_j, \sigma^2 \mathbf{I}_{n_j}\right)
\]
\[
\mathbf{u}_j \mid \boldsymbol{\Sigma} \sim \mathrm{N}(\mathbf{0}, \boldsymbol{\Sigma})
\]

完全条件分布为：

\[
\boldsymbol{\gamma} \mid \mathbf{y}, \mathbf{u}, \sigma^2 \sim \mathrm{N}(\boldsymbol{\mu}_\gamma, \boldsymbol{\Sigma}_\gamma)
\]
\[
\mathbf{u}_j \mid \mathbf{y}, \boldsymbol{\gamma}, \sigma^2, \boldsymbol{\Sigma} \sim \mathrm{N}(\boldsymbol{\mu}_{u_j}, \boldsymbol{\Sigma}_{u_j})
\]
\[
\sigma^2 \mid \mathbf{y}, \boldsymbol{\gamma}, \mathbf{u}, \boldsymbol{\Sigma} \sim \text{Inv-}\chi^2(\nu_n, \sigma_n^2)
\]

其中后验均值和方差依赖于数据的分组结构：当组内样本量大时，后验均值更依赖组内数据；当组内样本量小时，后验均值更依赖先验（总体效应）。\footnote{完全条件分布的详细推导见附录 \cref{sec:appendix-hlm-conditional}。}

\subsection{分层先验的构建}
\label{sec:hierarchical-prior}

对超参数 $\boldsymbol{\Sigma}$ 的先验选择是分层模型的关键。几种常见选择：

\begin{enumerate}
  \item \textbf{逆 Wishart 先验}：$\boldsymbol{\Sigma} \sim \text{Inv-Wishart}(\nu_0, \boldsymbol{\Sigma}_0)$
  \item \textbf{独立逆 Gamma 先验}：对方差分量独立建模
  \item \textbf{弱信息先验}：使用半参数或数据驱动的先验
\end{enumerate}

当方差分量接近零时，逆 Wishart 先验可能导致后验分布出现拖尾问题。因此，实际中常使用对数尺度上的先验，例如：

\[
\log(\sigma_j^2) \sim \mathrm{N}(\mu_0, \tau_0^2)
\]

\section{预测}
\label{sec:prediction-hlm}

分层模型的一个重要应用是预测新组或新观测的值。

\subsection{组均值预测}
\label{sec:group-mean-prediction}

给定新组 $j*$ 的协变量 $\mathbf{x}_{j*}$，我们希望预测其效应 $\boldsymbol{\beta}_{j*} = \mathbf{A}_{j*}\boldsymbol{\gamma} + \mathbf{u}_{j*}$。

后验预测分布为：

\[
\boldsymbol{\beta}_{j*} \mid \mathbf{y} \sim \mathrm{N}(\boldsymbol{\mu}_{\beta_{j*}}, \boldsymbol{\Sigma}_{\beta_{j*}})
\]

其中后验均值是总体效应 $\boldsymbol{\gamma}$ 的后验均值和组特异偏差 $\mathbf{u}_j$ 的后验均值的加权组合。当新组与已有组相似时，预测更依赖总体效应；当新组特征独特时，预测更依赖组内信息。\footnote{组均值预测的详细推导见附录 \cref{sec:appendix-hlm-prediction}。}

\subsection{新观测预测}
\label{sec:new-observation-prediction}

对于新组 $j*$ 的新观测 $y_{new}$，预测分布为：

\[
y_{new} \mid \mathbf{y} \sim \mathrm{N}\left(\mathbf{x}_{new}^\top \boldsymbol{\mu}_{\beta_{j*}}, \sigma^2 + \mathbf{x}_{new}^\top \boldsymbol{\Sigma}_{\beta_{j*}} \mathbf{x}_{new}\right)
\]

注意预测方差包括两部分：观测误差方差 $\sigma^2$ 和参数不确定性方差 $\mathbf{x}_{new}^\top \boldsymbol{\Sigma}_{\beta_{j*}} \mathbf{x}_{new}$。

\section{方差分量估计}
\label{sec:variance-components}

\subsection{方差分量模型}
\label{sec:variance-component-model}

最简单的分层线性模型是所谓的\textbf{方差分量模型}（variance components model）：

\[
y_{ij} = \mu + u_j + \epsilon_{ij}
\]

其中 $u_j \sim \mathrm{N}(0, \sigma_u^2)$ 是组效应，$\epsilon_{ij} \sim \mathrm{N}(0, \sigma_\epsilon^2)$ 是个体误差。组内相关系数（ICC）为：

\[
\mathrm{ICC} = \frac{\sigma_u^2}{\sigma_u^2 + \sigma_\epsilon^2}
\]

ICC 衡量了总变异中有多大比例来自组间变异。

\subsection{频率派与贝叶斯视角的差异}
\label{sec:frequentist-bayesian-variance}

对随机效应 $\mathbf{u}$ 进行边缘化后，似然函数涉及高维积分。使用 EM 算法可以方便地估计方差分量（频率派方法）。

E 步：计算随机效应的后验均值和协方差
\[
\widehat{\mathbf{u}}_j^{(t)} = \boldsymbol{\Sigma}_j \mathbf{Z}_j^\top \mathbf{V}_j^{-1} (\mathbf{y}_j - \mathbf{X}_j \widehat{\boldsymbol{\gamma}}^{(t)})
\]

M 步：更新参数估计
\[
\widehat{\boldsymbol{\gamma}}^{(t+1)} = \left(\sum_j \mathbf{X}_j^\top \mathbf{V}_j^{-1} \mathbf{X}_j \right)^{-1} \sum_j \mathbf{X}_j^\top \mathbf{V}_j^{-1} (\mathbf{y}_j - \mathbf{Z}_j \widehat{\mathbf{u}}_j^{(t)})
\]

方差分量的更新涉及更复杂的估计方程。\footnote{EM 算法对方差分量的更新公式见附录 \cref{sec:appendix-em-variance}。}

\begin{Remark}
在贝叶斯视角下，方差分量估计的"偏差"问题自然消失：我们对 $\sigma_u^2$ 指定先验，后验分布的平均值即为估计，无需像频率派那样讨论有偏/无偏问题。这是贝叶斯方法在分层模型中的一个重要优势。
\end{Remark}

\section{相关作业干预示例}
\label{sec:example-coaching}

考虑一个经典的教育研究问题：估计 $J = 8$ 所学校的培训效果。

\subsection{问题背景}
\label{sec:coaching-background}

每所学校报告了培训效果估计 $y_j$ 和标准误 $\sigma_j$。数据如下：

\begin{center}
\begin{tabular}{ccc}
\hline
学校 $j$ & 估计 $y_j$ & 标准误 $\sigma_j$ \\
\hline
1 & 28.4 & 14.9 \\
2 & 7.0 & 10.2 \\
3 & -4.4 & 16.3 \\
4 & 7.0 & 11.0 \\
5 & 1.9 & 9.4 \\
6 & 6.6 & 10.5 \\
7 & 10.8 & 17.6 \\
8 & 33.1 & 15.6 \\
\hline
\end{tabular}
\end{center}

观察数据，一个自然的问题是：学校 1（28.4）和学校 8（33.1）的估计远高于其他学校，这是真实差异还是抽样误差？如果我们假设所有学校效果相同（固定效应模型），学校 1 和 8 可能被过度拉伸；如果我们假设每所学校完全独立，学校 3（-4.4）这样的大负值可能只是噪声。

\subsection{模型设定}
\label{sec:coaching-model}

我们设定分层模型：

\[
y_j \mid \theta_j \sim \mathrm{N}(\theta_j, \sigma_j^2)
\]
\[
\theta_j \mid \mu, \tau \sim \mathrm{N}(\mu, \tau^2)
\]
\[
\mu \sim \mathrm{N}(0, 100^2), \quad \tau \sim \mathrm{Half-Cauchy}(0, 25)
\]

其中 $\theta_j$ 是学校 $j$ 的真实效果，$\mu$ 是总体效果，$\tau$ 是学校间标准差。

\subsection{后验推断}
\label{sec:coaching-posterior}

使用 MCMC 方法，我们可以得到 $\theta_j$ 的后验分布。关键发现是：

\begin{Remark}
后验均值是先验均值（0）和数据均值 $y_j$ 的加权平均，权重由精度（方差的倒数）给出。当 $\tau$ 小时，所有 $\theta_j$ 被强烈收缩到公共均值 $\mu$；当 $\tau$ 大时，每个 $\theta_j$ 主要由其自身数据决定。
\end{Remark}

这个例子展示了分层模型的核心优势：它提供了在\textbf{合并信息}和\textbf{保持特异性}之间的自然权衡。

\section{习题}
\label{sec:chapter15-exercises}

\begin{Exercise}{\ref{exr:15-1} 分层模型的收缩性质}
\label{exr:15-1}
考虑简化的方差分量模型 $y_j \sim \mathrm{N}(\theta_j, \sigma_j^2)$，$\theta_j \sim \mathrm{N}(\mu, \tau^2)$，其中 $\sigma_j^2$ 已知。

\begin{enumerate}
  \item 证明 $\theta_j$ 的后验均值为 $\widehat{\theta}_j = \frac{\frac{\mu}{\tau^2} + \frac{y_j}{\sigma_j^2}}{\frac{1}{\tau^2} + \frac{1}{\sigma_j^2}}$。
  \item 解释收缩因子 $\frac{1/\sigma_j^2}{1/\tau^2 + 1/\sigma_j^2}$ 的含义：当 $\tau \to 0$ 和 $\tau \to \infty$ 时分别发生什么？
  \item 如果有两所学校，A 学校的 $\sigma_A = 10$，B 学校的 $\sigma_B = 1$，在 $\tau = 5$ 的情况下，哪所学校会被收缩得更厉害？为什么？
\end{enumerate}
\end{Exercise}

\section{稳健推断}
\label{sec:robust-inference}

当数据存在异常值或偏离正态假设时，标准分层线性模型可能失效。

\subsection{稳健误差分布}
\label{sec:robust-error}

一种方法是使用 Student-$t$ 分布替代正态分布作为误差分布：

\[
\epsilon_{ij} \sim \text{Student-}t(\nu, 0, \sigma^2)
\]

当 $\nu \to \infty$ 时，Student-$t$ 分布趋近于正态分布；当 $\nu$ 较小时，分布的尾部更重，对异常值更稳健。

另一种方法是使用污染正态模型：

\[
\epsilon_{ij} \sim (1 - \epsilon) \mathrm{N}(0, \sigma^2) + \epsilon \mathrm{N}(0, k\sigma^2)
\]

其中 $\epsilon$ 是污染比例，$k$ 是方差放大因子。

\subsection{分层先验的稳健性}
\label{sec:robust-prior}

对方差分量使用弱信息先验可以提高稳健性。Half-t 先验和 Horseshoe 先验是常用的选择。

\begin{Remark}
选择先验时，应该考虑参数的尺度。例如，在 $\log(\tau)$ 上指定先验通常比在 $\tau$ 上更合理，因为前者保证 $\tau$ 始终为正。
\end{Remark}

\section{本章小结}
\label{sec:ch15-summary}

本章我们学习了分层线性模型的贝叶斯分析框架：

\begin{enumerate}
  \item \textbf{模型构建}：通过两层结构同时捕捉组内变异和组间变异
  \item \textbf{矩阵形式}：将分层模型表达为 $\mathbf{y} = \mathbf{X}\boldsymbol{\gamma} + \mathbf{Z}\mathbf{u} + \boldsymbol{\epsilon}$
  \item \textbf{后验推断}：利用完全条件分布进行 MCMC 抽样
  \item \textbf{预测}：提供组均值和新观测的预测分布
  \item \textbf{方差分量}：通过 ICC 等指标量化组间变异的重要性，理解贝叶斯与频率派视角的差异
  \item \textbf{稳健推断}：使用 Student-$t$ 误差或污染正态模型处理异常值
\end{enumerate}

分层线性模型是现代贝叶斯统计中最重要的工具之一，它的核心思想——在部分合并和完全分离之间寻求平衡——贯穿于整个统计学。

\section{下节预告}
\label{sec:ch15-next}

下一章我们将学习广义线性模型（Generalized Linear Models），这是将分层框架推广到非正态数据的关键步骤。

\endinput

% Chapter 16: [Title]
% Bayesian Data Analysis - Gelman et al.

\chapter{广义线性模型}

\section{本章导论}

\section{主要内容}

\section{本章小结}

\endinput

% Chapter 17: Models for robust inference
% Bayesian Data Analysis - Gelman et al.

\chapter{稳健推断模型}

\section{本章导论}

在前面的章节中，我们主要依赖正态分布、二项分布和泊松分布，以及这些分布的层次组合来建模数据和参数。然而，使用有限的分布类别会导致有限的、甚至不恰当的推断。许多问题超出了便捷模型的适用范围，模型的选择应该拟合底层科学和数据，而不是仅仅为了分析或计算上的方便。正如第 5 章所说明的，创建现实模型最有用的方法通常是层次化地工作，将简单的单变量模型组合起来。

如果我们为了方便而使用过于简化的模型，就必须回答一个关键问题：后验推断在何种程度上依赖于极端数据点和不可评估的模型假设？我们在第 6 章已经讨论过这个问题的后半部分，这本质上是敏感性分析的主题；现在我们用更先进的计算方法来更深入地探讨这个话题。

\textbf{历史脉络}。贝叶斯稳健推断的现代发展可以追溯到 Box 和 Tiao 的开创性工作（Box and Tiao, 1968, 1973），他们系统地研究了正态模型中异常值的贝叶斯处理方法。这一传统后来被 Berger 和同事们进一步发展，探索了贝叶斯稳健性的理论方面，包括最大稳健性先验族等。

\section{稳健性的各个方面}

\subsection{对异常值的稳健性}

基于正态分布的模型对异常值是出了名的``非稳健''——在某种意义上，一个异常数据点可以强烈影响模型中所有参数的推断，即使那些参数与该异常观测没有实质性关联。

例如，在第 5.5 节的学业测试例子中，我们对八个处理效应的估计是通过将各学校的样本均值向总体均值收缩（即向真实效应来自共同正态分布的先验信息收缩）而得到的，每所学校 $j$ 的收缩比例仅由其抽样误差 $\sigma_j$ 和学校间效应变异 $\tau$ 决定。

假设研究中第八所学校的观测值 $y_8$ （表 5.2 中报告的）本应是 100 而不是 12，那么八个观测值变为 28、8、−3、7、−1、1、18 和 100，标准误与表 5.2 中报告的相同。如果我们对这一数据集应用层次正态模型，后验分布会告诉我们 $\tau$ 有一个很高的值，因此每个估计值 $\hat{\theta}_j$ 将基本上等于其观测效应 $y_j$。

但这在实践中合理吗？毕竟，给定这些假设观测值，第八所学校似乎有一个极其有效的 SAT 辅导项目，或者 100 可能只是数据记录错误的结果。在这两种情况下，让单个观测值 $y_8$ 对我们估计 $\theta_1, \ldots, \theta_7$ 的方式产生如此强烈的影响，似乎是不合适的。

在贝叶斯框架中，我们可以通过将 $\theta_j$ 的正态总体模型替换为更长尾的分布族来减少异常第八观测值的影响。这里所说的长尾分布是指在远离中心的位置有相对较高的概率含量，其中``远离''的尺度是相对于包含分布 50\% 概率区域的直径来确定的。长尾分布的例子包括 $t$ 分布族（其中最极端的情况是柯西分布或 $t_1$），以及（有限）混合模型——后者通常使用简单分布（如正态分布）来拟合大部分值，但允许来自替代分布的观测或参数值有离散的将概率，该替代分布可以有不同的中心，通常有更大的散布。

在这个假设的 SAT 辅导例子修改中，使用长尾分布对 $\theta_j$ 进行分析会导致观测值 100 被解释为来自长尾分布的极端抽取，而不是正态分布效应具有高方差的证据。结果是，分析会将第八个观测值在一定程度上向其他观测值收缩，但相对于其与总体均值的距离而言，收缩程度远不如前七个相互之间的收缩程度那么大。

正如我们的假设例子所示，我们不必为了处理异常值而放弃贝叶斯原理。例如，像柯西分布甚至双组分混合模型这样的长尾模型仍然是 $(\theta_1, \ldots, \theta_8)$ 的可交换先验分布——这在没有什么先验信息区分八所学校时是恰当的。可交换先验模型的选择影响 $\theta_j$ 估计值的收缩方式，我们因此可以在不完全改变分析方式的情况下减少异常观测值的影响。（这不应该取代对数据的仔细检查和对异常值中可能存在的记录错误的查找。）有时会区分搜索异常值的方法（可能将其从分析中移除）和使推断不受异常值影响的稳健方法。在贝叶斯框架中，这两种方法不应该被区分对待。例如，使用混合模型（第 22 章中的有限混合模型或标准模型的过度分散版本）不仅能将极端观测分类为来自高方差混合成分（而不是简单地将其视为令人惊讶的``异常值''），还意味着这些点对总体均值和中位数等估计量的推断影响较小。

\subsection{敏感性分析}

除了补偿异常值外，稳健模型还可用于评估后验推断对模型假设的敏感性。例如，可以使用 $t$ 分布替代正态分布来进行敏感性分析，以评估对正态假设的敏感性，方法是将自由度从大变小。如第 6 章所讨论的，敏感性分析的基本思想是尝试各种不同的分布（对于似然和先验模型），并观察感兴趣估计量和预测量的后验推断如何变化。一旦已经从后验分布中抽取了样本，使用重要性重采样可以相当直接地从替代模型中抽取样本，精度足以检测模型间推断的主要差异。如果感兴趣估计量的后验分布对模型假设高度敏感，则可能需要迭代模拟方法来获得更精确的计算。

从某种意义上说，第 5.5 节中 SAT 辅导实验的大部分分析，特别是图 5.6 和 5.7，是一种敏感性分析，其中参数 $\tau$ 被允许从 0 变化到 $\infty$。如第 5.5 节所讨论的，观测数据实际上与所有相等效应模型（即 $\tau = 0$）一致，但该模型没有任何实质性意义，所以我们拟合允许 $\tau$ 为任意正值的模型。结果总结在 $\tau$ 的边际后验分布中（如图 5.5 所示），描述了数据支持的 $\tau$ 值范围。

\section{标准概率模型的过度分散版本}

有时使用标准模型之一——二项、正态、泊松、指数——看起来很自然，但数据的分散程度过大。例如，正态分布不应该用于拟合大样本中 10\% 的点距离中位数超过四分位距 1.5 倍的数据。在上一节的假设例子中，我们建议 $\theta_j$ 的先验或总体模型应该有比正态分布更长的尾巴。实际上，对于每个标准模型，都有一个自然的扩展，其中添加一个参数来允许过度分散。每个扩展模型都有作为混合分布的解释。

\textbf{关于过度分散的一个关键性质}。所有这些分布的一个特征是：基于混合表示，它们永远不可能欠分散。\footnote{推导见附录 \cref{sec:overdispersion-derivation}} 具体而言，对于正态-逆-$\chi^2$ 混合（产生 $t$ 分布）、Gamma-Poisson 混合（产生负二项分布）和 Beta-Binomial 混合，每种情况下的边际方差都等于或大于对应的原始分布的方差。这是混合结构的直接结果：条件均值的条件方差永远非负，因此边际方差 = 条件均值的无条件方差 + 条件方差的条件期望，前者贡献了额外的分散。相比之下，如果数据被认为相对于标准分布欠分散，则应使用不同的模型（如零膨胀模型或过分散的离散分布的反面）。

\subsection{$t$ 分布替代正态分布}

$t$ 分布比正态分布有更长的尾巴，可用于适应（1）数据分布中偶尔出现的异常观测值，或（2）先验分布或层次模型中偶尔出现的极端参数。$t$ 分布族——$t_\nu(\mu, \sigma^2)$——由三个参数表征：中心 $\mu$、尺度 $\sigma$ 和决定分布形状的``自由度''参数 $\nu$。$t$ 密度是对称的，$\nu$ 必须在 $(0, \infty)$ 范围内。当 $\nu = 1$ 时，$t$ 等价于柯西分布，其尾巴如此之长以至于均值和方差都不存在；当 $\nu \to \infty$ 时，$t$ 趋近于正态分布。如果 $t$ 分布是试图准确拟合长尾分布的概率模型的一部分，基于相当大量的数据，那么通常将自由度作为未知参数是恰当的。在选择 $t$ 作为正态分布的稳健替代方案的应用中，自由度可以固定在一个小值上以允许异常值，但不能小于先验理解所规定的值。例如，自由度为 1 或 2 的 $t$ 分布方差无限，在远尾处通常不现实。

\textbf{混合解释}。回想第 3.2 节和第 12.1 节，$t_\nu(\mu, \sigma^2)$ 分布可以解释为具有共同均值和方差服从逆-$\chi^2$ 分布的正态分布的混合。\footnote{推导见附录 \cref{sec:t-distribution-mixture-derivation}} 例如，模型 $y_i \sim t_\nu(\mu, \sigma^2)$ 等价于
\[
y_i \mid V_i \sim \mathrm{N}(\mu, V_i), \quad V_i \sim \mathrm{Inv-}\chi^2(\nu, \sigma^2).
\label{eq:t-distribution-mixture}
\]
这是我们在第 12.1 节中作为辅助变量和参数扩展的计算方法说明而引入的。从统计上讲，具有高方差的观测值可以被视为分布中的异常值。类似地，当建模可交换参数 $\theta_j$ 时也有类似的解释。

\textbf{不同 $\nu$ 值的直观理解}。为了建立直觉，考虑 $\nu$ 取不同值时 $t$ 分布的形状变化：$\nu = 1$（柯西）有极长尾；$\nu = 2$ 方差仍无限但尾巴收敛更快；$\nu = 4$ 是常用的稳健选择，接近 logistic 形状；$\nu = 10$ 已非常接近正态但仍有稍厚尾；$\nu \to \infty$ 精确等于正态分布。这种连续过渡使得敏感性分析可以系统地探测尾部行为的影响。

\subsection{负二项分布替代泊松分布}

在将泊松模型应用于数据时，一个常见困难是泊松模型要求方差等于均值；而在实践中，计数数据的分布通常过度分散，方差大于均值。我们已经在广义线性模型的背景下讨论过过度分散问题（第 16.1 节），第 16.4 节给出了层次正态模型用于过度分散泊松回归的例子。

另一种对过度分散计数数据进行建模的方法是使用负二项分布，这是一个允许均值和方差分别拟合的二参数族。\footnote{推导见附录 \cref{sec:negative-binomial-variance-derivation}} 服从 $\mathrm{Neg-bin}(\alpha, \beta)$ 分布的数据 $y_1, \ldots, y_n$ 可以被认为是具有均值 $\lambda_1, \ldots, \lambda_n$ 的泊松观测值，这些均值服从 $\mathrm{Gamma}(\alpha, \beta)$ 分布。负二项分布的方差是
\[
\mathrm{var}(y) = \beta + \frac{\beta^2}{\alpha}.
\label{eq:negative-binomial-variance}
\]
总是大于均值 $\beta$；这与泊松分布的方差总是等于其均值形成对比。当 $\alpha \to \infty$ 时，底层 gamma 分布趋近于一个确定性的尖峰，负二项分布退化为泊松分布。

\subsection{贝塔-二项分布替代二项分布}

类似地，二项离散数据模型的实际局限在于只有一个自由参数，这意味着方差由均值决定。标准的稳健替代方案是贝塔-二项分布（beta-binomial），顾名思义，它是二项分布的贝塔混合。贝塔-二项分布用于建模教育测试数据等，其中``成功''是正确答案，个人正确回答的概率差异很大。这里，数据 $y_i$——每个人 $i = 1, \ldots, n$ 的正确答案数——用 $\mathrm{Beta\text{-}bin}(m, \alpha, \beta)$ 分布建模，被认为是具有共同试验次数 $m$ 和不等概率 $\pi_1, \ldots, \pi_n$ 的二项观测值，这些概率服从 $\mathrm{Beta}(\alpha, \beta)$ 分布。贝塔-二项分布的方差大于具有相同均值的二项分布的方差，额外因子为 $\frac{\alpha+\beta+m}{\alpha+\beta+1}$。

\subsection{logistic 和 probit 回归的 $t$ 分布替代}

logistic 和 probit 回归可以说是不稳健的，因为线性预测器 $X\beta$ 的绝对值较大时，逆 logit 或 probit 变换给出的概率接近 0 或 1。可以通过允许对 $X\beta$ 的大值偶尔出现错误预测来使这些模型更加稳健。这种稳健性不是针对数据 $y$（在二元回归中等于 0 或 1）来定义的，而是针对预测变量 $X$ 定义的。一个更稳健的模型允许离散回归模型拟合大部分数据，同时偶尔产生孤立误差。

一个稳健模型——robit 回归——可以使用离散数据回归模型的潜变量公式（第 408 页）来实现，将潜在连续数据 $u$ 的 logistic 或正态分布替换为模型 $u_i \sim t_\nu((X\beta)_i, 1)$。在现实设置中，从数据中估计 $\nu$ 是不切实际的——因为潜数据 $u_i$ 从未被直接观测到，因此基本上不可能形成关于其连续底层分布形状的推断——所以将其设定为一个低值以确保稳健性。设定 $\nu = 4$ 得到的分布接近 logistic，当 $\nu \to \infty$ 时，模型趋近于 probit。

\subsection{为什么还要使用非稳健模型？}

$t$ 族包括正态分布作为特例，那么为什么我们还要使用正态分布，或者二项、泊松或其他标准模型呢？首先，每个标准模型都有一个逻辑状态，使其对许多应用问题是合理的。二项和多项分布适用于独立同分布结果的离散计数，具有固定的计数总数。泊松和指数分布拟合泊松过程的事件数和等待时间，这是按时间索引的独立离散事件的自然模型。最后，中心极限定理告诉我们，正态分布是作为大量独立分量之和形成的数据的适当模型。

即使它们不是自然由问题结构暗示的，标准模型在计算上也是方便的，因为共轭先验分布通常允许直接计算后验均值和方差以及容易进行模拟。这就是为什么对学业测试例子中的 $\theta_j$ 拟合正态总体模型很容易，以及为什么通常对全正值数据的对数或限制在 0 到 1 之间数据的 logit 拟合正态模型。这里需要做一个重要区分：当标准模型在结构上合理时，使用它们是正当的；但当模型是为了计算方便而任意选择时，敏感性分析就成为必要的步骤——而不是表明稳健模型不重要。在后一种情况下，稳健模型的使用与其说是``替代''，不如说是对假设的``压力测试''。当模型检查（第六章）显示标准模型拟合良好时，我们可以更有信心地使用它；但即使如此，探索稳健替代方案也能帮助我们理解推断对尾部假设的敏感性。

\section{后验推断与计算}

一如既往，我们可以使用第三部分描述的方法从后验分布（在敏感性分析情况下是一个或多个分布）中进行抽样。在本节中，我们简要描述在混合公式下稳健模型的吉布斯抽样使用。该方法在第 17.4 节的层次正态-$t$ 模型中进行了说明。然而，在扩展模型时，我们有可能使用不那么耗时的近似方法：我们可以用原始后验分布的抽取作为新模型模拟的起点。在本节中，我们还描述了可用于稳健模型和敏感性分析的两种技术：用于计算边际后验密度的重要性加权，以及用于近似稳健分析的重要性重采样。

\subsection{稳健模型作为简单模型扩展的符号}

我们使用 $p_0(\theta \mid y)$ 表示原始模型的后验分布，假设已经拟合到数据，使用 $\varphi$ 表示表征用于稳健性或敏感性分析的扩展模型的超参数。我们的目标是从
\[
p(\theta \mid \varphi, y) \propto p(\theta \mid \varphi) \, p(y \mid \theta, \varphi).
\label{eq:robust-extension-posterior}
\]
中抽样，使用 $\varphi$ 的预指定值（如稳健 $t$ 模型的 $\nu = 4$），或者对于一系列 $\varphi$ 值。在后一种情况下，我们还希望计算敏感性分析参数的边际后验分布 $p(\varphi \mid y)$。

稳健分布族可以通过 $p(\theta \mid \varphi)$ 或 $p(y \mid \theta, \varphi)$ 进入模型。例如，第 17.2 节专注于稳健数据分布，而我们第 17.4 节对 SAT 辅导实验的重新分析使用稳健分布建模模型参数。我们必须建立一个联合先验分布 $p(\theta, \varphi)$，这需要一些注意，因为它捕捉了 $\theta$ 和 $\varphi$ 之间的先验依赖关系。

\subsection{使用混合公式的吉布斯抽样}

马尔可夫链模拟可用于从后验分布 $p(\theta \mid \varphi, y)$ 中抽取。这可以使用混合公式完成，方法是从 $\theta$ 和额外未观测尺度参数（$t$ 模型中的 $V_i$、负二项模型中的 $\lambda_i$、贝塔-二项模型中的 $\pi_i$）的联合后验分布中抽样。

一个简单的例子，考虑将 $t_\nu(\mu, \sigma^2)$ 分布拟合到数据 $y_1, \ldots, y_n$，其中 $\mu$ 和 $\sigma$ 未知。给定 $\nu$，我们已经在第 12.1 节讨论了如何将吉布斯抽样器编程为涉及 $\mu$、$\sigma^2$、$V_1, \ldots, V_n$ 的参数化形式。如果 $\nu$ 本身未知，吉布斯抽样器必须扩展到包括从 $\nu$ 的条件后验分布中抽样的步骤。没有简单的方法来进行这一步，但可以使用 Metropolis 步骤。另一个复杂之处在于这类模型通常有多峰后验密度，不同的模式对应于 $t$ 分布尾巴中不同观测值的模式，这意味着在初始阶段需要额外的工作来搜索模式并在模拟中在模式之间跳跃，例如使用模拟退火（见第 12.3 节）。

\section{八校例子中的稳健推断与敏感性分析}

考虑基于第 5.5 节表 5.2 中数据的 SAT 辅导效应的层次模型。鉴于八个原始实验的大样本量，假设数据模型 $y_j \mid \theta_j, \sigma_j^2 \sim \mathrm{N}(\theta_j, \sigma_j^2)$（方差 $\sigma_j^2$ 已知）应该没有太大问题。总体模型 $\theta_j \mid \mu, \tau^2 \sim \mathrm{N}(\mu, \tau^2)$ 更难证明，尽管第 6.5 节的模型检查表明它对于获得学校效应的后验区间是足够的。然而，一般来说，后验推断可能对假设模型高度敏感，即使模型对观测数据提供了良好的拟合。为了说明稳健推断和敏感性分析方法，我们探索一个替代模型族，将 $t$ 分布拟合到学校效应总体：
\[
\theta_j \mid \nu, \mu, \tau \sim t_\nu(\mu, \tau^2), \quad j = 1, \ldots, 8.
\label{eq:t-population-model}
\]

我们使用符号 $p(\theta, \mu, \tau \mid \nu, y) \propto p(\theta, \mu, \tau \mid \nu) \, p(y \mid \theta, \mu, \tau, \nu)$ 表示在 $t_\nu$ 模型下的后验分布，$p_0(\theta, \mu, \tau \mid y) \equiv p(\theta, \mu, \tau \mid \nu = \infty, y)$ 表示第 5.5 节评估的正态模型下的后验分布。

\subsection{基于 $t_4$ 总体分布的稳健推断}

正如本章开头所讨论的，人们可能会担心正态总体模型导致最极端估计的学校效应被过多地向总体均值收缩。也许某个学校的辅导项目与其他学校足够不同，其估计值不应该被过多地向平均值收缩。一个相关担忧是最大观测效应可能对总体方差 $\tau^2$ 的估计产生不当影响，进而影响其他效应的贝叶斯估计。从建模角度来看，各种 SAT 辅导项目差异很大，其效应的总体可能更好地用长尾分布来拟合。为了评估这些担忧的重要性，我们执行稳健分析，用 $t$ 模型（\cref{eq:t-population-model}）替换正态总体分布，设定 $\nu = 4$，其他模型保持不变。

\textbf{吉布斯抽样}。我们使用第 12.1 节描述的方法进行吉布斯抽样，设定 $\nu = 4$。\footnote{详细计算步骤见附录 \cref{sec:eight-schools-gibbs-derivation}} 结果推断（基于从后验分布中抽取的 2500 次模拟——五个链的最后一半，每个链长 1000）汇总在表 17.1 中。结果与第 5.3 节正态模型下获得的推断\textbf{在中心趋势（后验中位数、均值）上基本相同}，但稳健模型对极端学校（如学校 A）的收缩程度稍小——这正是稳健分析的预期行为：减少对异常值的过度收缩。

\textbf{使用重要性重采样的计算}。尽管我们已经完成了马尔可夫链模拟，我们简要讨论如何应用重要性重采样来近似 $\nu = 4$ 的后验分布。首先，我们从 $p_0(\theta, \mu, \tau \mid y)$（正态模型下的后验分布）中抽取 5000 个 $(\theta, \mu, \tau)$ 的抽取值。接下来，我们计算每个抽取值的重要性比率：\footnote{推导见附录 \cref{sec:importance-ratio-derivation}}
\[
\frac{p(\theta, \mu, \tau \mid \nu, y)}{p_0(\theta, \mu, \tau \mid y)} = \frac{p(\mu, \tau \mid \nu) \, p(\theta \mid \mu, \tau, \nu) \, p(y \mid \theta, \mu, \tau, \nu)}{p_0(\mu, \tau) \, p_0(\theta \mid \mu, \tau) \, p_0(y \mid \theta, \mu, \tau)} \propto \prod_{j=1}^8 \frac{t_\nu(\theta_j \mid \mu, \tau^2)}{\mathrm{N}(\theta_j \mid \mu, \tau^2)}.
\label{eq:importance-ratio}
\]

然后我们使用重要性重采样，从 5000 个抽取值中无放回地抽取 500 个 $(\theta, \mu, \tau)$ 的子样本。这种近似对于评估稳健性可能已足够，但重要性比率对数的长尾分布（未显示）确实表明使用重要性重采样获得准确推断存在严重问题。

\subsection{基于不同 $\nu$ 值的 $t_\nu$ 分布的敏感性分析}

与稳健性略有不同的担忧是对正态总体分布先验假设的后验推断敏感性。为了研究敏感性，我们现在拟合一系列 $t$ 分布，自由度分别为 1、2、3、5、10 和 30。我们已经拟合了无限自由度（正态模型）和 4 个自由度（上述稳健模型）。

对于每个 $\nu$ 值，我们执行马尔可夫链模拟以从 $p(\theta, \mu, \tau \mid \nu, y)$ 中获取抽取值。我们通过每个学校效应 $\theta_j$ 的后验均值和标准差来总结结果。\cref{fig:nu-sensitivity-plot} 显示了作为 $\nu$ 函数的这些结果。在 $\nu_1$ 而非 $\nu$ 参数化方面有优势，因为包括正态分布在 $\nu_1 = 0$ 处，并将从正态到柯西分布的整个范围包含在有限区间 $[0, 1]$ 中。图中有一些变化，但没有显示出对超参数的系统性敏感性。

\subsection{将 $\nu$ 作为未知参数处理}

最后，我们将敏感性分析参数 $\nu$ 视为未知量，并在后验分布中对其进行平均。一般来说，这是关键的一步，因为我们通常只关心被数据支持的模型的敏感性。在这个特定例子中，推断对 $\nu$ 不敏感，因此计算边际后验分布是不必要的；我们在这里将其作为通用方法的说明。

\textbf{先验分布}。在计算 $\nu$ 的后验分布之前，我们必须为其指定先验分布。我们尝试在 $\nu_1$ 上使用均匀密度，范围为 $[0, 1]$（即从正态到柯西分布）。这个先验分布偏向长尾模型，一半的先验概率落在 $t_1$（柯西）和 $t_2$ 分布之间。

\textbf{后验推断}。为了将 $\nu$ 作为未知参数处理，我们修改稳健分析中使用的吉布斯抽样模拟，包括一个 Metropolis 步骤用于从 $\nu_1$ 的条件分布中抽样。\cref{fig:nu-posterior-histogram} 显示了 $\nu_1$ 模拟的直方图。

\textbf{讨论}。稳健性和对建模假设的敏感性取决于所研究的估计量。在 SAT 辅导例子中，八个学校效应的后验中位数、50\% 和 95\% 区间对正态总体分布假设不敏感（至少与 $t$ 族相比）。相比之下，99.9\% 区间可能强烈依赖于分布的尾部，并对 $t$ 分布中的自由度敏感——幸运的是，这些极端尾部在这个例子中不太可能有实质性意义。这一发现的意义在于：我们不需要担心正态假设会严重影响我们对这个特定数据集的中心推断；但如果我们对极尾部（如某效应超过 50 分的概率）感兴趣，那么选择哪个 $t$ 分布就变得关键了。

\section{使用 $t$ 分布误差的稳健回归}

与基于正态分布的其他模型一样，第 14 章正态线性回归模型下的推断对异常值或极端值敏感。稳健回归分析通过考虑回归误差的正态分布的稳健替代方案（如自由度数较少的 $t$ 分布）来获得。稳健误差分布（如方差较少的 $t$ 分布）将远离回归线的观测视为高方差观测，得到的结果类似于对异常值降权的结果。

\subsection{迭代加权线性回归和 EM 算法}

为了说明稳健回归计算，我们考虑将 $t_\nu$ 回归模型作为正态线性回归模型的替代方案，固定自由度参数。给定解释变量向量 $X_i$，个体响应变量 $y_i$ 的条件分布是 $p(y_i \mid X_i \beta, \sigma^2) = t_\nu(y_i \mid X_i \beta, \sigma^2)$。$t_\nu$ 分布可以表示为混合形式。$t$ 模型下 $p(\beta, \sigma^2 \mid \nu, y)$ 后验分布的模式可以直接使用牛顿法（第 13.1 节）或其他模式寻找技术获得。\footnote{详细推导见附录 \cref{sec:robust-regression-em-derivation}} 或者，我们可以利用 $t$ 模型的混合形式，使用 EM 算法，将方差 $V_i$ 视为``缺失数据''（即要平均的参数）来处理。

E 步计算在当前参数估计值给定的情况下，充分统计量的期望值。\footnote{详见附录 \cref{sec:robust-regression-em-derivation}} M 步是具有对角权重矩阵 $W$ 的加权线性回归，其中权重矩阵对角线包含 $1/V_i$ 的条件期望。EM 算法的迭代等价于迭代加权最小二乘算法中执行的迭代。给定回归参数的初始估计，计算每个案例的权重，残差较大的案例权重较小。然后通过加权线性回归获得回归参数的改进估计。

当自由度参数 $\nu$ 被视为未知时，可以应用 ECME 算法，添加一个额外步骤来更新自由度。

\textbf{其他稳健模型}。迭代加权线性回归或等效的 EM 算法可用于获得许多稳健替代模型的后验模式。更改用于观测方差 $V_i$ 的概率模型会产生替代稳健模型。例如，可以使用两点分布来建模具有污染误差的回归。

\subsection{吉布斯抽样和 Metropolis 算法}

可以使用吉布斯抽样和 Metropolis 算法从稳健回归模型的后验分布中获取样本。使用 $t_\nu$ 分布的混合参数化，我们可以通过交替从 $p(\beta, \sigma^2 \mid V_1, \ldots, V_n, \nu, y)$（使用加权线性回归的通常后验分布）和从 $p(V_1, \ldots, V_n \mid \beta, \sigma^2, \nu, y)$（一组独立的缩放逆-$\chi^2$ 分布）中抽取来获得后验样本。

\section{本章小结}

本章我们学习了：

\begin{enumerate}
  \item \textbf{稳健性的概念}：后验推断对异常值和模型假设的敏感性，以及 Box-Tiao 的历史贡献
  \item \textbf{标准模型的稳健替代}：
  \begin{enumerate}
    \item $t$ 分布替代正态分布（可使用正态-逆-$\chi^2$ 混合表示）
    \item 负二项分布替代泊松分布（Gamma-Poisson 混合）
    \item 贝塔-二项分布替代二项分布
    \item robit 回归（$t$ 分布替代 logistic/probit）
  \end{enumerate}
  \item \textbf{后验计算方法}：
  \begin{enumerate}
    \item 使用混合公式的吉布斯抽样
    \item 重要性加权和重要性重采样
  \end{enumerate}
  \item \textbf{敏感性分析}：通过将 $\nu$ 作为自由度参数变化来评估对假设的敏感性
  \item \textbf{稳健回归}：使用 $t$ 分布误差的稳健线性回归
\end{enumerate}

\section{下节预告}

下一章我们将学习缺失数据模型，这是贝叶斯分析中一个重要的实际问题，涉及如何对观测不完整的数据进行建模和推断。

\endinput

% Chapter 18: Models for missing data
% Bayesian Data Analysis - Gelman et al.

\chapter{缺失数据模型}
\label{ch:missing-data}

\section{本章导论}

在现实世界的统计研究中，缺失数据几乎无处不在——调查中的无响应、医学试验中的受试者退出、工业生产中的零件损坏。简单地丢弃不完整的观测会导致严重的偏倚，并低估推断的不确定性。

传统的统计方法往往假设数据是"完美"采集的，即每个应观测的变量都被记录下来。这一假设在许多应用中显然不成立。忽视缺失数据的存在，以"完全案例分析"（complete-case analysis）为例，只使用没有缺失值的观测，会导致估计量的偏倚。更隐蔽的是，完全案例分析通常会低估方差，因为它错误地假设样本量等于实际分析的观测数。

1976 年，Donald Rubin 首次系统性地提出了缺失机制的分类框架 \cite{Rubin1976}，将缺失数据分为完全随机缺失、随机缺失和非随机缺失三类。这一框架至今仍是该领域的标准，被广泛应用于生物统计、社会科学和计量经济学等领域。Rubin 的核心洞见是：缺失数据的处理方法应该依赖于缺失机制的性质，而正确识别缺失机制对于获得无偏估计至关重要。

本章我们将学习如何用概率模型描述数据的缺失机制，以及如何在贝叶斯框架下进行正确的推断。核心思想是：将缺失数据视为潜在的随机变量，为其指定概率模型，然后在完全贝叶斯分析中同时推断所有未知量。

\section{缺失数据机制}
\label{sec:missing-data-mechanisms}

\subsection{为什么需要描述缺失机制？}
\label{sec:why-describe-mechanism}

当我们面对缺失数据时，一个关键的问题是：哪些数据缺失了？为什么它们缺失了？

如果某些变量值的缺失与这些变量本身的值有关——例如，高收入人群更不愿意报告他们的工资——那么简单的补齐方法（如用均值填充）会产生严重偏倚。缺失机制（missing data mechanism）描述了缺失值与其对应变量值之间的关系。

统计学家 Rubin \cite{Rubin1976} 将缺失机制分为三类，这一分类至今仍是该领域的标准框架。

\subsection{三种缺失机制}

\begin{enumerate}
  \item \textbf{完全随机缺失（MCAR）}：缺失与观测数据或缺失数据都无关。例如，由于设备故障导致某些测量值丢失。
  \item \textbf{随机缺失（MAR）}：缺失只依赖于观测数据，而与缺失数据本身无关。例如，老年受试者更可能错过随访，但他们是否存活与是否参与随访无关（给定年龄）。
  \item \textbf{非随机缺失（MNAR）}：缺失依赖于缺失数据本身。例如，抑郁程度更严重的患者更可能退出治疗研究。
\end{enumerate}

\begin{Remark}
  MCAR 是最强的假设，意味着缺失是"真正随机"的。MAR 是许多多重插补方法的基础假设，但 MAR 假设本身无法从数据中验证——我们无法仅凭观测数据判断 MAR 是否成立，这正是敏感度分析存在的原因。MNAR 最难处理，因为它要求对缺失机制进行明确的建模。
\end{Remark}
\label{rem:mar-verifiability}

在贝叶斯分析中，我们通常假设缺失机制是\textbf{可忽略的（ignorable）}，这要求两个条件：
\begin{enumerate}
  \item \textbf{MAR（随机缺失）}：给定观测数据，缺失与否与缺失数据条件独立
  \item \textbf{区分性（Distinctness）}：完整数据模型的参数与缺失机制模型的参数是独立的
\end{enumerate}

\section{贝叶斯推断中的可忽略性}
\label{sec:ignorability}

\subsection{可忽略性的含义}

当缺失数据机制是可忽略的时候，我们可以在贝叶斯推断中"忽略"它——即不需要显式地建模缺失机制。后验分布只需基于观测数据的边际分布：

\[
p(\theta | y_{\text{obs}}) \propto p(\theta) p(y_{\text{obs}} | \theta)
\label{eq:ignorability-posterior}
\]

如果缺失机制不可忽略，后验分布需要包含缺失机制模型：

\[
p(\theta, \psi | y_{\text{obs}}) \propto p(\theta) p(\psi) p(y_{\text{obs}}, r | \theta, \psi)
\label{eq:non-ignorability-posterior}
\]

其中 $\psi$ 是缺失机制模型的参数，$r$ 是缺失指示向量。

\begin{Remark}
  区分性条件在实践中通常容易满足——我们通常假设数据模型和缺失机制模型的参数是不同的、独立的。例如，在正态均值估计中，均值参数 $\mu$ 与缺失概率参数通常是独立的。
\end{Remark}

\subsection{什么时候可忽略性成立？}
\label{sec:when-ignorability-holds}

可忽略的贝叶斯推断（ignorability）在以下情况成立：

\begin{enumerate}
  \item \textbf{MCAR 情况}：缺失完全随机，任何基于观测数据的推断都是无偏的（但可能低效）
  \item \textbf{MAR + 正确模型}：如果正确指定了条件缺失模型，且参数与数据模型区分
  \item \textbf{参数先验独立}：完整数据参数与缺失机制参数的先验分布相互独立
\end{enumerate}

\section{联合建模方法}
\label{sec:joint-modeling}

\subsection{完全数据模型}

联合建模方法（joint modeling approach）同时指定：

\begin{enumerate}
  \item \textbf{完全数据模型} $p(y_{\text{obs}}, y_{\text{mis}} | \theta)$
  \item \textbf{缺失机制模型} $p(r | y_{\text{obs}}, y_{\text{mis}}, \psi)$
\end{enumerate}

联合后验分布为：

\[
p(\theta, \psi, y_{\text{mis}} | y_{\text{obs}}, r) \propto p(\theta) p(\psi) p(y_{\text{obs}}, y_{\text{mis}} | \theta) p(r | y_{\text{obs}}, y_{\text{mis}}, \psi)
\label{eq:joint-posterior}
\]

\begin{Example}[二元数据的缺失机制]\label{ex:binary-missing-mechanism}
  假设 $y_i \in \{0, 1\}$ 是二项结果，缺失指示 $r_i$ 服从 logistic 模型：

  \[
  \mathrm{logit}(\mathbb{P}(r_i = 1)) = \psi_0 + \psi_1 y_i
  \label{eq:logistic-missing-model}
  \]

  当 $\psi_1 \neq 0$ 时，缺失依赖于潜在结果 $y_i$，这是 MNAR 情形。当 $\psi_1 = 0$ 时，缺失是 MCAR。

  我们可以在联合框架下估计 $\psi_1$，检验缺失机制是否依赖于潜在结果。
\end{Example}
推导见附录 \cref{sec:derivation-logistic-missing}。

联合建模的优点是提供了缺失机制的完整描述；缺点是需要对缺失机制做出明确的参数化假设，而这些假设往往难以验证。

\section{全贝叶斯方法与多重插补}
\label{sec:full-bayes-multiple-imputation}

\subsection{将缺失数据作为参数}

在全贝叶斯方法中，我们将缺失数据本身视为需要推断的参数。与潜在变量一样，缺失值 $y_{\text{mis}}$ 被赋予先验分布，然后通过 MCMC 同时抽样 $\theta$ 和 $y_{\text{mis}}$。

完整的分层结构为：

\begin{enumerate}
  \item \textbf{数据层}：$y_{\text{obs}}, y_{\text{mis}} | \theta \sim p(y | \theta)$
  \item \textbf{参数层}：$\theta \sim p(\theta)$
  \item \textbf{缺失机制层}：$r | y_{\text{obs}}, y_{\text{mis}}, \psi \sim p(r | y, \psi)$
\end{enumerate}

在 MAR 和可忽略性条件下，我们可以忽略缺失机制，只关注：

\[
p(\theta, y_{\text{mis}} | y_{\text{obs}}) \propto p(\theta) p(y_{\text{mis}} | \theta) p(y_{\text{obs}} | \theta, y_{\text{mis}})
\label{eq:bayesian-missing-posterior}
\]

\subsection{多重插补}

多重插补（multiple imputation）是一种实用的缺失数据处理技术，通过创建多个（通常是 $m = 5$ 到 $m = 50$）完整数据集来反映缺失数据的不确定性。\footnote{Rubin 组合规则的完整推导见附录 \cref{sec:derivation-rubin-combination}。}

\textbf{插补阶段}：给定观测数据和当前参数值，从后验分布中抽取缺失值：

\[
y_{\text{mis}}^{(l)} \sim p(y_{\text{mis}} | y_{\text{obs}}, \theta^{(l)})
\label{eq:imputation-step}
\]

\textbf{分析阶段}：对每个插补数据集应用标准的完全数据分析方法。

\textbf{合并阶段}：使用 Rubin 组合规则将 $m$ 个分析结果合并。

对于点估计 $\hat{Q}_l$ 和方差估计 $U_l$（$l = 1, \ldots, m$），合并的估计为：

\[
\bar{Q} = \frac{1}{m} \sum_{l=1}^m \hat{Q}_l
\label{eq:rubin-pooled-estimate}
\]

Rubin 组合规则的核心洞见是：总方差 $T$ 由两部分组成——来自数据分析的抽样方差 $\bar{U}$ 和来自缺失数据不确定性的插补间方差 $B$：

\[
T = \bar{U} + \left(1 + \frac{1}{m}\right) B
\label{eq:rubin-total-variance}
\]

其中 $\bar{U} = \frac{1}{m} \sum_{l=1}^m U_l$ 是 within-imputation 方差，$B = \frac{1}{m-1} \sum_{l=1}^m (\hat{Q}_l - \bar{Q})^2$ 是 between-imputation 方差。

自由度为：

\[
\nu = (m-1)\left(1 + \frac{\bar{U}}{(1 + 1/m)B}\right)^2
\label{eq:rubin-degrees-freedom}
\]

\begin{Remark}
  多重插补的关键洞见是：总方差 $T$ 包括两个部分——来自数据分析的抽样方差 $\bar{U}$ 和来自缺失数据不确定性的插补间方差 $B$。忽视后者会严重低估总体不确定性。
\end{Remark}

\section{敏感度分析}
\label{sec:sensitivity-analysis}

\subsection{为什么需要敏感度分析？}
\label{sec:why-sensitivity-analysis}

在 MNAR 情况下，缺失机制是不可忽略的，分析结果严重依赖于我们对缺失机制的假设。敏感度分析（sensitivity analysis）系统地检验：当缺失机制的假设发生变化时，结论如何变化？

即使在 MAR 情况下，如果模型误设（model misspecification），结论也可能对模型假设敏感。

\subsection{敏感度分析的实施方式}
\label{sec:sensitivity-methods}

敏感度分析有多种实施方式：

\begin{enumerate}
  \item \textbf{情景规划（Scenario-based）}：定义几个关于缺失机制的合理情景，比较不同情景下的推断
  \item \textbf{扰动分析（Perturbation）}：在似然函数或先验上引入小扰动，观察后验的稳定性
  \item \textbf{选择模型（Selection model）} vs \textbf{模式混合模型（pattern-mixture model）}：两种互补的建模框架
\end{enumerate}

\textbf{选择模型方法}将联合分布分解为：

\[
p(y, r) = p(y) p(r | y)
\label{eq:selection-model}
\]

\textbf{模式混合模型}则分解为：

\[
p(y, r) = \sum_{k} p(r = k) p(y | r = k)
\label{eq:pattern-mixture-model}
\]

其中 $k$ 表示不同的缺失模式（完全观测、部分缺失、完全缺失等）。

\begin{Example}[基于 Pattern-Mixture 的敏感度分析]\label{ex:pattern-mixture-sensitivity}
  假设二分之一的数据是完全观测的，其余二分之一在处理后缺失。模式混合方法对每种模式分别建模：

  \[
  y_{\text{obs},1} | \theta \sim \mathrm{N}(\mu, \sigma^2)
  \label{eq:pattern-mixture-complete}
  \]
  \[
  y_{\text{obs},2} | \theta \sim \mathrm{N}(\mu + \Delta, \sigma^2)
  \label{eq:pattern-mixture-missing}
  \]

  其中 $\Delta$ 衡量"缺失值与观测值的平均差异"。通过设定不同的 $\Delta$ 值（如 $\Delta = 0, \pm 0.5\sigma, \pm \sigma$），我们可以评估结论对 $\Delta$ 的敏感程度。
\end{Example}
推导见附录 \cref{sec:derivation-pattern-mixture}。

\section{缺失数据的实际处理策略}
\label{sec:practical-strategies}

\subsection{完整案例分析}

仅使用所有变量都被观测到的案例。优点是简单，缺点是：

\begin{enumerate}
  \item 如果 MCAR 不成立，会产生偏倚
  \item 浪费了部分观测的信息
  \item 方差估计基于不正确的样本量
\end{enumerate}

\subsection{单一插补}

用单一值（如均值、条件均值）填充缺失值。问题在于：

\begin{enumerate}
  \item 忽略了缺失数据的不确定性
  \item 通常会低估方差
  \item 可能扭曲数据的分布
\end{enumerate}

\subsection{逆概率加权}

在 MAR 假设下，通过对观测案例赋予权重来校正缺失偏倚。权重与缺失概率的倒数成正比：

\[
w_i = \frac{1}{\mathbb{P}(r_i = 1 | x_i)}
\label{eq:inverse-probability-weight}
\]

缺点是当缺失概率接近 0 时权重会变得极大，导致不稳定。

\begin{Remark}
  在实际应用中，多重插补通常是处理 MAR 数据的最稳健选择，因为它既保留了观测数据的信息，又通过多个插补反映了缺失的不确定性。
\end{Remark}

\section{本章小结}

本章我们学习了：

\begin{enumerate}
  \item 三种缺失数据机制：MCAR、MAR、MNAR \cite{Rubin1976}
  \item 可忽略性条件：MAR + 区分性
  \item 联合建模方法：同时建模完全数据和缺失机制
  \item 全贝叶斯方法：将缺失值作为参数进行 MCMC 推断
  \item 多重插补：Rubin 组合规则和方差的分解
  \item 敏感度分析：检验结论对缺失机制假设的依赖性
  \item 实际处理策略的优缺点比较
\end{enumerate}

缺失数据处理没有"一刀切"的解决方案。统计学家需要根据数据的具体缺失模式、缺失机制的合理性假设、以及计算资源，选择最合适的方法。在任何情况下，透明报告缺失数据的数量、模式和处理的假设，都是负责任的统计实践的基本要求。

\section{下节预告}

下一章我们将学习参数非线性模型（Parametric Nonlinear Models），这是对线性模型和广义线性模型的自然推广，允许更灵活的函数关系建模。

\endinput


% Part V: Nonlinear and Nonparametric Models
\part{非参数模型}

% Chapter 19: Parametric Nonlinear Models
% Bayesian Data Analysis - Gelman et al.

\chapter{参数非线性模型}

\section{本章导论}

在前几章中，我们主要关注线性模型——那些将参数以线性组合形式出现在均值函数中的模型。线性模型之所以被广泛使用，部分原因是它们易于解释：系数可以直接解读为"当自变量增加一个单位时，因变量期望值的变化"。然而，现实世界中的许多现象本质上是非线性的。

考虑几个简单的问题：光合作用如何随光照强度增加而变化？药物在体内的浓度如何随时间衰减？死亡率如何随年龄增长而变化？这些问题都无法用简单的线性关系来描述——它们涉及饱和效应、指数衰减、S形曲线等非线性现象。这就引出了本章的核心问题：\textbf{如何对这类非线性现象进行贝叶斯建模和推断？}

参数非线性模型为这个问题提供了一个优雅的解决方案。与非参数方法不同，我们仍然使用有限个参数来描述数据的生成机制；但与线性模型不同的是，这些参数以非线性方式出现在均值函数中。这种介于"完全参数化"与"完全非参数化"之间的方法，在许多科学领域有着悠久的应用历史，并在贝叶斯框架下获得了新的生命力。

\section{非线性模型的实例}

\subsection{光合作用与饱和效应}

在生态学中，研究光合作用速率如何随光照强度变化是一个经典问题。1975年，Jassby和Platform提出了以下模型来描述这一关系：

\begin{equation}
f(x) = \alpha \left(1 - \exp\left(-\frac{x}{\alpha}\right)\right), \label{eq:photosynthesis-model}
\end{equation}

其中 $x$ 是光照强度，$f(x)$ 是对应的光合作用速率，$\alpha$ 是最大光合作用速率。\footnote{光合作用模型的推导及生理学解释见附录 \cref{sec:nonlinear-photosynthesis-derivation}。}

这个模型捕捉了一个关键的非线性特征：\textbf{饱和效应}。当光照较弱时，光合作用速率近似线性增长；但随着光照强度的增加，增长速度逐渐放缓，最终趋于一个渐近值 $\alpha$。线性模型无法捕捉这种"边际效用递减"的特征——无论光照多强，线性模型总是预测线性增长。\footnote{线性模型与饱和效应的对比分析见附录 \cref{sec:linear-vs-saturation}。}

这种"饱和"现象在许多领域都有出现：药物疗效随剂量增加而饱和、学习曲线随练习增加而趋于平缓、化学反应速率随底物浓度增加而达到Plateau等。

\subsection{药物动力学：一室模型}

在药理学中，描述药物浓度随时间变化的模型是临床用药的基础。考虑最简单的一室模型：药物通过静脉注射进入血液，然后以固定速率从体内清除。设 $y$ 表示在时间 $t$ 测量的药物浓度，模型为

\begin{equation}
y = \frac{D}{V} \exp\left(-\frac{t}{\tau}\right) + \epsilon, \label{eq:pharmacokinetics-one-compartment}
\end{equation}

其中 $D$ 是给药剂量，$V$ 是表观分布容积，$\tau = V/k$ 是平均滞留时间（$k$ 是清除率），$\epsilon \sim \mathrm{N}(0, \sigma^2)$ 是测量误差。\footnote{一室模型的药代动力学推导见附录 \cref{sec:pharmacokinetics-derivation}。}

这个模型的非线性性来源于指数函数中的 $\tau$ 参数。指数衰减是非线性的：参数 $\tau$ 出现在指数的指数位置，而不是作为简单的乘法系数。

\subsection{Bioassay 实验}\label{sec:bioassay}

在毒理学中，生物测定（bioassay）是一个经典的应用场景。实验者对不同剂量 $x_i$ 的受试动物给予某种化合物，然后观察存活或死亡的比例。表19.1展示了一个典型的Bioassay数据集。

\begin{Example}[Bioassay 实验]\label{ex:bioassay}
设 $x_i$ 是第 $i$ 组的剂量（单位：gm/kg），$n_i$ 是该组的动物数量，$y_i$ 是存活的动物数量。假设每只动物的存活概率相互独立，则 $(y_i | \pi_i) \sim \mathrm{Bin}(n_i, \pi_i)$，其中 $\pi_i$ 是存活的概率。

我们使用逻辑斯谛模型描述剂量与存活概率的关系：

\begin{equation}
\pi_i = \mathrm{logit}^{-1}(\alpha + \beta x_i) = \frac{1}{1 + \exp(-(\alpha + \beta x_i))}. \label{eq:bioassay-logistic}
\end{equation}

参数 $\alpha$ 控制曲线的位置（中点），参数 $\beta$ 控制曲线的陡峭程度（斜率）。

\textbf{为何线性概率模型失效？}如果使用线性模型 $\pi_i = \alpha + \beta x_i$，当 $x_i$ 足够大时，预测值会超过1；当 $x_i$ 足够小时，预测值会小于0。这在概率意义上是不合理的。逻辑斯谛函数将整个实数轴映射到 $(0,1)$ 区间，自然地解决了这个问题。

表19.2展示了经典估计与贝叶斯估计的结果对比。\footnote{Bioassay 实验的完整数据分析包括 Tables 19.1 和 19.2，见附录 \cref{sec:bioassay-analysis}。}
\end{Example}

图 19.1 展示了后验分布的散点图以及对应的剂量-反应曲线（详见附录 \cref{sec:bioassay-rcode}）。

\subsection{种群动力学模型}

生物学中描述种群增长的经典模型也具有非线性结构。

\textbf{逻辑斯谛模型（Logistic growth）}假设种群增长率随种群规模增加而下降：

\begin{equation}
\frac{dN}{dt} = rN\left(1 - \frac{N}{K}\right), \label{eq:logistic-growth}
\end{equation}

其中 $N(t)$ 是时间 $t$ 的种群规模，$r$ 是固有增长率，$K$ 是环境容纳量。解这个微分方程得到

\begin{equation}
N(t) = \frac{K}{1 + \exp(-rt + \phi)}, \label{eq:logistic-solution}
\end{equation}

其中 $\phi$ 是与初始条件相关的参数。

\textbf{Ricker模型}是另一个经典的离散时间种群模型：

\begin{equation}
N_{t+1} = N_t \exp\left(r\left(1 - \frac{N_t}{K}\right)\right), \label{eq:ricker-model}
\end{equation}

这个模型用于描述鱼类等水产资源的种群动态，其中产卵量与亲代数量呈指数关系，但环境阻力导致幼鱼存活率下降。

\textbf{Lotka-Volterra 捕食者-猎物模型}描述两个物种的相互作用：

\begin{align}
\frac{dx}{dt} &= \alpha x - \beta xy, \label{eq:lotka-prey} \\
\frac{dy}{dt} &= \delta xy - \gamma y, \label{eq:lotka-predator}
\end{align}

其中 $x$ 是猎物数量，$y$ 是捕食者数量。\footnote{Lotka-Volterra 模型的详细推导及平衡点分析见附录 \cref{sec:lotka-volterra-derivation}。}

\section{非线性最小二乘估计}

\subsection{最小二乘框架}

对于非线性模型，估计参数的标准方法是最小二乘法。给定观测数据 $(t_i, y_i)$，$i = 1, \ldots, n$，我们最小化残差平方和

\begin{equation}
\mathrm{SSE}(\theta) = \sum_{i=1}^{n} [y_i - f(t_i, \theta)]^2. \label{eq:nonlinear-sse}
\end{equation}

其中 $f(t, \theta)$ 是非线性均值函数，$\theta$ 是 $k$ 维参数向量。

\textbf{最小二乘与最大似然的联系}：在正态误差假设下，即 $y_i \sim \mathrm{N}(f(t_i, \theta), \sigma^2)$，对数似然函数为

\[
\log p(y | \theta, \sigma) = -\frac{n}{2}\log(2\pi\sigma^2) - \frac{1}{2\sigma^2}\mathrm{SSE}(\theta).
\]

由此可见，\textbf{最小化 $\mathrm{SSE}(\theta)$ 等价于最大化正态似然}——非线性最小二乘估计正是正态误差假设下的 MLE。\footnote{最小二乘与 MLE 联系的详细推导见附录 \cref{sec:ls-mle-derivation}。}

与线性回归不同，非线性最小二乘估计没有闭式解，必须通过迭代优化算法求解。\footnote{非线性最小二乘估计为何没有闭式解的解释见附录 \cref{sec:nonlinear-no-closed-form}。}常用的算法包括：

\begin{enumerate}
  \item \textbf{高斯-牛顿法（Gauss-Newton）}：使用一阶导数近似 Hessian 矩阵
    \[
    H \approx J^T J, \quad \text{其中 } J_{ij} = \frac{\partial f(t_i, \theta)}{\partial \theta_j}. \label{eq:gauss-newton}
    \]
  \item \textbf{Levenberg-Marquardt法}：在高斯-牛顿法与最速下降法之间插值，通过调节参数 $\lambda$ 实现
  \item \textbf{牛顿法}：使用二阶导数，更精确但计算成本更高
\end{enumerate}

\textbf{迭代加权最小二乘（IRLS）}是拟合逻辑斯谛回归等广义线性模型的常用算法。对于逻辑斯谛回归，每次迭代求解：

\begin{equation}
\theta^{(t+1)} = (X^T W^{(t)} X)^{-1} X^T W^{(t)} z^{(t)}, \label{eq:irls-logistic}
\end{equation}

其中 $X$ 是设计矩阵（第1列为截距1，第2列为协变量 $x_i$），$z^{(t)}$ 是调整后的响应变量，$W^{(t)}$ 是权重矩阵。\footnote{IRLS 算法的详细推导见附录 \cref{sec:irls-derivation}。}

在 R 中，可以使用 \texttt{nls()} 函数实现这些算法。

\subsection{初始值与迭代收敛}

非线性最小二乘估计的一个关键挑战是\textbf{初始值的选择}。与线性模型不同，非线性优化通常不是凸的，可能存在多个局部极小值。因此，初始值的选择对最终估计有重要影响。

在实际应用中，初始值通常通过以下方式确定：

\begin{enumerate}
  \item \textbf{物理或生物学直觉}：基于对问题域的理解，估计参数的大致范围
  \item \textbf{数据可视化}：绘制数据图，估算渐近值、增长率等参数
  \item \textbf{简化模型拟合}：先拟合简化（线性化）的模型，获取参数初始估计
  \item \textbf{网格搜索}：在参数空间中搜索，找到使 SSE 全局最小的区域
\end{enumerate}

在 R 中实现带 bootstrap 的非线性最小二乘：

\begin{verbatim}
# R 代码示例：光合作用模型的非线性最小二乘拟合
library(nls)

# 数据
x <- c(0, 50, 100, 150, 200, 250, 300, 350, 400, 450, 500)
y <- c(4.2, 11.1, 19.5, 25.8, 30.2, 33.5, 35.8, 37.2, 38.1, 38.9, 39.2)

# 非线性最小二乘拟合
fit <- nls(y ~ alpha * (1 - exp(-x/alpha)),
           start = list(alpha = 50),
           algorithm = "port")

# Bootstrap 获取标准误
boot.se <- function(data, indices) {
  d <- data[indices, ]
  tryCatch(
    coef(nls(y ~ alpha * (1 - exp(-x/alpha)),
           data = d, start = list(alpha = 50))),
    error = function(e) NA
  )
}
\end{verbatim}

\section{非线性模型的贝叶斯推断}

\subsection{后验分布的计算}

在贝叶斯框架下，非线性模型的后验分布为

\begin{equation}
p(\theta | y) \propto p(\theta) p(y | \theta), \label{eq:nonlinear-posterior}
\end{equation}

其中 $p(\theta)$ 是先验分布，$p(y | \theta)$ 是似然函数，$y$ 是观测数据。

对于非线性模型，后验分布通常没有闭式表达。原因在于：非线性均值函数 $f(t, \theta)$ 使得似然函数 $p(y | \theta)$ 关于 $\theta$ 不是共轭的，也无法通过积分获得解析解。这就导致了后验分布的复杂非正态形状。\footnote{为何非线性模型后验没有闭式解的直觉解释见附录 \cref{sec:nonlinear-posterior-intuition}。}

我们使用马尔科夫链蒙特卡洛方法（MCMC）进行后验抽样。对于本章中的模型，一个有效的策略是使用 Metropolis 算法结合 Gibbs 抽样的混合方法。\footnote{Metropolis-Hastings 算法的详细推导见附录 \cref{sec:metropolis-hastings-derivation}。}

具体而言，我们可以采用以下步骤：

\begin{enumerate}
  \item 对每个参数 $\theta_j$，给定当前其他参数值 $\theta_{-j}$ 和数据 $y$，计算条件后验分布 $p(\theta_j | \theta_{-j}, y)$
  \item 使用 Metropolis-Hastings 步骤从条件后验分布中抽样
  \item 重复直到收敛，获得后验样本
\end{enumerate}

\subsection{Bioassay 模型的贝叶斯分析}\label{sec:bioassay-bayesian}

我们用 Bioassay 模型来说明贝叶斯推断的实际应用。设

\begin{align}
y_i &\sim \mathrm{Bin}(n_i, \pi_i), \label{eq:bioassay-likelihood} \\
\pi_i &= \mathrm{logit}^{-1}(\alpha + \beta x_i). \label{eq:bioassay-pi}
\end{align}

我们使用无信息先验：

\begin{equation}
p(\alpha, \beta) \propto 1. \label{eq:bioassay-prior}
\end{equation}

后验分布为

\begin{equation}
p(\alpha, \beta | y) \propto \prod_{i=1}^{n} \left(\frac{1}{1 + \exp(-(\alpha + \beta x_i))}\right)^{y_i} \left(\frac{1}{1 + \exp(\alpha + \beta x_i)}\right)^{n_i - y_i}. \label{eq:bioassay-posterior}
\end{equation}

在 R 中使用 JAGS 进行贝叶斯估计：

\begin{verbatim}
# R 代码示例：Bioassay 模型的贝叶斯分析
library(rjags)

model_string <- "
model {
  for (i in 1:n) {
    y[i] ~ dbin(p[i], n[i])
    logit(p[i]) <- alpha + beta * x[i]
  }
  alpha ~ dnorm(0, 0.0001)
  beta ~ dnorm(0, 0.0001)
}
"

# 准备数据
data <- list(y = c(0, 1, 3, 5, 6, 6),
              n = c(6, 6, 6, 6, 6, 6),
              x = c(-0.86, -0.30, -0.05, 0.73, 1.17, 1.80),
              n = 6)

# MCMC 抽样
model <- jags.model(textConnection(model_string), data = data)
samples <- coda.samples(model, c("alpha", "beta"), n.iter = 5000)
\end{verbatim}

\section{边缘似然与模型选择}\label{sec:marginal-likelihood}

\subsection{边缘似然的定义}

在贝叶斯框架下，模型比较的核心工具是边缘似然（marginal likelihood）。给定模型 $\mathcal{M}$ 和数据 $y$，边缘似然定义为

\begin{equation}
p(y | \mathcal{M}) = \int p(\theta | \mathcal{M}) p(y | \theta, \mathcal{M}) \, d\theta, \label{eq:marginal-likelihood}
\end{equation}

其中 $p(\theta | \mathcal{M})$ 是模型 $\mathcal{M}$ 的先验分布，$p(y | \theta, \mathcal{M})$ 是似然函数。

边缘似然衡量的是\textbf{在模型 $\mathcal{M}$ 下观察到数据 $y$ 的概率}。它自动实现了奥卡姆剃刀原则：复杂度较高的模型会被先验和似然的惩罚所平衡。

\textbf{为何边缘似然重要？}在非线性模型中，不同的参数化方式可能对应不同的物理或生物学意义。边缘似然提供了一种客观比较这些模型的方法。

\subsection{贝叶斯因子}

两个模型 $\mathcal{M}_1$ 和 $\mathcal{M}_2$ 的贝叶斯因子定义为

\begin{equation}
\mathrm{BF}_{12} = \frac{p(y | \mathcal{M}_1)}{p(y | \mathcal{M}_2)}. \label{eq:bayes-factor}
\end{equation}

当 $\mathrm{BF}_{12} > 1$ 时，数据更支持模型 $\mathcal{M}_1$。Jeffreys 提出了以下解释标准：

\begin{enumerate}
  \item $\log_{10} \mathrm{BF}_{12} \in (0, 0.5)$：微弱证据
  \item $\log_{10} \mathrm{BF}_{12} \in (0.5, 1)$：中等证据
  \item $\log_{10} \mathrm{BF}_{12} \in (1, 2)$：强证据
  \item $\log_{10} \mathrm{BF}_{12} > 2$：决定性证据
\end{enumerate}

对于非线性模型，边缘似然的计算通常需要使用变分近似或 MCMC 方法。Chib (1995) 提出的方法利用 MCMC 样本计算精确的边缘似然：\footnote{Chib (1995) 边缘似然计算方法的详细推导见附录 \cref{sec:chib-marginal-likelihood}。}

\begin{equation}
\log p(y) = \log p(y | \theta^*) + \log p(\theta^*) - \log p(\theta^* | y), \label{eq:chib-method}
\end{equation}

其中 $\theta^*$ 是后验众数或后验均值。

\subsection{嵌套模型检验}

在许多应用中，我们需要检验一个模型是否是另一个模型的特例。例如，在 Bioassay 模型中，检验 $\beta = 0$ 相当于检验剂量是否对存活率有影响。

\textbf{精确贝叶斯检验}：对于 Bioassay 模型，检验 $H_0: \beta = 0$（剂量无效应）的后验概率为

\begin{equation}
\mathbb{P}(\beta > 0 | y) = \frac{1}{B} \sum_{b=1}^{B} \mathbb{I}(\beta^{(b)} > 0), \label{eq:bioassay-exact-test}
\end{equation}

其中 $\beta^{(b)}$ 是从后验分布中抽取的样本。\footnote{嵌套模型检验的详细讨论见 Gelman et al. (2003) 第19.5节。}

在 R 中计算：

\begin{verbatim}
# 计算 P(beta > 0 | y)
library(rjags)
# ... (接上面的 JAGS 代码)

# 从后验样本中计算
beta_samples <- as.matrix(samples[, "beta"])
prob_positive <- mean(beta_samples > 0)
cat("P(beta > 0 | y) =", prob_positive, "\n")
\end{verbatim}

\section{模型不确定性}\label{sec:model-uncertainty}

在实际应用中，我们通常面临\textbf{模型不确定性}：我们不知道哪个模型真正生成了数据。这种不确定性可以通过模型平均来量化。

\subsection{过度拟合与偏差-方差权衡}

非线性模型的一个核心问题是\textbf{过度拟合（overfitting）}。当模型过于复杂时，它会拟合数据中的噪声而非真实信号。

对于非线性模型，偏差-方差分解仍然成立：

\begin{equation}
\mathbb{E}[( \hat{y} - y )^2] = \text{bias}^2 + \text{variance} + \sigma^2. \label{eq:bias-variance-nonlinear}
\end{equation}

在非线性模型中，偏差来源于模型形式本身的局限性（例如，用指数模型拟合逻辑斯谛型数据），而方差则随模型复杂度的增加而增加。

\subsection{模拟交叉验证}

评估模型预测能力的一种方法是\textbf{模拟交叉验证（simulation cross-validation）}：

\begin{enumerate}
  \item 将数据分为训练集和测试集
  \item 在训练集上拟合模型
  \item 在测试集上评估预测性能
  \item 重复多次，取平均
\end{enumerate}

对于非线性模型，由于参数估计的不确定性，预测区间应包含参数不确定性：

\begin{verbatim}
# 模拟交叉验证示例
n_folds <- 10
pred_errors <- numeric(n_folds)

for (k in 1:n_folds) {
  # 分割数据
  test_idx <- sample(1:n, n * 0.2)
  train_data <- data[-test_idx, ]
  test_data <- data[test_idx, ]

  # 在训练集上拟合
  fit <- nls(y ~ alpha * (1 - exp(-x/alpha)),
             data = train_data, start = list(alpha = 50))

  # 在测试集上预测
  y_pred <- predict(fit, newdata = test_data)
  pred_errors[k] <- mean((test_data$y - y_pred)^2)
}

cat("平均预测误差:", mean(pred_errors), "\n")
cat("预测误差标准误:", sd(pred_errors), "\n")
\end{verbatim}

\section{逻辑斯谛回归作为非线性模型}

逻辑斯谛回归是广义线性模型的一个特例，但也可以从非线性参数模型的角度来理解。设 $y_i \in \{0, 1\}$ 是二元结果，$x_i$ 是协变量。逻辑斯谛模型假设

\begin{equation}
\mathbb{P}(y_i = 1 | x_i) = \mathrm{logit}^{-1}(\beta_0 + \beta_1 x_i)
= \frac{1}{1 + \exp(-(\beta_0 + \beta_1 x_i))}. \label{eq:logistic-model}
\end{equation}

虽然这个模型通常通过最大似然方法拟合，但它确实具有非线性结构：参数出现在 logistic 函数的非线性变换中。

从贝叶斯角度，逻辑斯谛回归的后验分布没有闭式表达，需要使用 Metropolis 算法或拉普拉斯近似进行推断。\footnote{逻辑斯谛回归的贝叶斯推断细节见附录 \cref{sec:logistic-bayesian-derivation}。}

在 R 中使用 \texttt{glm()} 拟合逻辑斯谛回归：

\begin{verbatim}
# 逻辑斯谛回归示例
fit <- glm(y ~ x, data = data, family = binomial)
summary(fit)

# 置信区间
confint(fit)
\end{verbatim}

\section{非线性模型的选择与比较}

在非线性模型中，模型选择比线性模型更加复杂，因为不同的参数化方式可能对应不同的物理或生物学意义。常用的模型比较方法包括：

\begin{enumerate}
  \item \textbf{残差分析}：检查拟合残差的模式，识别模型偏差
  \item \textbf{信息准则}：AIC、BIC 等，考虑拟合优度与模型复杂度的平衡
  \item \textbf{交叉验证}：留出部分数据用于验证模型的预测能力
  \item \textbf{贝叶斯因子}：比较两个模型的边缘似然比
\end{enumerate}

\begin{Remark}
  在使用信息准则比较非线性模型时，必须注意参数数量的计算。对于非线性模型，"有效参数数量"可能小于显式参数数量，因为某些参数对拟合的贡献可能受到函数非线性结构的限制。
\end{Remark}

\section{非线性模型的诊断}

非线性模型拟合相关的诊断包括：

\begin{enumerate}
  \item \textbf{残差图}：绘制残差 vs 拟合值、残差 vs 协变量的图形
  \item \textbf{影响点分析}：识别对参数估计影响较大的异常观测
  \item \textbf{非线性性检验}：检验模型是否真正需要非线性结构
  \item \textbf{异方差性检验}：检查误差方差是否随均值变化
\end{enumerate}

\section{本章小结}

本章我们学习了参数非线性模型的贝叶斯分析方法：

\begin{enumerate}
  \item 非线性模型用于描述具有饱和效应、指数衰减、S形曲线等非线性特征的数据
  \item 常见的非线性模型包括光合作用模型、药物动力学模型、种群动力学模型、Bioassay 模型等
  \item Bioassay 示例展示了逻辑斯谛模型与贝叶斯分析的实际应用
  \item 非线性最小二乘估计通过迭代优化完成，需要合理选择初始值
  \item 在正态误差假设下，最小二乘估计等价于 MLE
  \item 边缘似然和贝叶斯因子是模型比较的核心工具
  \item 模型不确定性通过模型平均来量化
  \item 非线性模型的选择和诊断需要特别小心
\end{enumerate}

\section{下节预告}

下一章我们将学习基函数模型（Basis Function Models），这是另一类灵活的非参数化方法，通过线性组合基函数来近似复杂的函数关系。这类方法在函数估计、信号处理等领域有着重要应用。

\endinput

% Chapter 20: Basis Function Models
% Bayesian Data Analysis - Gelman et al.

\chapter{基函数模型}

\section{本章导论}

在前几章中，我们讨论了回归模型的参数估计和贝叶斯推断方法。然而，经典线性回归假设响应变量与预测变量之间存在线性关系——这一假设在许多实际应用中过于严格。例如，在金融时间序列分析中，资产收益率与宏观经济变量之间的关系往往是非线动的；在生物统计中，剂量反应曲线通常呈现 S 形而非直线。

\textbf{基函数方法}提供了一种优雅的解决方案：将非线性函数 $f(x)$ 表示为一组基函数的线性组合
\begin{equation}
f(x) = \sum_{j=1}^{J} h_j(x) \beta_j,
\label{eq:basis-function-decomposition}
\end{equation}
其中 $h_1(x), h_2(x), \ldots, h_J(x)$ 是选定的基函数，$\beta_1, \ldots, \beta_J$ 是对应的系数。通过贝叶斯推断，我们可以同时估计这些系数及其不确定程度，从而获得对非线性关系的灵活且有原则的估计。

本章我们将学习四种最重要的基函数方法：
\begin{enumerate}
\item \textbf{样条函数（Splines）}：通过分段多项式构建光滑的非线性函数；
\item \textbf{径向基函数（RBF）}：通过径向函数构造适用于多元函数的基函数；
\item \textbf{Sobolev 空间}：从函数空间的角度理解非参数推断的理论基础；
\item \textbf{小波（Wavelets）}：通过多尺度分解捕捉函数在不同分辨率下的特征。
\end{enumerate}

这四种方法虽然在构造基函数的方式上有所不同，但共享相同的核心思想：通过精心选择的基函数族，将无限维的函数估计问题转化为有限维的线性回归问题。

\section{样条函数模型}\label{sec:spline-models}

\subsection{从多项式回归到样条}

\textbf{多项式回归}是最直接的非线性建模方法——将预测变量的高次项作为额外的预测变量纳入模型：
\begin{equation}
y_i = \beta_0 + \beta_1 x_i + \beta_2 x_i^2 + \cdots + \beta_K x_i^K + \epsilon_i.
\label{eq:polynomial-regression}
\end{equation}
尽管多项式回归在概念上简单，但高次多项式在数据边界处往往表现不稳定——系数估计的方差随阶数增加而急剧膨胀，导致过拟合问题。

\textbf{样条函数}通过分段多项式克服了这些困难。在样条模型中，我们选择一个或多个\textbf{结点（knots）}$\xi_1 < \xi_2 < \cdots < \xi_K$，将预测变量域分割成多个区间，在每个区间内用低次多项式（如三次多项式）拟合，同时在结点处施加光滑性约束（通常要求函数值、一阶导数、二阶导数连续）。

具体而言，\textbf{三次样条（cubic spline）}可以表示为一组\textbf{截断幂基函数（truncated power basis）}：
\begin{align}
h_1(x) &= 1, & h_2(x) &= x, & h_3(x) &= x^2, & h_4(x) &= x^3, \label{eq:truncated-power-basis} \\
h_{4+k}(x) &= (x - \xi_k)_+^3,\quad k = 1,\ldots, K, \label{eq:knot-basis}
\end{align}
其中 $(x - \xi)_+ = \max(0, x - \xi)$ 是\textbf{正部（positive part）}函数。这 $4+K$ 个基函数张成了所有满足二阶导数连续的三次样条空间。

然而，在实际应用中，直接使用截断幂基函数在数值计算上不够稳定。一种更常用的方法是\textbf{B-样条（B-splines）}基函数——它具有局部支撑性（每个 B-样条仅在少数几个相邻区间内非零），这使得系数的解释更加直观，计算也更加高效。

\subsection{回归样条的贝叶斯推断}\label{sec:bayesian-spline-inference}

给定 $n$ 个观测点 $(x_i, y_i)$，样条回归模型为
\begin{equation}
y_i = \sum_{j=1}^{J} h_j(x_i) \beta_j + \epsilon_i,\quad \epsilon_i \sim \mathrm{N}(0, \sigma^2),
\label{eq:spline-model}
\end{equation}
其中 $J$ 是基函数的个数（对于三次样条，$J = 4 + K$，$K$ 为结点数）。

写成矩阵形式，设 $\mathbf{y} = (y_1, \ldots, y_n)^{\mathsf{T}}$，设计矩阵 $H$ 的第 $i$ 行第 $j$ 列为 $h_j(x_i)$，误差向量 $\boldsymbol{\epsilon} \sim \mathrm{N}(0, \sigma^2 I)$，则
\begin{equation}
\mathbf{y} \sim \mathrm{N}(H\boldsymbol{\beta}, \sigma^2 I).
\label{eq:spline-linear-model}
\end{equation}

\textbf{回归样条的贝叶斯推断}需要在系数上施加先验分布以实现光滑化。最常用的先验是\textbf{粗糖先验（g priors）}的推广——设
\begin{equation}
\boldsymbol{\beta} \sim \mathrm{N}(0, \sigma^2 \Sigma_\beta),
\label{eq:spline-g-prior}
\end{equation}
其中 $\Sigma_\beta$ 是先验协方差矩阵，反映了函数的光滑性先验。当 $\Sigma_\beta = \tau^2 P^{-1}$ 时（即 $\Sigma_\beta^{-1} = \tau^{-2} P$），后验均值中的惩罚项 $\Sigma_\beta^{-1}$ 恰好对应二阶导数惩罚 $\int |f''(x)|^2\, dx$——这建立了光滑性先验与惩罚泛函之间的对应关系。

具体而言，若使用截断幂基函数 \eqref{eq:truncated-power-basis}--\eqref{eq:knot-basis}，样条的光滑性可以用
\begin{equation}
\int |f''(x)|^2\, dx = \boldsymbol{\beta}^{\mathsf{T}} P \boldsymbol{\beta}
\label{eq:spline-roughness-penalty}
\end{equation}
衡量，其中 $P$ 是由基函数二阶导数构造的惩罚矩阵。

施加先验 \eqref{eq:spline-g-prior} 后的后验分布为
\begin{equation}
\boldsymbol{\beta} \mid \mathbf{y}, \sigma^2 \sim \mathrm{N}\left((H^{\mathsf{T}} H + \Sigma_\beta^{-1})^{-1} H^{\mathsf{T}} \mathbf{y},\ \sigma^2 (H^{\mathsf{T}} H + \Sigma_\beta^{-1})^{-1}\right).
\label{eq:spline-posterior}
\end{equation}
后验均值 $\hat{\boldsymbol{\beta}} = (H^{\mathsf{T}} H + \Sigma_\beta^{-1})^{-1} H^{\mathsf{T}} \mathbf{y}$ 可以看作是对最小二乘估计的收缩——这与岭回归的精神一致，只是惩罚矩阵 $\Sigma_\beta^{-1}$ 来自光滑性先验。

相应地，后验预测分布为
\begin{equation}
y_{\mathrm{new}} \mid x_{\mathrm{new}}, \mathbf{y} \sim \mathrm{N}\left(h(x_{\mathrm{new}})^{\mathsf{T}} \hat{\boldsymbol{\beta}},\ \sigma^2 \left(1 + h(x_{\mathrm{new}})^{\mathsf{T}} (H^{\mathsf{T}} H + \Sigma_\beta^{-1})^{-1} h(x_{\mathrm{new}})\right)\right).
\label{eq:spline-posterior-predictive}
\end{equation}
后验方差中的额外项 $h(x_{\mathrm{new}})^{\mathsf{T}} (H^{\mathsf{T}} H + \Sigma_\beta^{-1})^{-1} h(x_{\mathrm{new}})$ 反映了在未知新点处的不确定程度。

\begin{Example}[三次样条回归]\label{ex:cubic-spline-example}
考虑对非线性函数 $f(x) = \sin(2\pi x) + 0.5 \sin(4\pi x)$ 加上噪声观测生成的数据拟合三次样条。

设结点数 $K = 5$，等距分布在 $[0, 1]$ 区间。使用 B-样条基函数，模型为 $y_i = \sum_{j=1}^{J} h_j(x_i) \beta_j + \epsilon_i$。

先验协方差矩阵 $\Sigma_\beta$ 的构造使得相邻结点的系数变化受到惩罚，从而保证拟合曲线的光滑性。

后验均值的曲线将呈现为一条光滑曲线，既能捕捉到正弦函数的波动特征，又不会被噪声过度拟合。
\end{Example}
\footnote{三次样条回归的完整计算步骤，包括 B-样条基函数的构造和先验协方差矩阵的组装，见附录 \cref{sec:derivation-cubic-spline}。}

\subsection{平滑样条与等效核}\label{sec:smoothing-splines}

\textbf{平滑样条（smoothing splines）}提供了一种数据驱动的结点选择方法。与回归样条（需要预先指定结点）不同，平滑样条通过极小化带惩罚的残差平方和来同时选择拟合函数：
\begin{equation}
\sum_{i=1}^{n} (y_i - f(x_i))^2 + \lambda \int |f''(t)|^2\, dt,
\label{eq:smoothing-spline-objective}
\end{equation}
其中 $\lambda \geq 0$ 是光滑参数，$\int |f''(t)|^2\, dt$ 是函数二阶导数平方的积分，度量函数的光滑程度。

目标函数 \eqref{eq:smoothing-spline-objective} 的解是唯一的三次样条函数，其结点恰好位于所有唯一的 $x_i$ 处。这个解可以写成
\begin{equation}
\hat{f}(x) = \sum_{i=1}^{n} S_\lambda(x, x_i) y_i,
\label{eq:smoothing-spline-solution}
\end{equation}
其中 $S_\lambda(x, x')$ 称为\textbf{等效核（equivalent kernel）}\footnote{等效核的详细推导及其与核密度估计的联系，见附录 \cref{sec:derivation-equivalent-kernel}。}。

等效核具有以下重要性质：它是正定的，即 $\int S_\lambda(x, u) S_\lambda(u, x')\, du = S_\lambda(x, x')$，并且在不同 $x$ 处等效核的形状不同（这与标准核密度估计不同），这使得平滑样条在边界处具有自适应光滑的特性。

\textbf{自由度（degrees of freedom）}是量化平滑样条灵活性的重要指标。设 $S$ 为平滑矩阵（满足 $\hat{\mathbf{y}} = S \mathbf{y}$），则自由度定义为
\begin{equation}
\mathrm{df}(\lambda) = \mathrm{tr}(S).
\label{eq:smoothing-spline-df}
\end{equation}
当 $\lambda \to 0$ 时，$\mathrm{df} \to n$，拟合曲线接近插值（过拟合）；当 $\lambda \to \infty$ 时，$\mathrm{df} \to 1$，拟合曲线退化为线性回归（欠拟合）。因此，$\lambda$ 的选择控制了拟合曲线的光滑程度。

在贝叶斯框架下，$\lambda$ 可以通过边际似然或交叉验证来估计：
\begin{equation}
p(\lambda \mid \mathbf{y}) \propto p(\mathbf{y} \mid \lambda) p(\lambda).
\label{eq:smoothing-spline-lambda-posterior}
\end{equation}
计算 $p(\mathbf{y} \mid \lambda)$ 需要对系数 $\boldsymbol{\beta}$ 积分，这可以通过析因分解或 MCMC 完成。

\begin{Example}[等效核的可视化]\label{ex:equivalent-kernel-visualization}
对于均匀设计 $x_i = i/n$，等效核 $S_\lambda(x, \cdot)$ 在中心点 $x = 0.5$ 处近似高斯核的形状；但在边界点 $x = 0.05$ 处，等效核被截断并重新标准化以满足边界条件。

这种自适应的边界处理是平滑样条相对于传统核方法的优势之一。
\end{Example}
\footnote{等效核在边界点处行为的详细分析，包括 Mercer 定理和特征分解，见附录 \cref{sec:derivation-equivalent-kernel-boundary}。}

\subsection{基函数的构造：B-样条递推公式}\label{sec:b-spline-construction}

B-样条基函数是样条计算的核心工具。本节给出 de Boor 递推公式的完整描述。

设结点序列 $\xi_1 < \cdots < \xi_K$ 分布在 $[a, b]$，扩展结点为 $\xi_0 = \xi_{-1} = \cdots = a$ 和 $\xi_{K+1} = \xi_{K+2} = \cdots = b$。B-样条基函数 $B_{i,k}(x)$（次数为 $k-1$）由 de Boor 递推公式定义：
\begin{align}
B_{i,1}(x) &= \begin{cases} 1 & \xi_i \leq x < \xi_{i+1}, \\ 0 & \text{otherwise}, \end{cases} \label{eq:de-boor-base} \\
B_{i,k}(x) &= \frac{x - \xi_i}{\xi_{i+k-1} - \xi_i} B_{i,k-1}(x) + \frac{\xi_{i+k} - x}{\xi_{i+k} - \xi_{i+1}} B_{i+1,k-1}(x). \label{eq:de-boor-recursion}
\end{align}

\textbf{局部支撑性}：每个 $B_{i,k}$ 仅在区间 $[\xi_i, \xi_{i+k}]$ 上非零，这意味着改变一个结点只影响 $k$ 个相邻基函数，使得计算高效且系数的几何解释直观。

\textbf{单位分解性}：在每个内点区间 $[\xi_j, \xi_{j+1}]$ 上，恰好有 $k$ 个 $k-1$ 次 B-样条非零，且它们的和为 1：$\sum_{i=j-k+1}^{j} B_{i,k}(x) = 1$ for $x \in [\xi_j, \xi_{j+1}]$。

\textbf{例子（$k=4$，三次 B-样条）}：对于结点 $\xi_1 < \xi_2 < \xi_3$，第一个三次 B-样条 $B_{1,4}(x)$ 在 $[\xi_0, \xi_4]$ 上有定义，在 $[\xi_0, \xi_1]$ 和 $[\xi_3, \xi_4]$ 上分别为三次多项式，在 $[\xi_1, \xi_3]$ 上为四次多项式的组合。

样条函数 $f(x) = \sum_i \beta_i B_{i,k}(x)$ 在每个区间 $[\xi_j, \xi_{j+1}]$ 上是 $k-1$ 次多项式，且在整个 $[a, b]$ 上有 $k-2$ 阶连续导数——这正是样条"光滑拼接"的几何意义。

\section{径向基函数}\label{sec:radial-basis-functions}

\textbf{径向基函数（Radial Basis Functions, RBF）}提供了另一种重要的基函数构造方法。与样条基于多项式展开不同，RBF 的基函数是关于某个中心点的径向函数：
\begin{equation}
f(x) = \sum_{i=1}^{m} c_i \phi(\|x - c_i\|),
\label{eq:rbf-expansion}
\end{equation}
其中 $\phi: [0, \infty) \to \mathbb{R}$ 是\textbf{径向基函数}，$c_1, \ldots, c_m$ 是\textbf{中心点（centers）}，$\|\cdot\|$ 是欧氏距离。

RBF 方法的核心思想是：将多元函数分解为一系列以数据点为中心的局部径向函数的叠加。这种构造在空间维度上没有限制——无论是低维还是高维空间，RBF 都适用。

\textbf{常用径向基函数}包括：

\begin{enumerate}
\item \textbf{高斯径向基函数（Gaussian RBF）}：
\begin{equation}
\phi(r) = \exp(-\gamma r^2),\quad \gamma > 0.
\label{eq:gaussian-rbf}
\end{equation}
高斯 RBF 是正定的，其核矩阵始终是正定的。这使得高斯 RBF 在理论上和计算上都易于处理。

\item \textbf{多谐波样条（Polyharmonic splines）}：
\begin{equation}
\phi(r) = \begin{cases} r^k, & k = 1, 3, 5, \ldots, \\ r^k \log r, & k = 2, 4, 6, \ldots. \end{cases}
\label{eq:polyharmonic-spline}
\end{equation}
多谐波样条在插值问题中特别有用——当中心点互异时，插值矩阵是非奇异的。
\end{enumerate}

\textbf{高斯 RBF 的贝叶斯推断}等价于在函数空间上施加高斯过程先验。设中心点为 $c_1, \ldots, c_m$，则 $f(x) \sim \mathrm{GP}(0, K)$，其中核函数 $K(x, x') = \sum_i \tau_i^2 \phi_\gamma(\|x - c_i\|) \phi_\gamma(\|x' - c_i\|)$。从这一点看，RBF 是高斯过程的离散近似——这将在下一章（\cref{sec:basis-functions-outlook}）中详细展开。

\begin{Example}[高斯 RBF 回归]\label{ex:rbf-regression}
考虑对二维函数 $f(x_1, x_2) = \sin(x_1) \cos(x_2)$ 在网格点上加噪声观测的数据拟合高斯 RBF 模型。

设中心点为 $c_1, \ldots, c_m$，均匀分布在 $[0, 2\pi] \times [0, 2\pi]$ 上。模型为 $y_i = \sum_{j=1}^{m} c_j \exp(-\gamma \|x_i - c_j\|^2) + \epsilon_i$。

参数 $\gamma$ 控制基函数的宽度——$\gamma$ 越小，基函数越宽，模型越光滑；$\gamma$ 越大，基函数越窄，模型越灵活。$\gamma$ 和系数 $c_j$ 可以通过边际似然或交叉验证联合估计。
\end{Example}

RBF 方法的主要优势在于其\textbf{维数无关性}——即使在高维空间中，只要选择合适的中心点，RBF 也能有效逼近复杂函数。然而，当数据点很多时，RBF 的计算成本为 $O(n^2)$（核矩阵的构造和求逆），这限制了其在大规模数据中的应用。

\section{Sobolev 空间与函数空间视角}\label{sec:sobolev-spaces}

从更一般的角度看，样条函数属于一类特殊的函数空间——\textbf{Sobolev 空间}。Sobolev 空间 $W_2^k([0, 1])$ 定义为所有 $k$ 阶导数平方可积的函数构成的空间：
\begin{equation}
W_2^k([0, 1]) = \left\{ f : [0, 1] \to \mathbb{R} \mid f^{(j)} \in L^2([0, 1]),\ j = 0, 1, \ldots, k \right\}.
\label{eq:sobolev-space-definition}
\end{equation}
对于 $k=2$，三次样条恰好是 Sobolev 空间 $W_2^2([0, 1])$ 中使得泛函 $\int |f''(x)|^2\, dx$ 最小的函数类。

\textbf{为什么选择 $W_2^k$？}这里的上标 2 表示在 $L^2$ 空间（平方可积函数空间）中测量导数。Sobolev 选择 $L^2$ 而非 $L^1$ 或 $L^\infty$，是因为 $L^2$ 空间是 Hilbert 空间——具有内积结构，这使得泛函分析和谱理论工具得以应用。具体而言，$W_2^k$ 是一个再生核 Hilbert 空间（RKHS），这为核方法提供了几何基础。19 世纪末，Sobolev（以及更早的 Aronszajn）在研究偏微分方程时发展了这一理论，他们发现 $L^2$ 的选择使得许多分析问题在技术处理上比 $L^1$ 更加便利。

\textbf{Sobolev 嵌入定理}：当 $k > \ell + 1/2$ 时，$W_2^k([0, 1])$ 紧嵌入到 $C^\ell([0, 1]$（$\ell$ 阶连续可导函数空间）。这意味着，在 $W_2^2$ 中的函数（$f, f' \in L^2$）自动是连续的——这正是样条函数具有连续导数的数学保证。

这一函数空间视角揭示了样条方法与再生核 Hilbert 空间（RKHS）之间的深层联系。设 $K(x, x')$ 是满足再生性质的核函数，即 $\langle K(\cdot, x), K(\cdot, x') \rangle_H = K(x, x')$，则 $H = \{f : f(x) = \sum_i c_i K(x, x_i)\}$ 是 $L^2$ 的一个子空间。

\textbf{核函数的统计意义}在于：RKHS 中的范数 $\|f\|_H^2$ 恰好度量了函数的光滑程度。对于高斯过程回归（我们将在下一章详细讨论），核函数 $K(x, x')$ 定义了过程的协方差函数；而对于样条惩罚 \eqref{eq:smoothing-spline-objective}，其解恰好位于由惩罚核生成的 RKHS 中。

具体而言，平滑样条问题 \eqref{eq:smoothing-spline-objective} 的解空间可以表示为
\begin{equation}
\hat{f}(x) = \sum_{j=1}^{J} h_j(x) \beta_j,\quad \boldsymbol{\beta} \sim \mathrm{N}(0, \sigma^2 \lambda P^{-1}),
\label{eq:sobolev-rkhs-representation}
\end{equation}
其中 $P_{ij} = \int h_i''(x) h_j''(x)\, dx$ 是惩罚矩阵。这表明从贝叶斯视角看，样条回归等价于在 RKHS 上施加正态先验。

\textbf{核与惩罚的对偶性}是理解非参数贝叶斯方法的关键：高斯过程先验 $f \sim \mathrm{GP}(0, K)$ 等价于在 RKHS 上施加先验；惩罚泛函 $\int |f''(x)|^2\, dx$ 则对应于特定的核选择。这种对偶性使得我们可以在函数空间视角和系数空间视角之间自由切换——有时在系数空间更容易计算，有时在函数空间更容易理解。

\section{小波方法}\label{sec:wavelet-methods}

小波方法为非参数回归提供了另一种强大的工具。与样条函数通过分段多项式构造基函数不同，小波通过\textbf{伸缩（dilation）和平移（translation）}从单个母小波 $\psi$ 构造一族基函数：
\begin{equation}
\psi_{j,k}(x) = 2^{j/2} \psi(2^j x - k),\quad j, k \in \mathbb{Z}.
\label{eq:wavelet-basis}
\end{equation}
其中 $j$ 控制频率（尺度），$k$ 控制位置（平移）。

小波基函数具有两个显著特点：
\begin{enumerate}
\item \textbf{多分辨率性}：不同 $j$ 对应不同的空间分辨率。低 $j$ （大尺度）捕捉函数的整体轮廓；高 $j$ （小尺度）捕捉局部细节。
\item \textbf{局部支撑性}：大多数小波（如 Daubechies 小波）在时域和频域都是局部化的，使得小波展开能够有效捕捉局部特征。
\end{enumerate}

小波回归模型为
\begin{equation}
y_i = f(x_i) + \epsilon_i = \sum_{j,k} \theta_{jk} \psi_{jk}(x_i) + \epsilon_i,
\label{eq:wavelet-regression-model}
\end{equation}
其中 $\theta_{jk} = \langle f, \psi_{jk} \rangle$ 是小波系数。

由于小波基函数是正交的，小波系数 $\theta_{jk}$ 的最小二乘估计为 $\hat{\theta}_{jk} = \langle \mathbf{y}, \psi_{jk} \rangle$——正交投影得到的系数之间是独立的。然而，在实际应用中，我们通常使用\textbf{阈值法（thresholding）}来获得更好的统计性质，这涉及跨系数的全局决策（哪些系数保留、哪些收缩）。

\subsection{小波阈值法}\label{sec:wavelet-thresholding}

Donoho 和 Johnstone 的著名结果表明，在一定条件下，\textbf{小波阈值估计器}能够同时达到最优的收敛速率。两种最常用的阈值规则是：

\textbf{硬阈值（hard thresholding）}：
\begin{equation}
\hat{\theta}_{jk}^{\mathrm{hard}} = \hat{\theta}_{jk} \cdot \mathbb{I}\left(|\hat{\theta}_{jk}| > \lambda\right),
\label{eq:hard-thresholding}
\end{equation}
\textbf{软阈值（soft thresholding）}：
\begin{equation}
\hat{\theta}_{jk}^{\mathrm{soft}} = \mathrm{sgn}(\hat{\theta}_{jk}) \cdot (|\hat{\theta}_{jk}| - \lambda)_+,
\label{eq:soft-thresholding}
\end{equation}
其中 $\lambda$ 是阈值参数，$\mathbb{I}(\cdot)$ 是示性函数，$\mathrm{sgn}(\cdot)$ 是符号函数。

阈值参数 $\lambda$ 的选择至关重要。Donoho 和 Johnstone 建议使用\textbf{通用阈值} $\lambda = \sigma \sqrt{2\log n}$，其中 $\sigma$ 是噪声标准差，$n$ 是样本量。这个阈值在大样本下能够保证估计器的一致性，但可能存在过度收缩的问题。

实践中更常用的方法是\textbf{交叉验证阈值}或\textbf{Bayes 阈值}——后者假设小波系数服从双指数先验，通过边际似然最大化来选择 $\lambda$。

从贝叶斯视角看，\textbf{贝叶斯阈值估计器}等价于假设小波系数服从一个分层先验：
\begin{align}
\theta_{jk} \mid \sigma_j^2 &\sim \mathrm{N}(0, \sigma_j^2), \label{eq:wavelet-hierarchical-prior1} \\
\sigma_j^2 &\sim \text{Inverse-Gamma}(a, b), \label{eq:wavelet-hierarchical-prior2}
\end{align}
其中 $\sigma_j^2$ 是第 $j$ 个尺度上所有系数的共同方差。这种先验设计体现了小波系数的"稀疏性"先验——在大多数尺度上，只有少数系数有显著的非零值。

后验均值可以通过以下近似公式计算：
\begin{equation}
\mathbb{E}(\theta_{jk} \mid \hat{\theta}_{jk}) \approx \frac{\hat{\theta}_{jk}}{1 + \lambda_{\mathrm{Bayes}} / \hat{\theta}_{jk}^2},
\label{eq:wavelet-bayes-estimate}
\end{equation}
其中 $\lambda_{\mathrm{Bayes}}$ 是由先验超参数决定的收缩量。

\begin{Example}[Haar 小波的简单例子]\label{ex:haar-wavelet}
Haar 小波是最简单的正交小波，定义为
\begin{equation}
\psi(x) = \begin{cases} 1 & 0 \leq x < 1/2, \\ -1 & 1/2 \leq x < 1, \\ 0 & \text{otherwise}. \end{cases}
\label{eq:haar-wavelet}
\end{equation}
在区间 $[0, 1]$ 上，Haar 小波系数的计算是直观的：$\theta_{00} = \int_0^1 f(x)\, dx$（尺度系数），$\theta_{jk} = 2^{j/2} \int_{k/2^j}^{(k+1/2)/2^j} f(x)\, dx - 2^{j/2} \int_{(k+1/2)/2^j}^{(k+1)/2^j} f(x)\, dx$（小波系数）。

对于分段常数函数，Haar 小波给出了紧凑的表示；对于光滑函数，高尺度的小波系数趋于零，这正是小波压缩性的体现。
\end{Example}
\footnote{Haar 小波系数计算的详细推导，以及 Daubechies 小波的构造原理，见附录 \cref{sec:derivation-wavelet-coefficient}。}

\subsection{小波与样条的比较}\label{sec:wavelet-vs-spline}

小波方法和样条方法各有优劣，适用于不同的场景：

\textbf{样条的优势}：
\begin{itemize}
\item 结点的显式几何解释；
\item 在边界处自适应良好的表现；
\item 与高斯过程（见下一章）的紧密联系。
\end{itemize}

\textbf{小波的优势}：
\begin{itemize}
\item 压缩性：对于光滑函数，仅需少量小波系数即可精确近似；
\item 多分辨率分析：能够同时捕捉函数在不同尺度上的特征；
\item 阈值估计的最优收敛速率。
\end{itemize}

在实际应用中，选择样条还是小波取决于数据的特征和分析目标。对于边界效应不明显的区域型数据，小波往往更高效；对于需要光滑解释的曲线型数据，样条更合适。

\section{贝叶斯非参数推断的展望}\label{sec:basis-functions-outlook}

本章介绍的基函数方法——样条、径向基函数、Sobolev 空间和小波——代表了贝叶斯非参数统计的四大支柱。它们共同的核心思想是：通过精心选择的基函数族，将无限维的函数估计问题转化为有限维的线性回归问题，然后利用贝叶斯推断同时估计系数和光滑参数。

在下一章中，我们将看到这一思想的最优雅实现——\textbf{高斯过程（Gaussian Process）}回归。在高斯过程框架下，基函数的选择被隐含在核函数的设计中，而边际似然提供了自动选择光滑参数的原则性方法。高斯过程既是样条函数的贝叶斯极限，也能包容小波型的稀疏先验——它为我们提供了一个统一而灵活的非参数推断框架。

\section{本章小结}

本章我们学习了三种基函数方法：

\begin{enumerate}
\item \textbf{样条函数}（\cref{sec:spline-models}）：通过分段多项式和光滑性惩罚实现非线性建模。回归样条通过指定结点位置和个数构造基函数；平滑样条通过极小化带惩罚的残差平方和自动选择拟合函数。等效核和自由度是量化平滑样条性质的重要工具。B-样条的 de Boor 递推公式提供了局部支撑、高效计算的基函数构造方法（\cref{sec:b-spline-construction}）。
\item \textbf{径向基函数}（\cref{sec:radial-basis-functions}）：通过径向函数 $\phi(\|x - c_i\|)$ 构造基函数，适用于多元函数逼近。高斯 RBF 和多谐波样条是两种最重要的径向基函数。RBF 方法等价于高斯过程先验的离散近似。
\item \textbf{Sobolev 空间}（\cref{sec:sobolev-spaces}）：提供了理解非参数推断的函数空间视角。$W_2^k$ 空间的选择是因为 $L^2$ 是 Hilbert 空间，具有内积结构，便于应用泛函分析工具。样条空间是 Sobolev 空间的特例，惩罚泛函对应于 RKHS 的范数，核与惩罚之间存在对偶关系。
\item \textbf{小波方法}（\cref{sec:wavelet-methods}）：通过伸缩和平移从母小波构造正交基，实现多分辨率分析。阈值法是小波估计的核心技术，硬阈值和软阈值各有优劣。贝叶斯分层先验为阈值选择提供了原则性框架。
\end{enumerate}

这四种方法为下一章的高斯过程模型奠定了概念基础。

\section{附录：公式推导}\label{sec:appendix-chapter20}

\subsection{三次样条回归的完整计算}\label{sec:derivation-cubic-spline}

\textbf{背景}：在 \cref{ex:cubic-spline-example} 中，我们简要介绍了三次样条回归。本节给出 B-样条基函数的构造方法和先验协方差矩阵的组装细节。

\textbf{参数定义}：设 $K$ 个结点 $\xi_1 < \cdots < \xi_K$ 分布在 $[a, b]$ 区间，扩展结点为 $\xi_0 = \xi_{-1} = \cdots = a$ 和 $\xi_{K+1} = \xi_{K+2} = \cdots = b$。

\textbf{B-样条基函数的递推定义}：B-样条基函数 $B_{i,k}(x)$（次数为 $k-1$）由 de Boor 递推公式给出：
\begin{align}
B_{i,1}(x) &= \begin{cases} 1 & \xi_i \leq x < \xi_{i+1}, \\ 0 & \text{otherwise}, \end{cases} \label{eq:de-boor-base} \\
B_{i,k}(x) &= \frac{x - \xi_i}{\xi_{i+k-1} - \xi_i} B_{i,k-1}(x) + \frac{\xi_{i+k} - x}{\xi_{i+k} - \xi_{i+1}} B_{i+1,k-1}(x). \label{eq:de-boor-recursion}
\end{align}

\textbf{光滑性惩罚矩阵}：对于三次样条（$k=4$），惩罚矩阵 $P$ 的元素为
\begin{equation}
P_{ij} = \int_a^b B_{i,4}''(x) B_{j,4}''(x)\, dx.
\label{eq:spline-penalty-matrix}
\end{equation}
由于 B-样条基函数具有局部支撑性，$P$ 是一个带状矩阵，非零元素集中在对角线附近的三对角块中。

\textbf{先验协方差}：设 $\Sigma_\beta = \tau^2 P^{-1}$，则先验 \eqref{eq:spline-g-prior} 等价于在惩罚泛函 $\int |f''(x)|^2\, dx$ 上施加 Inverse-Gamma 先验。

\textbf{后验计算}：给定观测数据 $(x_i, y_i)$，后验分布 \eqref{eq:spline-posterior} 可以通过 Cholesky 分解高效计算。具体而言，设 $L$ 为 $H^{\mathsf{T}} H + \Sigma_\beta^{-1}$ 的 Cholesky 因子，则后验抽样可以通过前置-后向代换完成：
\begin{equation}
\boldsymbol{\beta}^* = L^{-\mathsf{T}} \mathbf{z} + L^{-1} H^{\mathsf{T}} \mathbf{y},
\label{eq:spline-posterior-sampling}
\end{equation}
其中 $\mathbf{z} \sim \mathrm{N}(0, I)$。

\subsection{等效核的推导}\label{sec:derivation-equivalent-kernel}

\textbf{背景}：平滑样条的解 \eqref{eq:smoothing-spline-solution} 可以看作是对观测值 $y_i$ 的加权线性组合，权重由等效核 $S_\lambda(x, x_i)$ 给出。本节推导等效核的显式形式。

\textbf{目标}：求满足
\begin{equation}
\hat{f}(x) = \int_0^1 S_\lambda(x, u) y(u)\, du
\label{eq:equivalent-kernel-goal}
\end{equation}
的核函数 $S_\lambda(x, u)$。

\textbf{特征分解方法}：设 $\{\phi_k, \lambda_k\}$ 为惩罚核 $P(x, u) = \sum_j h_j''(x) h_j''(u)$ 的特征函数和特征值，则等效核可以表示为
\begin{equation}
S_\lambda(x, u) = \sum_k \frac{\phi_k(x) \phi_k(u)}{\lambda_k + \lambda}.
\label{eq:equivalent-kernel-eigenfunction}
\end{equation}
将 \eqref{eq:equivalent-kernel-eigenfunction} 代入目标函数并施加极小化条件，即可验证其正确性。

\textbf{边界行为}：在边界点 $x \to 0$ 或 $x \to 1$ 处，等效核被截断并重新标准化，以确保 $\int_0^1 S_\lambda(x, u)\, du = 1$（无偏性）。这解释了平滑样条在边界处的自适应光滑行为。

\subsection{等效核的边界行为}\label{sec:derivation-equivalent-kernel-boundary}

\textbf{背景}：在 \cref{ex:equivalent-kernel-visualization} 中，我们指出等效核在边界点处形状不同。本节分析这一现象的数学机制。

\textbf{Mercer 定理}：惩罚核 $P(x, u)$ 在 $[0, 1]^2$ 上是连续、对称、正定的，因此可以展开为
\begin{equation}
P(x, u) = \sum_{k=1}^{\infty} \mu_k \phi_k(x) \phi_k(u),
\label{eq:penalty-kernel-mercer}
\end{equation}
其中 $\mu_k$ 是特征值，$\phi_k$ 是对应的特征函数，它们满足 $\int_0^1 P(x, u) \phi_k(u)\, du = \mu_k \phi_k(x)$。

\textbf{边界效应的机制}：对于均匀设计 $x_i = i/n$，经验核矩阵 $S_{ij} = S_\lambda(x_i, x_j)$ 是 $P(x, u)$ 的离散近似。在边界点 $i=1$ 或 $i=n$ 处，$S$ 的行向量被截断，导致有效邻域减小，等效核的方差增大。

这与核密度估计中的边界效应本质相同，但平滑样条通过数据驱动的 $\lambda$ 选择部分缓解了这个问题——在边界处 $\lambda$ 自动增大，使得估计更加保守。

\subsection{Haar 小波系数计算}\label{sec:derivation-wavelet-coefficient}

\textbf{背景}：在 \cref{ex:haar-wavelet} 中，我们给出了 Haar 小波系数的直观描述。本节给出严格的数学推导。

\textbf{Haar 尺度函数和小波}：
\begin{align}
\phi(x) &= \mathbb{I}_{[0,1]}(x), \label{eq:haar-scaling} \\
\psi(x) &= \mathbb{I}_{[0,1/2)}(x) - \mathbb{I}_{[1/2,1)}(x). \label{eq:haar-wavelet-function}
\end{align}
在尺度 $j$ 上，伸缩平移后的基函数为
\begin{align}
\phi_{j,k}(x) &= 2^{j/2} \phi(2^j x - k), \label{eq:haar-scaled-scaling} \\
\psi_{j,k}(x) &= 2^{j/2} \psi(2^j x - k), \label{eq:haar-scaled-wavelet}
\end{align}
其中 $k = 0, 1, \ldots, 2^{-j} - 1$。

\textbf{正交性验证}：利用 $\phi$ 和 $\psi$ 的定义，直接计算内积可以验证
\begin{equation}
\langle \phi_{j,k}, \phi_{j',k'} \rangle = \delta_{jj'} \delta_{kk'},\quad
\langle \psi_{j,k}, \psi_{j',k'} \rangle = \delta_{jj'} \delta_{kk'},\quad
\langle \phi_{j,k}, \psi_{j',k'} \rangle = 0.
\label{eq:haar-orthogonality}
\end{equation}

\textbf{系数计算}：对于任意 $L^2[0,1]$ 函数 $f$，其小波展开为
\begin{equation}
f(x) = \sum_k \alpha_{0,k} \phi_{0,k}(x) + \sum_{j=0}^{\infty} \sum_k \beta_{jk} \psi_{jk}(x),
\label{eq:wavelet-expansion}
\end{equation}
其中系数由内积给出：
\begin{align}
\alpha_{0,k} &= \int_0^1 f(x) \phi_{0,k}(x)\, dx, \label{eq:haar-scale-coefficient} \\
\beta_{jk} &= \int_0^1 f(x) \psi_{jk}(x)\, dx. \label{eq:haar-wavelet-coefficient}
\end{align}
对于分段常数函数，这些系数可以直接计算；对于一般光滑函数，高尺度系数 $\beta_{jk}$ 随 $j$ 增大而快速衰减。

\textbf{Daubechies 小波的构造原理}：Daubechies 小波通过要求消失矩条件来提高光滑性。设 $\psi$ 具有 $N$ 阶消失矩，即 $\int x^m \psi(x)\, dx = 0$ for $m = 0, 1, \ldots, N-1$，则对应的尺度函数需要满足更复杂的两尺度关系，其支撑长度随 $N$ 增加而增大。

\endinput

% Chapter 21: Gaussian Process Models
% Bayesian Data Analysis - Gelman et al.

\chapter{高斯过程模型}

\section{本章导论}

在前两章中，我们分别探讨了参数非线性模型（Chapter 19）和基函数模型（Chapter 20）。基函数模型通过线性组合一组固定的基函数（如样条基、小波基）来近似未知函数。这种方法的核心假设是：函数可以用少数基函数的加权线性组合很好地逼近。然而，基函数的选择——无论是样条节点的位置还是小波基的层次——本质上仍然是一种参数化决定。当数据的特征尺度事先未知时，选择恰当的基函数往往并不容易。

高斯过程（Gaussian Process，简称 GP）提供了一种原则上不同的思路：\textbf{我们不再指定基函数的形式，而是直接对函数空间本身施加概率结构}。换言之，我们把整个函数 $f(x)$ 视为一个随机对象，假设它的值（在任意有限多个点上）服从联合正态分布。这种假设的自然之处在于：正态分布是概率论中最自然的"默认"分布——在只知道均值和方差的条件下，正态分布是在所有分布中熵最大的分布，因此代表了最大的不确定性。

从历史脉络看，高斯过程的研究可追溯到Kolmogorov（1940年代）在平稳过程方面的工作，但在统计学中的广泛应用主要始于1990年代。O'Hagan（1978）首次在贝叶斯分析中系统使用高斯过程；Neal（1997）揭示了高斯过程与神经网络的深层联系；Rasmussen 和 Williams（2006）的专著将高斯过程塑造为机器学习的核心工具之一。在贝叶斯数据分析的框架下，高斯过程既是灵活的先验分布，又是可计算的后验推断对象——这是本章要揭示的核心主题。

\section{高斯过程回归}

\subsection{高斯过程的定义}

设 $f(x)$ 是定义在输入空间 $\mathcal{X}$ 上的函数。称 $f$ 服从高斯过程，如果对于任意有限的输入点 $x_1, \ldots, x_n \in \mathcal{X}$，向量
\[
\bigl(f(x_1), f(x_2), \ldots, f(x_n)\bigr)^\top
\]
的联合分布是多元正态分布。这个定义的核心是：\textbf{我们不对函数形式做参数化假设，而是对函数值之间的协方差结构做假设}。

高斯过程由两个函数完全确定：

\begin{enumerate}
  \item \textbf{均值函数（Mean function）}：$m(x) = \mathbb{E}[f(x)]$
  \item \textbf{协方差函数（Covariance function / Kernel）}：$k(x, x') = \mathrm{cov}(f(x), f(x'))$
\end{enumerate}

我们记作 $f(x) \sim \mathrm{GP}(m(x), k(x, x'))$。

协方差函数 $k(x, x')$ 编码了关于函数的重要先验假设。最常用的协方差函数是\textbf{平方指数（Squared Exponential，SE）}核：

\begin{equation}
k_{\mathrm{SE}}(x, x') = \sigma^2 \exp\left(-\frac{(x - x')^2}{2\ell^2}\right), \label{eq:se-kernel}
\end{equation}

其中 $\sigma^2$ 是信号的方差幅度（amplitude），$\ell$ 是长度尺度（length scale），控制函数变化的特征距离。\footnote{平方指数核的数学性质及与其他核函数的比较见附录 \cref{sec:gp-kernel-properties}。}

\subsection{回归模型与先验}

给定观测数据 $(x_i, y_i)$，$i = 1, \ldots, n$，我们假设

\begin{equation}
y_i = f(x_i) + \epsilon_i, \quad \epsilon_i \sim \mathrm{N}(0, \sigma_y^2), \label{eq:gp-observation-model}
\end{equation}

其中 $f(x) \sim \mathrm{GP}(m(x), k(x, x'))$ 是潜在函数，$\epsilon_i$ 是独立测量误差。

对于先验均值函数，我们通常取 $m(x) = 0$（中心化先验），这对应于假设 $f$ 在零附近波动。更信息化的先验可以通过设置非零均值函数来实现。

\textbf{超参数}：协方差函数中的参数 $\theta = (\sigma^2, \ell, \sigma_y^2)$ 称为高斯过程的超参数。它们的取值直接影响后验函数的平滑性和尺度。在贝叶斯框架下，我们既可以对这些超参数指定先验并同时推断，也可以在经验贝叶斯框架下通过边缘似然最大化来选择它们的值。

\subsection{后验预测}

设 $f_* = f(x_*)$ 是新输入点 $x_*$ 处的函数值，$y_*$ 是对应的观测。给定训练数据 $\mathcal{D} = \{(x_i, y_i)\}_{i=1}^n$，后验预测分布可以通过解析计算得到——这是高斯过程最重要的计算优势。

设 $K$ 是 $n \times n$ 协方差矩阵，$K_i = k(x_i, x_*)$ 是训练点与新点之间的协方差向量。则后验预测分布为

\begin{equation}
f_* \mid \mathcal{D}, x_* \sim \mathrm{N}\bigl(\mu_*(x_*), \sigma_*^2(x_*)\bigr), \label{eq:gp-posterior-predictive}
\end{equation}

其中

\begin{align}
\mu_*(x_*) &= K_*^\top (K + \sigma_y^2 I)^{-1} y, \label{eq:gp-posterior-mean} \\
\sigma_*^2(x_*) &= k(x_*, x_*) - K_*^\top (K + \sigma_y^2 I)^{-1} K_*. \label{eq:gp-posterior-variance}
\end{align}

后验均值 $\mu_*(x_*)$ 给出了在 $x_*$ 处对函数值的预测；后验方差 $\sigma_*^2(x_*)$ 量化了预测的不确定性——注意它由两项构成：$k(x_*, x_*)$ 是先验方差（$x_*$ 远离所有训练点时的最大不确定性），减去一项由数据带来的信息。\footnote{后验预测分布的完整推导见附录 \cref{sec:gp-posterior-derivation}。}

预测均值的计算涉及一个 $n \times n$ 矩阵的求逆，计算复杂度为 $O(n^3)$。这在高斯过程文献中称为"exact GP"方法，是精度与计算成本之间的基本权衡。

\begin{Remark}
  后验预测方差的解析表达式 \eqref{eq:gp-posterior-variance} 有一个值得注意的性质：即使我们只观测了 $y$ 而非 $f$ 本身，方差公式中的减项（即数据带来的信息）也自动反映了数据对 $f_*$ 预测精度的提升。换言之，\textbf{后验方差的衰减速度由核函数的局部相关性结构决定}，而不是由观测值本身决定。
\end{Remark}

\section{生日与出生日期的例子}\label{sec:gp-birthday-example}

我们用一个人口统计学例子来说明高斯过程回归的实际应用。

\subsection{数据背景}

考虑美国1981-1988年间每日出生人数的数据。设 $y_t$ 是在第 $t$ 天的新生儿数量，$t$ 是从1981年1月1日起的天数。这个数据集有若干有趣的特征：

\begin{enumerate}
  \item \textbf{年度周期}：夏季出生人数通常高于冬季
  \item \textbf{周内周期}：周末出生人数显著低于工作日
  \item \textbf{长期趋势}：出生人数有轻微的逐年下降趋势
  \item \textbf{异常日期}：某些日期（如圣诞节、新年）出生人数异常低
\end{enumerate}

传统的参数化方法需要为上述每个特征单独建模。高斯过程提供了一种更灵活的选择：我们可以选择合适的协方差函数来同时捕捉这些结构。

\subsection{周期协方差函数}

为了捕捉年度周期，我们可以构造一个周期协方差函数：

\begin{equation}
k_{\mathrm{periodic}}(d, d') = \sigma_p^2 \exp\left(-\frac{2 \sin^2(\pi |d - d'| / 365)}{\ell_p^2}\right), \label{eq:periodic-kernel}
\end{equation}

其中 $d$ 是日期（天数），$\ell_p$ 控制周期效应的衰减速度，$\sigma_p^2$ 是周期成分的方差。\footnote{周期协方差函数的构造及与傅里叶分析的关联见附录 \cref{sec:periodic-kernel-derivation}。}

将周期核与其他协方差函数（如SE核捕捉长期趋势、线性核捕捉线性趋势）组合，可以得到一个灵活的多成分协方差函数：

\begin{equation}
k(x, x') = k_{\mathrm{SE}}(x, x') + k_{\mathrm{periodic}}(x, x') + k_{\mathrm{linear}}(x, x') + \sigma_y^2 \delta_{xx'}, \label{eq:gp-composite-kernel}
\end{equation}

其中 $k_{\mathrm{linear}}(x, x') = \sigma_b^2 (x - \bar{x})(x' - \bar{x})$ 是一个线性核，捕捉线性趋势；$\delta_{xx'}$ 是克罗内克 delta 项，对应于观测噪声。

\subsection{后验推断与结果}

通过最大化边缘似然或使用MCMC，我们可以同时估计所有超参数和潜在函数值 $f(t)$。后验预测给出了每日潜在出生趋势的平滑估计，其后验方差在数据稀疏区域（如年末）显著增大。

\begin{Example}[生日数据的验证]\label{ex:gp-birthday-results}
  在上述模型下，圣诞节前后的预测出生人数显著低于其他日期——这反映在潜在函数 $f(t)$ 的负偏移上。后验置信区间自动包含了这一不确定性。周期成分的最大后验估计对应的周期长度约为365天，与经验观察一致。
\end{Example}

\section{潜变量高斯过程模型}\label{sec:latent-gp-models}

高斯过程的一个重要推广是将它作为\textbf{潜变量（latent variable）}嵌入更复杂的模型结构中。这种思路在分类、计数数据和空间建模中有着广泛应用。

\subsection{分类问题中的应用}

设 $y_i \in \{0, 1\}$ 是二元分类结果，$x_i$ 是 $p$ 维协变量。我们可以构造以下潜变量模型：

\begin{align}
f_i &= f(x_i) \sim \mathrm{GP}(0, k(x_i, x_j)), \label{eq:gp-class-latent} \\
\mathbb{P}(y_i = 1) &= \mathrm{logit}^{-1}(f_i), \label{eq:gp-class-link}
\end{align}

其中 $f(x)$ 是潜在高斯过程，logistic 变换将实数值 $f_i$ 映射到概率。这个模型称为\textbf{高斯过程分类（Gaussian Process Classification，GPC）}。

由于 logistic 变换的非线性性，后验分布 $p(f | y)$ 不再有解析形式。近似方法包括：

\begin{enumerate}
  \item \textbf{拉普拉斯近似（Laplace approximation）}：用高斯分布近似后验
  \item \textbf{期望传播（Expectation Propagation，EP）}：迭代地近似每个数据点的似然贡献
  \item \textbf{马尔科夫链蒙特卡洛（MCMC）}：对潜变量和超参数进行后验抽样
\end{enumerate}

\textbf{后验预测}：给定新点 $x_*$，我们需要计算
\[
\mathbb{P}(y_* = 1 | y) = \int \mathrm{logit}^{-1}(f_*) p(f_* | f, y) p(f | y) \, df_* \, df.
\]
这个积分同样没有解析解，需要使用近似方法。

\begin{Remark}
  GPC 的一个关键计算瓶颈是：每次超参数更新都需要重新计算协方差矩阵 $K$ 和它的逆。对于大规模数据集，$O(n^3)$ 的计算复杂度成为瓶颈。稀疏近似（sparse approximation）通过引入 $m \ll n$ 个诱导点（inducing points）将计算复杂度降为 $O(n m^2)$。详见附录 \cref{sec:sparse-gp-derivation}。
\end{Remark}

\subsection{计数数据的应用}

对于计数数据 $y_i \in \{0, 1, 2, \ldots\}$，我们可以使用对数线性模型：

\begin{align}
\lambda_i &= \exp(f(x_i)), \quad f \sim \mathrm{GP}(0, k), \label{eq:gp-poisson-link} \\
y_i &\sim \mathrm{Poisson}(\lambda_i). \label{eq:gp-poisson-likelihood}
\end{align}

这个模型等价于将高斯过程置于对数-link 函数之后。在生态学中，这称为\textbf{对数高斯 Cox 过程（Log-Gaussian Cox Process）}，用于空间点格局（spatial point pattern）的建模。

\section{函数数据分析}\label{sec:functional-data-analysis}

高斯过程为函数数据分析（Functional Data Analysis，FDA）提供了一种自然的贝叶斯框架。在 FDA 中，每个观测单位（如一个患者）记录的是整个函数曲线 $y_i(t)$，而非单个标量值。

\subsection{函数数据的层级模型}

设 $y_i(t)$ 是第 $i$ 个个体的函数型观测，$t$ 通常是时间或空间坐标。我们假设

\begin{align}
y_i(t) &= f_i(t) + \epsilon_i(t), \label{eq:fda-observation} \\
f_i(t) &= \mu(t) + g_i(t), \label{eq:fda-mean-decomposition} \\
g_i(t) &\sim \mathrm{GP}(0, k_g), \quad \epsilon_i(t) \sim \mathrm{GP}(0, \sigma_\epsilon^2 \delta), \label{eq:fda-gp-prior}
\end{align}

其中 $\mu(t)$ 是总体均值函数，$g_i(t)$ 是个体随机效应（捕获个体间的变异），$\epsilon_i(t)$ 是测量误差。

这个模型的优势在于：\textbf{高斯过程先验同时提供了函数空间的平滑性约束和数据驱动的平滑度估计}。个体效应 $g_i(t)$ 的后验均值给出了每个个体的平滑轨迹估计，其后验方差自动反映了样本量（观测时间点数量）对估计精度的影响。

\subsection{主成分分析与高斯过程的联系}

在经典的 FDA 中，主成分分析（PCA）用于降维——将无限维函数投影到少数主成分方向上。有趣的是，\textbf{高斯过程协方差函数的主成分分解与经典 PCA 之间存在深层对应关系}。

具体而言，协方差算子 $K(s, t) = \mathrm{cov}(f(s), f(t))$ 的特征函数给出了函数空间中最主要的变异模式。这在函数型数据的降维表示中有着直接应用。\footnote{高斯过程与函数型PCA的数学联系见附录 \cref{sec:gp-fda-connection}。}

\section{密度估计与回归}\label{sec:gp-density-estimation}

高斯过程的灵活性还体现在密度估计和回归问题中。通过对潜在函数施加适当的约束，高斯过程可以用于建模非负函数（密度）和单调函数。

\subsection{变换高斯过程}

设我们希望估计一个未知的一维密度函数 $p(x)$。一个自然的想法是先对 $\log p(x)$ 建模：

\[
\log p(x) = f(x), \quad f \sim \mathrm{GP}(m(x), k(x, x')).
\]

为保证 $p(x) \geq 0$ 且 $\int p(x) dx = 1$，我们需要对 $f$ 施加约束。一种方法是在有限网格上离散化，然后用 softmax 变换确保归一化：

\[
p_j = \frac{\exp(f_j)}{\sum_{k} \exp(f_k)}.
\]

在贝叶斯框架下，我们可以对 $f$ 指定高斯过程先验，然后使用 MCMC 抽样后验分布。\footnote{变换高斯过程密度估计的完整推断算法见附录 \cref{sec:gp-density-estimation-derivation}。}

另一种更直接的方法是对累积分布函数（CDF）建模：对 $F(x) = \mathbb{P}(X \leq x)$ 假设高斯过程先验（保证 $F$ 单调和取值于 $[0,1]$），然后通过微分得到密度 $p(x) = F'(x)$。

\subsection{高斯过程回归的边界约束}

在某些应用中，我们需要对回归函数施加单调性约束（$f'(x) \geq 0$）或其他边界约束。高斯过程的条件约束可以通过\textbf{镜像技术（reflection method）}或\textbf{回归方法（regression method）}实现。\footnote{单调高斯过程的构造与计算见附录 \cref{sec:monotone-gp-derivation}。}

\section{高斯过程的计算复杂度}\label{sec:gp-computation}

高斯过程的核心计算瓶颈是协方差矩阵的求逆。对于 $n$ 个训练点，精确计算需要 $O(n^3)$ 时间和 $O(n^2)$ 存储。

\subsection{大规模数据的稀疏近似}

当 $n$ 达到数万甚至更多时，精确高斯过程变得不可行。稀疏近似方法的核心思想是用 $m$ 个代表性诱导点（inducing points）来近似完整协方差矩阵，从而将 $O(n^3)$ 降为 $O(n m^2)$。具体而言，我们引入诱导输入点 $Z = \{z_j\}_{j=1}^m$ 和对应的函数值 $u_j = f(z_j)$，然后通过 $p(f_* | y) \approx \int p(f_* | u) q(u | y) \, du$ 来近似后验预测。\footnote{稀疏高斯过程近似的详细推导见附录 \cref{sec:sparse-gp-derivation}。}

代表性方法包括 FITC（Fully Independent Training Conditional）、PITC（Partially Independent Training Conditional）和 Nystr\"om 近似。

\subsection{梯度优化与超参数估计}

超参数 $\theta$ 的边缘似然（log-marginal likelihood）为

\begin{equation}
\log p(y | \theta) = -\frac{1}{2} y^\top (K_\theta + \sigma_y^2 I)^{-1} y - \frac{1}{2} \log |K_\theta + \sigma_y^2 I| - \frac{n}{2} \log 2\pi. \label{eq:gp-marginal-likelihood}
\end{equation}

边缘似然对超参数 $\theta$ 的梯度可以通过自动微分或有限差分计算。最大化边缘似然等价于经验贝叶斯估计。\footnote{边缘似然的梯度推导见附录 \cref{sec:gp-marginal-gradient}。}

\section{高斯过程与神经网络的联系}\label{sec:gp-nn-connection}

一个引人注目的理论结果是：\textbf{宽度趋于无穷的神经网络（neural network）收敛于一个高斯过程}。这一发现由Neal（1997）首先证明，随后在无限宽极限与神经网络的贝叶斯后验之间建立了精确对应。

具体而言，设一个全连接神经网络的权重和偏置独立同分布，当网络宽度（每层神经元数）趋于无穷时，网络输出的先验分布趋向于高斯过程。这意味着：

\begin{enumerate}
  \item 无限宽神经网络的函数空间先验等价于某个高斯过程先验
  \item 网络权重的贝叶斯后验预测在无限宽极限下等价于对应 GP 的预测
  \item GP 的协方差函数由网络的激活函数和权重初始化方差决定
\end{enumerate}

这个联系为理解神经网络的函数空间动力学提供了概率论框架，也为使用 GP 作为神经网络不确定性的替代模型提供了理论基础。

\section{本章小结}

本章我们系统介绍了高斯过程模型：

\begin{enumerate}
  \item 高斯过程是对函数空间的正态分布建模，为回归、分类、密度估计等问题提供了灵活的先验分布
  \item 协方差函数（核函数）编码了关于函数平滑性、周期性和尺度等重要先验假设
  \item 高斯过程回归的后验预测分布有解析形式，这是其最重要的计算优势
  \item 周期核、SE 核等可以组合成复合协方差函数，捕捉数据中的多尺度结构
  \item 潜变量高斯过程模型可以推广到分类、计数等非高斯数据
  \item 函数数据分析可以从高斯过程的视角统一理解
  \item 变换高斯过程可以用于密度估计和单调回归
  \item 计算复杂度 $O(n^3)$ 是精确高斯过程的主要瓶颈，稀疏近似方法在大规模数据中广泛应用
  \item 高斯过程与神经网络的深层联系揭示了无限宽网络的函数空间结构
\end{enumerate}

\section{下节预告}

下一章我们将学习有限混合模型（Finite Mixture Models），这是另一类灵活的非参数化方法。混合模型通过将分布表示为若干成分的加权混合来捕捉数据的异质性——这与高斯过程关注函数空间的连续建模形成互补。

\section{习题}\label{sec:chapter21-exercises}

\begin{Exercise}{\ref{exr:21-1} 高斯过程后验}\label{exr:21-1}
设 $f \sim \mathrm{GP}(0, k)$，观测模型为 $y = f + \epsilon$，其中 $\epsilon \sim \mathrm{N}(0, \sigma_y^2 I)$。证明后验预测分布的均值和方差满足 \cref{eq:gp-posterior-mean} 和 \cref{eq:gp-posterior-variance}。
\end{Exercise}

\begin{Exercise}{\ref{exr:21-2} 协方差函数的性质}\label{exr:21-2}
证明平方指数核 $k_{\mathrm{SE}}(x, x')$ 是正定的，即对任意实数 $a_1, \ldots, a_n$ 和点 $x_1, \ldots, x_n$，有 $\sum_{i,j} a_i a_j k_{\mathrm{SE}}(x_i, x_j) \geq 0$。
\end{Exercise}

\begin{Exercise}{\ref{exr:21-3} 周期核的构造}\label{exr:21-3}
设 $k_{\mathrm{base}}(x, x')$ 是一个协方差函数。证明 $k_{\mathrm{periodic}}(x, x') = k_{\mathrm{base}}(\sin(\pi (x - x') / T), \cos(\pi (x - x') / T))$ 产生一个以 $T$ 为周期的协方差函数。
\end{Exercise}

\begin{Exercise}{\ref{exr:21-4} 边缘似然的梯度}\label{exr:21-4}
推导边缘似然 \cref{eq:gp-marginal-likelihood} 对超参数 $\theta$ 的梯度公式，并说明如何在实际计算中利用协方差矩阵的 Cholesky 分解提高数值稳定性。
\end{Exercise}

\begin{Exercise}{\ref{exr:21-5} 稀疏近似}\label{exr:21-5}
描述诱导点（inducing points）方法的直觉，并比较 FITC 和 PITC 两种稀疏近似策略在计算复杂度和近似精度上的差异。
\end{Exercise}

\endinput

% Chapter 22: Finite Mixture Models
% Bayesian Data Analysis - Gelman et al.

\chapter{有限混合模型}

\section{本章导论}

在前面各章中，我们学习的模型大多假设数据来自某个特定的参数分布族——例如正态分布、二项分布或泊松分布。然而，现实世界的数据往往远比这复杂：一个总体可能包含多个异质性子群体，而每个子群体内部的行为相对简单。这种\textbf{异质性}（heterogeneity）在社会科学、生物统计学、金融和机器学习中无处不在。

如何为这类异质性数据构建统计模型？一个自然的想法是：将总体视为若干个简单分布的混合，每个简单分布代表一个潜在的子群体。这就是\textbf{混合模型}（mixture model）的核心思想。

本章研究有限混合模型（finite mixture models），即假设总体由有限个分量（component）混合而成的模型。我们将讨论：
\begin{enumerate}
  \item 混合模型的动机与应用场景
  \item 有限混合分布的数学形式
  \item 贝叶斯框架下的推断方法
  \item 标签切换问题及其处理
  \item 实际应用案例
\end{enumerate}

\section{异质性数据与混合模型的引入}

\subsection{为什么需要混合模型？}

考虑以下实际场景：

一个天文数据库中记录了星系的退行速度（redshift）。物理学家知道星系形成于不同的演化阶段，因此观测到的速度分布可能不是单峰的正态分布，而是多个峰值——每个峰值对应一类星系。

或者，考虑医学研究中的生存时间数据：患者在接受治疗后，有些人对治疗反应良好（生存时间较长），有些人对治疗反应不佳（生存时间较短）。如果我们忽略这种异质性，用单一分布建模，可能会得到误导性的结论。

在这些场景中，数据的异质性来源于\textbf{未观测到的分组结构}（unobserved grouping structure）。混合模型提供了一种原则性的方法来刻画这种结构。

\subsection{从潜在类别到混合模型}

设我们有一组数据 $y = (y_1, \ldots, y_n)$。假设每个观测 $y_i$ 属于某个潜在的类别 $z_i \in \{1, 2, \ldots, K\}$，其中 $K$ 是类别总数。类别 $z_i$ 是未被观测到的（latent）。

给定类别 $z_i = k$ 时，$y_i$ 的分布为某个简单的参数分布：
\[
y_i \mid z_i = k \sim f_k(y_i \mid \theta_k)
\]
其中 $f_k$ 是第 $k$ 个分量的密度函数，$\theta_k$ 是对应的参数。

类别变量 $z_i$ 本身服从多项分布：
\[
z_i \sim \text{Categorical}(\pi_1, \ldots, \pi_K)
\]
其中 $\pi_k = \mathbb{P}(z_i = k)$ 是第 $k$ 个类别的先验概率。

将潜在类别边缘化，我们得到观测数据的边际分布：
\[
p(y_i) = \sum_{k=1}^{K} \pi_k \cdot f_k(y_i \mid \theta_k)
\label{eq:finite-mixture-marginal}
\]

这就是\textbf{有限混合分布}（finite mixture distribution）的数学形式。

\section{有限混合分布的形式与性质}

\subsection{混合分布的基本形式}

有限混合模型的核心是\cref{eq:finite-mixture-marginal} 所示的加权求和形式。权重 $\pi_k$ 满足 $\pi_k \geq 0$ 且 $\sum_{k=1}^{K} \pi_k = 1$，因此构成一个概率分布。

\begin{Definition}[有限混合分布]
  设 $f_1(y), \ldots, f_K(y)$ 为 $K$ 个密度函数（或质量函数），$\theta = (\pi_1, \ldots, \pi_K, \theta_1, \ldots, \theta_K)$ 为参数向量，其中 $\pi_k \geq 0$，$\sum_{k=1}^K \pi_k = 1$。则由
  \[
  p(y \mid \theta) = \sum_{k=1}^{K} \pi_k \cdot f_k(y \mid \theta_k)
  \]
  定义的分布称为\textbf{有限混合分布}（finite mixture distribution），$K$ 称为混合分量数（number of components）。
\end{Definition}

混合分布的矩可以通过各分量的矩加权平均得到。例如，混合分布的均值为
\[
\mathbb{E}(y) = \sum_{k=1}^{K} \pi_k \cdot \mathbb{E}(y \mid z=k) = \sum_{k=1}^{K} \pi_k \cdot \mu_k
\]
其中 $\mu_k$ 是第 $k$ 个分量的均值。

混合分布的方差更为复杂，需要考虑组内方差和组间方差的叠加。

\subsection{混合正态分布}

最常用的混合分布是\textbf{混合正态分布}（Gaussian mixture model）：
\[
p(y \mid \theta) = \sum_{k=1}^{K} \pi_k \cdot \mathrm{N}(y \mid \mu_k, \sigma_k^2)
\]
其中 $\theta_k = (\mu_k, \sigma_k^2)$。

混合正态分布具有以下性质：
\begin{enumerate}
  \item 任何平滑的密度函数都可以用足够多分量的混合正态分布任意逼近（根据 Stone-Weierstrass 逼近定理）
  \item 混合正态分布本身不再是正态分布，但其对数是凹函数（当 $K > 1$ 时）
  \item 混合正态分布可以表示多峰、偏态等复杂形态
\end{enumerate}

\begin{Example}[双正态混合]
  设 $K=2$，$\pi_1 = \pi_2 = 0.5$，$\mu_1 = -2$，$\sigma_1^2 = 1$，$\mu_2 = 2$，$\sigma_2^2 = 1$。则混合分布为
  \[
  p(y) = 0.5 \cdot \mathrm{N}(y \mid -2, 1) + 0.5 \cdot \mathrm{N}(y \mid 2, 1)
  \]
  这个分布是双峰的，两个峰分别位于 $y = -2$ 和 $y = 2$ 附近。
\end{Example}

混合正态分布在聚类分析（clustering）中应用广泛：每个分量对应一个潜在的类别，我们可以用贝叶斯推断来同时估计类别参数和每个观测的类别归属概率。

\section{EM算法与混合模型的推断}

\subsection{潜在变量的引入}

混合模型的难点在于：潜在类别 $z_i$ 是未被观测到的。如果 $z_i$ 被观测到（即我们已知每个观测属于哪个分量），则问题将大大简化——我们可以分别对每个分量进行推断。

EM算法（Expectation-Maximization）由 Dempster, Laird 和 Rubin \cite{Dempster1977} 于 1977 年提出，正是利用了这种思想：通过迭代地估计潜在变量（E步）和更新参数（M步），来最大化观测数据的似然。

设 $\theta = (\pi, \theta_1, \ldots, \theta_K)$ 为全部参数，其中 $\pi = (\pi_1, \ldots, \pi_K)$。观测数据的对数似然为
\[
\ell(\theta) = \sum_{i=1}^{n} \log \left( \sum_{k=1}^{K} \pi_k \cdot f_k(y_i \mid \theta_k) \right)
\]

这个对数似然包含对潜在变量的求和，导致直接最大化困难。EM算法通过引入潜在变量的后验分布来解决这个问题。

\subsection{E步：计算潜在变量的后验}

在第 $t$ 次迭代，假设当前参数估计为 $\theta^{(t)}$。E步计算给定观测数据和当前参数下，每个观测属于各分量的后验概率：
\[
\gamma_{ik}^{(t)} = \mathbb{P}(z_i = k \mid y_i, \theta^{(t)}) = \frac{\pi_k^{(t)} \cdot f_k(y_i \mid \theta_k^{(t)})}{\sum_{j=1}^{K} \pi_j^{(t)} \cdot f_j(y_i \mid \theta_j^{(t)})}
\label{eq:mixture-responsibility}
\]

$\gamma_{ik}$ 被称为"责任"（responsibility）：第 $k$ 个分量对观测 $y_i$ 的责任程度。

从贝叶斯的角度，$\gamma_{ik}$ 就是给定 $y_i$ 后，$z_i = k$ 的后验概率。

\subsection{M步：更新参数估计}

M步在"完全数据"（观测数据 + 潜在变量）上最大化期望对数似然。完全数据的对数似然为
\[
\ell_c(\theta) = \sum_{i=1}^{n} \sum_{k=1}^{K} \mathbb{I}(z_i = k) \cdot \left[ \log \pi_k + \log f_k(y_i \mid \theta_k) \right]
\]

取期望（关于 $z_i$ 的后验分布），得到期望完全数据对数似然：
\[
Q(\theta, \theta^{(t)}) = \sum_{i=1}^{n} \sum_{k=1}^{K} \gamma_{ik}^{(t)} \cdot \left[ \log \pi_k + \log f_k(y_i \mid \theta_k) \right]
\]

M步通过对 $Q$ 关于 $\theta$ 最大化来更新参数。对于混合正态模型，更新公式为：
\[
\pi_k^{(t+1)} = \frac{1}{n} \sum_{i=1}^{n} \gamma_{ik}^{(t)}
\]
\[
\mu_k^{(t+1)} = \frac{\sum_{i=1}^{n} \gamma_{ik}^{(t)} \cdot y_i}{\sum_{i=1}^{n} \gamma_{ik}^{(t)}}
\]
\[
(\sigma_k^2)^{(t+1)} = \frac{\sum_{i=1}^{n} \gamma_{ik}^{(t)} \cdot (y_i - \mu_k^{(t+1)})^2}{\sum_{i=1}^{n} \gamma_{ik}^{(t)}}
\label{eq:mixture-variance-update}
\]\footnote{方差更新公式的推导见附录 \cref{sec:derivation-em-updates}。}

\subsection{EM算法的收敛性}

EM算法保证观测数据似然的单调递增：$\ell(\theta^{(t+1)}) \geq \ell(\theta^{(t)})$。但收敛到的点可能是局部极大值，而非全局极大值。因此，EM算法对初始值敏感，通常需要多次随机初始化。

对于混合模型的贝叶斯推断，我们通常关注后验分布 $p(\theta, z \mid y)$，而不仅仅是点估计。这需要使用MCMC方法（如Gibbs采样），我们将在下一节讨论。

\section{贝叶斯混合模型推断}

\subsection{完全条件分布}

在贝叶斯框架下，我们为所有参数指定先验分布 $p(\theta)$，然后推断后验分布 $p(\theta, z \mid y)$。

对于混合正态模型，一个方便的选择是共轭先验：
\[
\pi \sim \text{Dirichlet}(\alpha_1, \ldots, \alpha_K)
\]
\[
\mu_k \mid \sigma_k^2 \sim \mathrm{N}(\mu_0, \sigma_k^2 / n_0)
\]
\[
\sigma_k^2 \sim \text{Inverse-Gamma}(\nu_0/2, \nu_0 \sigma_0^2/2)
\]

给定这些先验，可以推导出各参数的完全条件分布（Gibbs采样所需）。

第 $k$ 个分量的混合比例 $\pi_k$ 的完全条件分布为：
\[
\pi \mid y, z, \theta_{-\pi} \sim \text{Dirichlet}(\alpha_1 + n_1, \ldots, \alpha_K + n_K)
\]
其中 $n_k = \sum_{i=1}^{n} \mathbb{I}(z_i = k)$ 是属于第 $k$ 个分量的观测数。

第 $k$ 个分量的均值和方差的完全条件分布为正态-逆Gamma分布。

潜在类别变量 $z_i$ 的完全条件分布为：
\[
\mathbb{P}(z_i = k \mid y_i, \theta) \propto \pi_k \cdot \mathrm{N}(y_i \mid \mu_k, \sigma_k^2)
\label{eq:z-complete-conditional}
\]

这就是\cref{eq:mixture-responsibility} 中的后验概率——在贝叶斯版本中，它是完全条件分布，而非EM算法中的点估计。

\subsection{Gibbs采样算法}

基于完全条件分布，我们可以设计Gibbs采样器来从后验分布中抽样：

\begin{enumerate}
  \item 初始化参数 $\theta^{(0)}$
  \item 对 $t = 1, 2, \ldots$ 迭代：
  \begin{enumerate}
    \item 对每个 $i = 1, \ldots, n$，从 \cref{eq:z-complete-conditional} 抽取 $z_i^{(t)}$
    \item 从 $\pi$ 的完全条件分布抽取 $\pi^{(t)}$
    \item 对每个 $k = 1, \ldots, K$，从 $\mu_k, \sigma_k^2$ 的完全条件分布抽取
  \end{enumerate}
\end{enumerate}

Gibbs采样器的收敛性可以通过监控多链的 Gelman-Rubin 统计量 \cite{GelmanRubin1992} 来诊断。

\section{标签切换问题}

\subsection{问题的本质}

混合模型的一个核心问题是\textbf{标签切换}（label switching）：由于似然函数在分量标签的置换下保持不变，后验分布是关于分量标签对称的。

具体来说，如果 $(z_i, \pi_k, \theta_k)$ 是一个合法的后验样本，那么对任意置换 $\sigma$，$(\sigma(z_i), \pi_{\sigma(k)}, \theta_{\sigma(k)})$ 同样是合法的后验样本。

这导致两个问题：
\begin{enumerate}
  \item 后验均值失去意义：$K$ 个分量的后验均值相互混合
  \item 迭代收敛后，每次迭代的标签可能不一致
\end{enumerate}

\subsection{标签切换的解决}

解决标签切换问题的方法主要有：

\begin{enumerate}
  \item \textbf{排序约束}：通过对参数施加排序约束（如 $\mu_1 < \mu_2 < \cdots < \mu_K$）来打破对称性。这是最简单的方法，但可能不适合所有场景。
  \item \textbf{标签切换检测与对齐}：在MCMC采样后，检测标签切换事件，并通过对齐链来恢复因果顺序。
  \item \textbf{后处理}：使用递归分区或聚类方法将后验样本分类到各个分量。
\end{enumerate}

排序约束是最常用的方法，但需要谨慎使用，因为人为的约束可能影响推断结果。

\section{应用实例：星系速度数据}

本节我们用混合正态模型分析星系退行速度数据。\footnote{数据来源见 \cite{ROBERT1995}。}

设 $y_i$ 为第 $i$ 个星系的退行速度（单位：km/s）。我们用双正态混合模型：
\[
y_i \sim \pi \cdot \mathrm{N}(\mu_1, \sigma_1^2) + (1-\pi) \cdot \mathrm{N}(\mu_2, \sigma_2^2)
\]

贝叶斯推断的目标是：
\begin{enumerate}
  \item 估计两个星族（population）的速度分布参数
  \item 计算每个星系属于各星族的后验概率
  \item 评估混合比例 $\pi$
\end{enumerate}

使用Gibbs采样，我们可以从后验分布中抽取样本，然后计算各参数的后验均值和置信区间。

\begin{Remark}
  混合模型的一个优势是它提供了\textbf{软聚类}（soft clustering）：与硬聚类（每个观测必须属于某一类）不同，混合模型给出每个观测属于各类的概率。这对于边界情况特别有价值。
\end{Remark}

\section{混合模型的其他应用}

混合模型在统计学和机器学习中有广泛的应用：

\begin{enumerate}
  \item \textbf{密度估计}：混合模型可以用于非参数密度估计，用有限个参数的混合分布逼近任意复杂的密度函数。
  \item \textbf{聚类分析}：如上所述，混合模型提供了一种基于概率的聚类方法。
  \item \textbf{缺失数据插补}：在存在缺失值的问题中，可以将缺失数据的插补值视为潜在变量，用混合模型进行多重插补。
  \item \textbf{生存分析}：考虑未观测到的异质性来源（如未观测到的风险因素），混合模型可以用于建模生存数据中的非比例风险。
  \item \textbf{时间序列}：隐马尔可夫模型（Hidden Markov Model）是混合模型在时间序列中的推广。
\end{enumerate}

\section{混合模型的扩展}

有限混合模型是更广泛模型类别的特例：

\begin{enumerate}
  \item \textbf{无限混合模型}（Infinite mixture models）：当分量数 $K \to \infty$ 时，得到非参数贝叶斯方法，如狄利克雷过程混合（Dirichlet process mixture）。这将在第23章讨论。
  \item \textbf{半参数混合模型}：放松对分量分布的假设，允许非参数形式的分量分布。
  \item \textbf{高斯过程混合}：在函数空间中进行混合，用于非参数回归和分类。
\end{enumerate}

\section{本章小结}

本章我们学习了有限混合模型：

\begin{enumerate}
  \item 混合模型用于建模具有未观测到的异质性的数据
  \item 有限混合分布是各分量分布的加权平均
  \item EM算法提供了一种寻找最大似然估计的迭代方法
  \item 贝叶斯推断使用Gibbs采样从后验分布中抽样
  \item 标签切换问题需要特殊处理
  \item 混合模型在聚类、密度估计、生存分析等领域有广泛应用
\end{enumerate}

\section{延伸阅读}

关于混合模型的更多信息，推荐以下参考文献：

\begin{itemize}
  \item \cite{McLachlan2007}：混合模型的经典教材
  \item \cite{Fruwirth2006}：有限混合模型的统计推断
  \item \cite{Robert1995}：贝叶斯混合模型的方法与应用
\end{itemize}

\endinput

% Chapter 23: Dirichlet Process Models
% Bayesian Data Analysis - Gelman et al.

\chapter{狄利克雷过程模型}

\section{本章导论}

在前几章中，我们探讨了有限混合模型——用有限个组件的加权组合来刻画异质数据。在实践中，分析者必须事先指定成分的数量 $K$，这往往是一个困难的选择。Dirichlet 过程（Dirichlet Process，DP）提供了一种优雅的解决方案：允许无限多个潜在成分，让数据自适应地确定有效成分的个数。

Dirichlet 过程是贝叶斯非参数统计的核心工具。它由 Ferguson (1973) 引入，是一种随机过程，其样本路径almost surely（以概率 1）是离散概率分布。这一看似特殊的性质——连续底空间上的随机概率分布却以概率 1 取离散值——正是 DP 成为非参数贝叶斯推断强大先验的根源。

本章的结构如下：先从有限混合模型与 DP 的联系入手，介绍 DP 的三种等价构造（stick-breaking、Pólya urn、Chinese restaurant process），然后讨论 DP 作为未知分布的先验及其后验推断，最后介绍 DP 混合模型的计算方法。

\section{从有限混合到无限混合}

在有限混合模型中，我们假设数据来自 $K$ 个成分的混合：
\[
p(y_i | \theta_1, \ldots, \theta_K, \pi) = \sum_{k=1}^K \pi_k \, p(y_i | \theta_k),
\]
其中 $\pi = (\pi_1, \ldots, \pi_K)$ 是混合权重，满足 $\pi_k \ge 0$，$\sum_{k=1}^K \pi_k = 1$。

一个自然的问题是：$K$ 本身能否从数据中学习，而不是预先固定？一个激进的思路是直接令 $K \to \infty$，但这在参数模型框架下并不容易处理。DP 提供了一种优雅的解决方案：先验性地将混合成分的数量设为无限，然后通过后验推断自动确定有效成分个数。

Dirichlet 过程先验的核心思想是：不去直接参数化混合权重 $\pi_k$，而是假设混合权重本身来自一个随机过程，使得"有多少个有效成分"由数据决定。

\section{狄利克雷过程的定义与构造}

\subsection{Stick-breaking 构造}

设 $G_0$ 为一个基分布（base measure），$\alpha > 0$ 为浓度参数（concentration parameter）。Dirichlet 过程 $G \sim \DP(\alpha, G_0)$ 可以通过以下 stick-breaking 构造（Sethuraman, 1994）生成：

\begin{Definition}[Stick-breaking 构造]
令 $V_1, V_2, \ldots$ 为独立同分布的 Beta 随机变量，$V_k \sim \mathrm{Beta}(1, \alpha)$。定义
\[
w_k = V_k \prod_{j=1}^{k-1} (1 - V_j), \quad k = 1, 2, \ldots
\]
则
\[
G = \sum_{k=1}^{\infty} w_k \, \delta_{\theta_k},
\]
其中 $\theta_1, \theta_2, \ldots$ 独立同分布自 $G_0$，且与 $(V_1, V_2, \ldots)$ 相互独立。
\end{Definition}

直观理解：取一根单位长度的棍子（stick），第一步折下长度为 $V_1$ 的第一段，分配给atom $\theta_1$；剩余长度为 $1-V_1$ 的棍子上，折下长度为 $(1-V_1)V_2$ 的第二段，分配给 $\theta_2$；依此类推。由于 $\prod_{j=1}^{\infty}(1-V_j) = 0$ with probability 1，这个过程永远不会终止，因此产生了无限多个atoms，但权重之和为 1。

权重 $w_k$ 的期望为 $\mathbb{E}w_k = \frac{1}{1+\alpha}\left(\frac{\alpha}{1+\alpha}\right)^{k-1}$，随 $k$ 指数衰减。$\alpha$ 越大，衰减越慢，有效成分个数越多。

\subsection{Pólya urn 构造}

Stick-breaking 构造从"折棍子"的角度直观理解了 DP 权重的来源。另一种等价的概率视角来自 Pólya urn scheme——一种经典的随机 urn 过程。设我们有颜色（atoms）集合，urn 中每次加入一个球：

\begin{enumerate}
  \item 以概率 $\alpha/(\alpha + n-1)$ 加入一个新颜色的球（颜色为从 $G_0$ 中独立抽取的 $\theta$）；
  \item 以概率 $(n_k)/( \alpha + n - 1)$ 加入一个颜色为 $k$ 的已有球（$n_k$ 为 urn 中颜色 $k$ 的球数）。
\end{enumerate}

这个 urn 过程产生的颜色分布恰好与 stick-breaking 构造等价。

\subsection{Chinese restaurant process（中国餐厅过程）}

第三种构造与中国餐厅过程的隐喻密切相关。

想象一个拥有无限多张桌子的餐厅，第一位顾客随意选择一张桌子就座。第 $n$ 位顾客到来时：
\begin{itemize}
  \item 以概率 $\alpha/(\alpha + n - 1)$ 选择一张新桌子；
  \item 以概率 $n_k/(\alpha + n - 1)$ 选择已有桌子 $k$（$n_k$ 为该桌已坐的顾客数）。
\end{itemize}

顾客坐在不同桌子上的概率分布恰好描述了 DP 样本中atoms的分配结构。同一桌子上的顾客可以理解为来自同一个聚类（cluster）。这就是著名的 Chinese restaurant process（CRP）。

CRP 的魅力在于它将离散概率分布的随机结构编码为一个sequential partition 过程——这与统计学中聚类分析的精神高度一致：数据自然地聚集成簇，而簇的数量本身由数据决定。

这三种构造实际上是等价的：stick-breaking 给出了权重的显式分布，urn 过程描述了 atoms 的依次分配，而 CRP 则是 urn 过程在聚类视角下的另一种表述。Ferguson (1973) 最初通过有限 Dirichlet 分布的类比构造了 DP；Sethuraman (1994) 的 stick-breaking 形式和 Blackwell-MacQueen (1973) 的 urn 形式则给出了更便于计算的等价表述。

\begin{Remark}
CRP 中的参数 $\alpha$ 控制聚类个数的"先验期望"。$\alpha$ 越小，倾向于更少的聚类；$\alpha$ 越大，倾向于更多的聚类。这与直觉一致：$\alpha$ 作为浓度参数，在 stick-breaking 中控制了权重衰减的速度。
\end{Remark}

\section{DP 作为未知分布的先验}

DP 的核心应用是作为未知分布 $F$ 的先验。设我们希望从分布 $F$ 中抽取独立样本 $x_1, \ldots, x_n$，但对 $F$ 的形式一无所知。

一个非参数的思路是：对 $F$ 赋予先验，然后在后验下进行推断。DP 正是这样一个先验：
\[
F \sim \DP(\alpha, G_0).
\]

具体而言，我们假设：
\[
\theta_i | F \stackrel{\mathrm{iid}}{\sim} F, \quad F \sim \DP(\alpha, G_0),
\]
其中 $G_0$ 是一个参数化的分布（如正态分布），$\alpha$ 控制先验的灵活性。

由于 $G$ 的样本路径是离散分布，我们可以将模型等价地写为：
\[
\theta_i \stackrel{\mathrm{iid}}{\sim} G, \quad G \sim \DP(\alpha, G_0).
\]
但更常见的是直接积分掉 $G$，得到 $\theta_i$ 的边缘分布——这正是下一节要讨论的折棍（stick-breaking）表示的核心应用。

\section{后验推断}

给定观测数据 $\mathbf{x} = (x_1, \ldots, x_n)$，DP 先验下的后验推断有一个优雅的形式。

假设基分布 $G_0$ 是某个参数 $\phi$ 的分布，即 $x_i | \phi \sim p(x | \phi)$，且 $\phi \sim G_0$。那么后验有如下形式：\footnote{推导见附录 \cref{sec:dp-posterior-derivation}}
\[
G | \mathbf{x} \sim \DP\!\left(\alpha + n, \frac{\alpha}{\alpha + n} G_0 + \frac{n}{\alpha + n} \cdot \frac{1}{n}\sum_{i=1}^n \delta_{x_i}\right).
\]

这一结果表明：后验等价于在先验 $G_0$ 的基础上"添加"了经验分布的 contribution。具体而言：
\begin{itemize}
  \item 有效样本量从 $\alpha$ 增加到 $\alpha + n$；
  \item 后验基分布是 $G_0$ 与经验分布 $\frac{1}{n}\sum_i \delta_{x_i}$ 的加权平均，权重分别为 $\frac{\alpha}{\alpha+n}$ 和 $\frac{n}{\alpha+n}$。
\end{itemize}

预测分布也有简洁形式。设新观测为 $\tilde{x}$，则其后验预测分布为：\footnote{推导见附录 \cref{sec:dp-predictive-derivation}}
\[
p(\tilde{x} | \mathbf{x}) = \frac{\alpha}{\alpha + n} \, G_0(\tilde{x}) + \sum_{k=1}^{K_n} \frac{n_k}{\alpha + n} \, \delta_{\theta_k}(\tilde{x}),
\]
其中 $K_n$ 是有效聚类个数，$n_k$ 是第 $k$ 个聚类中的点数，$\theta_k$ 是对应的聚类参数。

也就是说，新数据点有概率 $\alpha/(\alpha+n)$ 来自基分布 $G_0$（产生一个新聚类），有概率 $n_k/(\alpha+n)$ 落在第 $k$ 个已有聚类中（与该聚类的数据共享参数 $\theta_k$）。这正是 CRP 的 predictive probability。

\section{DP 混合模型}

DP 混合模型（Dirichlet process mixture model, DPM）是 DP 先验在密度估计中的主要应用形式。模型设定为：
\[
x_i | \theta_i \sim p(x | \theta_i), \quad \theta_i | G \stackrel{\mathrm{iid}}{\sim} G, \quad G \sim \DP(\alpha, G_0).
\]

与有限混合模型相比，DPM 的关键区别在于：成分的数量是无限的（由 stick-breaking 过程产生），而有效成分的个数由数据通过后验推断自动确定。

从聚类角度看，DPM 自然地产生了一个 partition——数据被划分为若干聚类，每个聚类内的数据共享同一参数 $\theta_k$。这个 partition 是随机且自适应的，与 CRP 的 sequential partition 精神一致。

\section{计算方法}

DP 混合模型的计算核心挑战在于：stick-breaking 构造涉及无限多个 atoms，而后验中 atoms 的个数同样是无限的。MCMC 方法是处理这一问题的主要工具。

\subsection{Blocked Gibbs 采样器}

一种直接的方法是引入隐含的聚类分配指标 $\mathbf{c} = (c_1, \ldots, c_n)$，其中 $c_i = k$ 表示 $x_i$ 属于第 $k$ 个聚类。

全条件后验分布为：
\[
p(c_i = k | \mathbf{c}_{-i}, \mathbf{x}, \alpha) \propto
\begin{cases}
n_{-i,k} \cdot p(x_i | \theta_k, \mathbf{x}_{-i}) & \text{if } k \text{ is an existing cluster}, \\
\alpha \cdot p(x_i | \theta_{\mathrm{new}}, G_0) & \text{if } k \text{ is a new cluster}.
\end{cases}
\]
其中 $n_{-i,k}$ 是除 $i$ 外属于第 $k$ 个聚类的数据点数。

然后从各自的后验分布中采样 $\theta_k$。这个过程在 $c_i$ 和 $\theta_k$ 之间交替进行，即 blocked Gibbs 采样。

\subsection{Chinese Restaurant Process 采样}

CRP 的视角给出了一种更直觉的采样方案。已知完整数据下的 CRP predictive probability，Gibbs 采样可以直接对聚类分配进行更新，而无需显式处理无限多个 atoms。

具体而言，对于数据点 $i$，计算其属于现有聚类 $k$ 和属于新聚类的概率，然后按概率重新分配。这个方法在高维数据和中等规模数据集上通常效率较高。

\section{应用举例}

\begin{Example}[基于 DP 聚类的密度估计]
假设我们有一组来自混合分布的观测数据，数据中存在两个自然聚类，但分析者事先并不知道聚类的个数。

应用 DPM 模型，设成分分布为正态分布 $N(\mu_k, \sigma^2_k)$，基分布 $G_0$ 取为 Normal-Inverse-Gamma 共轭先验。MCMC 采样将自动：
\begin{itemize}
  \item 在后验中探索不同的聚类个数；
  \item 收敛到数据中真实存在的聚类结构；
  \item 提供每个聚类的参数后验分布。
\end{itemize}
\end{Example}

DP 混合模型在以下场景中尤为有用：
\begin{itemize}
  \item 密度估计：当数据具有未知数量的峰（modes）时；
  \item 聚类分析：当聚类个数事先未知，且聚类结构可能随协变量变化时；
  \item 随机过程建模：如隐马尔可夫模型的状态转移矩阵非参数化推广；
  \item 函数逼近：如 DP 基函数回归（DP mixture of regressions）。
\end{itemize}

\section{与其他非参数方法的联系}

DP 并非唯一的贝叶斯非参数工具。在密度估计中，高斯过程（见 Chapter 21）提供了函数空间上的先验，Dirichlet 过程则提供了分布空间上的先验。

从数学上看，DP 是概率分布空间 $\mathcal{P}$ 上的先验（其样本是离散分布），而高斯过程是函数空间上的先验（其样本是函数曲线）。两者在无限维参数空间中都实现了"让数据自适应确定有效复杂度"的思想，但在应用层面：DP 更适合聚类、密度估计和混合模型；高斯过程更适合函数回归、时空建模。

一个有用的视角是：DP 是有限维 Dirichlet 分布 $\mathrm{Dir}(\alpha_1, \ldots, \alpha_K)$ 在 $K \to \infty$ 时的推广。有限维 Dirichlet 分布是定义在 $K$ 维单纯形上的分布，其样本是 $K$ 维向量 $(w_1, \ldots, w_K)$；DP 则是其无限维推广，样本是无限维的离散分布。

历史上，Ferguson (1973) 正是通过有限 Dirichlet 分布的类比，构造了无穷维 Dirichlet 过程。这一构造在数学上优雅，在应用中有效，是 20 世纪贝叶斯统计最重要的进展之一。

\section{本章小结}

本章我们学习了：
\begin{enumerate}
  \item Dirichlet 过程是定义在概率分布空间上的随机过程，其样本路径几乎 surely 是离散分布；
  \item DP 有三种等价的构造方式：stick-breaking 构造、Pólya urn 构造和中国餐厅过程（CRP）；
  \item 作为未知分布的先验，DP 的后验有简洁的闭式形式；
  \item DP 混合模型（DPM）将 DP 先验与参数化成分分布结合，实现了无限混合的自然贝叶斯非参数建模；
  \item 主要的计算方法包括 blocked Gibbs 采样和 CRP 采样；
  \item DP 是有限 Dirichlet 分布在无限维的自然推广，由 Ferguson (1973) 开创。
\end{enumerate}

\section{下节预告}

下一章我们将转向更深入的贝叶斯计算方法，讨论如何在大规模问题中高效地进行贝叶斯推断。

\endinput


% Appendix
\appendix
% Appendix: Formula Derivations
% Complete derivations for Chapter 1-2

\appendix
\chapter{公式推导}

\section{概率与推断（Chapter 1）推导}

\subsection{贝叶斯定理的推导}\label{sec:bayes-theorem-derivation}

\textbf{背景（Background）}：
贝叶斯定理是概率论中最基本的公式之一，它描述了如何根据新的证据更新概率估计。这个推导基于条件概率的定义和基本运算规则。

\textbf{参数定义（Parameter Definitions）}：
\begin{itemize}
  \item $A$：某个事件
  \item $B$：另一个事件（通常代表观测数据或证据）
  \item $p(A)$：事件 $A$ 的先验概率（prior probability）
  \item $p(A|B)$：在事件 $B$ 发生的条件下，事件 $A$ 的后验概率（posterior probability）
  \item $p(B|A)$：在事件 $A$ 发生的条件下，事件 $B$ 的似然（likelihood）
  \item $p(B)$：事件 $B$ 的边际概率（marginal probability）
\end{itemize}

\textbf{已知条件（Given）}：
\begin{itemize}
  \item 条件概率定义：$p(A|B) = \dfrac{p(A \cap B)}{p(B)}$，其中 $p(B) > 0$
  \item 乘法法则：$p(A \cap B) = p(A) \cdot p(B|A) = p(B) \cdot p(A|B)$
\end{itemize}

\textbf{目标（Goal）}：
从条件概率的基本定义推导出贝叶斯定理：
\[
p(A|B) = \frac{p(A) \cdot p(B|A)}{p(B)}
\]

\textbf{推导步骤（Derivation Steps）}：

\textbf{步骤 1}：从条件概率定义出发
\[
p(A|B) = \frac{p(A \cap B)}{p(B)}
\tag*{条件概率定义}

\textbf{步骤 2}：用乘法法则替换联合概率 $p(A \cap B)$
\[
p(A \cap B) = p(A) \cdot p(B|A)
\tag*{乘法法则}

\textbf{步骤 3}：将步骤 2 代入步骤 1
\[
p(A|B) = \frac{p(A) \cdot p(B|A)}{p(B)}
\tag*{代入即得}

\textbf{步骤 4}：对 $p(B)$ 进行边际化（如果 $A_1, A_2, \ldots, A_k$ 是互斥且完备的事件族）
\[
p(B) = \sum_{i=1}^{k} p(A_i) \cdot p(B|A_i)
\tag*{全概率公式}

\textbf{步骤 5}：将步骤 4 代入贝叶斯定理的分母，得到完整形式
\[
p(A_j|B) = \frac{p(A_j) \cdot p(B|A_j)}{\sum_{i=1}^{k} p(A_i) \cdot p(B|A_i)}
\tag*{贝叶斯定理（完整形式）}

\textbf{结论}：贝叶斯定理建立了先验概率与后验概率之间的桥梁——我们通过似然 $p(B|A)$ 将新的观测 $B$ 与先验知识 $p(A)$ 结合起来。

\clearpage

\subsection{先验预测分布与后验预测分布的推导}\label{sec:predictive-distributions-derivation}

\textbf{背景（Background）}：
在贝叶斯统计中，我们不仅要推断当前参数 $\theta$ 的后验分布，还需要预测未来观测 $\tilde{y}$ 的分布。这涉及对未知参数进行积分（边际化）。

\textbf{参数定义（Parameter Definitions）}：
\begin{itemize}
  \item $\theta$：不可观测的参数
  \item $y$：已观测的数据
  \item $\tilde{y}$：未来或未观测的数据
  \item $p(\theta)$：参数的先验分布
  \item $p(y|\theta)$：数据的似然函数
  \item $p(\theta|y)$：参数的后验分布
\end{itemize}

\textbf{已知条件（Given）}：
\begin{itemize}
  \item 贝叶斯定理：$p(\theta|y) \propto p(\theta) \cdot p(y|\theta)$
  \item 边际概率定义：$p(y) = \int p(\theta) \cdot p(y|\theta) \, d\theta$
  \item 条件独立性：给定 $\theta$，未来观测 $\tilde{y}$ 与当前观测 $y$ 条件独立
\end{itemize}

\textbf{目标（Goal）}：
\begin{enumerate}
  \item 推导先验预测分布 $p(y)$
  \item 推导后验预测分布 $p(\tilde{y}|y)$
\end{enumerate}

\textbf{推导步骤（Derivation Steps）}：

\subsubsection*{先验预测分布的推导}

\textbf{步骤 1}：从边际概率的定义出发
\[
p(y) = \int p(\theta) \cdot p(y|\theta) \, d\theta
\tag*{边际概率定义}

\textbf{步骤 2}：解释其含义——这是对所有可能的 $\theta$ 值加权平均，其中权重是先验概率

\textbf{步骤 3}：先验预测分布表示在我们观测数据之前，所有可能观测结果的分布

\clearpage

\subsubsection*{后验预测分布的推导}

\textbf{步骤 1}：从条件概率定义出发
\[
p(\tilde{y}|y) = \int p(\tilde{y}, \theta|y) \, d\theta
\tag*{条件概率的边际化}

\textbf{步骤 2}：将联合分布分解
\[
p(\tilde{y}, \theta|y) = p(\tilde{y}|\theta, y) \cdot p(\theta|y)
\tag*{链式法则}

\textbf{步骤 3}：利用条件独立性假设（给定 $\theta$，$\tilde{y}$ 与 $y$ 独立）
\[
p(\tilde{y}|\theta, y) = p(\tilde{y}|\theta)
\tag*{条件独立性}

\textbf{步骤 4}：代入得
\[
p(\tilde{y}|y) = \int p(\tilde{y}|\theta) \cdot p(\theta|y) \, d\theta
\tag*{后验预测分布}

\textbf{结论}：后验预测分布是\textbf{在参数后验分布上对条件预测进行平均}。这体现了贝叶斯推断的核心思想：不确定性在两个层次上都存在（参数的不确定性和预测的随机性），我们需要对两者都进行边际化。

\clearpage

\subsection{后验优势比的推导}\label{sec:posterior-odds-derivation}

\textbf{背景（Background）}：
优势比（odds ratio）是概率的另一种表示方式，在许多应用场景下比概率更直观。贝叶斯定理在优势比形式下有一个特别简洁的表达。

\textbf{参数定义（Parameter Definitions）}：
\begin{itemize}
  \item $\theta_1$：某个参数值
  \item $\theta_2$：另一个参数值
  \item $p(\theta_1)$：$\theta_1$ 的先验概率
  \item $p(\theta_1|y)$：$\theta_1$ 的后验概率
  \item $p(y|\theta_1)$：似然函数在 $\theta_1$ 处的值
  \item $O(\theta_1,\theta_2) = \dfrac{p(\theta_1)}{p(\theta_2)}$：先验优势比
  \item $O(\theta_1,\theta_2|y) = \dfrac{p(\theta_1|y)}{p(\theta_2|y)}$：后验优势比
\end{itemize}

\textbf{已知条件（Given）}：
贝叶斯定理：
\[
p(\theta|y) = \frac{p(\theta) \cdot p(y|\theta)}{p(y)}
\]

\textbf{目标（Goal）}：
证明后验优势比 = 先验优势比 × 似然比

\textbf{推导步骤（Derivation Steps）}：

\textbf{步骤 1}：写出后验优势比的定义
\[
O(\theta_1,\theta_2|y) = \frac{p(\theta_1|y)}{p(\theta_2|y)}
\tag*{后验优势比定义}

\textbf{步骤 2}：对分子和分母分别应用贝叶斯定理
\[
= \frac{\dfrac{p(\theta_1) \cdot p(y|\theta_1)}{p(y)}}{\dfrac{p(\theta_2) \cdot p(y|\theta_2)}{p(y)}}
\tag*{代入贝叶斯定理}

\textbf{步骤 3}：约去 $p(y)$
\[
= \frac{p(\theta_1) \cdot p(y|\theta_1)}{p(\theta_2) \cdot p(y|\theta_2)}
\tag*{约去边际概率}

\textbf{步骤 4}：重新组合
\[
= \frac{p(\theta_1)}{p(\theta_2)} \times \frac{p(y|\theta_1)}{p(y|\theta_2)}
\tag*{即先验优势比乘以似然比}

\textbf{结论}：这个简洁的关系式告诉我们，\textbf{后验优势比 = 先验优势比 × 似然比}。数据通过似然比来更新优势比——如果 $p(y|\theta_1) > p(y|\theta_2)$，则数据支持 $\theta_1$，后验优势比增大。

\clearpage

\section{单参数模型（Chapter 2）推导}

\subsection{Beta-Binomial 共轭后验分布推导}\label{sec:beta-binomial-posterior-derivation}

\textbf{背景（Background）}：
在二项模型中，如果我们使用 Beta 先验分布来估计成功概率 $\theta$，则后验分布也是 Beta 分布。这种性质称为共轭性，是贝叶斯统计中最重要的概念之一。

\textbf{参数定义（Parameter Definitions）}：
\begin{itemize}
  \item $n$：伯努利试验的总次数
  \item $y$：$n$ 次试验中成功的次数
  \item $\theta$：每次试验的成功概率（未知参数）
  \item $\alpha$：Beta 先验的第一个参数
  \item $\beta$：Beta 先验的第二个参数
  \item $p(y|\theta) = \binom{n}{y} \theta^y (1-\theta)^{n-y}$：二项似然函数
  \item $p(\theta) = \dfrac{\Gamma(\alpha+\beta)}{\Gamma(\alpha)\Gamma(\beta)} \theta^{\alpha-1}(1-\theta)^{\beta-1}$：Beta 先验
\end{itemize}

\textbf{已知条件（Given）}：
\begin{itemize}
  \item Beta 函数定义：$B(\alpha, \beta) = \dfrac{\Gamma(\alpha)\Gamma(\beta)}{\Gamma(\alpha+\beta)}$
  \item 似然函数：$p(y|\theta) \propto \theta^y (1-\theta)^{n-y}$（忽略不依赖 $\theta$ 的常数 $\binom{n}{y}$）
  \item 先验分布：$p(\theta) \propto \theta^{\alpha-1}(1-\theta)^{\beta-1}$
\end{itemize}

\textbf{目标（Goal）}：
证明后验分布 $p(\theta|y)$ 也是 Beta 分布，并确定其后验参数。

\textbf{推导步骤（Derivation Steps）}：

\textbf{步骤 1}：应用贝叶斯定理写出非归一化后验分布
\[
p(\theta|y) \propto p(\theta) \cdot p(y|\theta)
\tag*{贝叶斯定理（非归一化形式）}

\textbf{步骤 2}：代入先验和似然
\[
p(\theta|y) \propto \theta^{\alpha-1}(1-\theta)^{\beta-1} \cdot \theta^y (1-\theta)^{n-y}
\tag*{代入先验和似然}

\textbf{步骤 3}：合并同底数指数
\[
= \theta^{(\alpha-1)+y} \cdot (1-\theta)^{(\beta-1)+(n-y)}
= \theta^{\alpha+y-1} \cdot (1-\theta)^{\beta+n-y-1}
\tag*{合并指数}

\textbf{步骤 4}：识别分布形式
这正是 Beta 分布的核（kernel），即
\[
p(\theta|y) \propto \theta^{\alpha_{\text{post}}-1}(1-\theta)^{\beta_{\text{post}}-1}
\]
其中
\[
\alpha_{\text{post}} = \alpha + y, \quad \beta_{\text{post}} = \beta + n - y
\tag*{后验参数}

\textbf{步骤 5}：写出完整的后验分布
\[
\theta|y \sim \text{Beta}(\alpha + y, \beta + n - y)
\tag*{后验分布}

\textbf{结论}：后验分布与先验分布属于同一形式（Beta 分布），这就是\textbf{共轭性}。先验参数 $\alpha$ 可以解释为 $\alpha - 1$ 次"先验成功"，$\beta$ 解释为 $\beta - 1$ 次"先验失败"。

\clearpage

\subsection{Beta-Binomial 后验均值与方差推导}\label{sec:beta-binomial-moments-derivation}

\textbf{背景（Background）}：
在共轭框架下，我们可以推导出后验分布的均值和方差的闭式表达式。这些表达式展示了先验信息与数据信息如何"折中"。

\textbf{参数定义（Parameter Definitions）}：
\begin{itemize}
  \item $\alpha_{\text{post}} = \alpha + y$：后验第一个参数
  \item $\beta_{\text{post}} = \beta + n - y$：后验第二个参数
  \item $\mathbb{E}(\theta|y)$：后验均值
  \item $\text{Var}(\theta|y)$：后验方差
\end{itemize}

\textbf{已知条件（Given）}：
\begin{itemize}
  \item 后验分布：$\theta|y \sim \text{Beta}(\alpha_{\text{post}}, \beta_{\text{post}})$
  \item Beta 分布均值公式：$\mathbb{E}(X) = \dfrac{\alpha}{\alpha+\beta}$，其中 $X \sim \text{Beta}(\alpha, \beta)$
  \item Beta 分布方差公式：$\text{Var}(X) = \dfrac{\alpha\beta}{(\alpha+\beta)^2(\alpha+\beta+1)}$
\end{itemize}

\textbf{目标（Goal）}：
\begin{enumerate}
  \item 求后验均值 $\mathbb{E}(\theta|y)$
  \item 求后验方差 $\text{Var}(\theta|y)$
\end{enumerate}

\clearpage

\subsubsection*{后验均值推导}

\textbf{步骤 1}：应用 Beta 分布均值公式
\[
\mathbb{E}(\theta|y) = \frac{\alpha_{\text{post}}}{\alpha_{\text{post}} + \beta_{\text{post}}}
\tag*{Beta 分布均值公式}

\textbf{步骤 2}：代入后验参数
\[
= \frac{\alpha + y}{\alpha + y + \beta + n - y}
= \frac{\alpha + y}{\alpha + \beta + n}
\tag*{代入参数}

\textbf{步骤 3}：改写为加权平均形式
\[
= \frac{\alpha + y}{\alpha + \beta + n}
= \frac{\alpha + \beta}{\alpha + \beta + n} \cdot \frac{\alpha}{\alpha + \beta}
  + \frac{n}{\alpha + \beta + n} \cdot \frac{y}{n}
\tag*{加权平均分解}

\textbf{结论 1}：后验均值是先验均值 $\dfrac{\alpha}{\alpha+\beta}$ 和样本比例 $\dfrac{y}{n}$ 的加权平均，权重与样本量和先验强度相关。

\clearpage

\subsubsection*{后验方差推导}

\textbf{步骤 1}：应用 Beta 分布方差公式
\[
\text{Var}(\theta|y) = \frac{\alpha_{\text{post}} \cdot \beta_{\text{post}}}{(\alpha_{\text{post}} + \beta_{\text{post}})^2 (\alpha_{\text{post}} + \beta_{\text{post}} + 1)}
\tag*{Beta 分布方差公式}

\textbf{步骤 2}：代入后验参数
\[
= \frac{(\alpha + y)(\beta + n - y)}{(\alpha + \beta + n)^2 (\alpha + \beta + n + 1)}
\tag*{代入参数}

\textbf{步骤 3}：大样本近似（当 $y$ 和 $n-y$ 较大时）
当 $n$ 很大而 $\alpha, \beta$ 固定时，
\[
\text{Var}(\theta|y) \approx \frac{1}{n} \cdot \frac{y}{n}\left(1 - \frac{y}{n}\right)
\tag*{大样本近似}

\textbf{结论 2}：后验方差以 $1/n$ 的速率趋近于零，表明随着样本量增大，后验分布越来越集中在真值附近。

\clearpage

\subsection{拉普拉斯继承法则推导}\label{sec:laplace-succession-derivation}

\textbf{背景（Background）}：
拉普拉斯在 18 世纪研究"明天太阳会升起"这类问题时，提出了一个著名的法则。在均匀先验下，这个法则预测未来成功的概率为 $\dfrac{y+1}{n+2}$。

\textbf{参数定义（Parameter Definitions）}：
\begin{itemize}
  \item $n$：已完成的试验次数
  \item $y$：成功的次数
  \item $\tilde{y}$：下一次试验的结果（$\tilde{y} = 1$ 表示成功）
  \item $\theta$：成功概率
  \item $p(\theta|y) = \text{Beta}(y+1, n-y+1)$：均匀先验下的后验分布
\end{itemize}

\textbf{已知条件（Given）}：
\begin{itemize}
  \item 均匀先验：$\theta \sim \text{Uniform}(0,1) = \text{Beta}(1,1)$
  \item 后验分布：$\theta|y \sim \text{Beta}(y+1, n-y+1)$
  \item 后验预测分布：$p(\tilde{y}=1|y) = \int_0^1 \theta \cdot p(\theta|y) \, d\theta$
\end{itemize}

\textbf{目标（Goal）}：
求下一次试验成功的概率 $p(\tilde{y}=1|y)$。

\textbf{推导步骤（Derivation Steps）}：

\textbf{步骤 1}：写出后验预测分布的定义
\[
p(\tilde{y}=1|y) = \int_0^1 p(\tilde{y}=1|\theta) \cdot p(\theta|y) \, d\theta
\tag*{后验预测分布定义}

\textbf{步骤 2}：代入 $p(\tilde{y}=1|\theta) = \theta$（给定 $\theta$，成功的条件概率就是 $\theta$）
\[
= \int_0^1 \theta \cdot p(\theta|y) \, d\theta
= \mathbb{E}(\theta|y)
\tag*{代入条件概率}

\textbf{步骤 3}：识别这是后验均值
\[
= \mathbb{E}(\theta|y)
\tag*{后验均值}

\textbf{步骤 4}：应用 Beta 分布均值公式
\[
= \frac{y+1}{(y+1) + (n-y+1)}
= \frac{y+1}{n+2}
\tag*{Beta 均值公式}

\textbf{结论}：拉普拉斯继承法则表明，在均匀先验下，已知 $n$ 次试验中 $y$ 次成功，下一次成功的概率是 $\dfrac{y+1}{n+2}$。这推广了"瑞士表匠"问题的直觉——即使观测到 $n$ 次成功，概率也不是 1（因为 $\theta$ 本身可能是变化的）。

\clearpage

\subsection{总期望与总方差公式推导}\label{sec:law-of-total-expectation-variance-derivation}

\textbf{背景（Background）}：
全期望公式和全方差公式是概率论中的两个基本恒等式，它们描述了如何通过条件分布来计算无条件矩。

\textbf{参数定义（Parameter Definitions）}：
\begin{itemize}
  \item $\theta$：随机变量（参数）
  \item $y$：随机变量（数据）
  \item $\mathbb{E}(\theta)$：$\theta$ 的无条件（先验）均值
  \item $\text{Var}(\theta)$：$\theta$ 的无条件（先验）方差
  \item $\mathbb{E}(\theta|y)$：给定 $y$ 后 $\theta$ 的条件均值
  \item $\text{Var}(\theta|y)$：给定 $y$ 后 $\theta$ 的条件方差
\end{itemize}

\textbf{已知条件（Given）}：
\begin{itemize}
  \item 条件期望定义：$\mathbb{E}(\theta|y)$ 是 $y$ 的函数
  \item 迭代期望法则：$\mathbb{E}[\mathbb{E}(\theta|y)] = \mathbb{E}(\theta)$
\end{itemize}

\textbf{目标（Goal）}：
\begin{enumerate}
  \item 证明全期望公式：$\mathbb{E}(\theta) = \mathbb{E}[\mathbb{E}(\theta|y)]$
  \item 证明全方差公式：$\text{Var}(\theta) = \mathbb{E}[\text{Var}(\theta|y)] + \text{Var}[\mathbb{E}(\theta|y)]$
\end{enumerate}

\clearpage

\subsubsection*{全期望公式推导}

\textbf{步骤 1}：从条件期望的定义出发
\[
\mathbb{E}[\mathbb{E}(\theta|y)] = \int \mathbb{E}(\theta|y) \cdot p(y) \, dy
\tag*{外层期望定义}

\textbf{步骤 2}：代入条件期望 $\mathbb{E}(\theta|y) = \int \theta \cdot p(\theta|y) \, d\theta$
\[
= \int \left( \int \theta \cdot p(\theta|y) \, d\theta \right) p(y) \, dy
\tag*{代入条件期望定义}

\textbf{步骤 3}：交换积分次序（由富比尼定理保证）
\[
= \int \theta \left( \int p(\theta|y) p(y) \, dy \right) d\theta
\tag*{交换积分次序}

\textbf{步骤 4}：识别内层积分 $p(\theta) = \int p(\theta|y) p(y) dy$（全概率）
\[
= \int \theta \cdot p(\theta) \, d\theta
= \mathbb{E}(\theta)
\tag*{边际分布定义}

\textbf{结论 1}：全期望公式表明，$\theta$ 的先验均值等于在所有可能数据上后验均值的加权平均。

\clearpage

\subsubsection*{全方差公式推导}

\textbf{步骤 1}：从方差定义出发
\[
\text{Var}(\theta) = \mathbb{E}(\theta^2) - [\mathbb{E}(\theta)]^2
\tag*{方差定义}

\textbf{步骤 2}：对 $\mathbb{E}(\theta^2)$ 应用全期望公式
\[
\mathbb{E}(\theta^2) = \mathbb{E}[\mathbb{E}(\theta^2|y)]
\tag*{全期望公式应用于 $\theta^2$}

\textbf{步骤 3}：在给定 $y$ 的条件下，方差可以写为
\[
\text{Var}(\theta|y) = \mathbb{E}(\theta^2|y) - [\mathbb{E}(\theta|y)]^2
\tag*{条件方差定义}

\textbf{步骤 4}：解出 $\mathbb{E}(\theta^2|y)$
\[
\mathbb{E}(\theta^2|y) = \text{Var}(\theta|y) + [\mathbb{E}(\theta|y)]^2
\tag*{移项}

\textbf{步骤 5}：代入步骤 2
\[
\mathbb{E}(\theta^2) = \mathbb{E}[\text{Var}(\theta|y)] + \mathbb{E}[[\mathbb{E}(\theta|y)]^2]
\tag*{代入}

\textbf{步骤 6}：对 $[\mathbb{E}(\theta|y)]^2$ 应用方差分解
\[
\mathbb{E}[[\mathbb{E}(\theta|y)]^2] = \text{Var}[\mathbb{E}(\theta|y)] + \{\mathbb{E}[\mathbb{E}(\theta|y)]\}^2
\tag*{方差与均值平方的关系}

\textbf{步骤 7}：代入全期望公式的结果 $\mathbb{E}[\mathbb{E}(\theta|y)] = \mathbb{E}(\theta)$
\[
= \text{Var}[\mathbb{E}(\theta|y)] + [\mathbb{E}(\theta)]^2
\tag*{简化}

\textbf{步骤 8}：代回方差的定义
\[
\text{Var}(\theta) = \mathbb{E}(\theta^2) - [\mathbb{E}(\theta)]^2
= \mathbb{E}[\text{Var}(\theta|y)] + \text{Var}[\mathbb{E}(\theta|y)]
\tag*{最终结果}

\textbf{结论 2}：全方差公式表明，\textbf{总方差 = 条件方差的均值 + 条件均值的方差}。这在贝叶斯语境下的解释是：参数的总不确定性来自两个方面——（1）给定数据后参数仍有的剩余不确定性（条件方差的均值），以及（2）后验均值随数据变化的不确定性（条件均值的方差）。

\clearpage

\subsection{正态均值估计的后验分布推导}\label{sec:normal-mean-posterior-derivation}

\textbf{背景（Background）}：
当我们用正态似然估计正态均值时，如果使用正态先验，则后验分布也是正态分布。这个推导展示了共轭先验在正态模型中的应用。

\textbf{参数定义（Parameter Definitions）}：
\begin{itemize}
  \item $y \sim \text{Normal}(\theta, \sigma^2)$：观测数据（已知方差 $\sigma^2$）
  \item $\theta \sim \text{Normal}(\mu_0, \tau_0^2)$：正态先验
  \item $\mu_0$：先验均值
  \item $\tau_0^2$：先验方差
  \item $\sigma^2$：已知的观测方差
  \item $\mu_1, \tau_1^2$：后验均值和方差
\end{itemize}

\textbf{已知条件（Given）}：
\begin{itemize}
  \item 似然函数：$p(y|\theta) \propto \exp\left(-\frac{(y-\theta)^2}{2\sigma^2}\right)$
  \item 先验密度：$p(\theta) \propto \exp\left(-\frac{(\theta-\mu_0)^2}{2\tau_0^2}\right)$
  \item 配方法（completing the square）技巧
\end{itemize}

\textbf{目标（Goal）}：
证明后验分布 $\theta|y \sim \text{Normal}(\mu_1, \tau_1^2)$，并求出 $\mu_1$ 和 $\tau_1^2$。

\textbf{推导步骤（Derivation Steps）}：

\textbf{步骤 1}：写出非归一化后验分布的对数
\[
\log p(\theta|y) \propto \log p(\theta) + \log p(y|\theta)
\tag*{对数后验（忽略常数）}

\textbf{步骤 2}：代入先验和似然的对数
\[
= -\frac{(\theta-\mu_0)^2}{2\tau_0^2} - \frac{(y-\theta)^2}{2\sigma^2}
\tag*{代入}

\textbf{步骤 3}：展开平方项
\[
= -\frac{1}{2}\left[\frac{(\theta^2 - 2\mu_0\theta + \mu_0^2)}{\tau_0^2} + \frac{(\theta^2 - 2y\theta + y^2)}{\sigma^2}\right]
\tag*{展开}

\textbf{步骤 4}：合并 $\theta^2$ 项的系数
\[
= -\frac{1}{2}\left[\theta^2\left(\frac{1}{\tau_0^2} + \frac{1}{\sigma^2}\right) - 2\theta\left(\frac{\mu_0}{\tau_0^2} + \frac{y}{\sigma^2}\right) + \text{常数}\right]
\tag*{合并}

\textbf{步骤 5}：识别后验精度（方差的倒数）
\[
\frac{1}{\tau_1^2} = \frac{1}{\tau_0^2} + \frac{1}{\sigma^2}
\tag*{后验精度}

\textbf{步骤 6}：由配方法，后验均值满足
\[
\theta \cdot \frac{1}{\tau_1^2} = \frac{\mu_0}{\tau_0^2} + \frac{y}{\sigma^2}
\tag*{配方法条件}

\textbf{步骤 7}：解出后验均值
\[
\mu_1 = \frac{\mu_0/\tau_0^2 + y/\sigma^2}{1/\tau_0^2 + 1/\sigma^2}
\tag*{后验均值}

\textbf{步骤 8}：写出完整的后验分布
\[
\theta|y \sim \text{Normal}\left(\mu_1, \tau_1^2\right)
\tag*{后验分布}

\textbf{结论}：
\begin{enumerate}
  \item \textbf{精度加法}：后验精度 = 先验精度 + 数据精度
  \item \textbf{后验均值}：是先验均值 $\mu_0$ 和观测值 $y$ 的加权平均，权重与精度成正比
\end{enumerate}

\clearpage

\subsection{Jeffreys 先验与 Fisher 信息推导}\label{sec:jeffreys-prior-derivation}

\textbf{背景（Background）}：
Jeffreys 提出了一种构建非信息先验的系统方法，基于 Fisher 信息的不变性原则。这个先验在参数的单调变换下保持不变。

\textbf{参数定义（Parameter Definitions）}：
\begin{itemize}
  \item $\theta$：参数
  \item $p(y|\theta)$：似然函数
  \item $J(\theta)$：Fisher 信息
  \item $L(\theta) = \log p(y|\theta)$：对数似然
  \item $I(\theta) = -\mathbb{E}\left[\frac{\partial^2 \log p(y|\theta)}{\partial \theta^2}\right]$：Fisher 信息（另一种定义）
  \item $\phi = h(\theta)$：参数的单调变换
\end{itemize}

\textbf{已知条件（Given）}：
\begin{itemize}
  \item Jeffreys 原则：先验应在参数的单调变换下保持不变
  \item Fisher 信息定义：$J(\theta) = -\mathbb{E}\left[\frac{\partial^2 \log p(y|\theta)}{\partial \theta^2}\right]$
  \item 变换公式：$p(\phi) = p(\theta) \cdot \left|\frac{d\theta}{d\phi}\right|$
\end{itemize}

\textbf{目标（Goal）}：
推导 Jeffreys 先验公式并验证其在变换下的不变性。

\textbf{推导步骤（Derivation Steps）}：

\subsubsection*{Jeffreys 先验公式推导}

\textbf{步骤 1}：Jeffreys 提出先验正比于 Fisher 信息的平方根
\[
p(\theta) \propto [J(\theta)]^{1/2}
\tag*{Jeffreys 先验}

\textbf{步骤 2}：验证在变换 $\phi = h(\theta)$ 下的不变性

在新参数 $\phi$ 下，Fisher 信息为
\[
J(\phi) = -\mathbb{E}\left[\frac{\partial^2 \log p(y|\phi)}{\partial \phi^2}\right]
\tag*{Fisher 信息在新参数下}

\textbf{步骤 3}：利用链式法则
\[
\frac{\partial \log p(y|\theta)}{\partial \phi}
= \frac{\partial \log p(y|\theta)}{\partial \theta} \cdot \frac{d\theta}{d\phi}
\tag*{链式法则}

\textbf{步骤 4}：计算二阶导数
\[
\frac{\partial^2 \log p(y|\phi)}{\partial \phi^2}
= \frac{\partial}{\partial \phi}\left(\frac{\partial \log p(y|\theta)}{\partial \theta} \cdot \frac{d\theta}{d\phi}\right)
= \frac{\partial^2 \log p(y|\theta)}{\partial \theta^2} \left(\frac{d\theta}{d\phi}\right)^2 + \frac{\partial \log p(y|\theta)}{\partial \theta} \cdot \frac{d^2\theta}{d\phi^2}
\tag*{求导}

\textbf{步骤 5}：取期望，利用 $\mathbb{E}\left[\frac{\partial \log p(y|\theta)}{\partial \theta}\right] = 0$（得分函数的期望为零）
\[
J(\phi) = \left(\frac{d\theta}{d\phi}\right)^2 J(\theta)
\tag*{Fisher 信息的变换公式}

\textbf{步骤 6}：因此
\[
[J(\phi)]^{1/2} = \left|\frac{d\theta}{d\phi}\right| [J(\theta)]^{1/2}
\tag*{平方根}

\textbf{步骤 7}：新先验为
\[
p(\phi) \propto [J(\phi)]^{1/2} \propto [J(\theta)]^{1/2} \left|\frac{d\theta}{d\phi}\right|
\tag*{变换后的 Jeffreys 先验}

\textbf{结论}：这正是概率密度在变量变换下的标准变换公式，因此 Jeffreys 先验在单调变换下是不变的。

\clearpage

\subsubsection*{二项模型的 Jeffreys 先验推导}

\textbf{步骤 1}：写出二项模型的对数似然
\[
\log p(y|\theta) = \log \binom{n}{y} + y\log\theta + (n-y)\log(1-\theta)
\tag*{对数似然}

\textbf{步骤 2}：计算一阶导数（得分函数）
\[
\frac{\partial}{\partial\theta} \log p(y|\theta) = \frac{y}{\theta} - \frac{n-y}{1-\theta}
\tag*{一阶导数}

\textbf{步骤 3}：计算二阶导数
\[
\frac{\partial^2}{\partial\theta^2} \log p(y|\theta) = -\frac{y}{\theta^2} - \frac{n-y}{(1-\theta)^2}
\tag*{二阶导数}

\textbf{步骤 4}：取负期望（注意 $\mathbb{E}(Y) = n\theta$）
\[
J(\theta) = -\mathbb{E}\left[\frac{\partial^2}{\partial\theta^2} \log p(y|\theta)\right]
= \frac{n\theta}{\theta^2} + \frac{n(1-\theta)}{(1-\theta)^2}
= \frac{n}{\theta(1-\theta)}
\tag*{Fisher 信息}

\textbf{步骤 5}：应用 Jeffreys 公式
\[
p(\theta) \propto \left[\frac{n}{\theta(1-\theta)}\right]^{1/2}
\propto \theta^{-1/2}(1-\theta)^{-1/2}
\tag*{Jeffreys 先验}

\textbf{步骤 6}：识别为 Beta 分布
\[
p(\theta) \propto \theta^{(1/2)-1}(1-\theta)^{(1/2)-1}
\Rightarrow \theta \sim \text{Beta}\left(\frac{1}{2}, \frac{1}{2}\right)
\tag*{Beta(1/2, 1/2)}

\textbf{结论}：二项模型的 Jeffreys 先验是 $\text{Beta}(1/2, 1/2)$，这是一个弱信息先验，比均匀先验 $\text{Beta}(1,1)$ 在边界处更轻。

\clearpage

\subsection{后验预测分布的均值与方差推导}\label{sec:predictive-mean-variance-derivation}

\textbf{背景（Background）}：
后验预测分布 $\tilde{y}|y$ 的均值和方差有两个组成部分——一个来自模型本身的随机性，另一个来自参数的后验不确定性。

\textbf{参数定义（Parameter Definitions）}：
\begin{itemize}
  \item $\tilde{y}$：未来观测
  \item $y$：当前观测
  \item $\theta$：参数
  \item $\mathbb{E}(\tilde{y}|y)$：后验预测均值
  \item $\text{Var}(\tilde{y}|y)$：后验预测方差
\end{itemize}

\textbf{已知条件（Given）}：
\begin{itemize}
  \item $\tilde{y}|\theta \sim \text{Normal}(\theta, \sigma^2)$
  \item $\theta|y \sim \text{Normal}(\mu_1, \tau_1^2)$
  \item 全期望公式：$\mathbb{E}(\tilde{y}|y) = \mathbb{E}[\mathbb{E}(\tilde{y}|\theta, y)]$
  \item 全方差公式：$\text{Var}(\tilde{y}|y) = \mathbb{E}[\text{Var}(\tilde{y}|\theta, y)] + \text{Var}[\mathbb{E}(\tilde{y}|\theta, y)]$
\end{itemize}

\textbf{目标（Goal）}：
求后验预测分布 $\tilde{y}|y$ 的均值和方差。

\clearpage

\subsubsection*{后验预测均值推导}

\textbf{步骤 1}：应用全期望公式
\[
\mathbb{E}(\tilde{y}|y) = \mathbb{E}[\mathbb{E}(\tilde{y}|\theta, y)]
\tag*{全期望公式}

\textbf{步骤 2}：给定 $\theta$，$\tilde{y}$ 与 $y$ 条件独立，且 $\mathbb{E}(\tilde{y}|\theta) = \theta$
\[
= \mathbb{E}(\theta|y)
= \mu_1
\tag*{后验均值}

\textbf{结论 1}：后验预测均值等于参数的后验均值。

\clearpage

\subsubsection*{后验预测方差推导}

\textbf{步骤 1}：应用全方差公式
\[
\text{Var}(\tilde{y}|y) = \mathbb{E}[\text{Var}(\tilde{y}|\theta, y)] + \text{Var}[\mathbb{E}(\tilde{y}|\theta, y)]
\tag*{全方差公式}

\textbf{步骤 2}：第一项——给定 $\theta$ 下的方差是 $\sigma^2$（模型方差）
\[
\mathbb{E}[\text{Var}(\tilde{y}|\theta, y)] = \mathbb{E}(\sigma^2|y) = \sigma^2
\tag*{模型方差}

\textbf{步骤 3}：第二项——条件均值的方差
\[
\text{Var}[\mathbb{E}(\tilde{y}|\theta, y)] = \text{Var}(\theta|y) = \tau_1^2
\tag*{参数不确定性}

\textbf{步骤 4}：合并
\[
\text{Var}(\tilde{y}|y) = \sigma^2 + \tau_1^2
\tag*{后验预测方差}

\textbf{结论 2}：后验预测方差 = \textbf{模型方差} + \textbf{参数后验方差}。这体现了贝叶斯推断的核心：预测的不确定性来自两个方面——（1）模型本身的随机性（$\sigma^2$），以及（2）参数估计的不确定性（$\tau_1^2$）。

% ================================================
% Chapter 11 推导：马尔可夫链模拟基础
% =============================================

\section{拒绝抽样的正确性证明}\label{sec:derivation-rejection-sampling}

\textbf{背景（Background）}：
拒绝抽样通过随机接受机制从目标分布中抽取样本。我们需要证明被接受的样本确实服从目标分布 $p(\theta)$。

\textbf{参数定义（Parameter Definitions）}：
\begin{itemize}
  \item $p(\theta)$：目标（未归一化）分布
  \item $g(\theta)$：提议分布
  \item $M$：覆盖常数，满足 $p(\theta) \leq M \cdot g(\theta)$
  \item $\theta^*$：候选样本
  \item $u^*$：从 $\mathrm{Uniform}(0,1)$ 抽取的随机数
\end{itemize}

\textbf{目标（Goal）}：
证明被接受的样本 $\theta_{\text{accept}}$ 服从目标分布 $p(\theta)$。

\textbf{推导步骤（Derivation Steps）}：

\textbf{步骤 1}：写出接受的联合概率
\[
\mathbb{P}(\text{接受} \cap \theta^* \in d\theta)
= \mathbb{P}(\theta^* \in d\theta) \cdot \mathbb{P}(\text{接受} | \theta^*)
= g(\theta^*) d\theta \cdot \frac{p(\theta^*)}{M g(\theta^*)}
= \frac{p(\theta^*)}{M} d\theta
\tag*{接受联合概率}
\]

\textbf{步骤 2}：计算总接受概率（归一化常数）
\[
\mathbb{P}(\text{接受}) = \int \frac{p(\theta^*)}{M} d\theta = \frac{1}{M}
\tag*{总接受概率}
\]

\textbf{步骤 3}：计算接受样本的条件分布
\[
\mathbb{P}(\theta_{\text{accept}} \in d\theta | \text{接受})
= \frac{\mathbb{P}(\text{接受} \cap \theta^* \in d\theta)}{\mathbb{P}(\text{接受})}
= \frac{p(\theta^*)/M}{1/M} d\theta = p(\theta^*) d\theta
\tag*{条件分布 = 目标分布}
\]

\textbf{结论}：被接受的样本的条件分布正是目标分布 $p(\theta)$。证明完毕。

\clearpage

\subsection{拒绝抽样期望接受率推导}\label{sec:derivation-rejection-acceptance}

\textbf{背景（Background）}：
拒绝抽样的期望接受率是衡量算法效率的关键指标。

\textbf{目标（Goal）}：
推导期望接受率的公式。

\textbf{推导步骤（Derivation Steps）}：

\textbf{步骤 1}：期望接受率定义为
\[
\text{accept rate} = \int \mathbb{P}(\text{接受} | \theta) \cdot g(\theta) \, d\theta
\tag*{期望定义}
\]

\textbf{步骤 2}：代入接受概率
\[
\mathbb{P}(\text{接受} | \theta) = \min\!\left(1, \frac{p(\theta)}{M g(\theta)}\right)
\tag*{接受概率}
\]

\textbf{步骤 3}：于是
\[
\text{accept rate} = \int \min\!\left(\frac{p(\theta)}{M g(\theta)}, 1\right) \cdot g(\theta) \, d\theta
= \frac{1}{M} \int \min\!\bigl(p(\theta), M g(\theta)\bigr) \, d\theta
\tag*{积分结果}
\]

\textbf{步骤 4}：当覆盖条件 $M g(\theta) \geq p(\theta)$ 对所有 $\theta$ 成立时
\[
\text{accept rate} = \frac{1}{M} \int p(\theta) \, d\theta = \frac{1}{M}
\tag*{简化为 $1/M$}
\]

\textbf{结论}：期望接受率与 $M$ 成反比，选择尽可能小的 $M$ 可提高接受率。

\clearpage

\subsection{细致平衡条件推导}\label{sec:derivation-detailed-balance}

\textbf{背景（Background）}：
细致平衡是构造满足特定平稳分布的马尔可夫链的关键条件。

\textbf{目标（Goal）}：
证明如果转移核 $P$ 满足细致平衡，则其平稳分布就是 $\pi$。

\textbf{推导步骤（Derivation Steps）}：

\textbf{步骤 1}：写出细致平衡条件
\[
P(\theta' | \theta) \cdot \pi(\theta) = P(\theta | \theta') \cdot \pi(\theta')
\tag*{细致平衡}
\]

\textbf{步骤 2}：验证 $\pi$ 是平稳分布
\[
\int P(\theta' | \theta) \cdot \pi(\theta) \, d\theta
= \int P(\theta | \theta') \cdot \pi(\theta') \, d\theta
= \pi(\theta') \int P(\theta | \theta') \, d\theta
= \pi(\theta') \cdot 1
= \pi(\theta')
\tag*{平稳分布验证}
\]

其中第三步用到 $\int P(\theta | \theta') d\theta = 1$（转移核的归一化性）。

\textbf{结论}：细致平衡是平稳分布的充分条件，但不是必要条件。

\clearpage

\subsection{Metropolis 算法细致平衡证明}\label{sec:derivation-metropolis}

\textbf{背景（Background）}：
Metropolis 算法的提议分布是对称的，我们需要证明其满足细致平衡。

\textbf{目标（Goal）}：
证明 Metropolis 算法的转移核以目标分布 $p(\theta)$ 为平稳分布。

\textbf{推导步骤（Derivation Steps）}：

\textbf{步骤 1}：Metropolis 转移核
\[
P(\theta' | \theta) = q(\theta' | \theta) \cdot \min\!\left(1, \frac{p(\theta')}{p(\theta)}\right)
\]
其中 $q(\theta' | \theta) = q(\theta | \theta')$（对称提议分布）。

\textbf{步骤 2}：验证细致平衡
\[
P(\theta' | \theta) \cdot p(\theta)
= q(\theta' | \theta) \cdot \min(1, \frac{p(\theta')}{p(\theta)}) \cdot p(\theta)
\]

分两种情况讨论：

\textbf{情况 1}：$p(\theta') \geq p(\theta)$
\[
\text{右边} = q(\theta' | \theta) \cdot 1 \cdot p(\theta)
= q(\theta' | \theta) \cdot p(\theta')
\]

\textbf{情况 2}：$p(\theta') < p(\theta)$
\[
\text{右边} = q(\theta' | \theta) \cdot \frac{p(\theta')}{p(\theta)} \cdot p(\theta)
= q(\theta' | \theta) \cdot p(\theta')
\]

由于 $q(\theta' | \theta) = q(\theta | \theta')$，两种情况都满足
\[
P(\theta' | \theta) \cdot p(\theta) = P(\theta | \theta') \cdot p(\theta')
\tag*{细致平衡成立}
\]

\textbf{结论}：Metropolis 算法的平稳分布正是目标分布 $p(\theta)$。证明完毕。

\clearpage

\subsection{Gibbs 抽样接受率推导}\label{sec:derivation-gibbs-acceptance}

\textbf{背景（Background）}：
Gibbs 抽样是 Metropolis-Hastings 的特例，其接受率恒为 1。我们来推导这个结论。

\textbf{参数定义（Parameter Definitions）}：
\begin{itemize}
  \item $\theta = (\theta_1, \ldots, \theta_d)$：完整参数向量
  \item $\theta_{-j} = (\theta_1, \ldots, \theta_{j-1}, \theta_{j+1}, \ldots, \theta_d)$：除 $\theta_j$ 外的所有参数
  \item $p^*(\theta)$：未归一化的目标分布
  \item $q(\theta_j' | \theta_{-j}) = p(\theta_j' | \theta_{-j}, y)$：提议分布（条件分布）
\end{itemize}

\textbf{目标（Goal）}：
证明 Gibbs 抽样的接受率为 1。

\textbf{推导步骤（Derivation Steps）}：

\textbf{步骤 1}：写出 MH 接受比
\[
r = \frac{p^*(\theta_j', \theta_{-j})}{p^*(\theta_j, \theta_{-j})}
\cdot \frac{q(\theta_j | \theta_{-j})}{q(\theta_j' | \theta_{-j})}
\tag*{MH 接受比}
\]

\textbf{步骤 2}：代入 Gibbs 的提议分布
\[
q(\theta_j | \theta_{-j}) = p(\theta_j | \theta_{-j}, y)
\]
\[
q(\theta_j' | \theta_{-j}) = p(\theta_j' | \theta_{-j}, y)
\tag*{提议分布 = 条件分布}
\]

\textbf{步骤 3}：利用链式法则分解联合分布
\[
p^*(\theta_j', \theta_{-j}) = p^*(\theta_{-j}) \cdot p^*(\theta_j' | \theta_{-j})
\]
\[
p^*(\theta_j, \theta_{-j}) = p^*(\theta_{-j}) \cdot p^*(\theta_j | \theta_{-j})
\tag*{链式法则}
\]

\textbf{步骤 4}：代入接受比
\[
r = \frac{p^*(\theta_{-j}) \cdot p^*(\theta_j' | \theta_{-j})}{p^*(\theta_{-j}) \cdot p^*(\theta_j | \theta_{-j})}
\cdot \frac{p(\theta_j | \theta_{-j}, y)}{p(\theta_j' | \theta_{-j}, y)}
= \frac{p^*(\theta_j' | \theta_{-j})}{p^*(\theta_j | \theta_{-j})}
\cdot \frac{p(\theta_j | \theta_{-j}, y)}{p(\theta_j' | \theta_{-j}, y)}
\tag*{代入}
\]

\textbf{步骤 5}：注意 $p^*(\theta_j | \theta_{-j}, y) \propto p(\theta_j | \theta_{-j}, y)$，因为 $p^*(\theta_{-j})$ 不依赖于 $\theta_j$。于是两个因子互为倒数，
\[
r = 1
\tag*{接受率为 1}
\]

\textbf{结论}：Gibbs 抽样的提议分布恰好是条件分布，因此接受率恒为 1，无需拒绝。

\clearpage

\subsection{分层模型条件分布推导}\label{sec:derivation-hierarchical-conditionals}

\textbf{背景（Background）}：
在分层模型 \cref{eq:hierarchical-model} 中，我们需要推导各个参数的条件后验分布。

\textbf{模型回顾}：
\[
y_j | \theta_j \sim \mathrm{N}(\theta_j, \sigma^2), \quad \theta_j | \mu, \tau \sim \mathrm{N}(\mu, \tau^2)
\]
其中 $\sigma^2$ 已知。

\textbf{目标（Goal）}：
推导 $p(\theta_j | \mu, \tau, y)$、$p(\mu | \theta, \tau, y)$ 和 $p(\tau | \theta, \mu, y)$。

\clearpage

\subsubsection*{$p(\theta_j | \mu, \tau, y)$ 推导}

\textbf{步骤 1}：写出条件分布的核
\[
p(\theta_j | \mu, \tau, y) \propto p(y_j | \theta_j) \cdot p(\theta_j | \mu, \tau)
\propto \exp\!\left(-\frac{(y_j - \theta_j)^2}{2\sigma^2}\right)
\cdot \exp\!\left(-\frac{(\theta_j - \mu)^2}{2\tau^2}\right)
\tag*{核}
\]

\textbf{步骤 2}：合并指数
\[
= \exp\!\left(-\frac{(y_j - \theta_j)^2}{2\sigma^2} - \frac{(\theta_j - \mu)^2}{2\tau^2}\right)
\tag*{合并}
\]

\textbf{步骤 3}：配方得到正态分布形式
\[
= \exp\!\left(-\frac{1}{2}\left[\frac{1}{\sigma^2} + \frac{1}{\tau^2}\right]
\left(\theta_j - \frac{y_j/\sigma^2 + \mu/\tau^2}{1/\sigma^2 + 1/\tau^2}\right)^2\right)
\tag*{配方}
\]

因此
\[
\theta_j | \mu, \tau, y \sim \mathrm{N}\!\left(\frac{y_j/\sigma^2 + \mu/\tau^2}{1/\sigma^2 + 1/\tau^2},
\left(\frac{1}{\sigma^2} + \frac{1}{\tau^2}\right)^{-1}\right)
\tag*{条件后验}
\]

\clearpage

\subsubsection*{$p(\mu | \theta, \tau, y)$ 推导}

\textbf{步骤 1}：$\mu$ 的条件分布（似然部分只依赖于 $\theta_j$，与 $\mu$ 通过 $\theta_j | \mu, \tau$ 连接）
\[
p(\mu | \theta, \tau) \propto p(\theta | \mu, \tau) = \prod_{j=1}^J \exp\!\left(-\frac{(\theta_j - \mu)^2}{2\tau^2}\right)
\propto \exp\!\left(-\frac{\sum_j (\theta_j - \mu)^2}{2\tau^2}\right)
\tag*{先验部分}
\]

\textbf{步骤 2}：展开平方
\[
\sum_j (\theta_j - \mu)^2 = \sum_j (\theta_j - \bar{\theta} + \bar{\theta} - \mu)^2
= \sum_j (\theta_j - \bar{\theta})^2 + J(\bar{\theta} - \mu)^2
\tag*{平方展开}
\]

\textbf{步骤 3}：识别为正态核
\[
p(\mu | \theta, \tau) \propto \exp\!\left(-\frac{J(\bar{\theta} - \mu)^2}{2\tau^2}\right)
\tag*{正态核}
\]

因此
\[
\mu | \theta, \tau \sim \mathrm{N}\!\left(\bar{\theta}, \frac{\tau^2}{J}\right)
\tag*{条件后验}
\]

其中 $\bar{\theta} = \frac{1}{J}\sum_j \theta_j$。

\clearpage

\subsubsection*{$p(\tau | \theta, \mu, y)$ 的非标准形式}

\textbf{说明}：
\[
p(\tau | \theta, \mu, y) \propto p(\theta | \mu, \tau) \cdot p(\tau)
= \prod_{j=1}^J \frac{1}{\tau} \exp\!\left(-\frac{(\theta_j - \mu)^2}{2\tau^2}\right) \cdot p(\tau)
\tag*{非标准形式}
\]

这个分布不是标准形式（不是常见的正态、Gamma 等分布），因此需要使用 Metropolis-Hastings 步骤来抽样。

\textbf{结论}：在 Gibbs 抽样框架下，对 $\tau$ 需要使用 MH 步骤，这正是"混合 MCMC"的典型例子。

% ================================================
% 问答与注解
% ================================================
% Appendix: Questions and Answers
% User annotations and interactive Q&A records

\section{阅读过程中的问题与解答}

% === 用户问答记录 ===
% 本节记录阅读过程中的问答内容

\section{什么是分布的核}\label{sec:qa:DistributionKernel}

\subsection{问题}
分布的核是什么概念？

\subsection{回答}
分布的\textbf{核（kernel）}是指概率密度/质量函数中\textbf{依赖于参数的核心部分}，不包括归一化常数。

\textbf{举例说明}：

Beta-Binomial 共轭中，后验分布：
\[
p(\theta|y) \propto \theta^{\alpha+y-1}(1-\theta)^{\beta+n-y-1}
\]

其中 $\theta^{\alpha+y-1}(1-\theta)^{\beta+n-y-1}$ 是\textbf{核}，而完整的 Beta 分布还需要归一化常数：
\[
p(\theta|y) = \frac{\Gamma(\alpha+\beta+n)}{\Gamma(\alpha+y)\Gamma(\beta+n-y)} \theta^{\alpha+y-1}(1-\theta)^{\beta+n-y-1}
\]

\textbf{为什么核重要}：

\begin{enumerate}
  \item \textbf{识别分布形式}：知道核是 $\theta^a(1-\theta)^b$ 形式，就能识别这是 Beta 分布
  \item \textbf{贝叶斯推断中省略 $p(y)$}：$p(y)$ 是归一化常数，不包含关于 $\theta$ 的信息，所以可以省略
  \item \textbf{简化推导}：在推导过程中只关注核，省去繁琐的归一化常数
\end{enumerate}

★ Insight ─────────────────────────────────────
这本质上是\textbf{信息局部性}思想：归一化常数的作用只是"让面积等于1"，它本身不携带关于参数的信息。真正"有意义"的概率质量都集中在核上。
─────────────────────────────────────────────────

% === QA 记录格式 ===
% 使用以下格式记录问答：
%
% \section{问题标题}
% \label{sec:qa:LabelKey}
%
% \subsection{问题}
% 用户的原始问题
%
% \subsection{回答}
% 回答内容...
%
% 如需详细推导，添加：
% \footnote{推导见附录 \cref{sec:xxx}}

\section{学习笔记}

\textit{在这里记录你的学习心得和思考。}

% === 学习笔记格式 ===
% \section{笔记标题}
% \label{sec:notes:LabelKey}
%
% 笔记内容...


\endinput


\end{document}
